\documentclass[11pt,openany,UTF8]{ctexbook}
\usepackage{amsmath,amssymb}
\usepackage{array,booktabs,longtable}
\usepackage{fancyvrb}
\usepackage{xcolor}
\usepackage[a4paper,top=2.4cm,bottom=2.4cm,left=2.3cm,right=2.3cm,headsep=0.6cm]{geometry}
\usepackage{fancyhdr}
\usepackage{needspace}
\usepackage{enumitem}
\usepackage[hidelinks,bookmarksnumbered=true]{hyperref}
\usepackage{graphicx}
\providecommand{\iddots}{\reflectbox{$\ddots$}}
\providecommand{\xlongequal}[1]{\xrightarrow{\ #1\ }}
\providecommand{\rank}{\operatorname{rank}}
\providecommand{\tr}{\operatorname{tr}}


% ---- fonts
\setCJKmainfont{Songti SC}[BoldFont=Heiti SC]
\xeCJKsetup{CJKmath=true}
\xeCJKDeclareCharClass{CJK}{"2500 -> "257F, "25A0 -> "25FF, "2160 -> "2183, "2190 -> "21FF}
\setCJKsansfont{Heiti SC}
\setCJKmonofont{Songti SC}
\setmainfont{Times New Roman}

% ---- colours (restrained)
\definecolor{accent}{HTML}{1A1A1A}
\definecolor{keybar}{HTML}{1F4E79}
\definecolor{warnbar}{HTML}{9C2A20}
\definecolor{exambar}{HTML}{4A4A4A}
\definecolor{notebar}{HTML}{2E6B4F}
\definecolor{codebg}{HTML}{F4F4F4}
\definecolor{rulesep}{HTML}{BFBFBF}

% ---- table column type
\newcolumntype{L}[1]{>{\raggedright\arraybackslash}p{#1}}
\newcommand{\ccaret}{\ifmmode\hat{}\else\textasciicircum{}\fi}
\newcommand{\circled}[1]{\textcircled{\scriptsize\,#1}}
\newcommand{\ctilde}{\ifmmode\tilde{}\else\textasciitilde{}\fi}

% ---- code blocks
\DefineVerbatimEnvironment{codebox}{Verbatim}%
  {frame=single,framerule=0.3pt,rulecolor=\color{rulesep},%
   framesep=6pt,fontsize=\small,baselinestretch=0.98,%
   xleftmargin=4pt,xrightmargin=4pt}

% ---- callout environments (breakable: only a bar + rule, no box)
\newcommand{\callouthead}[2]{%
  \par\addvspace{7pt}\noindent
  {\color{#1}\rule[-0.15em]{3pt}{1.05em}}\hspace{0.5em}%
  {\bfseries\color{#1}#2}\par\nobreak\vspace{-0.35em}%
  {\color{rulesep}\hrule height 0.4pt}\nobreak\vspace{0.45em}%
  \begingroup\small\setlength{\parindent}{0pt}%
  \leftskip=1.1em\rightskip=0.6em\noindent\ignorespaces}
\newcommand{\calloutend}{\par\endgroup\addvspace{7pt}}

\newenvironment{callkey}[1]{\callouthead{keybar}{#1}}{\calloutend}
\newenvironment{callwarn}[1]{\callouthead{warnbar}{#1}}{\calloutend}
\newenvironment{callexam}[1]{\callouthead{exambar}{#1}}{\calloutend}
\newenvironment{callnote}[1]{\callouthead{notebar}{#1}}{\calloutend}

% ---- headings
\ctexset{
  chapter = {format = \raggedright\Large\bfseries, name = {第,章}, number = \chinese{chapter},
             aftername = {\quad}, beforeskip = 0pt, afterskip = 18pt, pagestyle = plain},
  section = {format = \raggedright\large\bfseries, name = {,、}, number = \chinese{section},
             aftername = {}, beforeskip = 16pt, afterskip = 8pt},
  subsection = {format = \raggedright\normalsize\bfseries, name = {（,）}, number = \chinese{subsection},
             aftername = {}, beforeskip = 11pt, afterskip = 6pt},
  subsubsection = {format = \raggedright\normalsize\itshape, name = {,.}, number = \arabic{subsubsection},
             aftername = {}, beforeskip = 9pt, afterskip = 4pt},
}
\setcounter{secnumdepth}{3}
\setcounter{tocdepth}{2}
\setlist[itemize]{topsep=3pt,itemsep=1.5pt,parsep=0pt,leftmargin=1.5em}
\setlist[enumerate]{topsep=3pt,itemsep=1.5pt,parsep=0pt,leftmargin=1.8em}
\linespread{1.16}
\tolerance=3000
\emergencystretch=3em
\hbadness=10000
\vbadness=10000
\setlength{\parskip}{3.2pt}
\setlength{\parindent}{0pt}
\renewcommand{\arraystretch}{1.25}

% ---- header / footer
\pagestyle{fancy}
\fancyhf{}
\fancyhead[L]{\small\color{exambar}考研数学二 · 线性代数知识点手册}
\fancyhead[R]{\small\color{exambar}\leftmark}
\fancyfoot[C]{\small\thepage}
\renewcommand{\headrulewidth}{0.3pt}
\renewcommand{\headrule}{\hbox to\headwidth{\color{rulesep}\leaders\hrule height \headrulewidth\hfill}}
\fancypagestyle{plain}{\fancyhf{}\fancyfoot[C]{\small\thepage}\renewcommand{\headrulewidth}{0pt}}

\begin{document}


\begin{titlepage}
\centering
\vspace*{2.6cm}
{\Huge\bfseries 考研数学二\par}
\vspace{0.5cm}
{\Huge\bfseries 线性代数知识点手册\par}
\vspace{0.9cm}
{\Large 按《考研数学二考试大纲》编排 · 五星重难点分级\par}
\vspace{1.8cm}
{\normalsize
\begin{tabular}{rl}
考查科目 & 线性代数（数学二不考概率统计） \\
分值占比 & 约 20\%（150 分中约 30 分） \\
本册章节 & 6 章，覆盖大纲线性代数全部条目 \\
分级方式 & $\bigstar$ 五级：考频 $\times$ 难度 $\times$ 分值 \\
主要素材 & 线性代数大讲义 · 线代强化/基础讲义 \\
 & 考研数学大观知识点 · 线代公式 · 数学二真题
\end{tabular}\par}
\vfill
{\small 本手册严格按《考研数学二考试大纲》顺序编排，每个考点标注五星重难点等级，供复习检索与精力分配使用。\par}
\vspace{1.0cm}
\end{titlepage}


\frontmatter
\pagestyle{fancy}
\chapter*{使用说明}
\addcontentsline{toc}{chapter}{使用说明}
\markboth{使用说明}{使用说明}
\section*{一、这本手册是什么}

这是一本\emph{严格按《考研数学二考试大纲》顺序}编排的线性代数知识点手册。它不替代教材与习题册，
解决的是三件事：

\begin{enumerate}
  \item \textbf{对照考纲查漏}：行列式、矩阵、向量、线性方程组、特征值与特征向量、二次型六章，
        大纲上出现的每一条本册都写到；考纲外的线性空间、线性变换、若尔当标准形等一律不收。
  \item \textbf{按星级分配精力}：线性代数只有约 20\% 的分值，但题型固定、套路清晰，
        把$\bigstar\bigstar\bigstar\bigstar$以上的考点练熟，是整张卷子里最稳的得分点。
  \item \textbf{考前快速回顾}：线性代数的失分绝大多数不是"不会"，而是"记错结论"（秩的不等式、
        可逆的充要条件、相似与合同的区别）。本册把这些结论集中成块，便于考前扫读。
\end{enumerate}

\section*{二、五星分级标准}

星标取\textbf{考频 $\times$ 难度 $\times$ 分值}的综合：

\begin{center}
\begin{tabular}{@{}clll@{}}
\toprule
\textbf{星标} & \textbf{名称} & \textbf{含义} & \textbf{复习要求} \\
\midrule
$\bigstar\bigstar\bigstar\bigstar\bigstar$ & 决胜考点 & 几乎年年必考，且是难题/易错点 & 闭卷默写 + 限时算对 \\
$\bigstar\bigstar\bigstar\bigstar$ & 高频重点 & 反复考，套路固定但计算量不小 & 熟练到条件反射 \\
$\bigstar\bigstar\bigstar$ & 常考点 & 隔年出现或以小题形式考 & 理解 + 会用 \\
$\bigstar\bigstar$ & 一般考点 & 低频，或只是解题的基础工具 & 会算即可 \\
$\bigstar$ & 了解性考点 & 大纲要求"了解"，考到即送分 & 记住结论 \\
\bottomrule
\end{tabular}
\end{center}

\section*{三、线性代数的三条主线}

线性代数看起来零散，实际上是三条主线贯穿六章。复习时始终问自己"这道题在考哪条线"：

\begin{enumerate}
  \item \textbf{秩的主线}：矩阵的秩 $=$ 行向量组的秩 $=$ 列向量组的秩 $=$ 非零子式的最高阶数；
        $r(AB)\le\min\{r(A),r(B)\}$，可逆矩阵相乘秩不变；$Ax=0$ 的解空间维数 $=n-r(A)$。
  \item \textbf{方程组的主线}：$Ax=0$ 只有零解 $\iff r(A)=n$ $\iff$ 列向量组线性无关 $\iff |A|\ne0$
        $\iff A$ 可逆；非齐次有解 $\iff r(A)=r(\overline A)$。这条线把矩阵、向量、方程组三章串成一处。
  \item \textbf{特征值的主线}：$|\lambda E-A|=0$ $\Rightarrow$ 特征值与迹、行列式的关系 $\Rightarrow$
        相似对角化的条件 $\Rightarrow$ 实对称矩阵正交对角化 $\Rightarrow$ 二次型的标准形与正定性。
\end{enumerate}

\section*{四、试卷结构与线代的分值}

数学二满分 150 分，考试时间 180 分钟；内容结构为\textbf{高等数学约 80\%}、\textbf{线性代数约 20\%}
（不考概率统计）。题型为：单项选择题 10 小题共 50 分、填空题 6 小题共 30 分、解答题 6 小题共 70 分。
线性代数通常占 2 个选择/填空加 1 个解答题，\textbf{分值不大但几乎不丢分}，性价比极高。

\section*{五、数学二线代不考什么}

\begin{itemize}
  \item 线性空间、向量空间的定义与判定，基变换与坐标变换
  \item 线性变换及其矩阵表示
  \item 若尔当（Jordan）标准形与 $\lambda$-矩阵
  \item 广义特征向量、最小多项式等更深的内容
  \item 概率论与数理统计（数学二完全不考）
\end{itemize}

\section*{六、怎么用：三轮法}

\begin{enumerate}
  \item \textbf{第一轮}：按行列式 $\to$ 矩阵 $\to$ 向量 $\to$ 方程组 $\to$ 特征值 $\to$ 二次型的顺序通读，
        把"结论/公式"抄成一张 A4 速查表（矩阵的秩、可逆的充要条件、相似与合同、正定的判别）。
  \item \textbf{第二轮}：集中攻$\bigstar\bigstar\bigstar\bigstar\bigstar$四类题——抽象行列式与伴随矩阵、
        向量组线性相关性的证明、含参数方程组的解的讨论、相似对角化与二次型标准化。
  \item \textbf{第三轮}：只看每章"易错清单"，重点核对：$|kA|=k^n|A|$、$|A^*|=|A|^{n-1}$、
        $A$ 与 $A^T$ 特征值相同但特征向量不同、相似必合同（实对称）而合同未必相似。
\end{enumerate}

\section*{七、使用提醒}

\begin{itemize}
  \item 本册按大纲顺序编排，章号即大纲条目号（一、行列式 $\dots$ 六、二次型）。
  \item "真题考法"只描述出题角度，不编造年份与题号。
  \item 手册中所有结论都写明了\textbf{前提}（如"$A$ 可逆时""$A$ 为实对称矩阵时"），
        请务必连前提一起记，这是线性代数最常见的失分原因。
\end{itemize}

\cleardoublepage
\tableofcontents
\mainmatter

\chapter{行列式}
\markboth{行列式}{行列式}
\subsection*{本章考纲要求}
\addcontentsline{toc}{subsection}{本章考纲要求}

\begin{quote}
\textbf{考试内容}

\begin{enumerate}
\item 行列式的概念和基本性质
\item 行列式按行（列）展开定理
\end{enumerate}
\textbf{考试要求}

\begin{enumerate}
\item 了解行列式的概念，掌握行列式的性质。
\item 会应用行列式的性质和行列式按行（列）展开定理计算行列式。
\end{enumerate}
\end{quote}
\subsection*{本章星级总表}
\addcontentsline{toc}{subsection}{本章星级总表}

\begingroup\small\setlength{\tabcolsep}{4.0pt}
\begin{longtable}{@{}L{0.4240\linewidth}L{0.1175\linewidth}L{0.1175\linewidth}L{0.1175\linewidth}L{0.1175\linewidth}@{}}
\toprule
\textbf{考点} & \textbf{星级} & \textbf{考纲要求} & \textbf{难度} & \textbf{常见题型} \\
\midrule
\endhead
行列式的概念与基本性质 & $\bigstar$$\bigstar$ & 了解/掌握 & 基础 & 选择 \\
数字型行列式的计算（化三角形、降阶） & $\bigstar$$\bigstar$$\bigstar$$\bigstar$ & 会 & 中等 & 填空、解答 \\
按行（列）展开定理与代数余子式 & $\bigstar$$\bigstar$$\bigstar$$\bigstar$ & 掌握 & 中等 & 填空、解答 \\
特殊行列式：范德蒙德、爪型、"三对角"、两行（列）成比例 & $\bigstar$$\bigstar$$\bigstar$$\bigstar$ & 会 & 中等 & 填空 \\
抽象行列式：$\vert kA\vert$、$\vert AB\vert$、$\vert A^*\vert$、$\vert A^{-1}\vert$、$\vert A^T\vert$ & $\bigstar$$\bigstar$$\bigstar$$\bigstar$$\bigstar$ & 会 & 中等 & 选择、填空 \\
代数余子式的线性组合与行列式的构造 & $\bigstar$$\bigstar$$\bigstar$$\bigstar$ & 会 & 难 & 选择、填空 \\
\bottomrule
\end{longtable}
\endgroup

\section{行列式的概念与基本性质}

\subsection{行列式的定义与记号}

\subsubsection{行列式的概念与基本性质 $\bigstar$$\bigstar$}

\textbf{定义}：由 $n^2$ 个数 $a_{ij}\ (i,j=1,2,\cdots,n)$ 排成的 $n$ 阶数表，按一定法则对应一个数，记作

\[D=\begin{vmatrix}a_{11}&a_{12}&\cdots&a_{1n}\\a_{21}&a_{22}&\cdots&a_{2n}\\\vdots&\vdots&&\vdots\\a_{n1}&a_{n2}&\cdots&a_{nn}\end{vmatrix}=\det(a_{ij}),\]

称为 $n$ 阶行列式。$a_{ij}$ 称为元素，$i$ 为行标、$j$ 为列标。

\textbf{展开式定义}：$n$ 阶行列式是取自\textbf{不同行、不同列}的 $n$ 个元素之积，带上符号后的和：

\[D=\sum_{j_1j_2\cdots j_n}(-1)^{\tau(j_1j_2\cdots j_n)}a_{1j_1}a_{2j_2}\cdots a_{nj_n},\]

其中 $j_1j_2\cdots j_n$ 是 $1,2,\cdots,n$ 的某一个排列，$\tau(j_1j_2\cdots j_n)$ 是该排列的\textbf{逆序数}（排列中"前面数大于后面数"的数对总数），$(-1)^{\tau}$ 为该项的符号。

这个定义统考只要求"了解"，但它是理解行列式性质、也是求\textbf{含大量零元素}的稀疏行列式的有力工具：只要判断某个乘积的行标列标是否各取一次，就能确定它是不是展开式中的项。

\textbf{逆序数与排列的奇偶性}：把 $1,2,\cdots,n$ 排成一列 $j_1j_2\cdots j_n$，对每一对数 $(j_p,j_q)$，若 $p<q$ 而 $j_p>j_q$，就称它构成一个\textbf{逆序}。逆序的总数叫这个排列的\textbf{逆序数}，记作 $\tau(j_1j_2\cdots j_n)$。逆序数为偶数的排列叫\textbf{偶排列}，为奇数的叫\textbf{奇排列}。

\textbf{怎么数逆序数}：从左往右，依次数出每个数\textbf{右边比它小}的数的个数，再全部相加。

\textbf{例}：排列 $3\,1\,4\,2$ 的逆序数。$3$ 右边比它小的有 $1,2$ 两个；$1$ 右边没有比它小的；$4$ 右边比它小的有 $2$ 一个；$2$ 右边没有。故 $\tau(3142)=2+0+1+0=3$，是奇排列，符号为 $(-1)^3=-1$。

\textbf{由定义能立刻读出的两件事}：$\circled{1}$ 展开式的每一项都是取自\textbf{不同行、不同列}的 $n$ 个元素之积，只要位置违反"不同行不同列"，这一项就不出现；$\circled{2}$ 每一项的符号只由\textbf{列标排列的逆序数}决定（因为行标固定按自然序 $1,2,\cdots,n$ 排列）。这两条是判断"某个乘积是不是行列式的项、带什么符号"这类题目的全部依据。

\textbf{六个基本性质}（前两条最常用）：

\begin{itemize}
\item \textbf{性质 1（线性性）}：某一行（列）的元素都是两个数之和时，行列式可拆成两个行列式之和。例如按第 $k$ 行拆分：
\end{itemize}
\[\begin{vmatrix}\alpha_1\\\vdots\\\alpha_k+\beta_k\\\vdots\\\alpha_n\end{vmatrix}=\begin{vmatrix}\alpha_1\\\vdots\\\alpha_k\\\vdots\\\alpha_n\end{vmatrix}+\begin{vmatrix}\alpha_1\\\vdots\\\beta_k\\\vdots\\\alpha_n\end{vmatrix}.\]

更一般地，某一行有公因子 $k$ 可提到行列式外：$\begin{vmatrix}\alpha_1\\\vdots\\k\alpha_k\\\vdots\\\alpha_n\end{vmatrix}=k\begin{vmatrix}\alpha_1\\\vdots\\\alpha_k\\\vdots\\\alpha_n\end{vmatrix}$。

\begin{itemize}
\item \textbf{性质 2（反对称性）}：交换两行（列），行列式变号。
\item \textbf{性质 3}：$\vert E\vert=1$（单位矩阵的行列式为 $1$）。
\item \textbf{性质 4（乘积性）}：$\vert AB\vert=\vert A\vert\cdot\vert B\vert$，$A,B$ 为\textbf{同阶方阵}。
\item \textbf{性质 5（转置不变）}：$\vert A^T\vert=\vert A\vert$。这条性质说明行列式中"行"与"列"地位完全对等。
\item \textbf{性质 6（成比例零）}：某一行（列）全为 $0$，或两行（列）元素对应相等，或两行（列）元素对应成比例，则 $D=0$。
\end{itemize}
\textbf{最常用的两个推论}：

\begin{itemize}
\item \textbf{倍加不变}：把某一行（列）的 $k$ 倍加到另一行（列），行列式的值不变。这是"化三角形"的核心工具。
\item \textbf{化三角形}：经过若干次倍加，把行列式化为上（下）三角后，$D$ 就等于主对角线上 $n$ 个元素之积。
\end{itemize}
\begin{callwarn}{易错}
性质 2、性质 6 都只对\textbf{行与行之间、列与列之间}成立。\textbf{行与列之间}不能做"交换""成比例则为零"这类操作：$a_{11}$ 与 $a_{12}$ 之间没有这种关系。另外，"某行提出 $k$"与"矩阵数乘 $k$"是两回事——从行列式提出公因子只影响\textbf{一行}，而 $\vert kA\vert=k^n\vert A\vert$ 是每一行都提一个 $k$。

\end{callwarn}

\begin{callexam}{真题考法}
这一考点的常见考法是：给出一个行列式，用性质 1、2、6 判断它等于多少倍的原行列式（如"交换两行后再加上某行的 $k$ 倍"）；或者给出 $\vert A\vert$ 与若干次初等变换后的 $\vert B\vert$，求比例关系。还有一种考法是给含参数的行列式，要求用"两行成比例则行列式为 $0$"反解参数。\textbf{陷阱}：把"交换两行"漏掉符号，或把"某行提 $k$"错记成"提 $k^n$"。

\end{callexam}

\subsection{三角、对角与副对角行列式}

这几个式子是所有行列式计算的落脚点，必须做到不假思索地写出。

\textbf{（1）主对角行列式}：

\[\begin{vmatrix}d_1&&\\&\ddots&\\&&d_n\end{vmatrix}=d_1d_2\cdots d_n .\]

\textbf{（2）上（下）三角行列式}：主对角线上方（下方）全为 $0$ 的行列式，等于主对角线元素之积。

\[\begin{vmatrix}a_{11}&a_{12}&\cdots&a_{1n}\\0&a_{22}&\cdots&a_{2n}\\\vdots&\vdots&\ddots&\vdots\\0&0&\cdots&a_{nn}\end{vmatrix}=a_{11}a_{22}\cdots a_{nn}.\]

\textbf{（3）副对角行列式}（"反对角"上非零，即只有 $a_{1n},a_{2,n-1},\cdots,a_{n1}$ 所在的那条线非零）：

\[\begin{vmatrix}&&&a_{1n}\\&&a_{2,n-1}&\\&\iddots&&\\a_{n1}&&&\end{vmatrix}=(-1)^{\frac{n(n-1)}{2}}a_{1n}a_{2,n-1}\cdots a_{n1}.\]

\textbf{（4）副对角三角行列式}：只有副对角线及其一侧非零时，同样有

\[\begin{vmatrix}&&&a_{1n}\\&&a_{2,n-1}&a_{2n}\\\iddots&\iddots&\iddots&\vdots\\a_{n1}&a_{n2}&\cdots&a_{nn}\end{vmatrix}=(-1)^{\frac{n(n-1)}{2}}a_{1n}a_{2,n-1}\cdots a_{n1}.\]

\textbf{符号怎么来的}：副对角情形可以逐行交换（第 1 行与第 $n$ 行换、第 2 行与第 $n-1$ 行换……）化成主对角，共交换 $\dfrac{n(n-1)}{2}$ 次，每次变号，故符号为 $(-1)^{\frac{n(n-1)}{2}}$。$n=2$ 时符号为 $-1$：$\begin{vmatrix}0&a\\b&0\end{vmatrix}=-ab$；$n=3$ 时符号为 $(-1)^3=-1$；$n=4$ 时符号为 $(-1)^6=+1$。

\begin{callwarn}{易错}
副对角公式的符号 $(-1)^{\frac{n(n-1)}{2}}$ 是本章第一大坑。$n=2,3$ 时是负号，$n=4,5$ 时是正号，$n=6,7$ 时又是负号（周期为 $4$）。记不清时\textbf{现场数交换次数}，不要硬背。

\end{callwarn}

\section{数字型行列式的计算}

\subsection{化三角形与降阶}

\subsubsection{数字型行列式的计算（化三角形、降阶）$\bigstar$$\bigstar$$\bigstar$$\bigstar$}

这是行列式题的\textbf{主战场}。所有具体数字型行列式，先问自己两句话：能不能化成三角形？能不能化零降阶？

\textbf{套路一：化零降阶法}（最通用）

\begin{enumerate}
\item 观察哪一行（列）的 $0$ 最多，或者哪个元素最容易作为"主元"。
\item 用倍加变换 $r_i+kr_j\to r_i$，把这一行（列）除一个元素外全部消成 $0$。
\item 按这一行（列）展开，降低一阶。
\item 重复上述步骤，直到降到二阶或三阶。
\end{enumerate}
\textbf{套路二：化三角形法}

\begin{enumerate}
\item 用第 1 行的适当倍数消去第 2 行至第 $n$ 行第 1 列的元素。
\item 用第 2 行消去下方第 2 列的元素……逐列推进。
\item 化为上三角后，$D$ 等于主对角线之积。
\end{enumerate}
\textbf{例}：计算 $D=\begin{vmatrix}3&1&2\\0&4&3\\1&2&5\end{vmatrix}$。先把第 1 行与第 3 行交换（变号），再用第 1 行消去第 3 行：

\[D=-\begin{vmatrix}1&2&5\\0&4&3\\3&1&2\end{vmatrix}\xrightarrow{r_3-3r_1}-\begin{vmatrix}1&2&5\\0&4&3\\0&-5&-13\end{vmatrix}.\]

按第 1 列展开，得 $D=-\begin{vmatrix}4&3\\-5&-13\end{vmatrix}=-[4\times(-13)-3\times(-5)]=-(-52+15)=37$。

\textbf{常用小技巧}（能省大量计算）：

\begin{itemize}
\item \textbf{各行（列）元素之和相等}时，把各列全部加到第 1 列，则第 1 列元素全为同一个数，提出公因子后第 1 列全为 $1$，再用各行减去第 1 行即可降阶。
\item 数字较大或含公因子时，\textbf{先提公因子}再算。
\item 目标行列式是"爪型"结构时，不要用通用的化三角形，见下一节。
\item \textbf{阶梯形行列式}（形状像 $\begin{vmatrix}1&1&1&3&1\\0&0&2&0&0\\0&0&0&0&4\\0&5&0&0&0\\6&0&0&0&0\end{vmatrix}$ 这类每行只有一个非零元的）直接用定义数"取哪几个元素"，一次定位符号，不要展开。
\end{itemize}
\textbf{例（稀疏行列式用定义）}：计算 $D=\begin{vmatrix}0&1&0&0\\0&0&2&0\\0&0&0&3\\4&0&0&0\end{vmatrix}$。

展开式中只有一个乘积可能非零，即取 $a_{12}a_{23}a_{34}a_{41}=1\cdot2\cdot3\cdot4=24$，其列标排列为 $2,3,4,1$，$\tau(2341)=3$（$2$ 右边比它小的有 $1$；$3$ 右边有 $1$；$4$ 右边有 $1$），故符号为 $(-1)^3=-1$，$D=-24$。

\textbf{验证}：把第 1 行依次与下面各行交换（$3$ 次对换）可化为副对角形状，$D=(-1)^3\cdot(1\cdot2\cdot3\cdot4)=-24$，与定义法一致。

\textbf{例（行和相等型）}：计算 $D=\begin{vmatrix}3&1&1\\1&3&1\\1&1&3\end{vmatrix}$。

每行元素之和都是 $5$。把第 2、3 列都加到第 1 列，第 1 列变为全 $5$：

\[D=\begin{vmatrix}5&1&1\\5&3&1\\5&1&3\end{vmatrix}=5\begin{vmatrix}1&1&1\\1&3&1\\1&1&3\end{vmatrix}.\]

再用第 2、3 行分别减去第 1 行（倍加不变值）：

\[D=5\begin{vmatrix}1&1&1\\0&2&0\\0&0&2\end{vmatrix}=5\times1\times2\times2=20 .\]

\textbf{验证（另一条路）}：该矩阵可写成 $A=2E+J$（$J$ 为全 $1$ 矩阵，$J$ 的秩为 $1$，特征值为 $3,0,0$）。故 $A$ 的特征值为 $2+3=5$ 与 $2,2$，$\vert A\vert=5\times2\times2=20$，与上面一致。

\textbf{例（"对角线为 $x$、其余为 $a$"型）}：计算 $D_n=\begin{vmatrix}x&a&\cdots&a\\a&x&\cdots&a\\\vdots&\vdots&&\vdots\\a&a&\cdots&x\end{vmatrix}$。

每行元素之和都是 $x+(n-1)a$。把各列都加到第 1 列并提取公因子：

\[D_n=\left[x+(n-1)a\right]\begin{vmatrix}1&a&\cdots&a\\1&x&\cdots&a\\\vdots&\vdots&&\vdots\\1&a&\cdots&x\end{vmatrix}.\]

再用第 $i$ 行（$i\ge2$）减去第 1 行，第 1 列下方化为 $0$，得到一个上三角行列式，其主对角元为 $1,x-a,\cdots,x-a$（共 $n-1$ 个 $x-a$）。故

\[D_n=\left[x+(n-1)a\right](x-a)^{\,n-1}.\]

\textbf{验证}：取 $n=3$、$x=3$、$a=1$，公式给出 $(3+2)(3-1)^2=5\times4=20$，与上例结果一致。

\textbf{例（局部拆项）}：$\begin{vmatrix}1+2&3\\4+5&6\end{vmatrix}=\begin{vmatrix}1&3\\4&6\end{vmatrix}+\begin{vmatrix}2&3\\5&6\end{vmatrix}=(6-12)+(12-15)=-6-3=-9$，而直接算 $\begin{vmatrix}3&3\\9&6\end{vmatrix}=18-27=-9$，一致。\textbf{注意性质 1 只允许"一次拆一行（列）"}，若第 1 行与第 2 行同时是两数之和，需按行拆两次，共得 $4$ 个行列式。

\begin{callwarn}{易错}
化三角形时用倍加变换\textbf{不改变行列式的值}；但一旦用了"交换两行"就必须补符号，用了"某行乘 $k$"就要在最后除回来或在外面记一个 $1/k$。很多考生算到一半把符号丢了，答案正好差一个负号。

\end{callwarn}

\begin{callexam}{真题考法}
数字型行列式的直接考法多出现在填空题里，阶数通常为 $3\sim 5$ 阶；四阶常见于选择题。命题人偏爱两类结构：一类是\textbf{某行（列）已含较多 $0$}（直接展开），另一类是\textbf{各行（列）元素之和相等}（先加列再降阶）。陷阱在于：给的数字故意让主元为 $0$，需要先交换两行；或者故意设置成成比例的某两行，答案其实是 $0$，但考生埋头硬算。

\end{callexam}

\subsection{特殊结构的行列式}

\subsubsection{特殊行列式：范德蒙德、爪型、"三对角"、两行（列）成比例 $\bigstar$$\bigstar$$\bigstar$$\bigstar$}

\textbf{（1）范德蒙德行列式}

\[D_n=\begin{vmatrix}1&1&\cdots&1\\x_1&x_2&\cdots&x_n\\x_1^2&x_2^2&\cdots&x_n^2\\\vdots&\vdots&&\vdots\\x_1^{n-1}&x_2^{n-1}&\cdots&x_n^{n-1}\end{vmatrix}=\prod_{1\le j<i\le n}(x_i-x_j).\]

\textbf{适用条件与要点}：第 1 行全为 $1$（即各列的 $0$ 次幂），第 2 行是一次幂，逐行升幂直到第 $n$ 行的 $n-1$ 次幂；结果是所有"下标大的减下标小的"之差连乘，共 $\dfrac{n(n-1)}{2}$ 个因子。\textbf{推论}：$D_n\ne0$ 当且仅当 $x_1,x_2,\cdots,x_n$ 互不相同。

\textbf{例}：$\begin{vmatrix}1&1&1\\2&3&4\\4&9&16\end{vmatrix}=(3-2)(4-2)(4-3)=1\times2\times1=2$。

\textbf{缺行缺列的还原}：考研常给"不完整的范德蒙德"，手法有三种——$\circled{1}$ \textbf{补行}：添上一行 $1$ 和一行适当幂次，凑成全范德蒙德；$\circled{2}$ \textbf{提公因子}：某列都有公因子 $x_i^{k}$ 时，先把 $x_i^k$ 提到行列式外；$\circled{3}$ \textbf{转置/换列}：转置后仍是范德蒙德结构，换列顺序只影响符号。

\textbf{（2）爪型（箭形）行列式}

形如主对角线上元素全非零、而第 1 行与第 1 列（或最后一行与最后一列）另有非零元素的"十"字形状：

\[D=\begin{vmatrix}a_{11}&a_{12}&a_{13}&\cdots&a_{1n}\\a_{21}&a_{22}&0&\cdots&0\\a_{31}&0&a_{33}&\cdots&0\\\vdots&\vdots&\vdots&\ddots&\vdots\\a_{n1}&0&0&\cdots&a_{nn}\end{vmatrix}.\]

\textbf{套路}：用第 $i$ 行（$i\ge2$）乘 $-\dfrac{a_{1i}}{a_{ii}}$ 加到第 1 行，把第 1 行从第 2 列起的元素全部消成 $0$，于是化为上三角，$D=\left(a_{11}-\sum_{i=2}^{n}\dfrac{a_{1i}a_{i1}}{a_{ii}}\right)a_{22}a_{33}\cdots a_{nn}$。\textbf{口诀：爪形一律化上三角，用"两杠"上的元素把"星"消掉。}

\textbf{（3）行（列）和相等型与加边法}

\textbf{行和相等}（每行元素之和都是同一个数）时，把各列全都加到第 1 列，第 1 列元素就都等于这个和，提出公因子后第 1 列全为 $1$，再用各行减去第 1 行，即可降阶。

\textbf{加边法}：给 $n$ 阶行列式添加一行一列并保证值不变（新添的行按 $1,0,\cdots,0$ 构造，新添的列第 1 个元素为 $1$、其余为 $0$），把"除主对角线外各行（列）同一位置元素相同"的行列式升成 $n+1$ 阶，再用爪型手法化简。\textbf{关键}：加边必须让新行列式的值与原行列式相等——按新添的第 1 列展开即知这一点。

\textbf{例（加边法）}：计算 $D_n=\begin{vmatrix}a_1+b_1&a_2&\cdots&a_n\\a_1&a_2+b_2&\cdots&a_n\\\vdots&\vdots&&\vdots\\a_1&a_2&\cdots&a_n+b_n\end{vmatrix}$（$b_i\ne0$）。

这是"除主对角线外，每列的其他元素都分别相同"的行列同数型行列式，用加边法。构造 $n+1$ 阶行列式

\[D_{n+1}=\begin{vmatrix}1&a_1&a_2&\cdots&a_n\\0&a_1+b_1&a_2&\cdots&a_n\\0&a_1&a_2+b_2&\cdots&a_n\\\vdots&\vdots&\vdots&&\vdots\\0&a_1&a_2&\cdots&a_n+b_n\end{vmatrix}.\]

按第 1 列展开知 $D_{n+1}=D_n$，值不变，故可改算这个新行列式。用\textbf{第 1 行乘 $-1$ 加到其余各行}，第 2 行起依次变为 $(0,\,b_1,\,0,\cdots,0)$、$(0,\,0,\,b_2,\cdots,0)$、……、$(0,0,\cdots,b_n)$——这就是爪型。再用第 $i$ 行（$i\ge2$）乘 $-\dfrac{a_{i-1}}{b_{i-1}}$ 加到第 1 行，把第 1 行第 2 列起的元素全部消成 $0$，于是

\[D_n=\left(1+\sum_{i=1}^{n}\frac{a_i}{b_i}\right)b_1b_2\cdots b_n .\]

\textbf{验证}：$n=1$ 时公式给出 $\left(1+\dfrac{a_1}{b_1}\right)b_1=a_1+b_1=D_1$；$n=2$ 时公式给出 $\left(1+\dfrac{a_1}{b_1}+\dfrac{a_2}{b_2}\right)b_1b_2=a_1b_2+a_2b_1+b_1b_2$，而直接计算 $\begin{vmatrix}a_1+b_1&a_2\\a_1&a_2+b_2\end{vmatrix}=(a_1+b_1)(a_2+b_2)-a_1a_2=a_1b_2+a_2b_1+b_1b_2$，两者一致。

\textbf{（4）三对角行列式与递推法}

形如

\[D_n=\begin{vmatrix}a&b&&&\\c&a&b&&\\&c&a&\ddots&\\&&\ddots&\ddots&b\\&&&c&a\end{vmatrix}\]

的三对角（或爪型）行列式，\textbf{零元素多，不要刻意消零}。做法是\textbf{按第 1 行（列）展开}，得到 $D_n=aD_{n-1}-bc\,D_{n-2}$ 这样的\textbf{递推关系}，再求出通项。常见的二阶线性递推可用特征方程法求通项，或用数学归纳法验证结论。

\textbf{常用特例}：$\begin{vmatrix}2a&1&&\\a^2&2a&\ddots&\\&\ddots&\ddots&1\\&&a^2&2a\end{vmatrix}=(n+1)a^n$（$n$ 阶）。

\textbf{递推法的完整流程（四步）}

\begin{enumerate}
\item \textbf{按第 1 行（列）展开}，把 $D_n$ 表示成 $D_{n-1}$ 与 $D_{n-2}$ 的组合（通常形如 $D_n=aD_{n-1}-\beta D_{n-2}$）。
\item \textbf{独立算出 $D_1$、$D_2$}（不能由递推式代出，必须手算）。
\item \textbf{解二阶线性递推}：由特征方程 $t^2-at+\beta=0$ 得两根 $t_1,t_2$，通项为 $D_n=At_1^{n-1}+Bt_2^{n-1}$（$t_1\ne t_2$ 时；$t_1=t_2$ 时为 $(A+Bn)t_1^{n-1}$），用 $D_1,D_2$ 定 $A,B$。
\item \textbf{用归纳法或代入 $n=1,2,3$ 验证}。
\end{enumerate}
\textbf{例}：计算 $n$ 阶行列式 $D_n=\begin{vmatrix}2&1&&&\\1&2&1&&\\&1&2&\ddots&\\&&\ddots&\ddots&1\\&&&1&2\end{vmatrix}$。

按第 1 行展开：$D_n=2D_{n-1}-\begin{vmatrix}1&1&\\1&2&1\\&\ddots&\ddots\end{vmatrix}$，其中第二个行列式再按第 1 列展开得 $D_{n-2}$，故

\[D_n=2D_{n-1}-D_{n-2}.\]

直接算起始项：$D_1=2$，$D_2=\begin{vmatrix}2&1\\1&2\end{vmatrix}=3$。特征方程 $t^2-2t+1=0$ 有二重根 $t=1$，故通项形如 $D_n=A+Bn$。代入 $n=1$ 得 $A+B=2$；代入 $n=2$ 得 $A+2B=3$。解得 $A=1,B=1$，于是

\[D_n=n+1 .\]

\textbf{验证}：$n=3$ 时 $D_3=2\times3-2=4=n+1$，成立。

\textbf{例（"$2a$ 与 $a^2$"型）}：$D_n=\begin{vmatrix}2a&1&&\\a^2&2a&\ddots&\\&\ddots&\ddots&1\\&&a^2&2a\end{vmatrix}$。同样按第 1 行展开得 $D_n=2aD_{n-1}-a^2D_{n-2}$，$D_1=2a$，$D_2=\begin{vmatrix}2a&1\\a^2&2a\end{vmatrix}=4a^2-a^2=3a^2$。特征方程 $t^2-2at+a^2=0$ 有二重根 $t=a$，故 $D_n=(A+Bn)a^{\,n-1}$。代入 $D_1=2a$ 得 $A=2$；代入 $D_2=3a^2$ 得 $(2+B)a=3a^2$，即 $B=1$（$a\ne0$ 时）。故

\[D_n=(n+1)a^n ,\]

与前面给出的常用特例一致。

\textbf{例（"$a$ 与 $b$"型）}：$D_n=\begin{vmatrix}a&b&&\\b&a&\ddots&\\&\ddots&\ddots&b\\&&b&a\end{vmatrix}$ 满足 $D_n=aD_{n-1}-b^2D_{n-2}$，$D_1=a$，$D_2=a^2-b^2$。特征方程 $t^2-at+b^2=0$ 的两根记为 $t_1,t_2$，则 $t_1+t_2=a$，$t_1t_2=b^2$，且

\[D_n=\frac{t_1^{\,n+1}-t_2^{\,n+1}}{t_1-t_2}\qquad(t_1\ne t_2\ \text{时}).\]

\textbf{验证}：有理化分式，$t_1^{n+1}-t_2^{n+1}=(t_1-t_2)\left(t_1^{n}+t_1^{n-1}t_2+\cdots+t_2^{n}\right)$，故 $D_n=t_1^n+t_1^{n-1}t_2+\cdots+t_2^n$（共 $n+1$ 项）。代入 $n=1$ 得 $D_1=t_1+t_2=a$，与直接计算一致；代入 $n=2$ 得 $D_2=t_1^2+t_1t_2+t_2^2=(t_1+t_2)^2-t_1t_2=a^2-b^2$，也与直接计算一致。公式正确。\textbf{这个"$a$ 与 $b$"型的三对角行列式是考研最常考的递推母题}，其特例 $a=2,b=1$ 就是上面的 $D_n=n+1$。

\textbf{（5）两行（列）成比例}

若行列式中有两行（列）元素对应成比例（包括相等、某行为另一行的 $k$ 倍、含全零行），则 $D=0$。这条性质常与含参数的题目结合：给定行列式，问参数取何值时行列式为 $0$。

\textbf{（6）行列式与高数"贴贴"型}

考研常把行列式与多项式零点结合：给出 $f(x)$ 是一个含 $x$ 的三阶或四阶行列式，问 $f(x)=0$ 有几个根、$x^3$ 的系数是多少。\textbf{核心观察}：$f(x)$ 是关于 $x$ 的多项式，其次数不超过"含 $x$ 的最大行数"；若能找出若干个 $x$ 的取值使\textbf{两行（列）相同或成比例}，则这些 $x$ 就是 $f(x)$ 的根。这类题的通用套路是：

\begin{enumerate}
\item 数清楚有几行含 $x$、每行含 $x$ 的最高次数，定出 $\deg f$ 的上界；
\item 找出使某两行（列）相同或成比例的 $x$ 值，即得根；
\item 用根的数量与次数上界比较，判断是否还有别的根（必要时用余下的一个点定系数）。
\end{enumerate}
\textbf{例}：$f(x)=\begin{vmatrix}1&1&1\\x&2&3\\x^2&4&9\end{vmatrix}$。这是范德蒙德结构（三列分别是 $x_1=x$、$x_2=2$、$x_3=3$ 的幂），故

\[f(x)=(2-x)(3-x)(3-2)=(2-x)(3-x),\]

零点为 $x=2$ 与 $x=3$，共 $2$ 个根。\textbf{验证}：$f(x)$ 按第 1 列看是关于 $x$ 的二次多项式，得到两个根正好用满次数，确实只有这两个根。

\textbf{（7）分块三角与拉普拉斯展开}

设 $A$ 为 $m$ 阶方阵、$B$ 为 $n$ 阶方阵，$C$ 为任意 $n\times m$ 矩阵，则

\[\begin{vmatrix}A&C\\O&B\end{vmatrix}=\begin{vmatrix}A&O\\C&B\end{vmatrix}=\vert A\vert\cdot\vert B\vert,\]

\[\begin{vmatrix}O&A\\B&O\end{vmatrix}=(-1)^{mn}\vert A\vert\cdot\vert B\vert .\]

\textbf{第二个式子的符号}：把 $O$ 右上角的 $A$ 所在的行块与 $B$ 所在的行块交换，共需 $m\cdot n$ 次行（列）对换，故符号为 $(-1)^{mn}$。\textbf{$m,n$ 分别为 $A,B$ 的阶数}，不要代错。

\textbf{拉普拉斯展开定理（按 $k$ 行展开）}：在 $n$ 阶行列式中任取 $k$ 行，则行列式等于由这 $k$ 行形成的\textbf{所有 $k$ 阶子式}与其对应\textbf{代数余子式}乘积之和。统考不要求熟练使用一般形式，但它正是上一条的来源：取前 $m$ 行时，唯一非零的 $m$ 阶子式是 $\vert A\vert$，其代数余子式带上 $(1+2+\cdots+m)+(1+2+\cdots+m)$ 的符号后恰好是 $\vert B\vert$，故直接得到 $\begin{vmatrix}A&C\\O&B\end{vmatrix}=\vert A\vert\vert B\vert$；而 $\begin{vmatrix}O&A\\B&O\end{vmatrix}$ 的非零 $m$ 阶子式位于第 $m+1$ 至 $m+n$ 列，行标和与列标和之差为 $mn$，故补上 $(-1)^{mn}$。

\textbf{例}：计算四阶行列式 $D=\begin{vmatrix}1&2&0&0\\3&4&0&0\\0&0&5&6\\0&0&7&8\end{vmatrix}$。

按前 $2$ 行取唯一非零的 $2$ 阶子式，得 $D=\begin{vmatrix}1&2\\3&4\end{vmatrix}\cdot\begin{vmatrix}5&6\\7&8\end{vmatrix}=(4-6)(40-42)=(-2)(-2)=4$。\textbf{验证}：对第 1 列做倍加消元，$c_2-2c_1$（把第 1 列的 $2$ 倍从第 2 列减去，列倍加不改变行列式的值），得 $D=\begin{vmatrix}1&0&0&0\\3&-2&0&0\\0&0&5&6\\0&0&7&8\end{vmatrix}$；再对右下角用 $c_4-\dfrac{6}{5}c_3$ 消元，得 $D=(1)\cdot(-2)\cdot\begin{vmatrix}5&0\\7&8-\frac{42}{5}\end{vmatrix}=(-2)\cdot5\cdot\left(-\dfrac25\right)=4$，与拉普拉斯结果一致。

\begin{callwarn}{易错}
$\circled{1}$ 范德蒙德公式中因子是 $x_i-x_j\ (i>j)$，不是 $x_j-x_i$，也不是 $x_i+x_j$；先写 $x_1,x_2,\cdots,x_n$ 的下标顺序，再确定符号。 $\circled{2}$ 副对角分块公式 $(-1)^{mn}$ 中，$m$ 是\textbf{左上角那个零块的阶数}（即 $A$ 的阶数），$n$ 是 $B$ 的阶数。四阶 $\begin{vmatrix}O&A\\B&O\end{vmatrix}$（$A,B$ 均二阶）符号为 $(-1)^{4}=+1$。 $\circled{3}$ 爪型消元要求"杠"上的元素\textbf{非零}（即 $a_{ii}\ne0$）。若某个 $a_{ii}=0$，需先交换两行（列）把它换到非零位置，并注意补符号。 $\circled{4}$ 递推法必须先算 $D_1,D_2$ 验证递推式的起始条件，否则通项会差一个常数。

\end{callwarn}

\begin{callexam}{真题考法}
这一考点的选择题、填空题最常出现三种形态：$\circled{1}$ 给出四阶的"分块反对角"型，考 $(-1)^{mn}$ 的符号（四个选项往往只差符号）；$\circled{2}$ 给出形如 $x_1,x_2,\cdots,x_n$ 的幂次表，识别为范德蒙德后直接代公式；$\circled{3}$ 给出含参数 $a,b$ 的"$x$ 在对角线、$a$ 在其余位置"的行列式，用行和相等的方法化为三角形。陷阱：拿"$\begin{vmatrix}a&b\\c&d\end{vmatrix}=ad-bc$"的肌肉记忆去套副对角，把符号漏掉。另外，考研常把行列式与高数结合（如 $f(x)$ 是一个行列式，问 $f(x)=0$ 的根的个数），此时要先观察结构，找出能让某两行（列）成比例或成零的 $x$ 值，再定根的个数。

\end{callexam}

\subsection{数字型行列式的计算流程总表}

拿到一个具体行列式，按下面的顺序判断，命中哪一条就做哪一条，\textbf{不要上来就硬展开}：

\begingroup\small\setlength{\tabcolsep}{4.0pt}
\begin{longtable}{@{}L{0.5009\linewidth}L{0.4448\linewidth}@{}}
\toprule
\textbf{结构特征} & \textbf{首选方法} \\
\midrule
\endhead
已有三角或对角形状、每行只有一个非零元 & 直接写主对角之积（或用定义数排列符号） \\
某行（列）含大量 $0$ & 按该行（列）展开，降阶 \\
各行（列）元素之和相等 & 各列加到第 1 列，提公因子后化三角 \\
主对角线非零、第 1 行与第 1 列另有非零元（爪型） & 用"两杠"消"星"，化上三角 \\
除主对角线外各行（列）同位置元素相同（行列同数型） & 加边法 $\to$ 爪型 $\to$ 三角 \\
三对角或爪型且 $n$ 较高 & 按行展开建立递推式，求通项 \\
幂次递增的 $x_1,\cdots,x_n$ 表 & 套范德蒙德公式 \\
分块三角或分块副对角 & 用 $\vert A\vert\vert B\vert$ 与 $(-1)^{mn}$ \\
含 $x$、问零点个数 & 找使两行（列）相同的 $x$ 值，定出根 \\
含字母参数且问行列式为 $0$ 的条件 & 展开成多项式，解方程 \\
\bottomrule
\end{longtable}
\endgroup

\textbf{一个统一的检查习惯}：算完后，用\textbf{另一种展开方式}（换一行或换一列）再算一遍关键步骤，或者代入一组简单数字检验。本章所有公式的正确性都能这样自查。

\subsection{用行列式判定方程组的解}

行列式是线性方程组的"理论探针"，本章只需掌握与行列式直接相关的两条结论，具体解法在第四章。

\textbf{克拉默法则}：设 $n$ 元线性方程组 $Ax=b$ 的系数行列式 $\vert A\vert\ne0$，则方程组有唯一解

\[x_1=\frac{\vert A_1\vert}{\vert A\vert},\quad x_2=\frac{\vert A_2\vert}{\vert A\vert},\quad\cdots,\quad x_n=\frac{\vert A_n\vert}{\vert A\vert},\]

其中 $A_j$ 是把 $A$ 的第 $j$ 列换成常数项列 $b$ 所得的矩阵。它的理论意义大于计算意义：\textbf{$\vert A\vert\ne0\iff$ 方程组有唯一解}，这是"可逆"与"唯一解"的桥梁（详见第二章、第四章）。

\textbf{齐次方程组的判据}：$n$ 元齐次方程组 $Ax=0$ \textbf{有非零解} $\iff\vert A\vert=0$；\textbf{只有零解} $\iff\vert A\vert\ne0$。这条结论在"求参数使方程组有非零解"的题目中就是解 $\vert A\vert=0$。

\section{按行（列）展开与代数余子式}

\subsection{展开定理}

\subsubsection{按行（列）展开定理与代数余子式 $\bigstar$$\bigstar$$\bigstar$$\bigstar$}

\textbf{余子式}：在 $n$ 阶行列式 $D$ 中，划去元素 $a_{ij}$ 所在的第 $i$ 行和第 $j$ 列，余下的 $(n-1)^2$ 个元素按原顺序构成的 $n-1$ 阶行列式，称为 $a_{ij}$ 的\textbf{余子式}，记作 $M_{ij}$。

\textbf{代数余子式}：$A_{ij}=(-1)^{i+j}M_{ij}$。

\textbf{展开定理}：$n$ 阶行列式等于它的\textbf{任意一行（列）}的各元素与其对应的代数余子式乘积之和，即对任意 $i,j$

\[D=a_{i1}A_{i1}+a_{i2}A_{i2}+\cdots+a_{in}A_{in}=\sum_{k=1}^{n}a_{ik}A_{ik}\quad(\text{按第 }i\text{ 行展开}),\]

\[D=a_{1j}A_{1j}+a_{2j}A_{2j}+\cdots+a_{nj}A_{nj}=\sum_{k=1}^{n}a_{kj}A_{kj}\quad(\text{按第 }j\text{ 列展开}).\]

\textbf{互交为零}（"一行的元素乘另一行的代数余子式，和为零"）：

\[\sum_{k=1}^{n}a_{ik}A_{jk}=0\quad(i\ne j),\qquad \sum_{k=1}^{n}a_{ki}A_{kj}=0\quad(i\ne j).\]

\textbf{记忆口诀}：\textbf{自己行乘自己的代数余子式（自交）得行列式；自己行乘别人的代数余子式（互交）得 $0$。}

\textbf{套路}：看到"求 $\sum a_{ij}A_{ij}$ 这类和"，不要逐个算余子式，而是\textbf{把待求行（列）的元素换成对应位置的系数，构造一个新的行列式}，用新行列式（往往结构更简单）算出结果。这是本章最高效的一招。

\textbf{例}：求 $A=\begin{pmatrix}1&2&3\\4&5&6\\7&8&10\end{pmatrix}$ 的 $A_{11}+2A_{12}+3A_{13}$。

这个和式不是"第 1 行元素乘第 1 行代数余子式"的形式（系数是 $1,2,3$，而元素本身也是 $1,2,3$，恰好相同），所以 $A_{11}+2A_{12}+3A_{13}=1\cdot A_{11}+2\cdot A_{12}+3\cdot A_{13}=\vert A\vert$。而按第 1 行展开可直接算：

\[\vert A\vert=\begin{vmatrix}1&2&3\\4&5&6\\7&8&10\end{vmatrix}=1\begin{vmatrix}5&6\\8&10\end{vmatrix}-2\begin{vmatrix}4&6\\7&10\end{vmatrix}+3\begin{vmatrix}4&5\\7&8\end{vmatrix}=1\cdot2-2\cdot(-2)+3\cdot(-3)=-3 .\]

若把系数改成 $1,2,3$ 以外的组合（例如求 $A_{21}+A_{22}+A_{23}$），就要构造"第 2 行全换成 $1$"的新行列式：

\[A_{21}+A_{22}+A_{23}=\begin{vmatrix}1&2&3\\1&1&1\\7&8&10\end{vmatrix}=1(10-8)-2(10-7)+3(8-7)=2-6+3=-1 .\]

\textbf{验证自洽}：按第 2 行展开原矩阵，$D_2=-4\begin{vmatrix}2&3\\8&10\end{vmatrix}+5\begin{vmatrix}1&3\\7&10\end{vmatrix}-6\begin{vmatrix}1&2\\7&8\end{vmatrix}=-4(-4)+5(-11)-6(-6)=16-55+36=-3$，与上面算出的 $\vert A\vert=-3$ 一致。

\textbf{符号的棋盘格分布}：$A_{ij}$ 前的符号 $(-1)^{i+j}$ 按位置排列成

\[\begin{pmatrix}+&-&+&\cdots\\-&+&-&\cdots\\+&-&+&\cdots\\\vdots&\vdots&\vdots&\end{pmatrix}.\]

\textbf{常用的是四种求和构造}（设 $A$ 为 $n$ 阶方阵）：

\[\sum_{k=1}^{n}A_{ik}=\begin{vmatrix}a_{11}&\cdots&a_{1n}\\\vdots&&\vdots\\1&\cdots&1\\\vdots&&\vdots\\a_{n1}&\cdots&a_{nn}\end{vmatrix}\ (\text{第 }i\text{ 行全换成 }1),\qquad \sum_{k=1}^{n}A_{kj}\ (\text{第 }j\text{ 列全换成 }1);\]

\[\sum_{i=1}^{n}A_{ij}=\sum_{k=1}^{n}A_{kl}\ \text{型的混合和}\ \Longrightarrow\ \text{把相应的行（列）换成对应的系数再算}.\]

\textbf{例（求全部代数余子式之和）}：设 $A=\begin{pmatrix}1&2&3\\4&5&6\\7&8&10\end{pmatrix}$，求 $\displaystyle\sum_{i=1}^{3}\sum_{j=1}^{3}A_{ij}$。

把三行都换成 $1$，得到的新行列式每行相同，故为 $0$，即 $\sum_i\sum_j A_{ij}=0$。\textbf{这也可以由 $A^*\mathbf{1}$ 的构造直接看出}：三行都是 $1$ 的行列式有两行相同。

\textbf{例（$a_{ij}+A_{ij}=0$ 型）}：设 $A=(a_{ij})_{3\times3}$ 是非零矩阵，且对一切 $i,j$ 有 $a_{ij}+A_{ij}=0$，求 $\vert A\vert$。

条件 $A_{ij}=-a_{ij}$ 意味着 $A^*=-A^T$（因为 $A^*$ 的第 $i$ 行第 $j$ 列是 $A_{ji}$，而 $A^T$ 的第 $i$ 行第 $j$ 列是 $a_{ji}$）。两边取行列式：

\[\vert A^*\vert=\vert -A^T\vert=(-1)^3\vert A^T\vert=-\vert A\vert .\]

而 $\vert A^*\vert=\vert A\vert^{\,3-1}=\vert A\vert^2$，故 $\vert A\vert^2=-\vert A\vert$，即 $\vert A\vert(\vert A\vert+1)=0$，得 $\vert A\vert=0$ 或 $\vert A\vert=-1$。

\textbf{再排除 $\vert A\vert=0$}：若 $\vert A\vert=0$，则由 $AA^*=\vert A\vert E=O$ 与 $A^*=-A^T$ 得 $-AA^T=O$，即 $AA^T=O$。$AA^T$ 的主对角元为 $\sum_j a_{ij}^2$（实矩阵），全为 $0$ 迫使 $A=O$，与"$A$ 非零"矛盾。故 $\vert A\vert=-1$。

\textbf{验证}：$n=3$ 时 $\vert -A^T\vert=(-1)^3\vert A^T\vert=-\vert A\vert$，符号确为负；答案 $\vert A\vert=-1$ 满足 $\vert A\vert^2=1=-\vert A\vert$，自洽。

\textbf{例（含参数的三角化）}：计算 $D=\begin{vmatrix}1&1&1\\1&2&a\\1&a&4\end{vmatrix}$，并讨论 $a$ 取何值时 $D=0$。

按第 1 行展开：

\[D=1\cdot\begin{vmatrix}2&a\\a&4\end{vmatrix}-1\cdot\begin{vmatrix}1&a\\1&4\end{vmatrix}+1\cdot\begin{vmatrix}1&2\\1&a\end{vmatrix}=(8-a^2)-(4-a)+(a-2)=2-a^2+2a .\]

\textbf{验证}：代入 $a=0$ 得 $D=2$；换用第 1 列展开：$D=1(8-a^2)-1(4-a)+1(a-2)=2-a^2+2a$，与上式相同，展开无误。

令 $D=0$ 即 $a^2-2a-2=0$，解得 $a=1\pm\sqrt3$。故当且仅当 $a=1+\sqrt3$ 或 $a=1-\sqrt3$ 时该行列式为 $0$。

\textbf{例（两种方法互验）}：计算 $D=\begin{vmatrix}2&1&3\\4&2&6\\1&5&0\end{vmatrix}$。

\textbf{方法一（看结构）}：第 2 行是第 1 行的 $2$ 倍（$4=2\times2$，$2=2\times1$，$6=2\times3$），两行成比例，故 $D=0$。

\textbf{方法二（硬算验证）}：按第 1 行展开，$D=2(0-30)-1(0-6)+3(20-2)=-60+6+54=0$，两法一致。\textbf{结论：动手前先扫一眼有没有成比例的行（列），能省一整个计算过程。}

\begin{callkey}{必背}
构造法的全部依据只有两条：$\sum_{k}a_{ik}A_{ik}=D$ 与 $\sum_{k}a_{ik}A_{jk}=0\ (i\ne j)$。因此：

\begin{itemize}
\item 求"某一行（列）的代数余子式之和 $\sum_j A_{ij}$"$\to$ 把第 $i$ 行全换成 $1$，求新行列式；
\item 求"带系数的组合 $\sum_j k_jA_{ij}$"$\to$ 把第 $i$ 行换成 $k_1,k_2,\cdots,k_n$，求新行列式；
\item 求"所有元素的代数余子式之和"$\to$ 把所有行都换成 $1$ 后逐行求和。
\end{itemize}
\end{callkey}

\textbf{"选零多的行（列）展开"的定量理由}

按第 $i$ 行展开时 $D=a_{i1}A_{i1}+a_{i2}A_{i2}+\cdots+a_{in}A_{in}$。若某个 $a_{ij}=0$，这一项就\textbf{不必去算 $A_{ij}$}，而算一个 $n-1$ 阶行列式正是整道题的主要工作量。因此有三条操作原则：

\begin{enumerate}
\item \textbf{优先选 $0$ 最多的行（列）展开}。若某行只有一个非零元，只需算\textbf{一个} $n-1$ 阶子式，工作量最小。
\item \textbf{先做倍加制造 $0$，再展开}。用几次倍加换掉"少算 $n-1$ 个余子式"，几乎总是划算的。
\item \textbf{不要对每一行都展开}。逐行展开只适合阶数很低（二阶、三阶）或结构特别规则（范德蒙德型、递推型）的场合；四阶以上的"杂乱"数字行列式，一定先化零再展开。
\end{enumerate}
\textbf{一个工作量估计}：$n$ 阶行列式按定义硬展开有 $n!$ 项，四阶 $24$ 项、五阶 $120$ 项，\textbf{绝不能硬展开}；而"化零降阶"每降一阶只需算一个降阶行列式，四阶问题两次展开即可解决。这就是化零降阶法成为主战场的原因。

\textbf{本考点的三则常见失误}：$\circled{1}$ 把 $a_{ij}$ 与 $A_{ij}$ 的下标写反（$A_{ij}$ 的 $i$ 是被划掉的行号，$j$ 是被划掉的列号）；$\circled{2}$ 展开时漏项或把 $+,-,+,\cdots$ 的符号交替写错；$\circled{3}$ 只熟"按行展开"，遇到明显"按列更好算"的题仍按行硬展开。

\begin{callwarn}{易错}
$\circled{1}$ 代数余子式的符号 $(-1)^{i+j}$ 绝不能漏，$n$ 阶行列式的符号呈棋盘格分布：$(1,1)$ 正、$(1,2)$ 负、$(2,1)$ 负、$(2,2)$ 正。 $\circled{2}$ 余子式 $M_{ij}$ 与代数余子式 $A_{ij}$ 差一个符号，题目问哪个就答哪个；求"余子式之和"时要先把每个 $M_{ij}$ 化成 $(-1)^{i+j}A_{ij}$，角标之和为奇数时添负号。 $\circled{3}$ $A_{ij}$ 的数值与 $a_{ij}$ 本身\textbf{无关}。例如 $\begin{vmatrix}1&2&3\\2&3&4\\3&4&5\end{vmatrix}$ 与 $\begin{vmatrix}1&2&3\\5&6&7\\3&4&5\end{vmatrix}$ 的第 1 行元素的代数余子式完全相同（因为划掉第 1 行后剩下的东西一样）。

\end{callwarn}

\begin{callexam}{真题考法}
这是填空题的高频考法，典型设问是"设 $A=(a_{ij})_{3\times3}$，$\vert A\vert=2$，求 $A_{11}+A_{21}+A_{31}$"或"$M_{11}+M_{21}+M_{31}+M_{41}$"。答题步骤固定两步：\textbf{先把余子式统一化成代数余子式（补符号），再把待求行（列）替换成对应系数，构造新行列式并计算}。陷阱：忘记"角标和为奇数添负号"；把待求的行余子式错当成列余子式；新行列式算错。另一类是"$a_{ij}+A_{ij}=0$"型条件反求 $\vert A\vert$——把条件翻译成 $A^*=-A^T$，两边取行列式得 $\vert A\vert^{n-1}=(-1)^n\vert A\vert$；\textbf{$n=3$ 时得 $\vert A\vert^2=-\vert A\vert$，故 $\vert A\vert=0$ 或 $\vert A\vert=-1$，必须再用"$A$ 非零"等条件排除其中一个}，否则会漏解或多解。

\end{callexam}

\subsection{抽象行列式}

\subsubsection{抽象行列式：$\vert kA\vert$、$\vert AB\vert$、$\vert A^*\vert$、$\vert A^{-1}\vert$、$\vert A^T\vert$ $\bigstar$$\bigstar$$\bigstar$$\bigstar$$\bigstar$}

设 $A,B$ 均为 $n$ 阶方阵，$k$ 为常数。以下是必须闭卷默写的全套公式：

\[\vert A^T\vert=\vert A\vert,\qquad \vert kA\vert=k^n\vert A\vert,\qquad \vert AB\vert=\vert A\vert\vert B\vert,\]

\[\vert A^{-1}\vert=\frac{1}{\vert A\vert}\ (\vert A\vert\ne0),\qquad \vert A^*\vert=\vert A\vert^{\,n-1},\qquad \vert A^k\vert=\vert A\vert^k .\]

\textbf{关键点说明}：

\begin{itemize}
\item $\vert kA\vert=k^n\vert A\vert$ 中的指数是\textbf{阶数} $n$，不是 $2$ 也不是 $1$。因为 $kA$ 相当于把 $A$ 的\textbf{每一行}都乘 $k$，共 $n$ 行，所以提出来 $n$ 个 $k$。
\item $\vert AB\vert=\vert A\vert\vert B\vert$ 要求 $A,B$ 是\textbf{同阶方阵}；它可推广到多个：$\vert A_1A_2\cdots A_m\vert=\vert A_1\vert\vert A_2\vert\cdots\vert A_m\vert$。由它立即得到\textbf{重要推论}：$\vert AB\vert=\vert BA\vert$（哪怕 $AB\ne BA$，它们的行列式也相等）。
\item $\vert A^*\vert=\vert A\vert^{n-1}$ 的推导：由 $AA^*=\vert A\vert E$ 两边取行列式，得 $\vert A\vert\vert A^*\vert=\vert A\vert^n$，当 $\vert A\vert\ne0$ 时约去 $\vert A\vert$ 即得。\textbf{$\vert A\vert=0$ 时该式给出 $\vert A^*\vert=0^{n-1}$}：$n\ge2$ 时这正是对的（此时 $r(A^*)\le1<n$，故 $\vert A^*\vert=0$），所以公式对 $n\ge2$ 一律成立。
\item \textbf{相似不变}：若 $A$ 与 $B$ 相似（$B=P^{-1}AP$），则 $\vert B\vert=\vert P^{-1}\vert\vert A\vert\vert P\vert=\vert A\vert$。更一般地，$\vert f(A)\vert$ 可由 $A$ 的特征值直接算出：若 $A$ 的特征值为 $\lambda_1,\cdots,\lambda_n$，则 $\vert f(A)\vert=\prod_{i=1}^{n}f(\lambda_i)$。
\end{itemize}
\textbf{套路（抽象行列式的固定三步）}

\begin{enumerate}
\item \textbf{统一阶数}：把所有矩阵化成同阶方阵的形式，看清每个符号的阶数 $n$。
\item \textbf{套公式化简}：优先用 $\vert AB\vert=\vert A\vert\vert B\vert$ 把乘积拆开，用 $\vert kA\vert=k^n\vert A\vert$ 把系数提到阶数次幂，用 $\vert A^*\vert=\vert A\vert^{n-1}$ 把伴随化成行列式的幂。
\item \textbf{由条件反解}：把已知的 $\vert A\vert$ 或矩阵等式代入。遇矩阵等式时先取行列式。
\end{enumerate}
\textbf{例}：设 $A$ 为 $3$ 阶方阵，$\vert A\vert=2$，求 $\vert (3A)^{-1}-2A^*\vert$。

先化简表达式。由 $(3A)^{-1}=\dfrac{1}{3}A^{-1}=\dfrac{1}{3\vert A\vert}A^*=\dfrac{1}{6}A^*$，所以

\[(3A)^{-1}-2A^*=\frac{1}{6}A^*-2A^*=-\frac{11}{6}A^*,\]

\[\left\vert (3A)^{-1}-2A^*\right\vert=\left(-\frac{11}{6}\right)^3\vert A^*\vert=-\frac{1331}{216}\cdot\vert A\vert^{2}=-\frac{1331}{216}\times4=-\frac{1331}{54}.\]

\textbf{验证关键步骤}：$n=3$ 时 $\vert A^*\vert=\vert A\vert^{3-1}=2^2=4$；系数提出时是 $3$ 次幂。这两处是本题唯一的两个坑。

\textbf{抽象行列式的六种常见题面}（先归类，再套对应动作）

\begin{enumerate}
\item \textbf{给 $\vert A\vert$ 与阶数，求若干个矩阵组合的行列式}：把 $A^*$ 换成 $\vert A\vert A^{-1}$、把 $kA$ 的系数按 $k^n$ 提出，最终都化成 $\vert A\vert$ 的幂。
\item \textbf{给向量组的行列式}：如 $A=(\alpha_1,\alpha_2,\alpha_3)$，$B=(\alpha_1+\alpha_2+\alpha_3,\ \alpha_1+2\alpha_2+4\alpha_3,\ \alpha_1+3\alpha_2+9\alpha_3)$，求 $\vert B\vert$。做法是\textbf{把 $B$ 写成 $A$ 右乘某个矩阵}：$B=A\begin{pmatrix}1&1&1\\1&2&3\\1&4&9\end{pmatrix}$，右乘矩阵恰是（转置的）范德蒙德行列式，其行列式为 $(2-1)(3-1)(3-2)=2$，故 $\vert B\vert=\vert A\vert\cdot2$。\textbf{这一步把所有"向量组重组"题都归约为"$\vert B\vert=\vert A\vert\cdot\vert C\vert$"。}
\item \textbf{矩阵方程型}：方程中含 $A^*$ 或 $A^{-1}$ 时，先统一换成 $A^{-1}$（用 $A^*=\vert A\vert A^{-1}$）再取行列式。
\item \textbf{相似型}：若 $A$ 与 $B$ 相似，则 $\vert A\vert=\vert B\vert$，且 $f(A)$ 与 $f(B)$ 相似、$\vert f(A)\vert=\vert f(B)\vert$。
\item \textbf{特征值型}：若 $A$ 的特征值为 $\lambda_1,\cdots,\lambda_n$，则 $\vert A\vert=\prod\lambda_i$，$\vert f(A)\vert=\prod f(\lambda_i)$——这是求抽象行列式最快的一条路。
\item \textbf{分块型}：把整体分成若干块后套 $\begin{vmatrix}A&O\\O&B\end{vmatrix}=\vert A\vert\vert B\vert$ 或 $\begin{vmatrix}O&A\\B&O\end{vmatrix}=(-1)^{mn}\vert A\vert\vert B\vert$。
\end{enumerate}
\textbf{例（向量组重组型）}：设 $\alpha_1,\alpha_2,\alpha_3$ 是三维列向量，$A=(\alpha_1,\alpha_2,\alpha_3)$，$\vert A\vert=1$，求 $\vert B\vert$，其中 $B=(\alpha_1+\alpha_2+\alpha_3,\ \alpha_1+2\alpha_2+4\alpha_3,\ \alpha_1+3\alpha_2+9\alpha_3)$。

把 $B$ 按列看成 $A$ 的列向量的线性组合，写成矩阵乘积：

\[B=A\begin{pmatrix}1&1&1\\1&2&3\\1&4&9\end{pmatrix}=AC .\]

$\vert B\vert=\vert A\vert\vert C\vert$，而 $C$ 是范德蒙德行列式（列分别是 $1,2,3$ 的幂），

\[\vert C\vert=(2-1)(3-1)(3-2)=1\times2\times1=2 .\]

故 $\vert B\vert=1\times2=2$。\textbf{验证}：$C$ 的第一行全为 $1$、第二行是 $1,2,3$、第三行是 $1,4,9$，确为三阶范德蒙德型，$\vert C\vert=(x_2-x_1)(x_3-x_1)(x_3-x_2)=2$，无误。

\textbf{例（相似型）}：设 $A$ 为 $3$ 阶矩阵，其特征值为 $1,-2,3$，求 $\vert A^2-2A+3E\vert$。

设 $f(x)=x^2-2x+3$。矩阵 $f(A)$ 的特征值为 $f(1)=1-2+3=2$、$f(-2)=4+4+3=11$、$f(3)=9-6+3=6$（这一条的证明见第五章：若 $A\alpha=\lambda\alpha$ 则 $f(A)\alpha=f(\lambda)\alpha$）。故

\[\vert A^2-2A+3E\vert=2\times11\times6=132 .\]

\textbf{验证}：$\vert A\vert=1\times(-2)\times3=-6$；若取 $A=\operatorname{diag}(1,-2,3)$ 这一最简相似代表，则 $f(A)=\operatorname{diag}(2,11,6)$，行列式 $=132$，与上式一致。

\textbf{例（含伴随矩阵）}：设 $A$ 为 $3$ 阶矩阵，$\vert A\vert=3$，求 $\left\vert (A^*)^{-1}\right\vert$ 与 $\left\vert (A^*)^*\right\vert$。

第一个用 $\vert A^*\vert=\vert A\vert^{n-1}$ 的倒数：

\[\left\vert (A^*)^{-1}\right\vert=\frac{1}{\vert A^*\vert}=\frac{1}{\vert A\vert^{\,n-1}}=\frac{1}{3^{2}}=\frac{1}{9}.\]

第二个要小心两层伴随。分两步：先由 $(A^*)^{-1}=\left(\vert A\vert A^{-1}\right)^{-1}=\dfrac{1}{\vert A\vert}A$ 与 $(A^*)^*=\vert A^*\vert(A^*)^{-1}$ 得

\[(A^*)^*=\vert A\vert^{\,n-1}\cdot\frac{1}{\vert A\vert}A=A\,\vert A\vert^{\,n-2}.\]

再取行列式。注意 $\vert A\vert^{\,n-2}$ 是一个\textbf{标量} $c$，而 $\vert cA\vert=c^{\,n}\vert A\vert$（这里的 $n$ 也要对 $c$ 起作用）：

\[\left\vert (A^*)^*\right\vert=\left\vert A\cdot\vert A\vert^{\,n-2}\right\vert=\left(\vert A\vert^{\,n-2}\right)^{n}\vert A\vert=\vert A\vert^{\,n(n-2)+1}=\vert A\vert^{\,(n-1)^2}.\]

代入 $n=3$、$\vert A\vert=3$ 得 $\left\vert (A^*)^*\right\vert=3^{\,4}=81$。

\textbf{验证}：换一条路。$(A^*)^*$ 是 $A^*$ 的伴随矩阵，$A^*$ 的阶数仍是 $n=3$，故可直接用"伴随的行列式"公式：

\[\left\vert (A^*)^*\right\vert=\vert A^*\vert^{\,n-1}=\left(\vert A\vert^{\,n-1}\right)^{\,n-1}=\vert A\vert^{\,(n-1)^2}=3^{\,4}=81,\]

两条路一致。\textbf{这就是本考点最容易翻车的一步}：标量提出来时忘记乘 $n$ 次幂，会得到错误的 $3^9=19683$ 或 $3^6=729$。

\begin{callkey}{必背}
全套结论写成一行：

\[\vert A^T\vert=\vert A\vert,\quad \vert A^{-1}\vert=\vert A\vert^{-1},\quad \vert A^*\vert=\vert A\vert^{n-1},\quad \vert kA\vert=k^n\vert A\vert,\quad \vert AB\vert=\vert A\vert\vert B\vert,\quad \vert A^k\vert=\vert A\vert^k .\]

再加一条分块公式：$\begin{vmatrix}A&O\\O&B\end{vmatrix}=\vert A\vert\vert B\vert$，$\begin{vmatrix}O&A\\B&O\end{vmatrix}=(-1)^{mn}\vert A\vert\vert B\vert$。

\end{callkey}

\begin{callwarn}{易错}
$\circled{1}$ $\vert A+B\vert\ne\vert A\vert+\vert B\vert$，也\textbf{没有任何}一般公式。遇到 $\vert A+B\vert$ 只能靠已有条件另想办法（如凑成已知矩阵的乘积、或用特征值）。 $\circled{2}$ $\vert kA\vert=k^n\vert A\vert$ 的 $n$ 是\textbf{矩阵的阶数}，与 $\vert A^*\vert=\vert A\vert^{n-1}$ 的 $n-1$ 不要串。 $\circled{3}$ $\vert A^*\vert=\vert A\vert^{n-1}$ 中若 $\vert A\vert=0$ 且 $n=1$，$0^0$ 无意义——但 $1$ 阶情形考研不考，$n\ge2$ 时上式都成立。 $\circled{4}$ 抽象行列式题中，\textbf{先看清楚 $A$ 是几阶}。题干说"$A$ 为 $3$ 阶矩阵"还是"$n$ 阶矩阵"，决定了答案是数字还是含 $n$ 的表达式。

\end{callwarn}

\begin{callexam}{真题考法}
这是线代选择题的必考题型。设问常为"已知 $\vert A\vert=2$，$A$ 为 $3$ 阶矩阵，求 $\vert (2A)^{-1}+A^*\vert$"或"$\vert A^*\vert=\underline{\quad}$"。解题路径固定：\textbf{先把待求式用 $A^{-1}=\frac{1}{\vert A\vert}A^*$ 统一成 $A^*$ 或 $A$ 的倍数，再提系数（记得是 $n$ 次幂），最后代 $\vert A\vert$}。另一个高频形态是分块矩阵的行列式，例如 $M=\begin{pmatrix}O&A\\B&O\end{pmatrix}$ 问 $\vert M\vert$，此时要用 $(-1)^{mn}$ 且注意 $m,n$ 分别是 $A,B$ 的阶数。陷阱：把 $\vert (2A)^{-1}\vert$ 算成 $\frac{1}{2}\vert A^{-1}\vert$（漏了 $2^{-n}$），或把 $(-1)^{mn}$ 写成 $(-1)^{m+n}$。

\end{callexam}

\subsubsection{代数余子式的线性组合与行列式的构造 $\bigstar$$\bigstar$$\bigstar$$\bigstar$}

这一考点是上一考点的"逆用"：题目给的不是 $a_{ij}$，而是 $A_{ij}$ 之间的关系（例如"所有代数余子式之和""$A_{ij}$ 与 $a_{ij}$ 满足某等式"），要求反过来求行列式或某组代数余子式之和。

\textbf{核心工具}：伴随矩阵。$A$ 的伴随矩阵定义为

\[A^*=\begin{pmatrix}A_{11}&A_{21}&\cdots&A_{n1}\\A_{12}&A_{22}&\cdots&A_{n2}\\\vdots&\vdots&&\vdots\\A_{1n}&A_{2n}&\cdots&A_{nn}\end{pmatrix},\]

即\textbf{把代数余子式按"行变列"的位置排列}（$A^*$ 的第 $i$ 行第 $j$ 列是 $A_{ji}$）。于是"求所有代数余子式之和"就是"求 $A^*$ 的所有元素之和"。

\textbf{四种常见题型与套路}

\textbf{题型一：求 $\sum_{i,j}A_{ij}$ 或 $\sum_j A_{ij}$}

用展开定理构造新行列式。特别是，若 $A$ 的\textbf{各行元素之和都等于 $k$}，则有一组漂亮的结论：

\begin{itemize}
\item $A\begin{pmatrix}1\\1\\\vdots\\1\end{pmatrix}=k\begin{pmatrix}1\\1\\\vdots\\1\end{pmatrix}$，即 $k$ 是 $A$ 的特征值，$\begin{pmatrix}1&1&\cdots&1\end{pmatrix}^T$ 是对应的特征向量。
\item 由 $AA^*=\vert A\vert E$，两边右乘全 $1$ 列向量：$A(A^*\mathbf{1})=\vert A\vert\mathbf{1}$，即 $A^*\mathbf{1}=\dfrac{\vert A\vert}{k}\mathbf{1}$（$\vert A\vert\ne0$ 时）。这说明 \textbf{$A^*$ 的各行元素之和都等于 $\dfrac{\vert A\vert}{k}$}。
\item 特别地，若 $\vert A\vert=1$，则 $A^*$ 的每行元素之和为 $\dfrac{1}{k}$。转置后可得"各列元素之和为 $k$"的版本：此时 $A^T$ 各行和为 $k$，故 $(A^T)^*=A^{*T}$ 各行和为 $\dfrac{\vert A\vert}{k}$，即 $A^*$ \textbf{各列}元素之和为 $\dfrac{\vert A\vert}{k}$。
\end{itemize}
\textbf{题型二：给出 $A_{ij}$ 与 $a_{ij}$ 的等式，求 $\vert A\vert$}

把 $A_{ij}=-a_{ij}$ 这类条件翻译成矩阵语言。若 $a_{ij}+A_{ij}=0$ 对所有 $i,j$ 成立，即 $A^*=-A^T$（因为 $A^*$ 的第 $i$ 行第 $j$ 列是 $A_{ji}$，而 $A^T$ 的第 $i$ 行第 $j$ 列是 $a_{ji}$）。两边取行列式：$\vert A^*\vert=\vert -A^T\vert=(-1)^n\vert A\vert$，而 $\vert A^*\vert=\vert A\vert^{n-1}$，于是

\[\vert A\vert^{n-1}=(-1)^n\vert A\vert .\]

$n=3$ 时得 $\vert A\vert^2=-\vert A\vert$，即 $\vert A\vert(\vert A\vert+1)=0$。若题设另有 $A\ne O$ 且 $A^*\ne O$ 等信息，可排掉 $\vert A\vert=0$ 得到 $\vert A\vert=-1$。

\textbf{题型三：给出新的矩阵 $B$，求 $B$ 的代数余子式之和}

先把 $B$ 用 $A$ 表示（如 $B=QAQ^{-1}$ 或 $B=A^T$、$B=kA$），再用 $\vert B\vert$、$B^*$ 与 $A^*$ 的关系换算。常用关系：

\[(kA)^*=k^{\,n-1}A^*,\qquad (A^T)^*=(A^*)^T,\qquad (AB)^*=B^*A^* .\]

\textbf{题型四：行列式与矩阵方程结合}

若题干给出含 $A^*$ 的矩阵方程（如 $ABA^*-BA^*=BAB^*$），先在方程两侧同乘或提取公因子，把 $A^*,B^*$ 换回 $A^{-1},B^{-1}$（用 $A^*=\vert A\vert A^{-1}$），把方程化成只含 $A,B$ 的形式，再取行列式解出 $\vert B\vert$ 或 $B^*$。

\textbf{例}：设 $A=(a_{ij})_{3\times3}$ 的每行元素之和均为 $2$，$\vert A\vert=6$，求 $A_{11}+A_{21}+A_{31}$。

因为每行和为 $2$，所以 $\vert A\vert=2$ 倍的"把第 1 列换成全 $1$ 后"……更直接地，$A_{11}+A_{21}+A_{31}$ 是 $A^*$ 第 1 行的和。由上面的结论，$A^*$ 每行元素之和为 $\dfrac{\vert A\vert}{k}=\dfrac{6}{2}=3$。故 $A_{11}+A_{21}+A_{31}=3$。

\textbf{直接验证}：由 $A^*\mathbf{1}=\dfrac{\vert A\vert}{k}\mathbf{1}=3\mathbf{1}$，第 1 个分量正是 $A_{11}\cdot1+A_{21}\cdot1+A_{31}\cdot1=3$。

\begin{callkey}{必背}
伴随矩阵的位置约定与三条换算：

\[A^*=(A_{ji})_{n\times n},\qquad AA^*=A^*A=\vert A\vert E,\qquad A^{-1}=\frac{1}{\vert A\vert}A^*\ (\vert A\vert\ne0),\]

\[\vert A^*\vert=\vert A\vert^{n-1},\qquad (kA)^*=k^{n-1}A^*,\qquad (A^*)^*=A\,\vert A\vert^{\,n-2}\ (n\ge2),\qquad (AB)^*=B^*A^* .\]

"各行和为 $k$"的翻译：$k$ 是特征值，全 $1$ 向量是特征向量，$A^*$ 各行和为 $\vert A\vert/k$。

\end{callkey}

\begin{callwarn}{易错}
$\circled{1}$ \textbf{伴随矩阵是"转置着放"的}：$A^*$ 的第 $i$ 行是 $A$ 第 $i$ 列对应的代数余子式，即 $A^*$ 第 $i$ 行第 $j$ 列 $=A_{ji}$。组装时排错位置是最常见的错误——写成 $A_{ij}$ 就得到了"代数余子式矩阵"而不是伴随矩阵，两者恰好互为转置。 $\circled{2}$ 用公式 $A^{-1}=\dfrac{1}{\vert A\vert}A^*$ 时，必须保证 $\vert A\vert\ne0$；$\vert A\vert=0$ 时 $A^{-1}$ 不存在，但 $A^*$ 仍可能非零（$\vert A\vert=0$ 且 $r(A)=n-1$ 时 $r(A^*)=1$）。 $\circled{3}$ "两行之和相等"与"两列之和相等"结论不同：\textbf{行和}对应 $A$ 右乘全 $1$ 向量，\textbf{列和}对应左乘全 $1$ 向量，换算到 $A^*$ 上方向正好相反，不要搞混。 $\circled{4}$ $(A^*)^*=A\vert A\vert^{n-2}$：$n=2$ 时 $\vert A\vert^{0}=1$，即 $(A^*)^*=A$，这是对的；但若 $\vert A\vert=0$ 且 $n=2$，公式给出 $O$，实际用 $A^*$ 的计算式验证更好。

\end{callwarn}

\begin{callexam}{真题考法}
常以填空题形式出现："设 $A=(a_{ij})_{3\times3}$，$A_{ij}$ 为 $a_{ij}$ 的代数余子式，且 $a_{ij}+A_{ij}=0$，$\vert A\vert=\underline{\quad}$"，或"设 $A$ 的每行元素之和为 $k$，$\vert A\vert=a$，则 $A_{11}+A_{12}+\cdots+A_{1n}=\underline{\quad}$"。第一类的解法是把条件翻译成 $A^*=-\alpha A^T$（$\alpha=1$）再取行列式得方程；第二类的解法是用 $AA^*=\vert A\vert E$ 右乘全 $1$ 向量。另有一类综合题把行列式与伴随矩阵方程、矩阵方程混在一起考，关键动作是\textbf{用 $A^*=\vert A\vert A^{-1}$ 把伴随换回可逆矩阵}。

\end{callexam}

\subsection*{本章易错清单}
\addcontentsline{toc}{subsection}{本章易错清单}

\begin{enumerate}
\item \textbf{副对角行列式的符号}：$\begin{vmatrix}0&a\\b&0\end{vmatrix}=-ab$，$n$ 阶副对角公式带 $(-1)^{\frac{n(n-1)}{2}}$。这个符号 $n=2,3$ 为负、$n=4,5$ 为正，周期为 $4$；记不住就现场数交换次数，\textbf{不要凭感觉写正号}。
\item \textbf{代数余子式漏符号}：$A_{ij}=(-1)^{i+j}M_{ij}$。题目问"余子式"还是"代数余子式"要看清楚；求"余子式之和"时必须逐个补上 $(-1)^{i+j}$，角标和为奇数就添负号。
\item \textbf{伴随矩阵位置排错}：$A^*$ 的第 $i$ 行第 $j$ 列是 $A_{ji}$，不是 $A_{ij}$。伴随 $=$ 代数余子式矩阵的\textbf{转置}，横写时列标要"前面"。
\item \textbf{$\vert kA\vert=k^n\vert A\vert$ 的指数}：是矩阵的阶数 $n$，不是 $2$，也不是 $1$。这条与 $\vert A^*\vert=\vert A\vert^{n-1}$ 的 $n-1$、$(kA)^*=k^{n-1}A^*$ 的 $n-1$ 极易混淆，务必分清"数乘提 $n$ 次、伴随提 $n-1$ 次"。
\item \textbf{误用 $\vert A+B\vert=\vert A\vert+\vert B\vert$}：行列式对加法\textbf{不成比例}，$\vert A+B\vert$ 没有一般公式，只能借条件（特征值、特征多项式）另算；也不要写成 $\vert A+B\vert=\vert A\vert+\vert B\vert$ 或从 $\vert A+B\vert=0$ 推出 $\vert A\vert=-\vert B\vert$。
\item \textbf{$\vert AB\vert=\vert A\vert\vert B\vert$ 的前提}：$A,B$ 必须是\textbf{同阶方阵}。只有在这个前提下才有 $\vert AB\vert=\vert A\vert\vert B\vert=\vert B\vert\vert A\vert=\vert BA\vert$。若 $A$ 为 $m\times n$、$B$ 为 $n\times m$ 且 $m\ne n$，$AB$ 与 $BA$ 阶数都不同，上式不适用（这类长方阵结论统考不考）。
\item \textbf{范德蒙德公式因子顺序}：结果是 $\prod_{1\le j<i\le n}(x_i-x_j)$，即"大下标减小下标"，且 $x_1,\cdots,x_n$ 的顺序就是行列式中列的顺序（第 1 列对应 $x_1$）。缺行缺列时先补行、提公因子或转置，不要硬算。
\item \textbf{分块公式 $(-1)^{mn}$ 中的 $m,n$}：分别是 $A$ 与 $B$ 的\textbf{阶数}（即两个对角块的阶数），不是位置序号。也常有考生把 $(-1)^{mn}$ 写成 $(-1)^{m+n}$——四阶时前者为 $+1$，后者为 $+1$，恰好撞对；五阶时前者为 $+1$、后者为 $-1$，就露馅了。
\item \textbf{化三角形时丢符号}：用"交换两行"和"某行乘非零数"都会改变行列式的值，必须补符号或补系数；只有"倍加"不变值。中途丢了符号是"算到最后差一个负号"的头号原因。
\item \textbf{递推法的起始项}：用递推关系求通项前，必须先独立算出 $D_1,D_2$，确认递推式在两个起始点上成立，否则通项会整体差一个常数或符号。递推式的"两次展开"（按第 1 行展开后再对其中一个子式按第 1 列展开）也要写清楚。
\item \textbf{$\vert A\vert=0$ 与 $A=O$ 混淆}：$\vert A\vert=0$ 只说明 $A$ 不可逆、行（列）向量组线性相关，\textbf{不说明} $A$ 是零矩阵；同理 $A^*=O$ 也不等价于 $A=O$（$A^*=O$ 只说明 $r(A)\le n-2$）。
\item \textbf{抽象题中阶数看漏}：题干里"$A$ 为 $3$ 阶矩阵"与"$A$ 为 $n$ 阶矩阵"答案形式完全不同，前者给具体数字，后者给含 $n$ 的表达式。做题第一件事就是把阶数圈出来，并核对每个符号的阶数。
\item \textbf{$\vert A\vert=0$ 时滥用 $\vert A^*\vert=\vert A\vert^{n-1}$ 的"倒数"形式}：不能由 $A^*=\vert A\vert A^{-1}$ 去求 $A^*$（$A^{-1}$ 不存在），要用 $AA^*=O$ 与秩的分类讨论来处理奇异情形。
\item \textbf{把"某一行元素之和"当成"行列式的值"}：行列式不等于任何一行（列）元素之和。只有当各行（列）元素之和相等时，才可以用"加列法"把这个和提取出来作为公因子，且\textbf{提取后还要乘以剩下的那个行列式}，不能直接把和当作答案。
\end{enumerate}
\chapter{矩阵}
\markboth{矩阵}{矩阵}
\subsection*{本章考纲要求}
\addcontentsline{toc}{subsection}{本章考纲要求}

\begin{quote}
\textbf{考试内容}

\begin{enumerate}
\item 矩阵的概念
\item 矩阵的线性运算
\item 矩阵的乘法
\item 方阵的幂
\item 方阵乘积的行列式
\item 矩阵的转置
\item 逆矩阵的概念和性质
\item 矩阵可逆的充分必要条件
\item 伴随矩阵
\item 矩阵的初等变换
\item 初等矩阵
\item 矩阵的秩
\item 矩阵的等价
\item 分块矩阵及其运算
\end{enumerate}
\textbf{考试要求}

\begin{enumerate}
\item 理解矩阵的概念，了解单位矩阵、数量矩阵、对角矩阵、三角矩阵、对称矩阵、反对称矩阵和正交矩阵以及它们的性质。
\item 掌握矩阵的线性运算、乘法、转置以及它们的运算规律，了解方阵的幂与方阵乘积的行列式的性质。
\item 理解逆矩阵的概念，掌握逆矩阵的性质以及矩阵可逆的充分必要条件。理解伴随矩阵的概念，会用伴随矩阵求逆矩阵。
\item 了解矩阵初等变换的概念，了解初等矩阵的性质和矩阵等价的概念，理解矩阵的秩的概念，掌握用初等变换求矩阵的秩和逆矩阵的方法。
\item 了解分块矩阵及其运算。
\end{enumerate}
\end{quote}
\subsection*{本章星级总表}
\addcontentsline{toc}{subsection}{本章星级总表}

\begingroup\small\setlength{\tabcolsep}{4.0pt}
\begin{longtable}{@{}L{0.4484\linewidth}L{0.1074\linewidth}L{0.1074\linewidth}L{0.1074\linewidth}L{0.1235\linewidth}@{}}
\toprule
\textbf{考点} & \textbf{星级} & \textbf{考纲要求} & \textbf{难度} & \textbf{常见题型} \\
\midrule
\endhead
矩阵的概念与特殊矩阵（单位、数量、对角、三角、对称、反对称、正交） & $\bigstar$$\bigstar$ & 理解/了解 & 基础 & 选择 \\
矩阵的线性运算、乘法、转置及其运算规律 & $\bigstar$$\bigstar$$\bigstar$ & 掌握 & 基础 & 选择、填空 \\
方阵的幂与方阵乘积的行列式 & $\bigstar$$\bigstar$$\bigstar$$\bigstar$ & 了解 & 中等 & 填空、解答 \\
逆矩阵的概念、性质与可逆的充分必要条件 & $\bigstar$$\bigstar$$\bigstar$$\bigstar$$\bigstar$ & 理解/掌握 & 中等 & 选择、解答 \\
伴随矩阵及其公式（$AA^*=\vert A\vert E$、$\vert A^*\vert$、$(A^*)^*$、秩的讨论） & $\bigstar$$\bigstar$$\bigstar$$\bigstar$$\bigstar$ & 理解/会 & 难 & 选择、填空、证明 \\
矩阵的初等变换与初等矩阵（左行右列） & $\bigstar$$\bigstar$$\bigstar$$\bigstar$ & 了解/理解 & 中等 & 选择、填空 \\
矩阵的秩（定义、求法、秩的不等式与常用结论） & $\bigstar$$\bigstar$$\bigstar$$\bigstar$$\bigstar$ & 理解/掌握 & 难 & 选择、填空、证明 \\
矩阵等价与分块矩阵及其运算 & $\bigstar$$\bigstar$ & 了解 & 中等 & 选择 \\
矩阵方程 $AX=B$、$XA=B$、$AXB=C$ 的求解 & $\bigstar$$\bigstar$$\bigstar$$\bigstar$ & 会 & 中等 & 解答 \\
\bottomrule
\end{longtable}
\endgroup

\section{矩阵的概念与运算}

\subsection{矩阵与特殊矩阵}

\subsubsection{矩阵的概念与特殊矩阵（单位、数量、对角、三角、对称、反对称、正交）$\bigstar$$\bigstar$}

\textbf{定义}：由 $m\times n$ 个数 $a_{ij}$ 排成的 $m$ 行 $n$ 列的矩形数表

\[A=\begin{pmatrix}a_{11}&a_{12}&\cdots&a_{1n}\\a_{21}&a_{22}&\cdots&a_{2n}\\\vdots&\vdots&&\vdots\\a_{m1}&a_{m2}&\cdots&a_{mn}\end{pmatrix}\]

称为 $m\times n$ 矩阵，记作 $A=(a_{ij})_{m\times n}$，$a_{ij}$ 是第 $i$ 行第 $j$ 列的元素。当 $m=n$ 时称为\textbf{方阵}，$n$ 称为阶数。只有方阵才谈行列式、伴随矩阵、逆矩阵、特征值。

\textbf{几类必须记住的特殊矩阵}

\begin{itemize}
\item \textbf{零矩阵} $O$：所有元素为 $0$。
\item \textbf{单位矩阵} $E$（也记作 $I$）：主对角元素全为 $1$、其余全为 $0$。性质：$AE=EA=A$。
\item \textbf{数量矩阵} $kE$：主对角元素全为同一个数 $k$。它与\textbf{任意同阶方阵都可交换}，即 $(kE)A=A(kE)=kA$——这是"矩阵可交换"最常见的来源。
\item \textbf{对角矩阵} $\Lambda=\operatorname{diag}(d_1,d_2,\cdots,d_n)$：非主对角元全为 $0$。两个同阶对角阵之积仍是对角阵，且 $\operatorname{diag}(a_i)\operatorname{diag}(b_i)=\operatorname{diag}(a_ib_i)$，\textbf{必可交换}。
\item \textbf{三角矩阵}：主对角线上方（下方）全为 $0$ 的方阵，分别称为上（下）三角矩阵。两个同阶同向三角阵之积仍是同向三角阵；三角阵的行列式等于主对角元之积。
\item \textbf{对称矩阵}：满足 $A^T=A$ 的方阵，即 $a_{ij}=a_{ji}$。
\item \textbf{反对称矩阵}：满足 $A^T=-A$ 的方阵，即 $a_{ij}=-a_{ji}$。\textbf{必然有主对角元全为 $0$}（取 $i=j$ 得 $a_{ii}=-a_{ii}$）。
\item \textbf{正交矩阵}：满足 $Q^TQ=E$（等价于 $Q^{-1}=Q^T$）的方阵。
\end{itemize}
\textbf{关于对称与反对称的三条常考结论}

\begin{enumerate}
\item 任意 $n$ 阶方阵 $A$ 都可唯一分解为一个对称矩阵与一个反对称矩阵之和：
\end{enumerate}
\[A=\underbrace{\frac{A+A^T}{2}}_{\text{对称}}+\underbrace{\frac{A-A^T}{2}}_{\text{反对称}} .\]

验证：$\left(\dfrac{A+A^T}{2}\right)^T=\dfrac{A^T+A}{2}$，故为对称阵；$\left(\dfrac{A-A^T}{2}\right)^T=\dfrac{A^T-A}{2}=-\dfrac{A-A^T}{2}$，故为反对称阵。

\begin{enumerate}
\item 对任意 $m\times n$ 矩阵 $B$，$B^TB$ 是 $n$ 阶对称矩阵，$BB^T$ 是 $m$ 阶对称矩阵。
\item \textbf{奇数阶反对称矩阵的行列式必为 $0$}：由 $A^T=-A$ 两边取行列式得 $\vert A\vert=\vert A^T\vert=\vert -A\vert=(-1)^n\vert A\vert$；$n$ 为奇数时得 $\vert A\vert=-\vert A\vert$，故 $\vert A\vert=0$。\textbf{推论}：奇数阶反对称矩阵必不可逆；以它为系数矩阵的齐次方程组必有非零解。
\end{enumerate}
\textbf{正交矩阵的性质}：$\vert Q\vert=\pm1$；$Q^T$、$Q^{-1}$、$Q^*$ 都是正交矩阵；两个同阶正交矩阵之积仍是正交矩阵；$Q$ 的列（行）向量组是标准正交向量组（两两正交且长度为 $1$）。

\begin{callwarn}{易错}
$\circled{1}$ 反对称矩阵的主对角元\textbf{全为 $0$}，这是常考小结论；但反过来，"主对角元全为 $0$"的矩阵不一定是反对称矩阵。 $\circled{2}$ 只有\textbf{方阵}才可能有对称、反对称、正交、可逆、伴随、特征值这些概念，看到长方阵不要往这上面想。 $\circled{3}$ 正交矩阵的定义是 $Q^TQ=E$。对\textbf{实方阵}而言它等价于 $QQ^T=E$（同阶方阵时 $AB=E\Rightarrow BA=E$），但 "$Q$ 是方阵"这一前提不能少。

\end{callwarn}

\begin{callexam}{真题考法}
选择题常给出一个含参数的矩阵，问它是对称、反对称还是正交阵，并求参数。答题动作：对称$\to$列 $a_{ij}=a_{ji}$ 的方程；反对称$\to$列 $a_{ij}=-a_{ji}$ 且 $a_{ii}=0$；正交$\to$列 $Q^TQ=E$（即列向量两两正交且长度为 $1$）。另一种考法是给"$A$ 为奇数阶反对称矩阵"问 $\vert A\vert$，或给"$A^T=-A$ 且 $n$ 为奇数"证明 $\vert A\vert=0$。

\end{callexam}

\subsection{矩阵的运算}

\subsubsection{矩阵的线性运算、乘法、转置及其运算规律 $\bigstar$$\bigstar$$\bigstar$}

\textbf{加法}：只有\textbf{同型}矩阵（行数、列数都相同）才能相加，对应位置元素相加。满足交换律、结合律，且 $A+O=A$。

\textbf{数乘}：$kA$ 表示 $A$ 的每个元素都乘 $k$。满足 $k(A+B)=kA+kB$、$(k+l)A=kA+lA$、$(kl)A=k(lA)$、$1\cdot A=A$。

\textbf{乘法}：设 $A=(a_{ij})_{m\times s}$、$B=(b_{ij})_{s\times n}$，则 $C=AB=(c_{ij})_{m\times n}$，其中

\[c_{ij}=\sum_{k=1}^{s}a_{ik}b_{kj}=a_{i1}b_{1j}+a_{i2}b_{2j}+\cdots+a_{is}b_{sj}\qquad(1\le i\le m,\ 1\le j\le n).\]

\textbf{可乘的条件}：\textbf{左矩阵的列数 $=$ 右矩阵的行数}；乘积的型为"左矩阵的行数 $\times$ 右矩阵的列数"。记作 $(m\times s)\cdot(s\times n)=(m\times n)$：\textbf{内标相同才能乘，外标就是结果型}。

\textbf{乘法的运算规律}

\begin{itemize}
\item \textbf{结合律}：$(AB)C=A(BC)$。
\item \textbf{分配律}：$A(B+C)=AB+AC$，$(B+C)A=BA+CA$。
\item 与数乘相容：$k(AB)=(kA)B=A(kB)$。
\item \textbf{交换律不成立}：一般 $AB\ne BA$。
\item \textbf{消去律不成立}：$AB=AC$ 且 $A\ne O$，\textbf{不能}推出 $B=C$。
\item \textbf{零因子存在}：$AB=O$ 且 $A\ne O,B\ne O$ 是可能的。
\item \textbf{没有除法}，但可逆矩阵有逆。
\end{itemize}
\textbf{矩阵乘法的三种理解方式}（考试中经常要切换视角）

\begin{enumerate}
\item \textbf{元素角度}：$c_{ij}$ 是 $A$ 的第 $i$ 行与 $B$ 的第 $j$ 列的内积。
\item \textbf{行、列角度}：把 $B$ 按列分块为 $B=(b_1,b_2,\cdots,b_n)$，则 $AB=(Ab_1,Ab_2,\cdots,Ab_n)$——\textbf{$AB$ 的列向量可由 $A$ 的列向量线性表示}；把 $A$ 按行分块，则 $AB$ 的行向量可由 $B$ 的行向量线性表示。
\item \textbf{对角阵角度}：$A\operatorname{diag}(d_1,\cdots,d_n)$ 相当于把 $A$ 的\textbf{第 $j$ 列}乘 $d_j$；$\operatorname{diag}(d_1,\cdots,d_m)A$ 相当于把 $A$ 的\textbf{第 $i$ 行}乘 $d_i$。这与初等变换的"左行右列"完全一致。
\end{enumerate}
\textbf{转置}：把 $A$ 的行与列互换得到 $A^T$。若 $A$ 为 $m\times n$，则 $A^T$ 为 $n\times m$，且 $(A^T)_{ij}=a_{ji}$。运算规律：

\[(A^T)^T=A,\qquad (A+B)^T=A^T+B^T,\qquad (kA)^T=kA^T,\qquad (AB)^T=B^TA^T .\]

最后一条最重要：\textbf{乘积转置要反序}，可推广为 $(A_1A_2\cdots A_m)^T=A_m^T\cdots A_2^TA_1^T$。

\textbf{例}：设 $A=\begin{pmatrix}1&0\\2&-1\end{pmatrix}$，$B=\begin{pmatrix}3&1\\0&4\end{pmatrix}$，求 $AB$ 与 $BA$。

\[AB=\begin{pmatrix}1\cdot3+0\cdot0&1\cdot1+0\cdot4\\2\cdot3+(-1)\cdot0&2\cdot1+(-1)\cdot4\end{pmatrix}=\begin{pmatrix}3&1\\6&-2\end{pmatrix},\]

\[BA=\begin{pmatrix}3\cdot1+1\cdot2&3\cdot0+1\cdot(-1)\\0\cdot1+4\cdot2&0\cdot0+4\cdot(-1)\end{pmatrix}=\begin{pmatrix}5&-1\\8&-4\end{pmatrix}\ne AB .\]

\textbf{验证}：$AB\ne BA$，交换律不成立；但 $\vert AB\vert=3\times(-2)-1\times6=-12$，$\vert BA\vert=5\times(-4)-(-1)\times8=-12$，两者相等——这正是 $\vert AB\vert=\vert A\vert\vert B\vert=\vert BA\vert$ 的体现。

\begin{callkey}{必背}
\textbf{可交换的常见来源}（见到 $AB=BA$ 就想到这几类）：$\circled{1}$ 其中一个为数量矩阵 $kE$；$\circled{2}$ 其中一个为另一个的多项式 $f(A)$；$\circled{3}$ 两者都是对角矩阵；$\circled{4}$ $A^*$、$A^{-1}$、$A^k$、$f(A)$ 都与 $A$ 可交换。

\textbf{$AB=O$ 的四种翻译}：$\circled{1}$ 秩：$r(A)+r(B)\le n$（$n$ 为 $A$ 的列数）；$\circled{2}$ 方程组：$B$ 的每一列都是 $Ax=0$ 的解；$\circled{3}$ 正交：$A$ 的行向量与 $B$ 的列向量正交；$\circled{4}$ 特征：$B$ 的非零列向量是 $A$ 的属于特征值 $0$ 的特征向量。

\end{callkey}

\textbf{矩阵等式化简的四板斧}（抽象矩阵题的全部招式）

统考中的抽象矩阵题，几乎都是"给一个矩阵等式，要求证明某结论或求某矩阵"。化简时只有四个动作，来回组合就够了：

\textbf{（1）同乘}：两边同时左乘（或右乘）同一个可逆矩阵。\textbf{注意左右顺序不可交换}：要消去左边的 $A$ 就左乘 $A^{-1}$，要消去右边的 $A$ 就右乘 $A^{-1}$。

\textbf{（2）提}：把相同的矩阵因子从同一侧提出来。$AX+BX=(A+B)X$（右提），$XA+XB=X(A+B)$（左提）。\textbf{同一侧才能提}，$AX+XB$ 提不出任何公因子。

\textbf{（3）转置}：对等式两边同时取转置，注意 $(AB)^T=B^TA^T$ 的顺序。常用于把 $AB=A^TB^T$ 型条件转化为可提公因子的形式。

\textbf{（4）$E$ 的妙用}：$E=AA^{-1}=A^{-1}A$ 可以在任意位置乘上去；若 $A$ 是正交矩阵则 $E=AA^T=A^TA$。\textbf{在等式两边同时"添 $E$"再分解}，是把 $AB=A+B$ 这类式子凑成 $(A-E)(B-E)=E$ 的关键动作。

\textbf{标准流程}：\textbf{补 $E$ $\to$ 分解 $E$ $\to$ 两边提公因子 $\to$ 读出结论。}

\textbf{例}：设 $A,B$ 为 $n$ 阶矩阵，$AB=A+B$，证明 $AB=BA$ 并求 $B$。

移项得 $AB-A-B=O$。\textbf{补 $E$}：两边加 $E$，得 $AB-A-B+E=E$。\textbf{分解}：左边恰好是 $(A-E)(B-E)$。于是

\[(A-E)(B-E)=E .\]

由此：$\circled{1}$ $A-E$ 与 $B-E$ 互为逆矩阵，都可逆；$\circled{2}$ $B=E+(A-E)^{-1}$；$\circled{3}$ 由 $(B-E)(A-E)=E$ 展开得 $BA-B-A+E=E$，与 $AB-A-B+E=E$ 对照即得 $AB=BA$。

\textbf{例}：设 $A$ 为 $n$ 阶矩阵且 $A^2=A$（幂等），证明 $A$ 可对角化（本处只用到化简技巧）。由 $A^2-A=O$ 得 $A(A-E)=O$；由此立即可得 $r(A)+r(A-E)\le n$（秩一章用），同时 $(A-E)A=O$ 也成立。这是"化简后发现两个积为零的矩阵关系"的典型。

\begin{callkey}{必背}
抽象矩阵等式的三个高频"配方"（背下来直接套）：

\[AB=A+B\ \Longrightarrow\ (A-E)(B-E)=E;\]

\[A^2-3A+2E=O\ \Longrightarrow\ A(A-3E)=-2E\ \Longrightarrow\ A^{-1}=\frac{1}{2}(3E-A);\]

\[A^3=O\ \Longrightarrow\ (E-A)(E+A+A^2)=E\ \Longrightarrow\ (E-A)^{-1}=E+A+A^2 .\]

\end{callkey}

\begin{callwarn}{易错}
$\circled{1}$ $AB=O\ \not\Rightarrow\ A=O$ 或 $B=O$。反例：$\begin{pmatrix}1&1\\1&1\end{pmatrix}\begin{pmatrix}1&-1\\-1&1\end{pmatrix}=O$。 $\circled{2}$ $AB=AC\ \not\Rightarrow\ B=C$，除非 $A$ 可逆（此时两边左乘 $A^{-1}$）。 $\circled{3}$ $(A+B)^2=A^2+2AB+B^2$ \textbf{不成立}，正确的是 $(A+B)^2=A^2+AB+BA+B^2$。同理 $(A-B)(A+B)=A^2+AB-BA-B^2\ne A^2-B^2$。 $\circled{4}$ $(AB)^k\ne A^kB^k$（$k\ge2$ 时一般不等），除非 $A,B$ 可交换。 $\circled{5}$ $(AB)^T=B^TA^T$ 的顺序不能丢，转置\textbf{必然反序}。 $\circled{6}$ 由 $A^2=A$ 不能推出 $A=O$ 或 $A=E$；由 $A^2=O$ 不能推出 $A=O$。

\end{callwarn}

\begin{callexam}{真题考法}
选择题的经典形态是"下列命题中正确的是"，四个选项分别考察交换律、消去律、$AB=O$ 的推论、$(AB)^T$ 的顺序。另一类考法是给抽象矩阵等式化简：已知 $A^2=A$（幂等）或 $A^2=E$（对合）求某个 $f(A)$ 的表达式——做法是把多项式作带余除法，用 $A^2=A$ 或 $A^2=E$ 降幂。还有一类是具体的 $AB$ 或 $(AB)^T$ 计算题，注意先检查可乘性。

\end{callexam}

\subsection{方阵的幂}

\subsubsection{方阵的幂与方阵乘积的行列式 $\bigstar$$\bigstar$$\bigstar$$\bigstar$}

\textbf{定义}：设 $A$ 为 $n$ 阶方阵，$k$ 为正整数，$A^k=\underbrace{A\cdot A\cdots A}_{k\ \text{个}}$，并规定 $A^0=E$。运算律：$A^mA^n=A^{m+n}$，$(A^m)^n=A^{mn}$。

\textbf{注意}：$(AB)^k\ne A^kB^k$，除非 $AB=BA$。

\textbf{求 $A^k$ 的五种固定套路}

\textbf{套路一：试算找规律}（最简单的情形）

先算 $A^2$、$A^3$，观察是否有周期性或简单形式。例如 $A=\begin{pmatrix}0&1\\1&0\end{pmatrix}$ 满足 $A^2=E$，于是 $A^{2025}=A$；$A=\begin{pmatrix}0&1\\-1&0\end{pmatrix}$ 满足 $A^2=-E$，于是 $A^4=E$，按 $k$ 除以 $4$ 的余数分四类讨论。\textbf{旋转矩阵} $C=\begin{pmatrix}\cos\theta&-\sin\theta\\\sin\theta&\cos\theta\end{pmatrix}$ 满足 $C^n=\begin{pmatrix}\cos n\theta&-\sin n\theta\\\sin n\theta&\cos n\theta\end{pmatrix}$。

\textbf{套路二：幂零矩阵}（存在 $k$ 使 $A^k=O$）

若 $A$ 是主对角元全为 $0$ 的严格三角阵，则 $A^n=O$。求 $A^k$ 时直接乘几次即可，或用二项式展开时只保留前几项。

\textbf{套路三：二项展开型} $A=\lambda E+N$，其中 $N$ 幂零

因为 $\lambda E$ 与任何矩阵都可交换，可用二项式定理：

\[A^k=(\lambda E+N)^k=\sum_{i=0}^{k}\binom{k}{i}\lambda^{k-i}N^i .\]

由于 $N$ 幂零（设 $N^m=O$），求和实际只有前 $m$ 项。最典型的是 $A=\begin{pmatrix}\lambda&1\\0&\lambda\end{pmatrix}=\lambda E+N$（$N^2=O$），得

\[A^k=\lambda^kE+k\lambda^{k-1}N=\begin{pmatrix}\lambda^k&k\lambda^{k-1}\\0&\lambda^k\end{pmatrix}.\]

\textbf{套路四：秩为 $1$ 的矩阵}

若 $r(A)=1$，则 $A$ 必可分解为一个非零列向量与一个非零行向量之积：$A=\alpha\beta^T$（$\alpha,\beta$ 均为 $n$ 维列向量）。此时

\[A^2=(\alpha\beta^T)(\alpha\beta^T)=\alpha(\beta^T\alpha)\beta^T=(\beta^T\alpha)A,\]

其中 $\beta^T\alpha=\alpha^T\beta=\operatorname{tr}(A)$ 是一个\textbf{数}。记 $l=\operatorname{tr}(A)$，则 $A^2=lA$，进而

\[A^k=l^{\,k-1}A=(\operatorname{tr}A)^{k-1}A .\]

\textbf{例}：$A=\begin{pmatrix}2&0&1\\4&0&2\\6&0&3\end{pmatrix}$。注意每列都是 $\begin{pmatrix}1\\2\\3\end{pmatrix}$ 的倍数，写成 $A=\begin{pmatrix}2\\4\\6\end{pmatrix}\begin{pmatrix}1&0&\frac12\end{pmatrix}$……为了结论干净，改用整数分解：$A=\begin{pmatrix}1\\2\\3\end{pmatrix}\begin{pmatrix}2&0&1\end{pmatrix}$。此时 $\operatorname{tr}(A)=2+0+3=5$，而 $\beta^T\alpha=1\times2+2\times0+3\times1=5$，两者一致。于是

\[A^{n}=5^{\,n-1}A=\begin{pmatrix}2\cdot5^{n-1}&0&5^{n-1}\\4\cdot5^{n-1}&0&2\cdot5^{n-1}\\6\cdot5^{n-1}&0&3\cdot5^{n-1}\end{pmatrix}.\]

\textbf{验证}：代入 $n=2$，得 $A^2=5A$。直接算 $A^2$：$A$ 的第 1 行乘第 1 列为 $2\cdot2+0\cdot4+1\cdot6=10=5\times2$；第 1 行乘第 3 列为 $2\cdot1+0\cdot2+1\cdot3=5=5\times1$，与 $5A$ 的对应元素一致，结论正确。

\textbf{套路五：相似对角化}

若 $A=P\Lambda P^{-1}$（$\Lambda$ 为对角阵），则 $A^k=P\Lambda^kP^{-1}$，其中 $\Lambda^k=\operatorname{diag}(\lambda_1^k,\cdots,\lambda_n^k)$。这是最通用的一条路，具体判定与求法在特征值一章。

\textbf{方阵乘积的行列式}

对同阶方阵 $A,B$：

\[\vert AB\vert=\vert A\vert\cdot\vert B\vert,\qquad \vert A^k\vert=\vert A\vert^k,\qquad \vert AB\vert=\vert BA\vert .\]

\textbf{重要推论}：$A,B$ 中只要有一个不可逆（行列式为 $0$），则 $AB$ 必不可逆；反过来，若 $AB$ 可逆，则 $A$ 与 $B$ 都可逆。

\begin{callwarn}{易错}
$\circled{1}$ $A^mA^n=A^{m+n}$ 只对\textbf{同一个矩阵} $A$ 成立；$(AB)^m=A^mB^m$ 一般\textbf{不成立}。 $\circled{2}$ 幂零矩阵的幂零指数不一定等于阶数：$n$ 阶严格上三角阵的 $n$ 次幂必为零矩阵，但可能更早为零。 $\circled{3}$ 秩为 $1$ 时 $A^k=l^{k-1}A$ 中的 $l=\operatorname{tr}(A)$ 是一个\textbf{数}；当 $l=0$ 时 $A^2=O$，此时 $A^k=O\ (k\ge2)$，且 $A$ 不能相似对角化（$A\ne O$ 时）。 $\circled{4}$ 用二项式展开的前提是两项\textbf{可交换}。写成 $A=B+C$ 时若 $BC\ne CB$，二项式展开无效。

\end{callwarn}

\begin{callexam}{真题考法}
填空题与解答题都常见，典型设问是"设 $A=\begin{pmatrix}1&0&1\\0&2&0\\1&0&1\end{pmatrix}$，求 $A^n$"或"设 $A=\begin{pmatrix}1&2\\0&1\end{pmatrix}$，求 $A^{2025}$"。答题第一步永远是\textbf{判断 $A$ 属于哪一类}：试算有周期？幂零？$\lambda E+$ 幂零？秩为 $1$？可对角化？判错类会白算半天。陷阱：把 $(AB)^k$ 拆成 $A^kB^k$；秩为 $1$ 时忘记 $A^k$ 中 $\operatorname{tr}A$ 的指数是 $k-1$ 而不是 $k$。

\end{callexam}

\section{逆矩阵与伴随矩阵}

\subsection{可逆的判定}

\subsubsection{逆矩阵的概念、性质与可逆的充分必要条件 $\bigstar$$\bigstar$$\bigstar$$\bigstar$$\bigstar$}

\textbf{定义}：设 $A$ 为 $n$ 阶方阵，若存在 $n$ 阶矩阵 $B$ 使得

\[AB=BA=E,\]

则称 $A$ \textbf{可逆}，$B$ 称为 $A$ 的\textbf{逆矩阵}，记作 $A^{-1}$。逆矩阵若存在必\textbf{唯一}。

\textbf{方阵情形的重要简化}：若 $A,B$ 为\textbf{同阶方阵}且 $AB=E$，则必有 $BA=E$，即只验证一侧即可。\textbf{注意}：对非方阵不成立。

\textbf{矩阵可逆的充分必要条件（必须能一口气说全）}

设 $A$ 为 $n$ 阶方阵，下列条件\textbf{互相等价}：

\begin{enumerate}
\item $\vert A\vert\ne0$；
\item $r(A)=n$（满秩）；
\item $A$ 可逆（存在 $B$ 使 $AB=BA=E$）；
\item $A$ 的行向量组线性无关（等价地，列向量组线性无关）；
\item $Ax=0$ 只有零解；
\item 对任意 $n$ 维列向量 $b$，$Ax=b$ 有唯一解；
\item $A$ 的特征值全不为 $0$；
\item $A$ 可表示为若干初等矩阵的乘积；
\item $A$ 与 $E$ 等价（即 $A$ 可经初等变换化为 $E$）。
\end{enumerate}
\textbf{口诀：秩满可逆不为零，无关唯一只零解。}

\textbf{逆矩阵的运算性质}

\[(A^{-1})^{-1}=A,\qquad (AB)^{-1}=B^{-1}A^{-1},\qquad (kA)^{-1}=\frac{1}{k}A^{-1}\ (k\ne0),\]

\[(A^T)^{-1}=(A^{-1})^T,\qquad (A^{-1})^T=(A^T)^{-1},\qquad \vert A^{-1}\vert=\vert A\vert^{-1},\qquad (A^*)^{-1}=(A^{-1})^* .\]

\textbf{求逆矩阵的四条路}

\textbf{（1）伴随矩阵法}（二阶必用，三阶可用）

\[A^{-1}=\frac{1}{\vert A\vert}A^*\qquad(\vert A\vert\ne0).\]

二阶矩阵有速记口诀——\textbf{主对角对调，副对角变号，再除以行列式}：

\[A=\begin{pmatrix}a&b\\c&d\end{pmatrix}\ \Longrightarrow\ A^{-1}=\frac{1}{ad-bc}\begin{pmatrix}d&-b\\-c&a\end{pmatrix}.\]

\textbf{（2）初等行变换法}（三阶及以上首选）

\[(A\mid E)\xrightarrow{\ \text{初等行变换}\ }(E\mid A^{-1}).\]

只在\textbf{左半边}做行变换，右半边跟着一起变；当左半边化为 $E$ 时，右半边就是 $A^{-1}$。若过程中左半边出现全零行，则 $A$ 不可逆。

\textbf{（3）定义法（凑因式法）}（抽象题首选）

要把 $A^{-1}$ 求出来，就是找出一个 $B$ 使 $AB=E$。手法是从已知的矩阵等式出发，\textbf{凑出 $A$ 与某个多项式的乘积等于 $E$}。

\textbf{例}：设 $A$ 满足 $A^2-3A+2E=O$。由 $A^2-3A=-2E$ 得 $A(A-3E)=-2E$，两边同除以 $-2$：

\[A\cdot\left[-\frac{1}{2}(A-3E)\right]=E\ \Longrightarrow\ A^{-1}=-\frac{1}{2}(A-3E)=\frac{1}{2}(3E-A).\]

\textbf{验证}：$A\cdot\dfrac{1}{2}(3E-A)=\dfrac{1}{2}(3A-A^2)=\dfrac{1}{2}(3A-(3A-2E))=E$，成立。

\textbf{求 $(A-\lambda E)^{-1}$ 的通用手法（带余除法）}：设 $A$ 满足多项式方程 $f(A)=O$，要找 $(A-\lambda E)^{-1}$，就用 $f(x)$ 除以 $x-\lambda$：

\[f(x)=(x-\lambda)q(x)+c\quad(c=f(\lambda)\ \text{为常数}).\]

代入 $A$ 得 $O=f(A)=(A-\lambda E)q(A)+f(\lambda)E$，若 $f(\lambda)\ne0$ 则

\[(A-\lambda E)\cdot\left[-\frac{1}{f(\lambda)}q(A)\right]=E\ \Longrightarrow\ (A-\lambda E)^{-1}=-\frac{1}{f(\lambda)}q(A).\]

\textbf{例}：$A^3-2A^2+3A-4E=O$，求 $A^{-1}$ 与 $(A-E)^{-1}$。

（i）由 $A^3-2A^2+3A=4E$，即 $A(A^2-2A+3E)=4E$，得

\[A^{-1}=\frac{1}{4}(A^2-2A+3E).\]

（ii）令 $f(x)=x^3-2x^2+3x-4$，$f(1)=1-2+3-4=-2\ne0$。作带余除法：$f(x)=(x-1)(x^2-x+2)-2$，故

\[(A-E)(A^2-A+2E)-2E=O\ \Longrightarrow\ (A-E)(A^2-A+2E)=2E,\]

\[(A-E)^{-1}=\frac{1}{2}(A^2-A+2E).\]

\textbf{验证（i）}：$A\cdot\dfrac14(A^2-2A+3E)=\dfrac14(A^3-2A^2+3A)=\dfrac14\cdot4E=E$，成立。

\textbf{（4）分块矩阵法}

\[\begin{pmatrix}A&O\\O&B\end{pmatrix}^{-1}=\begin{pmatrix}A^{-1}&O\\O&B^{-1}\end{pmatrix},\qquad \begin{pmatrix}A&O\\C&B\end{pmatrix}^{-1}=\begin{pmatrix}A^{-1}&O\\-B^{-1}CA^{-1}&B^{-1}\end{pmatrix},\]

\[\begin{pmatrix}O&A\\B&O\end{pmatrix}^{-1}=\begin{pmatrix}O&B^{-1}\\A^{-1}&O\end{pmatrix}.\]

\textbf{抽象矩阵可逆性的三种证明入口}

\begin{enumerate}
\item \textbf{定义法}：凑出 $AB=E$（只需一侧，因是同阶方阵），或解出 $B=A^{-1}$。
\item \textbf{行列式法}：证 $\vert A\vert\ne0$。常配合"$\vert A\vert\ne0\iff Ax=0$ 只有零解"，用\textbf{反证法}：若 $A$ 不可逆，则存在 $\xi\ne0$ 使 $A\xi=0$，推出矛盾。
\item \textbf{秩法}：证 $r(A)=n$，常用 $r(A)\le n$ 与 $r(A)\ge n$ 两侧夹逼。
\end{enumerate}
\textbf{例}：设 $A$ 为 $n$ 阶方阵且 $A^2=A$（幂等），$A\ne E$。证明 $A$ 不可逆。

\textbf{反证}：若 $A$ 可逆，在 $A^2=A$ 两边左乘 $A^{-1}$ 得 $A=E$，与 $A\ne E$ 矛盾，故 $A$ 不可逆。

\textbf{例}：设 $n$ 阶矩阵 $A$ 满足 $A^3=O$，证明 $E-A$ 可逆并求逆。

由 $E-A^3=E$ 且 $E-A^3=(E-A)(E+A+A^2)$（因为 $E-A^3=(E-A)(E+A+A^2)$ 是立方差分解，代入 $A^3=O$ 逐项验证亦成立），得

\[(E-A)^{-1}=E+A+A^2 .\]

\textbf{验证}：$(E-A)(E+A+A^2)=E+A+A^2-A-A^2-A^3=E-A^3=E$，成立。

\begin{callkey}{必背}
\textbf{求逆两把钥匙}：具体矩阵用初等行变换 $(A\mid E)\to(E\mid A^{-1})$；抽象矩阵用\textbf{定义法凑因式}，把已知等式因式分解出 $(A-\lambda E)$ 并配成 $E$（本质是带余除法）。

\end{callkey}

\begin{callwarn}{易错}
$\circled{1}$ $AB=E\Rightarrow BA=E$ 只在 $A,B$ \textbf{同阶方阵}时成立，长方阵不成立。 $\circled{2}$ $(AB)^{-1}=B^{-1}A^{-1}$，顺序必须颠倒。 $\circled{3}$ $(A+B)^{-1}\ne A^{-1}+B^{-1}$；一般地 $A+B$ 甚至可能不可逆（如 $B=-A$）。 $\circled{4}$ $(kA)^{-1}=\dfrac{1}{k}A^{-1}$，系数在逆中\textbf{不提幂次}（对比 $\vert kA\vert=k^n\vert A\vert$ 要提 $n$ 次幂）。 $\circled{5}$ 用初等变换法求逆时\textbf{只能做行变换}。若用列变换，应写成 $\begin{pmatrix}A\\E\end{pmatrix}\xrightarrow{\text{列变换}}\begin{pmatrix}E\\A^{-1}\end{pmatrix}$。 $\circled{6}$ 由 $A^2=O$ 可推出 $A$ 必不可逆（若 $A$ 可逆则 $A=O$，与 $A^2=O$ 且 $A\ne O$ 矛盾）；但由 $A^2=A$ 得不出可逆与否，要结合 $r(A)$、$A\ne E$ 等信息。

\end{callwarn}

\begin{callexam}{真题考法}
线代每年的必考内容。\textbf{选择/填空题}常问"下列矩阵中不可逆的是"，或给含参数矩阵问可逆的参数条件（用 $\vert A\vert=0$ 解方程）。\textbf{解答题}常以抽象矩阵等式出现，如"设 $A$ 满足 $A^3-2A^2+3A-4E=O$，求 $A^{-1}$ 与 $(A-E)^{-1}$"——解法固定：凑成 $A\cdot(\cdots)=E$ 得 $A^{-1}$；求 $(A-\lambda E)^{-1}$ 用带余除法。另一高频形态是"$A$ 为 $n$ 阶矩阵，且 $A,A-E$ 均可逆，化简 $A^{-1}-(A-E)^{-1}$"，用恒等式 $(A^{-1}-(A-E)^{-1})$ 化为 $-(A^2-A)^{-1}$ 之类。陷阱：忘记 $AB=E$ 只需验一侧；把 $(AB)^{-1}$ 写成 $A^{-1}B^{-1}$；分块求逆时负号位置放错。

\end{callexam}

\subsection{伴随矩阵}

\subsubsection{伴随矩阵及其公式（$AA^*=\vert A\vert E$、$\vert A^*\vert$、$(A^*)^*$、秩的讨论）$\bigstar$$\bigstar$$\bigstar$$\bigstar$$\bigstar$}

\textbf{定义}：设 $A=(a_{ij})_{n\times n}$，$A_{ij}$ 为元素 $a_{ij}$ 的代数余子式，则称矩阵

\[A^*=\begin{pmatrix}A_{11}&A_{21}&\cdots&A_{n1}\\A_{12}&A_{22}&\cdots&A_{n2}\\\vdots&\vdots&&\vdots\\A_{1n}&A_{2n}&\cdots&A_{nn}\end{pmatrix}\]

为 $A$ 的\textbf{伴随矩阵}。即 $A^*=(A_{ji})_{n\times n}$：$A^*$ 的第 $i$ 行第 $j$ 列元素是 $A_{ji}$。\textbf{一句话记住：伴随矩阵 $=$ 代数余子式矩阵的转置。}

\textbf{核心公式（必须闭卷默写）}

\[AA^*=A^*A=\vert A\vert E,\]

\[A^{-1}=\frac{1}{\vert A\vert}A^*\quad(\vert A\vert\ne0),\qquad A^*=\vert A\vert A^{-1}\quad(\vert A\vert\ne0),\]

\[\vert A^*\vert=\vert A\vert^{\,n-1},\]

\[(kA)^*=k^{\,n-1}A^*,\qquad (A^T)^*=(A^*)^T,\qquad (AB)^*=B^*A^*,\qquad (A^{-1})^*=(A^*)^{-1},\]

\[(A^*)^*=A\,\vert A\vert^{\,n-2}\qquad(n\ge2).\]

\textbf{推导要点}：$AA^*=\vert A\vert E$ 直接来自展开定理（主对角元是"自交"得 $\vert A\vert$，非对角元是"互交"得 $0$）。由此两边取行列式得 $\vert A\vert\vert A^*\vert=\vert A\vert^n$；$\vert A\vert\ne0$ 时约去 $\vert A\vert$ 即 $\vert A^*\vert=\vert A\vert^{n-1}$。再用 $A^*=\vert A\vert A^{-1}$：

\[(A^*)^*=\vert A^*\vert(A^*)^{-1}=\vert A\vert^{n-1}\cdot\frac{1}{\vert A\vert}A=A\,\vert A\vert^{\,n-2}.\]

\textbf{二阶矩阵的伴随矩阵}（速记：\textbf{主对调，副变号}）

\[A=\begin{pmatrix}a&b\\c&d\end{pmatrix}\ \Longrightarrow\ A^*=\begin{pmatrix}d&-b\\-c&a\end{pmatrix}.\]

\textbf{三阶矩阵求伴随的机械化操作}（省时高效，四步走）

\begin{enumerate}
\item \textbf{抄前两列、再抄前两行}：把原矩阵的第 $1,2$ 列接在右边、第 $1,2$ 行接在下边，把 $3\times3$ 扩充成 $5\times5$。
\item 对原矩阵的每个位置 $(i,j)$，\textbf{划掉第 $i$ 行第 $j$ 列}，在剩余的 $4\times4$ 区块中取"两条对角线相乘再相减"。
\item \textbf{竖着算、横着写}：每个位置的代数余子式算出来后，按 $A_{ji}$ 的位置摆放（第 $i$ 行第 $j$ 列放 $A_{ji}$）。
\item 最后除以 $\vert A\vert$ 即得逆矩阵。
\end{enumerate}
\textbf{伴随矩阵的秩（必考）}

设 $A$ 为 $n$ 阶方阵（$n\ge2$），则

\[r(A^*)=\begin{cases}n, & r(A)=n,\\1, & r(A)=n-1,\\0, & r(A)\le n-2.\end{cases}\]

\textbf{理由}：$r(A)=n$ 时 $A$ 可逆，由 $A^*=\vert A\vert A^{-1}$ 知 $A^*$ 可逆，$r(A^*)=n$；$r(A)=n-1$ 时 $A$ 存在 $n-1$ 阶非零子式，故至少一个 $A_{ij}\ne0$，即 $A^*\ne O$；又 $AA^*=\vert A\vert E=O$，由 $AB=O$ 的秩不等式得 $r(A)+r(A^*)\le n$，即 $r(A^*)\le1$，结合 $A^*\ne O$ 得 $r(A^*)=1$；$r(A)\le n-2$ 时所有 $n-1$ 阶子式全为 $0$，故 $A^*=O$。

\textbf{由秩看 $A^*$ 的形态}：$r(A^*)=1$ 时 $A^*$ 是秩 $1$ 矩阵，可写成 $\alpha\beta^T$ 的形式，从而 $(A^*)^2=\operatorname{tr}(A^*)A^*$，这一分解在求 $A^*$ 的幂、特征值时很有用。

\textbf{伴随矩阵的三大应用}

\begin{itemize}
\item \textbf{求逆}：$A^{-1}=\dfrac{1}{\vert A\vert}A^*$，二阶必用。
\item \textbf{克拉默法则的矩阵形式}：$Ax=b$（$\vert A\vert\ne0$）的解为 $x=\dfrac{1}{\vert A\vert}A^*b$，即 $x_j=\dfrac{\vert A_j\vert}{\vert A\vert}$，其中 $A_j$ 是把 $A$ 的第 $j$ 列换成 $b$ 所得的矩阵。
\item \textbf{"代数余子式的线性组合"问题}：全部归结为 $AA^*=\vert A\vert E$ 与展开定理。
\end{itemize}
\textbf{例}：设 $A=\begin{pmatrix}1&2\\3&4\end{pmatrix}$，求 $A^*$ 与 $A^{-1}$。

\[\vert A\vert=1\times4-2\times3=-2,\qquad A^*=\begin{pmatrix}4&-2\\-3&1\end{pmatrix},\qquad A^{-1}=\frac{1}{-2}\begin{pmatrix}4&-2\\-3&1\end{pmatrix}=\begin{pmatrix}-2&1\\\frac{3}{2}&-\frac{1}{2}\end{pmatrix}.\]

\textbf{验证}：$AA^{-1}=\begin{pmatrix}1&2\\3&4\end{pmatrix}\begin{pmatrix}-2&1\\\frac32&-\frac12\end{pmatrix}=\begin{pmatrix}-2+3&1-1\\-6+6&3-2\end{pmatrix}=\begin{pmatrix}1&0\\0&1\end{pmatrix}$，正确。

\textbf{例（秩的讨论）}：设 $A$ 为 $4$ 阶方阵，$r(A)=3$，求 $r(A^*)$ 与 $\vert A^*\vert$。

由 $r(A)=n-1=3$ 得 $r(A^*)=1$；又 $r(A)=3<4$ 故 $\vert A\vert=0$，从而 $\vert A^*\vert=\vert A\vert^{\,n-1}=0^3=0$。两个结论并不矛盾：$r(A^*)=1$ 表示 $A^*\ne O$，但它作为 $4$ 阶矩阵仍不可逆。

\begin{callkey}{必背}
\textbf{伴随矩阵六大公式一览}

\[AA^*=A^*A=\vert A\vert E,\qquad \vert A^*\vert=\vert A\vert^{n-1},\qquad (kA)^*=k^{n-1}A^*,\qquad (A^T)^*=(A^*)^T,\]

\[(AB)^*=B^*A^*,\qquad (A^*)^*=A\vert A\vert^{n-2}.\]

\textbf{秩三分类}：$r(A)=n\Rightarrow r(A^*)=n$；$r(A)=n-1\Rightarrow r(A^*)=1$；$r(A)\le n-2\Rightarrow A^*=O$。

\textbf{特征值关联}：若 $\lambda$ 是 $A$ 的特征值且 $\lambda\ne0$，则 $\dfrac{\vert A\vert}{\lambda}$ 是 $A^*$ 的特征值，且对应的特征向量不变。

\end{callkey}

\begin{callwarn}{易错}
$\circled{1}$ \textbf{位置不能排错}：$A^*$ 的元素是 $A_{ji}$（下标前后调换），直接按 $A_{ij}$ 摆会得到转置，这是最常见的失分点。 $\circled{2}$ $r(A)=n-1$ 时 \textbf{$A^*\ne O$}（这是 $r(A^*)=1$ 的一半理由），不能因为 $\vert A\vert=0$ 就断定 $A^*=O$。 $\circled{3}$ $n=2$ 时 $(A^*)^*=A\vert A\vert^{0}=A$，这个特例最容易被忽略，公式依然有效。 $\circled{4}$ $\vert A^*\vert=\vert A\vert^{n-1}$、$(kA)^*=k^{n-1}A^*$ 的指数都是 $n-1$，而 $\vert kA\vert=k^n\vert A\vert$ 的指数是 $n$。三个式子常同时出现，务必分清。 $\circled{5}$ $(AB)^*=B^*A^*$ 的顺序必须颠倒（与 $(AB)^{-1}=B^{-1}A^{-1}$ 一致）。 $\circled{6}$ 用 $A^*=\vert A\vert A^{-1}$ 时必须 $\vert A\vert\ne0$；$\vert A\vert=0$ 时该式无意义，但 $AA^*=\vert A\vert E=O$ 仍成立，这是处理奇异情形的关键。

\end{callwarn}

\begin{callexam}{真题考法}
伴随矩阵是线代命题的重灾区，题型非常固定。\textbf{选择题}常考：$\circled{1}$ 给 $\vert A\vert$ 与阶数求 $\vert (kA)^*\vert$、$\vert (A^*)^*\vert$、$r(A^*)$；$\circled{2}$ 给 $A$ 的秩求 $A^*$ 的秩；$\circled{3}$ 给分块矩阵 $\begin{pmatrix}O&A\\B&O\end{pmatrix}$ 求其伴随矩阵（四个选项只差系数位置与系数值）。\textbf{填空题}常给含 $A^*$ 的方程反求 $A$ 或 $\vert A\vert$，做法是用 $A^*=\vert A\vert A^{-1}$ 把伴随换成逆。\textbf{证明题}常用 $AA^*=\vert A\vert E$ 配合秩不等式讨论。经典杀招：正交矩阵满足 $A^*=\vert A\vert A^{-1}=\vert A\vert A^T$，若 $\vert A\vert=-1$ 则 $A^*=-A^T$——"伴随等于转置的倍数（或相反数）"，是选择题的速解入口。

\end{callexam}

\section{初等变换与矩阵的秩}

\subsection{初等变换与初等矩阵}

\subsubsection{矩阵的初等变换与初等矩阵（左行右列）$\bigstar$$\bigstar$$\bigstar$$\bigstar$}

\textbf{三类初等行变换}（对列有完全平行的一套）

\begin{enumerate}
\item \textbf{对换变换}：交换两行，记 $r_i\leftrightarrow r_j$。
\item \textbf{倍乘变换}：用非零常数 $k$ 乘某一行，记 $kr_i\to r_i$（$k\ne0$）。
\item \textbf{倍加变换}：把某一行的 $k$ 倍加到另一行，记 $r_i+kr_j\to r_i$（$i\ne j$）。
\end{enumerate}
\textbf{行阶梯形}：每个非零行的首非零元（主元）所在列标严格递增，且主元下方元素全为 $0$。\textbf{行最简形}：在行阶梯形基础上，每个主元都是 $1$，且主元所在列的其余元素都是 $0$。

\textbf{初等矩阵}：单位矩阵经\textbf{一次}初等变换所得的矩阵，分三类：

\begin{itemize}
\item 对换矩阵 $E_{ij}$：交换 $E$ 的第 $i,j$ 行（列）所得；
\item 倍乘矩阵 $E_i(k)$：把 $E$ 的第 $i$ 行（列）乘 $k\ne0$ 所得；
\item 倍加矩阵 $E_{ij}(k)$：把 $E$ 的第 $j$ 行的 $k$ 倍加到第 $i$ 行（或第 $i$ 列的 $k$ 倍加到第 $j$ 列）所得。
\end{itemize}
\textbf{三个初等矩阵的行列式与逆（必须记住）}

\[\vert E_{ij}\vert=-1,\qquad \vert E_i(k)\vert=k\ne0,\qquad \vert E_{ij}(k)\vert=1;\]

\[E_{ij}^{-1}=E_{ij},\qquad E_i(k)^{-1}=E_i\left(\frac{1}{k}\right),\qquad E_{ij}(k)^{-1}=E_{ij}(-k).\]

即：\textbf{初等矩阵都可逆，且逆仍是同类型的初等矩阵}。

\textbf{最核心的定理（"左行右列"）}

设 $A$ 为 $m\times n$ 矩阵：

\begin{itemize}
\item 对 $A$ 做一次\textbf{初等行变换} $\iff$ 在 $A$ 的\textbf{左边}乘上相应的 $m$ 阶初等矩阵；
\item 对 $A$ 做一次\textbf{初等列变换} $\iff$ 在 $A$ 的\textbf{右边}乘上相应的 $n$ 阶初等矩阵。
\end{itemize}
\textbf{口诀：左行右列。}

\textbf{例}：设 $A$ 为 $3$ 阶矩阵，$P=\begin{pmatrix}0&1&0\\1&0&0\\0&0&1\end{pmatrix}$，则 $PA$ 相当于把 $A$ 的第 $1,2$ 行交换（左乘、行变换）；$AP$ 相当于把 $A$ 的第 $1,2$ 列交换（右乘、列变换）。

\textbf{两条重要推论}

\begin{enumerate}
\item \textbf{初等变换不改变矩阵的秩}：因为初等矩阵可逆，而"左乘或右乘可逆矩阵不改变秩"，即 $r(A)=r(PA)=r(AQ)=r(PAQ)$（$P,Q$ 可逆）。
\item \textbf{可逆矩阵可表示为若干初等矩阵的乘积}：$A$ 可逆 $\iff A=E_1E_2\cdots E_k$（每个 $E_i$ 都是初等矩阵）。
\end{enumerate}
\textbf{等价标准形}：任意 $m\times n$ 矩阵 $A$ 总可经过初等变换化为标准形

\[\begin{pmatrix}E_r&O\\O&O\end{pmatrix}_{m\times n},\qquad r=r(A).\]

该标准形由 $m,n,r$ 完全确定，且唯一。

\textbf{用初等变换求逆（回顾）}：$(A\mid E)\xrightarrow{\ \text{行变换}\ }(E\mid A^{-1})$。其成立理由正是"左行"——每一步行变换相当于左乘一个初等矩阵，把所有初等矩阵之积记为 $P$，则 $PA=E$ 即 $P=A^{-1}$；这些变换作用在右半边的 $E$ 上便得到 $A^{-1}$。

\begin{callkey}{必背}
\textbf{判断"左乘还是右乘"的黄金流程}：先看变换是\textbf{行}的还是\textbf{列}的 $\to$ 行变换左乘、列变换右乘 $\to$ 再数清楚乘的是几阶初等矩阵（行变换用左矩阵的行数，列变换用左矩阵的列数）。

\textbf{"反求 $A$"的固定动作}：把题设每一步都写成矩阵等式。例如"把 $A$ 的第 2 列加到第 1 列得 $B$"即 $B=AP_1$；"再交换 $B$ 的第 2、3 行得 $E$"即 $E=P_2B$。于是 $P_2AP_1=E$，得 $A=P_2^{-1}P_1^{-1}=P_2P_1$（因为对换矩阵的逆是它自己）。

\end{callkey}

\begin{callwarn}{易错}
$\circled{1}$ \textbf{左行右列不能反}：行变换左乘、列变换右乘。写成"行变换右乘"是最高频的低级错误。 $\circled{2}$ 初等矩阵的阶数要匹配：对 $m\times n$ 矩阵做行变换，左乘 \textbf{$m$ 阶}初等矩阵；做列变换，右乘 \textbf{$n$ 阶}初等矩阵。 $\circled{3}$ $E_{ij}(k)^{-1}=E_{ij}(-k)$（符号相反），不是 $E_{ij}(1/k)$；而 $E_i(k)^{-1}=E_i(1/k)$（取倒数），不是 $E_i(-k)$。 $\circled{4}$ 只有\textbf{倍加}变换不改变行列式的值；对换变号，倍乘使行列式乘 $k$。若 $A$ 经 $s$ 次对换与若干次倍乘 $k_1,\cdots,k_m$ 化为 $B$，则 $\vert A\vert=(-1)^s\dfrac{1}{k_1\cdots k_m}\vert B\vert$。 $\circled{5}$ 用初等变换求逆时\textbf{只能做行变换}；求秩时可同时用行、列变换。 $\circled{6}$ "$A$ 与 $B$ 等价"指存在可逆 $P,Q$ 使 $B=PAQ$，不要与相似、合同混淆。

\end{callwarn}

\begin{callexam}{真题考法}
近几年数二的高频考点。典型设问：$\circled{1}$ 给"$A$ 经某列变换得 $B$，再经某行变换得 $E$"，并给出两个初等矩阵 $P_1,P_2$，问 $A$ 等于 $P_1P_2$ 还是 $P_2P_1$——解题动作是把\textbf{每一步都写成矩阵等式再反解}；$\circled{2}$ 给 $A$ 与可逆矩阵 $P$，问 $P^TAP$ 或 $PAP^{-1}$ 的效果，把 $P$ 分解成初等矩阵之积再逐步施加变换；$\circled{3}$ 求初等矩阵的逆或行列式；$\circled{4}$ 问"下列矩阵中能经若干初等行变换化为给定矩阵的是哪个"，做法是\textbf{比较秩与行最简形}（行变换保持行空间，两个矩阵行等价当且仅当行最简形相同）；$\circled{5}$ 给"$A$ 为可逆矩阵，交换 $A$ 的第 1、2 行得 $B$"，问 $A^*$ 与 $B^*$ 的关系——用 $B=E_{12}A$，$B^*=(E_{12}A)^*=A^*E_{12}^*=A^*E_{12}$ 得"$B^*$ 是 $A^*$ 交换第 1、2 列"，同时 $\vert B\vert=-\vert A\vert$，故也可写"交换 $A^*$ 的第 1、2 列并变号"（两种表述等价，注意选项中给的是哪一种）。陷阱：把列变换当行变换导致左右颠倒；忘记对换矩阵行列式为 $-1$。

\end{callexam}

\subsection{矩阵的秩}

\subsubsection{矩阵的秩（定义、求法、秩的不等式与常用结论）$\bigstar$$\bigstar$$\bigstar$$\bigstar$$\bigstar$}

\textbf{$k$ 阶子式}：在 $m\times n$ 矩阵 $A$ 中任取 $k$ 行、$k$ 列，位于交叉位置的 $k^2$ 个元素（保持原序）构成的 $k$ 阶行列式，称为 $A$ 的一个 $k$ 阶子式。

\textbf{秩的定义}：若 $A$ 中存在一个 $r$ 阶子式不为零，而所有 $r+1$ 阶子式（若存在）全为零，则称 $r$ 为矩阵 $A$ 的\textbf{秩}，记作 $r(A)=r$。即\textbf{矩阵的秩 $=$ 最高阶非零子式的阶数}。

\textbf{秩的等价说法}（对 $m\times n$ 矩阵 $A$）

\begin{itemize}
\item $r(A)=r$ $\iff$ $A$ 至少有一个 $r$ 阶子式非零，且所有高于 $r$ 阶的子式全为零；
\item $r(A)\ge r$ $\iff$ $A$ 至少有一个 $r$ 阶子式非零；
\item $r(A)\le r$ $\iff$ $A$ 所有高于 $r$ 阶的子式全为零；
\item $r(A)=$ 行阶梯形中非零行的行数 $=$ 等价标准形中 $E_r$ 的阶数 $=$ $A$ 的行秩 $=$ $A$ 的列秩。
\end{itemize}
\textbf{求秩的方法}

\textbf{（1）初等变换法（具体矩阵首选）}：用初等行变换把 $A$ 化为行阶梯形，\textbf{非零行的行数就是秩}。含参数时先化行阶梯形，再按主元是否为零分类讨论。

\textbf{（2）定义法（抽象矩阵）}：找一个 $r$ 阶非零子式，再论证所有 $r+1$ 阶子式为零。

\textbf{（3）性质法}：直接用下面的常用结论推断。

\textbf{常用结论（第一组：借助定义与初等变换）}

\begin{enumerate}
\item $0\le r(A)\le\min\{m,n\}$。
\item 对 $n$ 阶方阵：$r(A)=n\iff\vert A\vert\ne0$（\textbf{满秩}）；$r(A)<n\iff\vert A\vert=0$（\textbf{降秩}）。
\item $r(kA)=r(A)\ (k\ne0)$。
\item $r(A^T)=r(A)$。
\item $r(A)=r(PA)=r(AQ)=r(PAQ)$，其中 $P,Q$ 可逆。\textbf{即：乘可逆矩阵，秩不变。}
\item 初等变换不改变矩阵的秩。
\item $r(A)=r(A^TA)=r(AA^T)$。
\item 同型矩阵 $A,B$ 等价 $\iff r(A)=r(B)$。
\item $A$ 行满秩（$r(A)=m$，$A$ 为 $m\times n$）$\iff$ $A$ 的行向量组线性无关 $\iff$ $A$ 的等价标准形为 $(E_m\ O)$；$A$ 列满秩（$r(A)=n$）$\iff$ $A$ 的列向量组线性无关 $\iff$ $Ax=0$ 只有零解 $\iff$ $A$ 的等价标准形为 $\begin{pmatrix}E_n\\O\end{pmatrix}$。
\end{enumerate}
\textbf{常用结论（第二组：秩的不等式）}

\[r(A\pm B)\le r(A)+r(B);\]

\[\max\{r(A),\,r(B)\}\le r(A,B)\le r(A)+r(B)\qquad(\text{拼矩阵：越拼秩越大});\]

\[r(AB)\le\min\{r(A),\,r(B)\}\qquad(\text{越乘秩越小});\]

\[r(AB)\ge r(A)+r(B)-n\qquad(A\text{ 为 }m\times n,\ B\text{ 为 }n\times s,\ \text{即 Sylvester 秩不等式});\]

\[AB=O\ \Longrightarrow\ r(A)+r(B)\le n\qquad(n\text{ 为 }A\text{ 的列数});\]

\[r(ABC)\ge r(AB)+r(BC)-r(B)\qquad(\text{Frobenius 秩不等式}).\]

\textbf{第三组：分块矩阵的秩}

\[r\begin{pmatrix}A&O\\O&B\end{pmatrix}=r(A)+r(B),\qquad r\begin{pmatrix}O&A\\B&O\end{pmatrix}=r(A)+r(B),\]

\[r\begin{pmatrix}A&C\\O&B\end{pmatrix}=r\begin{pmatrix}A&O\\C&B\end{pmatrix}\ge r(A)+r(B).\]

\textbf{第四组：伴随矩阵的秩}（见"伴随矩阵"一节）：$r(A)=n\Rightarrow r(A^*)=n$；$r(A)=n-1\Rightarrow r(A^*)=1$；$r(A)\le n-2\Rightarrow A^*=O$。

\textbf{第五组：由特征值看秩}：若 $n$ 阶矩阵 $A$ 可相似对角化（特别地，$A$ 为实对称矩阵），则 $r(A)$ 等于 $A$ 的非零特征值的个数（重根按个数累加）。

\textbf{例（初等变换求秩）}：设 $A=\begin{pmatrix}1&2&3\\4&5&6\\7&8&9\end{pmatrix}$，求 $r(A)$。

\[A\xrightarrow[r_2-4r_1]{r_3-7r_1}\begin{pmatrix}1&2&3\\0&-3&-6\\0&-6&-12\end{pmatrix}\xrightarrow{r_3-2r_2}\begin{pmatrix}1&2&3\\0&-3&-6\\0&0&0\end{pmatrix}.\]

非零行有 $2$ 行，故 $r(A)=2$。

\textbf{验证}：$\vert A\vert=1\times(45-48)-2\times(36-42)+3\times(32-35)=-3+12-9=0$，说明 $r(A)\le2$；而二阶子式 $\begin{vmatrix}1&2\\4&5\end{vmatrix}=-3\ne0$，说明 $r(A)\ge2$。故 $r(A)=2$，与初等变换结果一致。

\textbf{例（$AB=O$ 型）}：设 $A$ 为 $3$ 阶非零矩阵且 $A^2=O$，求 $r(A)$。

由 $A^2=A\cdot A=O$ 与 $AB=O\Rightarrow r(A)+r(B)\le n$（此处 $n=3$）得 $2r(A)\le3$，故 $r(A)\le1$；又 $A\ne O$ 得 $r(A)\ge1$。所以 $r(A)=1$。

\textbf{例（可逆时可消去）}：设 $A,B$ 均为 $n$ 阶方阵，$A$ 可逆且 $AB=O$。由 $A$ 可逆立得 $B=A^{-1}O=O$。这说明\textbf{只有在可逆的前提下，$AB=O$ 才能消去 $A$}；一般情形只能得到秩的不等式。

\textbf{例（分块矩阵的秩）}：设 $C=\begin{pmatrix}E_n&E_n\\E_n&E_n\end{pmatrix}$，求 $r(C)$。

用分块初等变换：$\begin{pmatrix}E&E\\E&E\end{pmatrix}\xrightarrow{\text{把第 1 块行的 }-1\text{ 倍加到第 2 块行}}\begin{pmatrix}E&E\\O&O\end{pmatrix}$，即 $C$ 等价于 $\begin{pmatrix}E_n&E_n\\O&O\end{pmatrix}$，其秩为 $n$。故 $r(C)=n$。

\begin{callkey}{必背}
\textbf{五句口诀}：越乘秩越小（$r(AB)\le\min\{r(A),r(B)\}$）；越拼秩越大（$r(A,B)\le r(A)+r(B)$）；乘转置不变秩（$r(A)=r(A^T)=r(A^TA)$）；乘可逆不变秩（$r(PAQ)=r(A)$）；$AB=O$ 则相邻下标之和有界（$r(A)+r(B)\le n$）。

\textbf{Sylvester 不等式} $r(AB)\ge r(A)+r(B)-n$ 给出秩的\textbf{下界}，与"越乘秩越小"给出的上界 $r(AB)\le\min\{r(A),r(B)\}$ 合起来才能夹出确定值；当一侧满秩时上界可达。

\end{callkey}

\begin{callwarn}{易错}
$\circled{1}$ $r(A+B)\le r(A)+r(B)$ 是\textbf{不等式}，取等号需额外条件。把"$\le$"写成"$=$"是高频错误。 $\circled{2}$ $AB=O$ 给出的上界是 $r(A)+r(B)\le n$，其中 $n$ \textbf{是 $A$ 的列数（也是 $B$ 的行数）}，不是 $A$ 的阶数。若 $A$ 为 $m\times n$、$B$ 为 $n\times s$，界限就是 $n$。 $\circled{3}$ "乘可逆矩阵秩不变"只对\textbf{可逆}矩阵成立；对不可逆矩阵，$r(PA)$ 可能小于 $r(A)$。 $\circled{4}$ $r(A)=r(B)$ \textbf{只是}同型矩阵等价的充要条件，对相似、合同都\textbf{不充分}：秩相同不意味着相似或合同。 $\circled{5}$ 由 $r(A)<n$ 只能推出 $Ax=0$ 有非零解、$\vert A\vert=0$、$A$ 不可逆，推不出 $r(A)$ 的具体数值。 $\circled{6}$ 秩为 $1$ 的矩阵必可写成 $\alpha\beta^T$ 的形式，这一分解在求幂、求特征值时极有用，不要只用来判断秩。

\end{callwarn}

\begin{callexam}{真题考法}
这是线代分值最高、出题最灵活的内容之一。\textbf{选择题}常考秩的不等式取等条件、$r(A^*)$ 的分类讨论、分块矩阵的秩。\textbf{填空题}常给含参数矩阵问秩，需要分类讨论。\textbf{证明题}常要求证 $r(A)+r(B)\le n$（由 $AB=O$）或 Sylvester 不等式。\textbf{解答题}中秩常与方程组解的结构、向量组的相关性联动。近年频繁出现"由矩阵等式（如 $A(A-A^*)=O$）反求 $r(A)$"的题型——解法是\textbf{用秩不等式夹出 $r(A)$ 的可能值，再逐一代回验证}，注意别漏解（典型答案形如"$1$ 或 $2$"）。陷阱：把 $r(AB)=r(A)$ 当成普遍成立；把 $r(A,B)$ 与 $r\begin{pmatrix}A\\B\end{pmatrix}$ 混为一谈。

\end{callexam}

\subsection{矩阵等价与分块矩阵}

\subsubsection{矩阵等价与分块矩阵及其运算 $\bigstar$$\bigstar$}

\textbf{矩阵等价}

\textbf{定义}：设 $A,B$ 为同型（都是 $m\times n$）矩阵，若存在 $m$ 阶可逆矩阵 $P$ 与 $n$ 阶可逆矩阵 $Q$，使得

\[B=PAQ,\]

则称 $A$ 与 $B$ \textbf{等价}。

\textbf{充要条件}：同型矩阵 $A,B$ 等价 $\iff r(A)=r(B)$。

\textbf{等价标准形}：秩为 $r$ 的 $m\times n$ 矩阵必等价于 $\begin{pmatrix}E_r&O\\O&O\end{pmatrix}$，且该标准形唯一。全体 $m\times n$ 矩阵按等价关系恰好分成 $\min\{m,n\}+1$ 个等价类（由 $r=0,1,\cdots,\min\{m,n\}$ 确定）。\textbf{秩是矩阵在等价关系下的完全不变量。}

\textbf{三种关系必须分清}

\begingroup\small\setlength{\tabcolsep}{4.0pt}
\begin{longtable}{@{}L{0.1736\linewidth}L{0.1736\linewidth}L{0.3906\linewidth}L{0.1736\linewidth}@{}}
\toprule
\textbf{关系} & \textbf{定义式} & \textbf{不变量} & \textbf{前提} \\
\midrule
\endhead
等价 & $B=PAQ$，$P,Q$ 可逆 & 秩 & 同型矩阵 \\
相似 & $B=P^{-1}AP$，$P$ 可逆 & 特征值、特征多项式、秩、迹、行列式 & 同阶方阵 \\
合同 & $B=P^TAP$，$P$ 可逆 & 正、负惯性指数（对实对称矩阵） & 同阶方阵 \\
\bottomrule
\end{longtable}
\endgroup

\textbf{分块矩阵及其运算}

用若干条横线和竖线把矩阵分成若干小块，每一小块称为子块，以子块为元素的矩阵称为\textbf{分块矩阵}。

\textbf{（1）加法与数乘}：分法相同的分块矩阵可直接相加；数乘对每个子块乘。

\textbf{（2）乘法}：当 $A$ 按列的分法与 $B$ 按行的分法一致时，

\[\begin{pmatrix}A_{11}&A_{12}\\A_{21}&A_{22}\end{pmatrix}\begin{pmatrix}B_{11}&B_{12}\\B_{21}&B_{22}\end{pmatrix}=\begin{pmatrix}A_{11}B_{11}+A_{12}B_{21}&A_{11}B_{12}+A_{12}B_{22}\\A_{21}B_{11}+A_{22}B_{21}&A_{21}B_{12}+A_{22}B_{22}\end{pmatrix}.\]

即\textbf{分块矩阵的乘法规则与普通矩阵乘法完全一样}，前提是每个子块都满足可乘条件。

\textbf{（3）转置}："大转小也转"——子块的位置按转置规则调换，每个子块\textbf{各自再取转置}：

\[\begin{pmatrix}A&B\\C&D\end{pmatrix}^T=\begin{pmatrix}A^T&C^T\\B^T&D^T\end{pmatrix}.\]

\textbf{（4）分块对角阵的幂}：

\[\begin{pmatrix}A&O\\O&B\end{pmatrix}^k=\begin{pmatrix}A^k&O\\O&B^k\end{pmatrix}.\]

\textbf{（5）分块矩阵的逆}（设 $A,B$ 可逆）：

\[\begin{pmatrix}A&O\\O&B\end{pmatrix}^{-1}=\begin{pmatrix}A^{-1}&O\\O&B^{-1}\end{pmatrix},\qquad \begin{pmatrix}O&A\\B&O\end{pmatrix}^{-1}=\begin{pmatrix}O&B^{-1}\\A^{-1}&O\end{pmatrix},\]

\[\begin{pmatrix}A&O\\C&B\end{pmatrix}^{-1}=\begin{pmatrix}A^{-1}&O\\-B^{-1}CA^{-1}&B^{-1}\end{pmatrix},\qquad \begin{pmatrix}A&C\\O&B\end{pmatrix}^{-1}=\begin{pmatrix}A^{-1}&-A^{-1}CB^{-1}\\O&B^{-1}\end{pmatrix}.\]

\textbf{口诀}：主对角型——主对角块"不用换序，直接添逆"，关键块写负号；副对角型——"先添逆，再换序"。统一记作：\textbf{分块副对角阵求逆时，两条对角线上的块都要互换位置}。

\textbf{（6）分块矩阵的行列式}（设 $A$ 为 $m$ 阶、$B$ 为 $n$ 阶）

\[\begin{vmatrix}A&O\\O&B\end{vmatrix}=\begin{vmatrix}A&C\\O&B\end{vmatrix}=\begin{vmatrix}A&O\\C&B\end{vmatrix}=\vert A\vert\vert B\vert,\qquad \begin{vmatrix}O&A\\B&O\end{vmatrix}=(-1)^{mn}\vert A\vert\vert B\vert .\]

更一般地（\textbf{打洞法/舒尔补}）：若 $A$ 可逆，则

\[\begin{vmatrix}A&B\\C&D\end{vmatrix}=\vert A\vert\cdot\vert D-CA^{-1}B\vert .\]

\textbf{（7）分块伴随}

\[\begin{pmatrix}A&O\\O&B\end{pmatrix}^*=\begin{pmatrix}\vert B\vert A^*&O\\O&\vert A\vert B^*\end{pmatrix}.\]

即分块对角阵求伴随时，\textbf{两个对角块"调换"后再各乘上对方块的行列式}。

\textbf{（8）分块矩阵的广义初等变换（左行右列）}

用可逆的分块矩阵左乘相当于做"分块行变换"，右乘相当于做"分块列变换"：

\[\begin{pmatrix}E&M\\O&E\end{pmatrix}\begin{pmatrix}A&B\\C&D\end{pmatrix}=\begin{pmatrix}A+MC&B+MD\\C&D\end{pmatrix},\]

\[\begin{pmatrix}A&B\\C&D\end{pmatrix}\begin{pmatrix}E&O\\M&E\end{pmatrix}=\begin{pmatrix}A+BM&B\\C+DM&D\end{pmatrix}.\]

\textbf{注意}：分块倍加时 $M$ 必须与相邻块\textbf{可乘}；分块倍乘时若乘的矩阵不可逆，则秩可能变小。

\textbf{例}：设 $A=\begin{pmatrix}1&2\\0&1\end{pmatrix}$，$B=E_2$，$C=\begin{pmatrix}0&1\\0&0\end{pmatrix}$，求 $M=\begin{pmatrix}A&O\\C&B\end{pmatrix}$ 的逆。

由公式 $\begin{pmatrix}A&O\\C&B\end{pmatrix}^{-1}=\begin{pmatrix}A^{-1}&O\\-B^{-1}CA^{-1}&B^{-1}\end{pmatrix}$。先算 $A^{-1}=\begin{pmatrix}1&-2\\0&1\end{pmatrix}$，$B^{-1}=B=E_2$，则

\[-B^{-1}CA^{-1}=-\begin{pmatrix}0&1\\0&0\end{pmatrix}\begin{pmatrix}1&-2\\0&1\end{pmatrix}=-\begin{pmatrix}0&1\\0&0\end{pmatrix}=\begin{pmatrix}0&-1\\0&0\end{pmatrix},\]

\[M^{-1}=\begin{pmatrix}1&-2&0&0\\0&1&0&0\\0&-1&1&0\\0&0&0&1\end{pmatrix}.\]

\textbf{验证}：$M$ 的左下 $2\times2$ 块为 $C=\begin{pmatrix}0&1\\0&0\end{pmatrix}$。验算 $M\cdot M^{-1}$ 的左下块：$C\cdot A^{-1}+B\cdot(-B^{-1}CA^{-1})=CA^{-1}-CA^{-1}=O$；左上块：$A\cdot A^{-1}+O\cdot C=E$；其余块显然正确。故 $M^{-1}$ 无误。

\begin{callwarn}{易错}
$\circled{1}$ 分块乘法时\textbf{左矩阵列的分法与右矩阵行的分法必须一致}，否则子块不可乘。 $\circled{2}$ 分块转置是"\textbf{大转小也转}"：既要把子块的位置转置，每个子块自身也要转置。漏掉子块自身的转置是常见错误。 $\circled{3}$ 分块副对角矩阵 $\begin{pmatrix}O&A\\B&O\end{pmatrix}$ \textbf{求幂没有} $\begin{pmatrix}A^k&O\\O&B^k\end{pmatrix}$ 这样的规律。 $\circled{4}$ $\begin{vmatrix}A&B\\C&D\end{vmatrix}\ne\vert AD-BC\vert$。分块矩阵的"行列式乘法"一般不成立，只有 $A$ 可逆时才有 $\begin{vmatrix}A&B\\C&D\end{vmatrix}=\vert A\vert\vert D-CA^{-1}B\vert$。而 $\begin{vmatrix}A&C\\O&B\end{vmatrix}=\vert A\vert\vert B\vert$ 要求 $A,B$ 都是\textbf{方阵}。 $\circled{5}$ 矩阵等价（$B=PAQ$）、相似（$B=P^{-1}AP$）、合同（$B=P^TAP$）三者不同，判定的不变量也不同。

\end{callwarn}

\begin{callexam}{真题考法}
矩阵等价多与秩联合出选择题："设 $A,B$ 为同型矩阵，则 $A$ 与 $B$ 等价是 $r(A)=r(B)$ 的什么条件"（充要）。分块矩阵的考法主要有：$\circled{1}$ 求分块对角阵（或副对角阵）的逆与伴随，四个选项差别在系数与位置；$\circled{2}$ 求分块矩阵的行列式（常与 $(-1)^{mn}$ 结合）；$\circled{3}$ 用分块矩阵配合广义初等变换证明秩的不等式；$\circled{4}$ 给分块矩阵 $\begin{pmatrix}A&O\\C&B\end{pmatrix}$ 问可逆的充要条件（$A,B$ 都可逆）。陷阱：分块求逆时负号位置放错；把 $\begin{vmatrix}A&B\\C&D\end{vmatrix}$ 错算成 $\vert AD-BC\vert$。

\end{callexam}

\section{矩阵方程的求解}

\subsection{三类矩阵方程}

\subsubsection{矩阵方程 $AX=B$、$XA=B$、$AXB=C$ 的求解 $\bigstar$$\bigstar$$\bigstar$$\bigstar$}

\textbf{核心思想}：矩阵方程就是"矩阵版的方程"，解法与代数方程完全平行——\textbf{能除就除（左乘或右乘逆矩阵），不能除就化行最简形}。关键是\textbf{左乘与右乘不能混}。

\textbf{（1）$AX=B$：左乘 $A^{-1}$}

若 $A$ 可逆，则 $X=A^{-1}B$。实际操作有两种：

\begin{itemize}
\item \textbf{初等行变换法}（推荐）：$(A\mid B)\xrightarrow{\ \text{行变换}\ }(E\mid A^{-1}B)$，右半边就是 $X$。这比先单独求 $A^{-1}$ 再乘 $B$ 快得多。
\item \textbf{先求逆再乘}：$X=A^{-1}B$。
\end{itemize}
若 $A$ \textbf{不可逆}（甚至是长方阵），则把 $X$ 按列分块：$X=(x_1,x_2,\cdots,x_s)$，$B=(b_1,b_2,\cdots,b_s)$，方程等价于 $s$ 个线性方程组 $Ax_j=b_j$，逐个用初等行变换解并判断有解性。

\textbf{（2）$XA=B$：右乘 $A^{-1}$}

若 $A$ 可逆，则 $X=BA^{-1}$。\textbf{注意是右乘，顺序不能颠倒成 $A^{-1}B$。} 操作上用\textbf{初等列变换}：$\begin{pmatrix}A\\B\end{pmatrix}\xrightarrow{\ \text{列变换}\ }\begin{pmatrix}E\\BA^{-1}\end{pmatrix}$，下边就是 $X$；也可以转置化成 $A^TX^T=B^T$ 的类型一处理。

\textbf{（3）$AXB=C$：解两次}

若 $A,B$ 都可逆，则 $X=A^{-1}CB^{-1}$。

\textbf{（4）含 $X$ 的复合方程：先提公因子}

\[AX+X=B\ \Longrightarrow\ (A+E)X=B\ \Longrightarrow\ X=(A+E)^{-1}B;\]

\[AX-X=B\ \Longrightarrow\ (A-E)X=B\ \Longrightarrow\ X=(A-E)^{-1}B;\]

\[XA+X=B\ \Longrightarrow\ X(A+E)=B\ \Longrightarrow\ X=B(A+E)^{-1};\]

\[X-XA=B\ \Longrightarrow\ X(E-A)=B\ \Longrightarrow\ X=B(E-A)^{-1}.\]

\textbf{提取公因子的方向必须统一}：$AX+X$ 中 $X$ 在右侧，故提\textbf{右}公因子得 $(A+E)X$；$XA+X$ 中 $X$ 在左侧，故提\textbf{左}公因子得 $X(A+E)$。

\textbf{（5）含 $A^*$、$A^{-1}$ 的方程}

先用 $A^*=\vert A\vert A^{-1}$ 把伴随换成逆，把方程化为只含 $A,B$ 及其逆的形式，再按上述套路求解。

\textbf{例}：解矩阵方程 $AX+X=B$，其中 $A=\begin{pmatrix}2&1\\0&3\end{pmatrix}$，$B=\begin{pmatrix}1&-1\\0&2\end{pmatrix}$。

先合并：$AX+X=(A+E)X=B$，其中 $A+E=\begin{pmatrix}3&1\\0&4\end{pmatrix}$。用初等行变换解 $(A+E\mid B)$：

\[\left(\begin{array}{cccc}3&1&1&-1\\0&4&0&2\end{array}\right)\xrightarrow{r_1\times\frac13}\left(\begin{array}{cccc}1&\frac13&\frac13&-\frac13\\0&1&0&\frac12\end{array}\right)\xrightarrow{r_1-\frac13r_2}\left(\begin{array}{cccc}1&0&\frac13&-\frac12\\0&1&0&\frac12\end{array}\right).\]

故 $X=\begin{pmatrix}\frac13&-\frac12\\0&\frac12\end{pmatrix}$。

\textbf{验证}：$(A+E)X=\begin{pmatrix}3&1\\0&4\end{pmatrix}\begin{pmatrix}\frac13&-\frac12\\0&\frac12\end{pmatrix}=\begin{pmatrix}1&-\frac32+\frac12\\0&2\end{pmatrix}=\begin{pmatrix}1&-1\\0&2\end{pmatrix}=B$，正确。

\textbf{例（$XA=B$ 型）}：解 $XA=B$，其中 $A=\begin{pmatrix}1&2\\0&1\end{pmatrix}$，$B=E_2$。

由 $A^{-1}=\begin{pmatrix}1&-2\\0&1\end{pmatrix}$，得 $X=BA^{-1}=\begin{pmatrix}1&-2\\0&1\end{pmatrix}$。

\textbf{验证}：$XA=\begin{pmatrix}1&-2\\0&1\end{pmatrix}\begin{pmatrix}1&2\\0&1\end{pmatrix}=\begin{pmatrix}1&0\\0&1\end{pmatrix}=B$，正确。

\textbf{例（抽象型：由 $AB=A+B$ 求 $B$ 与 $AB=BA$）}：设 $A,B$ 为 $n$ 阶矩阵，$AB=A+B$。

移项并凑因式：$AB-A-B=O$，两边加 $E$ 得 $AB-A-B+E=E$，即

\[(A-E)(B-E)=E .\]

由此立即得到三条结论：$\circled{1}$ $A-E$ 与 $B-E$ 都可逆，且 $(A-E)^{-1}=B-E$；$\circled{2}$ $B=E+(A-E)^{-1}$；$\circled{3}$ \textbf{$AB=BA$}：因为 $(B-E)(A-E)=E=(A-E)(B-E)$，展开得 $BA-B-A+E=AB-A-B+E$，故 $BA=AB$。

\begin{callkey}{必背}
\textbf{"$AX=B$ 还是 $XA=B$"的判断}：看 $X$ 在 $A$ 的\textbf{右边}还是\textbf{左边}。$X$ 在右边（$AX=B$）$\to$ 左乘 $A^{-1}$，$X=A^{-1}B$，用\textbf{行}变换；$X$ 在左边（$XA=B$）$\to$ 右乘 $A^{-1}$，$X=BA^{-1}$，用\textbf{列}变换。

\textbf{抽象矩阵等式的标准动作}：移项 $\to$ 两边加/减 $E$ $\to$ 分解成 $(A-\alpha E)(B-\beta E)=cE$ 的形式 $\to$ 读出可逆性与逆。

\end{callkey}

\begin{callwarn}{易错}
$\circled{1}$ \textbf{左右不能颠倒}：$AX=B$ 的解是 $A^{-1}B$，不是 $BA^{-1}$；$XA=B$ 的解是 $BA^{-1}$，不是 $A^{-1}B$。这是本节最大的坑。 $\circled{2}$ 合并同类项时\textbf{公因子方向要一致}：$AX+X=(A+E)X$，而 $XA+X=X(A+E)$，两者不能混。用错方向会得到完全错误的结果（除非 $A+E$ 与 $X$ 恰好可交换）。 $\circled{3}$ 解 $(A\mid B)$ 时\textbf{只能做行变换}；解 $\begin{pmatrix}A\\B\end{pmatrix}$ 时\textbf{只能做列变换}。 $\circled{4}$ 若 $A$ 不可逆，$X=A^{-1}B$ 型解法失效，必须按列分块化为线性方程组逐个解，并判断 $r(A)=r(A\mid b_j)$ 是否成立。 $\circled{5}$ 在 $AB=A+B$ 这类等式中，\textbf{不能直接"两边除以 $A$"}；正确做法是加减 $E$ 后因式分解。同时注意由此得到的可逆性结论是"$A-E$ 可逆"而\textbf{不是}"$A$ 可逆"。

\end{callwarn}

\begin{callexam}{真题考法}
解答题的高频题型，几乎每隔一两年出现一次。设问形式包括：$\circled{1}$ 直接给具体矩阵解 $AX=B$ 或 $XA=B$（考左右方向与变换类型）；$\circled{2}$ 给抽象等式如 $AB=A+B$，要求证明 $AB=BA$ 或求 $B$——做法是凑成 $(A-E)(B-E)=E$；$\circled{3}$ 给含 $A^*$ 或 $A^{-1}$ 的方程，先化简再求解；$\circled{4}$ 给"$A$ 可逆，$A^*X=B$"求 $X$，用 $A^*=\vert A\vert A^{-1}$ 转换。陷阱：$AX=X+B$ 之类移项后提公因子方向搞反；以及 $A$ 不可逆时强行求逆。

\end{callexam}

\subsection*{本章易错清单}
\addcontentsline{toc}{subsection}{本章易错清单}

\begin{enumerate}
\item \textbf{$(AB)^T=B^TA^T$ 与 $(AB)^{-1}=B^{-1}A^{-1}$ 都要反序}，写成 $A^TB^T$、$A^{-1}B^{-1}$ 都是错的。同理 $(AB)^*=B^*A^*$。
\item \textbf{消去律不成立}：$AB=AC$ 且 $A\ne O$ 推不出 $B=C$；只有 $A$ 可逆时才可两边左乘 $A^{-1}$。
\item \textbf{$AB=O$ 推不出 $A=O$ 或 $B=O$}；它能推出的是 $r(A)+r(B)\le n$（$n$ 为 $A$ 的列数）。
\item \textbf{$(A+B)^2\ne A^2+2AB+B^2$}，因为 $AB\ne BA$；正确的是 $A^2+AB+BA+B^2$。
\item \textbf{$(AB)^k\ne A^kB^k$}，只有 $AB=BA$ 时才可拆；二项式展开同理要求两项可交换。
\item \textbf{三个系数指数不能串}：$\vert kA\vert=k^n\vert A\vert$（指数 $n$）、$\vert A^*\vert=\vert A\vert^{n-1}$（指数 $n-1$）、$(kA)^*=k^{n-1}A^*$（指数 $n-1$）。
\item \textbf{伴随矩阵的位置}：$A^*$ 的第 $i$ 行第 $j$ 列是 $A_{ji}$，即"代数余子式矩阵的转置"；直接按 $A_{ij}$ 摆会写错。
\item \textbf{$r(A^*)$ 三分类的前提是 $n\ge2$}：$r(A)=n\Rightarrow n$；$r(A)=n-1\Rightarrow1$；$r(A)\le n-2\Rightarrow A^*=O$。别漏掉"$r(A)=n-1$ 时 $A^*\ne O$"这一半信息。
\item \textbf{初等变换"左行右列"}：行变换左乘、列变换右乘，且初等矩阵的阶数要匹配（行变换用 $m$ 阶、列变换用 $n$ 阶）。
\item \textbf{初等矩阵的逆}：$E_{ij}^{-1}=E_{ij}$，$E_i(k)^{-1}=E_i(1/k)$，$E_{ij}(k)^{-1}=E_{ij}(-k)$；不要写成 $E_{ij}(1/k)$。
\item \textbf{求逆只能用行变换}：$(A\mid E)\to(E\mid A^{-1})$ 全程只做行变换；混入列变换得到的结果没有意义。
\item \textbf{秩的不等式方向}：$r(A+B)\le r(A)+r(B)$，$r(AB)\le\min\{r(A),r(B)\}$，$AB=O\Rightarrow r(A)+r(B)\le n$，$r(AB)\ge r(A)+r(B)-n$。取等号都需额外条件，不能随意写等号。
\item \textbf{矩阵方程的方向}：$AX=B$ 的解是 $A^{-1}B$（左乘、行变换）；$XA=B$ 的解是 $BA^{-1}$（右乘、列变换）。提公因子时 $AX+X=(A+E)X$，而 $XA+X=X(A+E)$。
\item \textbf{分块矩阵}：转置要"大转小也转"；$\begin{vmatrix}A&B\\C&D\end{vmatrix}\ne\vert AD-BC\vert$；副对角分块求幂没有对角块各自求幂的规律。
\item \textbf{等价、相似、合同不能混}：等价看秩，相似看特征值（特征多项式），合同看正负惯性指数；同秩不等于相似或合同。
\end{enumerate}
\chapter{向量}
\markboth{向量}{向量}
\subsection*{本章考纲要求}
\addcontentsline{toc}{subsection}{本章考纲要求}

\begin{quote}
\textbf{考试内容}

\begin{enumerate}
\item 向量的概念
\item 向量的线性组合与线性表示
\item 向量组的线性相关与线性无关
\item 向量组的极大线性无关组
\item 等价向量组
\item 向量组的秩
\item 向量组的秩与矩阵的秩之间的关系
\item 向量的内积
\item 线性无关向量组的正交规范化方法
\end{enumerate}
\textbf{考试要求}

\begin{enumerate}
\item 理解 n 维向量、向量的线性组合与线性表示的概念。
\item 理解向量组线性相关、线性无关的概念，掌握向量组线性相关、线性无关的有关性质及判别法。
\item 理解向量组的极大线性无关组和向量组的秩的概念，会求向量组的极大线性无关组及秩。
\item 理解向量组等价的概念，理解矩阵的秩与其行(列)向量组的秩之间的关系。
\item 了解内积的概念，掌握线性无关向量组正交规范化的施密特(Schmidt)方法。
\end{enumerate}
\end{quote}
\subsection*{本章星级总表}
\addcontentsline{toc}{subsection}{本章星级总表}

\begingroup\small\setlength{\tabcolsep}{4.0pt}
\begin{longtable}{@{}L{0.4240\linewidth}L{0.1175\linewidth}L{0.1175\linewidth}L{0.1175\linewidth}L{0.1175\linewidth}@{}}
\toprule
\textbf{考点} & \textbf{星级} & \textbf{考纲要求} & \textbf{难度} & \textbf{常见题型} \\
\midrule
\endhead
n 维向量、线性组合与线性表示 & $\bigstar$$\bigstar$$\bigstar$$\bigstar$ & 理解 & 中等 & 选择、解答 \\
向量组线性相关与线性无关的概念、性质与判别法 & $\bigstar$$\bigstar$$\bigstar$$\bigstar$$\bigstar$ & 理解/掌握 & 难 & 选择、证明 \\
线性相关、线性无关的证明（抽象向量组） & $\bigstar$$\bigstar$$\bigstar$$\bigstar$$\bigstar$ & 掌握 & 难 & 证明 \\
极大线性无关组与向量组的秩 & $\bigstar$$\bigstar$$\bigstar$$\bigstar$ & 了解/会 & 中等 & 填空、解答 \\
等价向量组、矩阵的秩与向量组的秩的关系 & $\bigstar$$\bigstar$$\bigstar$ & 了解 & 中等 & 选择 \\
内积与施密特正交规范化 & $\bigstar$$\bigstar$ & 了解/掌握 & 中等 & 填空 \\
\bottomrule
\end{longtable}
\endgroup

\paragraph{学这一章要建立的两个直觉}\mbox{}{\par}

第一个直觉是"\textbf{向量组 = 一个整体}"。相关性、秩、极大无关组，谈的都不是某一个向量，而是若干向量凑在一起时呈现出的整体性质。做题时只要看到"向量组"三个字，就应当自动把它按列排成一个矩阵，接下来所有问题都能翻译成矩阵问题——这是本章与上一章（矩阵）之间最重要的通道。

第二个直觉是"\textbf{相关 = 有冗余，无关 = 无冗余}"。这个直觉能推出几乎所有判别法而不需要死记：有冗余的组，去掉冗余向量后表示能力不变（所以至少有一个向量能被其余的表示）；无冗余的组，任何一个向量去掉都会缩小表示能力（所以每个向量都不能被其余的表示）；冗余的组加上更多向量还是有冗余（部分相关则整体相关）；无冗余的组取一部分仍然无冗余（整体无关则部分无关）。把这些结论理解成"冗余"这一件事的不同侧面，比逐条背诵可靠得多。

\section{向量的概念与线性运算}

本章是线性代数的"语言基础"：后面所有关于方程组解的结构、特征向量、二次型的讨论，都要用向量组的语言来表述。数二试卷中，向量这一章的独立分值通常不大，但它是\textbf{证明题的策源地}——"抽象向量组线性无关的证明"几乎每年都以某种变形出现在线代大题的第一问里。

本章只讨论\textbf{实向量}。统考范围内向量默认指\textbf{列向量}，写成 $(a_1,a_2,\cdots,a_n)^T$ 的形式。

\subsection{向量的线性运算与内积}

\subsubsection{n 维向量、线性组合与线性表示 $\bigstar$$\bigstar$$\bigstar$$\bigstar$}

\textbf{定义（n 维向量）}：$n$ 个实数组成的有序数组

\[\alpha=\begin{pmatrix} a_1 \\ a_2 \\ \vdots \\ a_n \end{pmatrix} \quad \text{或} \quad (a_1,a_2,\cdots,a_n)^T\]

称为一个 \textbf{$n$ 维向量}，$a_i$ 称为 $\alpha$ 的第 $i$ 个分量。分量全为 $0$ 的向量称为\textbf{零向量}，记作 $\mathbf{0}$。若题目没有指明是行向量还是列向量，\textbf{统考一律按列向量理解}。

\textbf{向量的线性运算}：向量本质上是 $n\times 1$ 矩阵，加法和数乘完全照搬矩阵的运算规律：

\[\alpha+\beta=\begin{pmatrix} a_1+b_1 \\ a_2+b_2 \\ \vdots \\ a_n+b_n \end{pmatrix},\qquad k\alpha=\begin{pmatrix} ka_1 \\ ka_2 \\ \vdots \\ ka_n \end{pmatrix}\]

\textbf{定义（线性组合与线性表示）}：设 $\alpha_1,\alpha_2,\cdots,\alpha_s$ 是 $n$ 维向量，$k_1,k_2,\cdots,k_s$ 是数，则称

\[k_1\alpha_1+k_2\alpha_2+\cdots+k_s\alpha_s\]

为向量组 $\alpha_1,\alpha_2,\cdots,\alpha_s$ 的一个\textbf{线性组合}；若向量 $\beta$ 能写成这样的线性组合，即存在 $k_1,\cdots,k_s$ 使 $\beta=k_1\alpha_1+\cdots+k_s\alpha_s$，则称 $\beta$ \textbf{可由} $\alpha_1,\alpha_2,\cdots,\alpha_s$ \textbf{线性表示}。

\begin{callkey}{线性表示的三条等价判据}
设 $A=(\alpha_1,\alpha_2,\cdots,\alpha_s)$ 是以这些向量为列拼成的矩阵，则下面三个命题互相等价：

\[\beta \text{ 可由 } \alpha_1,\cdots,\alpha_s \text{ 线性表示}
\iff (\alpha_1,\cdots,\alpha_s)x=\beta \text{ 有解}
\iff r(\alpha_1,\cdots,\alpha_s)=r(\alpha_1,\cdots,\alpha_s,\beta)\]

若 $r(A)=r(\bar A)=r<s$（即方程组有无穷多解），则\textbf{表示方法不唯一}，通解中的自由参数就是所有可能的表示系数；若 $r(A)=r(\bar A)=s$，则\textbf{表示方法唯一}。

\textbf{把这条判据翻译成四种等价说法}（做题时要能在它们之间自由切换）：

\begingroup\small\setlength{\tabcolsep}{4.0pt}
\begin{longtable}{@{}L{0.2351\linewidth}L{0.2351\linewidth}L{0.4584\linewidth}@{}}
\toprule
\textbf{说法} & \textbf{数学表达} & \textbf{典型用法} \\
\midrule
\endhead
方程组语言 & $(\alpha_1,\cdots,\alpha_s)x=\beta$ 有解 & 求具体的表示系数 \\
秩语言 & $r(A)=r(\bar A)$ & 抽象判定、证明确实能表示 \\
向量语言 & $\beta$ 属于 $\alpha_1,\cdots,\alpha_s$ 张成的空间 & 概念辨析题 \\
秩的增量语言 & $r(\bar A)\ne r(A)+1$，即 $r(\bar A)=r(A)$ & 判定"不能表示"最快 \\
\bottomrule
\end{longtable}
\endgroup

其中最后一条最好用：\textbf{"不能表示"当且仅当添上 $\beta$ 之后秩恰好增加 $1$}（$r(\bar A)=r(A)+1$）。证明题里反复使用这一条。

\end{callkey}

\textbf{三条最常用的"送分"结论}：

\begin{itemize}
\item 零向量 $\mathbf{0}$ 可由\textbf{任何}向量组线性表示（取系数全为 $0$ 即可）。
\item 任一向量都可由\textbf{包含它自己的}向量组线性表示。
\item 若 $r(A)<s$，即 $\alpha_1,\cdots,\alpha_s$ 线性相关，则 $\beta$ 一旦能被表示，就有\textbf{无穷多种}表示方式。
\end{itemize}
\textbf{解题套路（判定 + 求系数）}：

\begin{enumerate}
\item 把向量组按列拼成 $A$，把 $\beta$ 添在最后一列得增广矩阵 $\bar A=(A\mid\beta)$；
\item 对 $\bar A$ 作\textbf{初等行变换}化为行阶梯形，比较 $r(A)$ 与 $r(\bar A)$：相等则有解，$r(\bar A)=r(A)+1$ 则不能表示；
\item 若有解，继续化为\textbf{行最简形}，回代即得表示系数。
\end{enumerate}
\textbf{例}：设 $\alpha_1=(1,2)^T,\alpha_2=(-1,1)^T,\beta=(3,3)^T$，判断 $\beta$ 能否由 $\alpha_1,\alpha_2$ 线性表示。

解方程组 $\begin{pmatrix} 1 & -1 \\ 2 & 1 \end{pmatrix}\begin{pmatrix} x_1 \\ x_2 \end{pmatrix}=\begin{pmatrix} 3 \\ 3 \end{pmatrix}$，由第一式 $x_1-x_2=3$，第二式 $2x_1+x_2=3$，解得 $x_1=2,x_2=-1$。故

\[\beta=2\alpha_1-\alpha_2\]

且因 $r(A)=2=s$，表示方式唯一。

\begin{callwarn}{易错点}
"$\beta$ 可由 $\alpha_1,\cdots,\alpha_s$ 表示"与"$\alpha_1,\cdots,\alpha_s$ 线性相关"是\textbf{两件事}：前者说的是组\textbf{外}的向量与这组向量的关系，后者说的是组\textbf{内}向量之间的关系。只有当 $\beta$ 能由线性相关的向量组表示时，才出现表示不唯一。

另一个高频陷阱：比较 $r(A)$ 与 $r(\bar A)$ 时必须用\textbf{行变换}（行阶梯形），因为行变换不改变线性相关性；如果对 $A$ 作列变换，就破坏了"列的线性关系"，结论全错。

\end{callwarn}

\begin{callexam}{真题考法}
这一考点常以两种面孔出现：$\circled{1}$ 给具体的数字向量，问"能否表示 / 求出表示式"，属于填空或选择题；$\circled{2}$ \textbf{抽象形式}："已知 $r(\alpha_1,\alpha_2,\alpha_3)=r(\alpha_1,\alpha_2,\alpha_3,\beta)=3$，问 $\beta$ 与 $\alpha_1,\alpha_2,\alpha_3$ 的关系"——考的就是秩判据。陷阱在于把 $r(\alpha_1,\alpha_2,\alpha_3,\beta)=3$ 与"$\alpha_1,\alpha_2,\alpha_3,\beta$ 线性相关"混为一谈（$4$ 个向量秩为 $3$，确实相关，但由此\textbf{能}推出 $\beta$ 可由 $\alpha_1,\alpha_2,\alpha_3$ 表示——这一步推理必须写清楚，因为 $\alpha_1,\alpha_2,\alpha_3$ 本身已构成极大无关组）。

\end{callexam}

\paragraph{施密特正交化的记忆方法与验算习惯}\mbox{}{\par}

公式看起来长，但结构非常统一：\textbf{只有第一项是"原向量"，后面每一项都是"原向量在已经造好的正交向量上的投影"}。于是记忆口诀是"\textbf{原向量减去它在所有前辈上的投影}"。

具体到第三个向量：$\beta_3=\alpha_3-(\alpha_3$ 在 $\beta_1$ 上的投影$)-(\alpha_3$ 在 $\beta_2$ 上的投影$)$，而投影的公式是 $\dfrac{(\alpha_3,\beta_k)}{(\beta_k,\beta_k)}\beta_k$。把"投影"这两个字换成公式，就得到标准写法。分子是内积、分母是该正交向量与自己的内积（即模长的平方），这一点绝不能记反。

\textbf{良好的验算习惯}：正交化只做到第二步时就有可以检查的地方——

\begin{itemize}
\item $\beta_1$ 与 $\beta_2$ 的内积应当为 $0$：$(\beta_2,\beta_1)=(\alpha_2,\beta_1)-\dfrac{(\alpha_2,\beta_1)}{(\beta_1,\beta_1)}(\beta_1,\beta_1)=0$；
\item 若算出 $(\beta_2,\beta_1)\ne0$，一定是分母用了 $(\alpha_1,\alpha_1)$ 或分子用错了向量。
\end{itemize}
全部做完后，把三个 (或更多) 结果两两作内积检查一遍，$\eta_i^T\eta_j=0$（$i\ne j$）且 $\eta_i^T\eta_i=1$。这是本节唯一可靠的检验手段，花不了两分钟，却能避免整道大题失分。

\textbf{为什么要正交化}：正交向量组一定是线性无关的（设 $\sum k_i\eta_i=\mathbf{0}$，两端与 $\eta_j$ 作内积即得 $k_j=0$），反过来线性无关的向量组未必正交。正交化的作用是在\textbf{不改变张成空间}的前提下，把一组"歪斜"的基换成一组"互相垂直"的基，从而让计算（投影、坐标、二次型的标准形）变得简单。第五章"求正交矩阵 $Q$ 使 $Q^{-1}AQ=\Lambda$"时的正交化与单位化，用的就是这里的方法。

\textbf{一个必须掌握的衔接点}：实对称矩阵的不同特征值对应的特征向量\textbf{自动正交}，不需要正交化；只有\textbf{同一个特征值}对应的多个线性无关特征向量之间才需要做施密特正交化。考试中漏做或错做这一步都很常见，判断依据就是"这几个特征向量是不是对应同一个特征值"。

\subsubsection{内积与施密特正交规范化 $\bigstar$$\bigstar$}

\textbf{定义（内积）}：设 $\alpha=(a_1,a_2,\cdots,a_n)^T$，$\beta=(b_1,b_2,\cdots,b_n)^T$，称

\[(\alpha,\beta)=\alpha^T\beta=\sum_{i=1}^{n}a_ib_i\]

为 $\alpha$ 与 $\beta$ 的\textbf{内积}。内积满足：$(\alpha,\beta)=(\beta,\alpha)$；$(k\alpha,\beta)=k(\alpha,\beta)$；$(\alpha+\beta,\gamma)=(\alpha,\gamma)+(\beta,\gamma)$；$(\alpha,\alpha)\ge 0$，且 $(\alpha,\alpha)=0\iff\alpha=\mathbf{0}$。

\textbf{由内积派生的三个概念}：

\begin{itemize}
\item \textbf{模（长度）}：$\Vert\alpha\Vert=\sqrt{(\alpha,\alpha)}=\sqrt{a_1^2+a_2^2+\cdots+a_n^2}$；$\Vert\alpha\Vert=1$ 时称 $\alpha$ 为\textbf{单位向量}，把非零向量 $\alpha$ 变成 $\dfrac{\alpha}{\Vert\alpha\Vert}$ 的过程叫\textbf{单位化}（标准化）。
\item \textbf{正交}：$(\alpha,\beta)=0$ 时称 $\alpha$ 与 $\beta$ \textbf{正交}，记作 $\alpha\perp\beta$。规定零向量与任何向量都正交。
\item \textbf{柯西不等式}：$(\alpha,\beta)^2\le(\alpha,\alpha)(\beta,\beta)$，即 $(\alpha,\beta)^2\le\Vert\alpha\Vert^2\Vert\beta\Vert^2$。由此可定义夹角 $\cos\theta=\dfrac{(\alpha,\beta)}{\Vert\alpha\Vert\,\Vert\beta\Vert}$。
\end{itemize}
\begin{callkey}{施密特（Schmidt）正交规范化方法}
设 $\alpha_1,\alpha_2,\alpha_3$ 线性无关，要把它变成\textbf{两两正交}的 $\beta_1,\beta_2,\beta_3$（正交化），再各自除以模长得到标准正交组 $\eta_1,\eta_2,\eta_3$（单位化）。

\textbf{第一步：正交化}

\[\beta_1=\alpha_1,\qquad
\beta_2=\alpha_2-\frac{(\alpha_2,\beta_1)}{(\beta_1,\beta_1)}\beta_1,\qquad
\beta_3=\alpha_3-\frac{(\alpha_3,\beta_1)}{(\beta_1,\beta_1)}\beta_1-\frac{(\alpha_3,\beta_2)}{(\beta_2,\beta_2)}\beta_2\]

\textbf{第二步：单位化}

\[\eta_1=\frac{\beta_1}{\Vert\beta_1\Vert},\qquad
\eta_2=\frac{\beta_2}{\Vert\beta_2\Vert},\qquad
\eta_3=\frac{\beta_3}{\Vert\beta_3\Vert}\]

则 $\eta_1,\eta_2,\eta_3$ 两两正交且都是单位向量，即\textbf{标准正交组}。

\textbf{适用前提}：原向量组必须线性无关。若向量组本身线性相关，正交化过程中会出现某个 $\beta_k=\mathbf{0}$（无法单位化），这不是公式出错，而是向量组"秩不够"的警报。

\end{callkey}

\textbf{例}：把 $\alpha_1=(1,1,0)^T,\alpha_2=(1,0,1)^T,\alpha_3=(0,1,1)^T$ 正交规范化。

\[\beta_1=\alpha_1=(1,1,0)^T,\qquad (\beta_1,\beta_1)=2,\quad (\alpha_2,\beta_1)=1\]

\[\beta_2=\alpha_2-\frac{1}{2}\beta_1=(1,0,1)^T-\left(\tfrac12,\tfrac12,0\right)^T=\left(\tfrac12,-\tfrac12,1\right)^T\]

\[(\alpha_3,\beta_1)=1,\qquad (\alpha_3,\beta_2)=\tfrac12\cdot 0-\tfrac12\cdot 1+1\cdot 1=\tfrac12,\qquad (\beta_2,\beta_2)=\tfrac14+\tfrac14+1=\tfrac32\]

\[\beta_3=\alpha_3-\frac{1}{2}\beta_1-\frac{1/2}{3/2}\beta_2=(0,1,1)^T-\left(\tfrac12,\tfrac12,0\right)^T-\tfrac13\left(\tfrac12,-\tfrac12,1\right)^T=\left(-\tfrac23,\tfrac23,\tfrac23\right)^T\]

单位化后得

\[\eta_1=\frac{1}{\sqrt2}\begin{pmatrix} 1 \\ 1 \\ 0 \end{pmatrix},\quad
\eta_2=\frac{1}{\sqrt6}\begin{pmatrix} 1 \\ -1 \\ 2 \end{pmatrix},\quad
\eta_3=\frac{1}{\sqrt3}\begin{pmatrix} -1 \\ 1 \\ 1 \end{pmatrix}\]

\textbf{验算}：$\eta_1^T\eta_2=\frac{1}{\sqrt{12}}(1-1+0)=0$，$\eta_1^T\eta_3=\frac{1}{\sqrt6}(-1+1+0)=0$，$\eta_2^T\eta_3=\frac{1}{\sqrt{18}}(-1-1+2)=0$。三向量两两正交且模长皆为 $1$。

\begin{callwarn}{易错点}
正交化公式里分母是 $(\beta_k,\beta_k)$，\textbf{不是} $(\alpha_k,\alpha_k)$，也不是 $\Vert\beta_k\Vert$ 的一次方。这是丢分率最高的一处。

单位化时忘记除以模长、或者把模长算成未开方的 $(\alpha,\alpha)$，都会让答案"差一个系数"。

正交化\textbf{必须按顺序}用已经生成的 $\beta_1,\beta_2,\cdots$ 逐个作投影减去；漏减某一项后，$\beta_3$ 与 $\beta_1$ 未必正交，务必用内积回代验算。

\end{callwarn}

\begin{callexam}{真题考法}
数二对施密特的考查通常不是"给三个具体向量让你正交化"（那是数一的老题型），而是\textbf{披着二次型外衣}出现：求正交矩阵 $Q$ 化二次型为标准形时，对同一特征值的多个线性无关特征向量必须先正交化再单位化；不同特征值的特征向量已经自动正交，无需再做。这一区别是每年的固定采分点。

\end{callexam}

\section{向量组的线性表示与线性相关性}

这一板块是本手册线代部分\textbf{分值密度最高}的地方：概念辨析题（选择）与抽象证明题（解答）都从这里出。两条主线务必分清——\textbf{组间关系（能否互相表示）} 与 \textbf{组内关系（相关还是无关）}。

\subsection{线性相关与线性无关}

\subsubsection{向量组线性相关与线性无关的概念、性质与判别法 $\bigstar$$\bigstar$$\bigstar$$\bigstar$$\bigstar$}

\textbf{定义}：设 $\alpha_1,\alpha_2,\cdots,\alpha_s$ 是 $n$ 维向量组。

\begin{itemize}
\item 若存在\textbf{不全为零}的数 $k_1,k_2,\cdots,k_s$，使 $k_1\alpha_1+k_2\alpha_2+\cdots+k_s\alpha_s=\mathbf{0}$，则称该向量组\textbf{线性相关}；
\item 若上式成立\textbf{当且仅当} $k_1=k_2=\cdots=k_s=0$，则称该向量组\textbf{线性无关}。
\end{itemize}
\textbf{同义改写（做题时按需要切换）}：以 $\alpha_1,\cdots,\alpha_s$ 为列构成矩阵 $A=(\alpha_1,\alpha_2,\cdots,\alpha_s)$，则

\[\alpha_1,\cdots,\alpha_s \text{ 线性相关}\iff Ax=\mathbf{0}\ \text{有非零解}\iff r(A)<s\]

\[\alpha_1,\cdots,\alpha_s \text{ 线性无关}\iff Ax=\mathbf{0}\ \text{只有零解}\iff r(A)=s\]

\textbf{组内关系与组间关系的桥梁}：当 $s\ge 2$ 时，

\[\alpha_1,\cdots,\alpha_s \text{ 线性相关}\iff \text{其中至少有一个向量可由其余 } s-1 \text{ 个向量线性表示}\]

\[\alpha_1,\cdots,\alpha_s \text{ 线性无关}\iff \text{其中任一向量都不能由其余 } s-1 \text{ 个向量线性表示}\]

\textbf{注意措辞的差别}：相关时是"\textbf{至少有一个}能表示"（不是每一个都能表示）；无关时是"\textbf{任意一个都不能}"。

\textbf{例}：判断 $\alpha_1=(1,2,3)^T,\alpha_2=(2,4,6)^T,\alpha_3=(1,1,1)^T$ 的线性相关性。

因 $\alpha_2=2\alpha_1$，存在不全为零的系数使 $2\alpha_1+(-1)\alpha_2+0\cdot\alpha_3=\mathbf{0}$，故线性相关。用秩判据：以它们为列作初等行变换，

\[A=\begin{pmatrix} 1 & 2 & 1 \\ 2 & 4 & 1 \\ 3 & 6 & 1 \end{pmatrix}\to\begin{pmatrix} 1 & 2 & 1 \\ 0 & 0 & -1 \\ 0 & 0 & -2 \end{pmatrix}\to\begin{pmatrix} 1 & 2 & 1 \\ 0 & 0 & 1 \\ 0 & 0 & 0 \end{pmatrix}\]

$r(A)=2<3=s$，故线性相关。

\paragraph{判别法一览与选择原则}\mbox{}{\par}

把上面所有判据集中成一张"工具表"，考试时按题目给的信息选择最快的路径。

\begingroup\small\setlength{\tabcolsep}{4.0pt}
\begin{longtable}{@{}L{0.4121\linewidth}L{0.2347\linewidth}L{0.2817\linewidth}@{}}
\toprule
\textbf{题目给出的信息} & \textbf{首选判据} & \textbf{理由} \\
\midrule
\endhead
具体的数字向量 & 秩判据（初等行变换） & 一定能算，且顺带得到极大无关组 \\
具体向量且个数等于维数 & 行列式判据 & 一个行列式解决问题 \\
抽象向量 + 线性组合关系 & 定义法或矩阵法 & 见下一考点的四招 \\
已知"某组无关"要推出"新组无关" & 矩阵法（$B=AC$） & 把抽象问题数字化 \\
判断"某个向量能否被表示" & 秩的增量判据 $r(A\mid\beta)=r(A)$ & 一步到位 \\
只知道秩的信息（如 $r(\alpha_1,\alpha_2,\alpha_3)=2$） & 秩与个数的关系 & 秩小于个数即相关 \\
\bottomrule
\end{longtable}
\endgroup

\textbf{从"秩"的角度重新理解相关与无关}。向量组的秩是这组向量"真实占用的维数"，向量个数是"名义上的维数"。于是

\begin{itemize}
\item 相关 $\iff$ 真实维数小于名义个数 $\iff$ 有向量是"多余"的，去掉它不损失任何表示能力；
\item 无关 $\iff$ 真实维数等于名义个数 $\iff$ 每一个向量都"各司其职"，去掉任何一个都会缩小张成的空间。
\end{itemize}
这个视角能解释很多结论。例如"部分相关则整体相关"：部分组已经有多余向量，整体组自然更多余。"整体无关则部分无关"：整体组每个向量都不可替代，抽出一部分当然也不可替代。"以少表多，多必相关"：被表示的向量组的所有表示能力都来自那少数几个向量，它的真实维数不会超过表出者的个数，所以名义个数多于表出者个数时必然有多余向量。

\textbf{从"方程组"的角度重新理解相关与无关}。$k_1\alpha_1+\cdots+k_s\alpha_s=\mathbf{0}$ 本身就是以 $k_1,\cdots,k_s$ 为未知量的齐次线性方程组，写成矩阵形式就是 $Ax=\mathbf{0}$，其中 $A$ 的\textbf{列}是 $\alpha_i$。于是"线性相关"就是"这个齐次方程组有非零解"，"线性无关"就是"只有零解"。这一翻译的好处是立刻可以用上方程组的全部工具：自由变量个数 $=s-r(A)$ 就是"有多少种本质不同的相关方式"。

\textbf{一个常考的几何直观}（帮助记结论，不必写进答案）：在平面里两个向量相关就是共线，在空间里三个向量相关就是共面；$n$ 维空间里最多只能有 $n$ 个互不"共超平面"的向量，这就是"向量个数大于维数必相关"的来源。

\subsubsection{线性相关、线性无关的证明（抽象向量组） $\bigstar$$\bigstar$$\bigstar$$\bigstar$$\bigstar$}

这是\textbf{数二线代压轴证明题的核心}。题目不给数字，只给"$\alpha_1,\cdots,\alpha_s$ 线性无关"和一组由它们线性组合出的新向量 $\beta_1,\cdots,\beta_t$，要求判断或证明新向量组的相关性。全部套路只有三招。

\begin{callkey}{三大招数（按优先级使用）}
\textbf{第 1 招：定义法（"设 $k_1\beta_1+\cdots+k_t\beta_t=\mathbf{0}$"）}

按定义写出 $k_1\beta_1+\cdots+k_t\beta_t=\mathbf{0}$，把每个 $\beta_j$ 用它关于 $\alpha_i$ 的表达式代入，\textbf{按 $\alpha_1,\cdots,\alpha_s$ 重新整理合并}，利用 $\alpha_1,\cdots,\alpha_s$ 线性无关得到关于 $k$ 的齐次方程组，解出 $k_1=\cdots=k_t=0$（或发现非零解）。

\textbf{第 2 招：矩阵法（"$B=AC$"）}

把 $\beta$ 组写成 $(\beta_1,\cdots,\beta_t)=(\alpha_1,\cdots,\alpha_s)C$，其中 $C$ 是 $s\times t$ 的\textbf{表示系数矩阵}。于是

\[\beta \text{ 组线性无关}\iff r(C)=t,\qquad \beta \text{ 组线性相关}\iff r(C)<t\]

特别地：当 $(\alpha_1,\cdots,\alpha_s)$ 列满秩（线性无关）时，$r(B)=r(C)$，于是\textbf{$C$ 的作用完全可以代替 $B$} 来判断。这一招把抽象向量问题变成了"求矩阵 $C$ 的行列式/秩"的具体计算。

\textbf{第 3 招："乘"与"重组"（强化技巧）}

在 $k_1\beta_1+\cdots+k_t\beta_t=\mathbf{0}$ 两端左乘某个矩阵 $A$（常用题目条件中现成的 $A^k$、$A^{k-1}$），或把已知关系 $A\alpha_i=\cdots$ 代入并重新分组，达到"缩短等式、逐项逼出 $k_i=0$"的目的。

\textbf{辅助手段}：反证法（假设相关导出矛盾）。若向量个数 $=$ 向量维数，可直接算行列式。

\end{callkey}

\textbf{例 1（第 2 招的典型）}：设 $\alpha_1,\alpha_2,\alpha_3$ 线性无关，$\beta_1=\alpha_1+\alpha_2$，$\beta_2=\alpha_2+\alpha_3$，$\beta_3=\alpha_3+\alpha_1$，判断 $\beta_1,\beta_2,\beta_3$ 的相关性。

写成矩阵形式：

\[(\beta_1,\beta_2,\beta_3)=(\alpha_1,\alpha_2,\alpha_3)\begin{pmatrix} 1 & 0 & 1 \\ 1 & 1 & 0 \\ 0 & 1 & 1 \end{pmatrix}=(\alpha_1,\alpha_2,\alpha_3)C\]

\[\vert C\vert=\begin{vmatrix} 1 & 0 & 1 \\ 1 & 1 & 0 \\ 0 & 1 & 1 \end{vmatrix}\]

按第一行展开：$1\cdot(1\cdot1-0\cdot1)+1\cdot(1\cdot1-1\cdot0)=1+1=2\ne 0$，故 $r(C)=3$，$\beta_1,\beta_2,\beta_3$ \textbf{线性无关}。

若改用定义法也可：设 $k_1\beta_1+k_2\beta_2+k_3\beta_3=\mathbf{0}$，代入得

\[(k_1+k_3)\alpha_1+(k_1+k_2)\alpha_2+(k_2+k_3)\alpha_3=\mathbf{0}\]

由 $\alpha_1,\alpha_2,\alpha_3$ 线性无关，系数必须全为 $0$：

\[\begin{cases} k_1+k_3=0 \\ k_1+k_2=0 \\ k_2+k_3=0 \end{cases}\]

三式相加得 $2(k_1+k_2+k_3)=0$，再结合第一式得 $k_2=0$，进而 $k_1=0,k_3=0$。故 $\beta_1,\beta_2,\beta_3$ 线性无关。

\textbf{例 2（第 3 招的典型）}：设 $A$ 为 $n$ 阶矩阵，$\alpha$ 满足 $A^{m}\alpha=\mathbf{0}$ 而 $A^{m-1}\alpha\ne\mathbf{0}$，证明 $\alpha,A\alpha,\cdots,A^{m-1}\alpha$ 线性无关。

设 $k_0\alpha+k_1A\alpha+\cdots+k_{m-1}A^{m-1}\alpha=\mathbf{0}$，两端\textbf{左乘} $A^{m-1}$：

\[k_0A^{m-1}\alpha+k_1A^{m}\alpha+\cdots+k_{m-1}A^{2m-2}\alpha=\mathbf{0}\]

因 $A^m\alpha=\mathbf{0}$，故 $A^{m+1}\alpha=A(A^m\alpha)=\mathbf{0}$，更高次幂亦为 $\mathbf{0}$，于是左端只剩 $k_0A^{m-1}\alpha$，即 $k_0A^{m-1}\alpha=\mathbf{0}$。又 $A^{m-1}\alpha\ne\mathbf{0}$，得 $k_0=0$。再左乘 $A^{m-2}$ 同法得 $k_1=0$，依次递推即得全部 $k_i=0$，故线性无关。

\begin{callwarn}{易错点}
\textbf{第 1 招最常见的失败}是代入后\textbf{没有按 $\alpha_1,\cdots,\alpha_s$ 归并同类项}就急于下结论，直接说"因为 $\alpha_i$ 无关所以 $k_i=0$"——必须先化成 $\sum(\cdots)\alpha_i=\mathbf{0}$ 的形式，才能用无关性推出系数为零。

\textbf{第 2 招最常见的失败}是搞错 $C$ 的形状与秩的方向：$(\beta_1,\cdots,\beta_t)=(\alpha_1,\cdots,\alpha_s)C$ 中 $C$ 是 $s\times t$ 的，$\beta$ 组无关 $\iff r(C)=t$；判据里出现的是 $t$（$\beta$ 的\textbf{个数}），不是 $s$。

\textbf{概念陷阱}：由 $(\beta_1,\cdots,\beta_t)=(\alpha_1,\cdots,\alpha_s)C$ 且 $\alpha$ 组线性无关，\textbf{不能}因为"$C$ 可逆"就自动断言 $\beta$ 组无关——必须先把 $C$ 算出来。反过来，若 $r(C)=s<t$，则 $\beta$ 组\textbf{必相关}（以少表多）。

\textbf{方向和量词是选择题的重灾区}：

\begin{itemize}
\item 正确：$\alpha$ 组无关 $\Rightarrow$ 其\textbf{任一}部分组无关；$\alpha$ 组某部分组相关 $\Rightarrow$ \textbf{整体}相关。
\item 错误（常被当成正确选项）："$\alpha_1,\alpha_2$ 无关，$\alpha_3,\alpha_4$ 都不能由 $\alpha_1,\alpha_2$ 表示 $\Rightarrow$ $\alpha_1,\alpha_2,\alpha_3,\alpha_4$ 无关"——反例：$\alpha_3,\alpha_4$ 之间可以相关。
\item 错误："某部分组线性无关 $\Rightarrow$ 整体线性无关"。
\item 错误："$\alpha_1,\alpha_2,\alpha_3$ 线性无关 $\Rightarrow$ $\alpha_1+\alpha_2,\alpha_2+\alpha_3,\alpha_3-\alpha_1$ 线性无关"——注意 $\alpha_3-\alpha_1=(\alpha_1+\alpha_2)-(\alpha_2+\alpha_3)+(\alpha_3+\alpha_1)-2\alpha_1$ 这类配对要具体算行列式，不能凭直觉。实际上 $(\alpha_1+\alpha_2)-(\alpha_2+\alpha_3)+(\alpha_3-\alpha_1)=\mathbf{0}$，故这三个向量\textbf{线性相关}。
\end{itemize}
\end{callwarn}

\begin{callexam}{真题考法}
固定出法有两种：$\circled{1}$ \textbf{选择题的命题判断}——给出四五条关于相关性/可表示性的命题，问正确命题的个数（这类题的正确率最低，因为每条命题都要独立验证）；$\circled{2}$ \textbf{解答/证明题的第一问}——"证明 $\beta_1,\cdots,\beta_t$ 线性无关"，标准答案就是上面第 1 招或第 2 招。

另一类高频变形是把矩阵关系翻译成向量语言：由 $AB=O$（$B$ 非零）可知 \textbf{$B$ 的每一列都是 $Ax=\mathbf{0}$ 的解}，从而 $B$ 的列向量组线性相关；由 $A\alpha_i=\beta_i$ 且 $A$ 可逆，可知 $\alpha$ 组与 $\beta$ 组\textbf{相关性相同}（因为可逆矩阵左乘不改变秩）。看到 $AB=O$ 就往"列向量是解"上想，看到"$A$ 可逆"就往"相关性不变"上想，这是命题人反复使用的两个开关。

\end{callexam}

\section{向量组的秩与等价}

\subsection{极大线性无关组与秩}

\subsubsection{极大线性无关组与向量组的秩 $\bigstar$$\bigstar$$\bigstar$$\bigstar$}

\textbf{定义（极大线性无关组）}：向量组 $\alpha_1,\cdots,\alpha_s$ 的一个部分组 $\alpha_{i_1},\cdots,\alpha_{i_r}$，若满足

\begin{enumerate}
\item $\alpha_{i_1},\cdots,\alpha_{i_r}$ \textbf{线性无关}；
\item 原向量组中\textbf{每个}向量都可由 $\alpha_{i_1},\cdots,\alpha_{i_r}$ 线性表示（等价说法：原组中任意 $r+1$ 个向量——若有的话——都线性相关），
\end{enumerate}
则称它是原向量组的一个\textbf{极大线性无关组}。

\textbf{定义（向量组的秩）}：极大线性无关组所含向量的个数 $r$ 称为向量组的\textbf{秩}，记作 $r(\alpha_1,\cdots,\alpha_s)$。规定只含零向量的向量组没有极大无关组，其秩为 $0$。

\begin{callkey}{极大无关组的核心性质}
\begin{enumerate}
\item 极大无关组\textbf{一般不唯一}，但\textbf{所含向量的个数（即秩）唯一确定}。
\item 向量组与其\textbf{任一}极大无关组\textbf{等价}；任意两个极大无关组\textbf{等价}。
\item 若 $S_1$ 可由 $S_2$ 线性表示，则 $r(S_1)\le r(S_2)$——\textbf{"被表示的秩不大"}。
\item 若 $S_1$ 可由 $S_2$ 表示且 $r(S_1)=r(S_2)$，则 $S_1$ 与 $S_2$ \textbf{等价}。
\item $0\le r(\alpha_1,\cdots,\alpha_s)\le\min\{s,n\}$（$s$ 为向量个数，$n$ 为向量维数）。
\item 单个向量 $\alpha$ 线性相关 $\iff\alpha=\mathbf{0}$；两个向量线性相关 $\iff$ 对应分量成比例。
\end{enumerate}
\end{callkey}

\textbf{求极大无关组与秩的标准步骤（初等行变换法）}：

\begin{enumerate}
\item 把向量组按\textbf{列}拼成矩阵 $A=(\alpha_1,\alpha_2,\cdots,\alpha_s)$；
\item 对 $A$ 只作\textbf{初等行变换}，化为行阶梯形 $B$（原理：初等行变换不改变列向量组的线性相关性）；
\item $B$ 的\textbf{非零行数}就是向量组的秩 $r$；
\item 在 $B$ 中每个台阶取一个主元，\textbf{记下主元所在的列号}；在\textbf{原矩阵 $A$} 中抽取这些列号的列，得到的部分组就是原向量组的一个极大无关组；
\item 若要把其余向量用极大无关组表示，继续把 $B$ 化为\textbf{行最简形}，再解方程组即得表示系数。
\end{enumerate}
\textbf{例}：求 $\alpha_1=(1,1,1)^T,\alpha_2=(1,2,3)^T,\alpha_3=(1,3,5)^T,\alpha_4=(2,3,4)^T$ 的秩与一个极大无关组。

\[A=\begin{pmatrix} 1 & 1 & 1 & 2 \\ 1 & 2 & 3 & 3 \\ 1 & 3 & 5 & 4 \end{pmatrix}
\to\begin{pmatrix} 1 & 1 & 1 & 2 \\ 0 & 1 & 2 & 1 \\ 0 & 2 & 4 & 2 \end{pmatrix}
\to\begin{pmatrix} 1 & 1 & 1 & 2 \\ 0 & 1 & 2 & 1 \\ 0 & 0 & 0 & 0 \end{pmatrix}\]

非零行数为 $2$，故 $r(\alpha_1,\alpha_2,\alpha_3,\alpha_4)=2$。主元在第 $1,2$ 列，故 $\alpha_1,\alpha_2$ 是一个极大无关组。继续化行最简形：

\[\begin{pmatrix} 1 & 1 & 1 & 2 \\ 0 & 1 & 2 & 1 \\ 0 & 0 & 0 & 0 \end{pmatrix}
\to\begin{pmatrix} 1 & 0 & -1 & 1 \\ 0 & 1 & 2 & 1 \\ 0 & 0 & 0 & 0 \end{pmatrix}\]

即第 $3,4$ 列的坐标分别为 $(-1,2,0)^T$ 与 $(1,1,0)^T$，于是

\[\alpha_3=-\alpha_1+2\alpha_2,\qquad \alpha_4=\alpha_1+\alpha_2\]

（可代回验证：$-\alpha_1+2\alpha_2=(-1+2,-1+4,-1+6)^T=(1,3,5)^T=\alpha_3$，$\alpha_1+\alpha_2=(2,3,4)^T=\alpha_4$。）

\begin{callwarn}{易错点}
\textbf{先把行最简形里主元列的列向量当成极大无关组}——这是最经典的错误。行变换\textbf{改变了列向量本身}，只有在\textbf{原矩阵}里按同样列号抽取，得到才是原向量组的极大无关组。

只有零向量的向量组没有极大无关组，其秩规定为 $0$，不要写成 $1$。

"极大无关组不唯一"不等于"随便凑 $r$ 个向量都行"：必须是\textbf{线性无关}的 $r$ 个向量，且要能从原组中抽出来。由 $r(\alpha_1,\alpha_2,\alpha_3)=2$ 推出 $\alpha_1,\alpha_3$ 与 $\alpha_2,\alpha_3$ 也都是极大无关组，但 $\alpha_1,2\alpha_1$ 就不是。

秩为 $r$ 的向量组中，任意 $r+1$ 个向量一定线性相关（这是定义中条件 2 的等价说法）；但"任意 $r$ 个向量都线性无关"是\textbf{错的}。反例：向量组 $(1,0)^T,(2,0)^T,(0,1)^T$ 的秩为 $2$，可是从中取出的 $(1,0)^T$ 与 $(2,0)^T$ 线性相关。正确的命题只能是"\textbf{存在} $r$ 个向量线性无关"（极大无关组的存在性），不能加强成"任意 $r$ 个"。

\end{callwarn}

\begin{callexam}{真题考法}
两大类题型：$\circled{1}$ \textbf{给具体向量求秩与极大无关组}，并"把其余向量用极大无关组线性表示"——这是解答题的固定问法，采分点在"回原矩阵取列"这一步；$\circled{2}$ \textbf{含参数}："当 $a$ 为何值时，$\alpha_1,\cdots,\alpha_4$ 线性相关？此时求一个极大无关组并将其余向量线性表出"——先算行列式或化阶梯形，按参数分类讨论秩，再抽列。

还有一种"\textbf{限定包含某向量}"的问法："求\textbf{包含 $\alpha_1,\alpha_5$} 的一个极大无关组"。此时不能直接从主元列抽列，标准做法是：先判断 $\alpha_1,\alpha_5$ 是否无关（无关则可作为起点），再把其余向量逐个添入，用秩是否增加来决定取舍；或对重排后的矩阵（把 $\alpha_1,\alpha_5$ 放在前两列）作行变换，看前两列是否同时成为主元列。

\end{callexam}

\paragraph{为什么"秩"是向量组唯一的不变量}\mbox{}{\par}

同一个向量组可以有多个极大无关组（例如上面例子里的 $\alpha_1,\alpha_2$ 与 $\alpha_1,\alpha_3$ 都是），它们的\textbf{向量不同}，但\textbf{个数相同}。这说明"向量组里有多少个向量是真正独立的"不是人为选择的结果，而是这组向量固有的属性。

理解这一点最有效的方式是把极大无关组理解成"\textbf{信息压缩}"：原向量组可能很臃肿（几十个向量），但它所携带的全部信息（即它能张成的那个空间）只需要 $r$ 个向量就能完整表达。多出来的向量都是这 $r$ 个向量的线性组合，不提供新信息。于是

\begin{itemize}
\item 秩 $=$ 这组向量携带的"独立信息量"；
\item 极大无关组 $=$ 一份"最省的完整存储"；
\item 相关 $\iff$ 存在冗余 $\iff$ 秩小于向量个数；
\item 两个向量组等价 $\iff$ 它们携带的信息完全相同 $\iff$ 它们张成同一个空间。
\end{itemize}
最后一条给出了等价性的几何本质：\textbf{等价就是"张成同一个空间"}。这也解释了为什么两个等价向量组的向量个数可以不同（一个空间可以用 3 个向量张成，也可以用 5 个向量张成），但秩必须相同（同一个空间的维数只有一个）。

\textbf{线性表示与秩大小关系的直观解读}。若 $B$ 可由 $A$ 线性表示，则 $B$ 的每个向量都"住在" $A$ 张成的空间里，$B$ 张成的空间是这个空间的子空间，所以 $r(B)\le r(A)$。写成口诀就是"\textbf{被表示的秩不大}"。反过来，若 $r(B)\le r(A)$ 并不能推出 $B$ 可由 $A$ 表示——因为 $B$ 完全可以张成 $A$ 的张成空间之外的另一个更小的空间。

\textbf{判断"能否相互表示"的实操建议}：先用秩做快速筛查（三秩判据），只有在秩的条件满足时才去解具体的表示方程组；若秩的条件已经不满足，立刻可以下"不等价"的结论，省下大量计算。

\subsubsection{等价向量组、矩阵的秩与向量组的秩的关系 $\bigstar$$\bigstar$$\bigstar$}

\textbf{定义（向量组等价）}：设向量组 $(I):\alpha_1,\cdots,\alpha_s$ 与 $(II):\beta_1,\cdots,\beta_t$。若 $(I)$ 中每个向量都可由 $(II)$ 线性表示，且 $(II)$ 中每个向量都可由 $(I)$ 线性表示，则称这两个向量组\textbf{等价}。

\textbf{等价关系的性质}：满足\textbf{反身性}、\textbf{对称性}、\textbf{传递性}（传递性在证明题里最常用："$(I)$ 可由 $(II)$ 表示，$(II)$ 可由 $(III)$ 表示，则 $(I)$ 可由 $(III)$ 表示"）。

\begin{callkey}{等价的三个常用充要条件（三秩相同）}
\[(I)\sim(II)\iff (I),(II)\text{ 可相互线性表示}\iff r(I)=r(II)=r(I,II)\]

其中 $r(I,II)$ 是把两组向量\textbf{放在一起}构成的向量组的秩。另有常用的单项版本：

\[(II)\text{ 可由 }(I)\text{ 线性表示，且 } r(I)=r(II)\implies (I)\sim(II)\]

\textbf{常用派生性质}：向量组与其任一极大无关组等价；任意两个极大无关组等价；等价向量组的极大无关组等价；等价向量组\textbf{秩必相等}。

\end{callkey}

\textbf{矩阵的秩与向量组的秩的关系（本考点的第二条主线）}：

设 $A=(\alpha_1,\alpha_2,\cdots,\alpha_s)$ 为 $m\times n$ 矩阵，则

\[r(A)=A\text{ 的列向量组的秩}=A\text{ 的行向量组的秩}\]

即\textbf{矩阵的秩 $=$ 行秩 $=$ 列秩}。由此得到几条高频结论：

\begin{itemize}
\item $A$ 的列向量组线性无关 $\iff r(A)=n$（列满秩）；$A$ 的行向量组线性无关 $\iff r(A)=m$（行满秩）。
\item 若 $AB=C$，则 $C$ 的\textbf{列}向量组可由 $A$ 的列向量组线性表示，$C$ 的\textbf{行}向量组可由 $B$ 的行向量组线性表示。
\item 若 $B$ 可逆，则 $AB$ 的列向量组与 $A$ 的列向量组\textbf{等价}；若 $A$ 可逆，则 $AB$ 的行向量组与 $B$ 的行向量组\textbf{等价}。
\item 若 $AB=O$，则 $B$ 的每一列都是 $Ax=\mathbf{0}$ 的解，从而 $r(A)+r(B)\le n$（$n$ 为 $A$ 的列数）。
\end{itemize}
\textbf{例}：设 $\alpha_1=(1,0,1)^T,\alpha_2=(0,1,1)^T$；$\beta_1=(1,1,2)^T,\beta_2=(1,1,0)^T$，判断两个向量组是否等价。

\textbf{第一步，算三秩}。$A=(\alpha_1,\alpha_2)$ 的前两行构成单位阵，故 $r(A)=2$；$B=(\beta_1,\beta_2)$ 的前两行也构成单位阵，故 $r(B)=2$。合并矩阵

\[(A,B)=\begin{pmatrix} 1 & 0 & 1 & 1 \\ 0 & 1 & 1 & 1 \\ 1 & 1 & 2 & 0 \end{pmatrix}
\to\begin{pmatrix} 1 & 0 & 1 & 1 \\ 0 & 1 & 1 & 1 \\ 0 & 0 & 0 & -2 \end{pmatrix}\]

（第三行减去第一、二行即得末行。）故 $r(A,B)=3$。

\textbf{第二步，用三秩判据}。$r(A)=r(B)=2$，而 $r(A,B)=3$，不满足"三秩相同"，故两个向量组\textbf{不等价}。

\textbf{第三步，找出是谁不能表示谁}。因为 $r(A,B)>r(A)$，说明 $\beta_1,\beta_2$ 中至少有一个不能由 $A$ 线性表示。检验 $\beta_1$：设 $(x_1,x_2,x_1+x_2)^T=(1,1,2)^T$，取 $x_1=x_2=1$ 满足全部三式，故 $\beta_1=\alpha_1+\alpha_2$。再检验 $\beta_2$：需要 $x_1=x_2=1$ 但同时 $x_1+x_2=0$，矛盾，\textbf{无解}。故是 $\beta_2$ 不能被 $A$ 表示，两向量组不等价。

用三秩判据一步定胜负，比逐个去解表示方程组快得多——这是含参数等价问题的标准打法。

\begin{callwarn}{易错点}
\textbf{"等价向量组"与"等价矩阵"是两个不同的概念}。矩阵等价指同型且秩相等（存在可逆 $P,Q$ 使 $PAQ=B$）；向量组等价要求\textbf{能相互线性表示}，是两个更严的条件。若两个矩阵等价，它们的行向量组不一定等价，列向量组也不一定等价。

\textbf{"秩相等"推不出"等价"}：$r(I)=r(II)$ 只说明两组"规模相当"，还需要一个方向的可表示性（或三秩相同）。典型反例是两个平面里各自的基向量组。

\textbf{"$r(B)\le r(A)$"的逆不成立}：$B$ 可由 $A$ 表示 $\Rightarrow r(B)\le r(A)$，但 $r(B)\le r(A)$ \textbf{不能}推出 $B$ 可由 $A$ 表示。

数向量个数时注意：\textbf{等价的两个向量组所含向量个数可以不同}（总个数不同而秩相同），这一点常被误判。

\end{callwarn}

\begin{callexam}{真题考法}
高频出法是\textbf{"含参等价"型选择题}："设 $\alpha_1,\alpha_2,\alpha_3$ 与 $\alpha_1,\alpha_2,\alpha_4$ 等价，求参数取值范围。"标准解法是用三秩判据 $r(I)=r(II)=r(I,II)$，先算 $r(I)$ 与 $r(II)$，把参数使秩下降的值（通常是某些行列式为零的值）找出来，再逐个回代检验是否真正等价——\textbf{千万不要只排除让行列式为零的值就下结论}，必须回代验证。

另一类是\textbf{命题判断题}："若 $AB=C$ 且 $B$ 可逆，则 $C$ 的列向量组与 $A$ 的列向量组等价"（正确）；"若矩阵 $A$ 与 $C$ 等价，则 $A$ 的行向量组与 $C$ 的行向量组等价"（错误，秩相等推不出向量组等价）。

\end{callexam}

\paragraph{专题：抽象向量组线性无关的证明套路（强化）}\mbox{}{\par}

这是数二线代解答题中最"值钱"的一类，单独抽出来系统训练。此类题目的\textbf{题目特征}是：条件里出现"$\alpha_1,\cdots,\alpha_s$ 线性无关"，结论要求证明由它们线性组合出的 $\beta_1,\cdots,\beta_t$ 线性无关（或者反过来说明相关）。

\#\#\#\#\#\# 四种标准打法与选择顺序

\textbf{打法一：矩阵法（首选，最省事）}

把 $\beta$ 组写成 $(\beta_1,\cdots,\beta_t)=(\alpha_1,\cdots,\alpha_s)C$。因为 $\alpha_1,\cdots,\alpha_s$ 线性无关，即 $(\alpha_1,\cdots,\alpha_s)$ 列满秩，故 $r(\beta_1,\cdots,\beta_t)=r(C)$。于是\textbf{把抽象问题完全转化为对具体数字矩阵 $C$ 的秩的计算}：$\beta$ 组线性无关 $\iff r(C)=t$。

$C$ 的列就是每个 $\beta_j$ 关于 $\alpha_1,\cdots,\alpha_s$ 的坐标。写 $C$ 的方法是：看 $\beta_j$ 中 $\alpha_i$ 的系数，填到 $C$ 的第 $i$ 行第 $j$ 列。

\textbf{适用场合}：$\beta_j$ 关于 $\alpha_i$ 的表达式已经明确给出，且 $\alpha$ 组线性无关。绝大多数真题都属于此类。

\textbf{打法二：定义法（万能，但计算量不定）}

设 $k_1\beta_1+\cdots+k_t\beta_t=\mathbf{0}$，逐个代入 $\beta_j$ 的表达式，\textbf{按 $\alpha_1,\cdots,\alpha_s$ 合并同类项}得到 $\sum_{i=1}^{s}\left(\sum_{j=1}^{t}c_{ij}k_j\right)\alpha_i=\mathbf{0}$，由 $\alpha_i$ 无关得齐次方程组

\[\sum_{j=1}^{t}c_{ij}k_j=0,\qquad i=1,\cdots,s\]

再解这个方程组：只有零解则无关，有非零解则相关。

\textbf{注意}：打法二与打法一本质上是同一件事（$C=(c_{ij})$ 就是打法一中的系数矩阵），但定义法把过程完整写出，证明题的格式更规范；矩阵法写起来快，适合选择题。

\textbf{打法三："乘"的技巧（用于条件里含矩阵运算）}

当题目条件中出现 $A\alpha_i=\cdots$、$A^k\alpha=\mathbf{0}$ 之类的关系时，在 $k_1\beta_1+\cdots+k_t\beta_t=\mathbf{0}$ 两端\textbf{左乘}某个合适的矩阵（常用 $A^{m-1}$、$A$、$A^T$），让等式"缩短"，逐项逼出 $k_i=0$。典型题：已知 $A^{m}\alpha=\mathbf{0}$、$A^{m-1}\alpha\neq\mathbf{0}$，证明 $\alpha,A\alpha,\cdots,A^{m-1}\alpha$ 线性无关——依次左乘 $A^{m-1},A^{m-2},\cdots$ 即可。

\textbf{打法四：反证法（配合前三种使用）}

假设 $\beta_1,\cdots,\beta_t$ 线性相关，则存在不全为零的系数使 $\sum k_j\beta_j=\mathbf{0}$，代入后由 $\alpha_i$ 无关得到"系数不全为零却使 $Ck=\mathbf{0}$"，与某个已知的无关性矛盾。反证法适合结论"线性无关"而难以正面算出 $k_i$ 的情形。

\#\#\#\#\#\# 三个必练的标准题型

\textbf{题型 1：初等变换型}。设 $\alpha_1,\cdots,\alpha_s$ 线性无关，$\beta_1=\alpha_1$、$\beta_2=\alpha_1+\alpha_2$、$\cdots$、$\beta_s=\alpha_1+\alpha_2+\cdots+\alpha_s$，证明 $\beta_1,\cdots,\beta_s$ 线性无关。

系数矩阵是下三角且对角元全为 $1$：

\[C=\begin{pmatrix} 1 & 1 & \cdots & 1 \\ 0 & 1 & \cdots & 1 \\ \vdots & & \ddots & \vdots \\ 0 & 0 & \cdots & 1 \end{pmatrix},\qquad \vert C\vert=1\ne0\]

故 $r(C)=s$，$\beta_1,\cdots,\beta_s$ 线性无关。用定义法同样可做：由 $\beta_j-\beta_{j-1}=\alpha_j$（约定 $\beta_0=\mathbf{0}$）知每个 $\alpha_i$ 都能由 $\beta$ 组表示，故两组等价、秩相同。

\textbf{题型 2：循环型}。设 $\alpha_1,\alpha_2,\alpha_3$ 线性无关，$\beta_1=\alpha_1+\alpha_2$、$\beta_2=\alpha_2+\alpha_3$、$\beta_3=\alpha_3+\alpha_1$，判断相关性。

\[C=\begin{pmatrix} 1 & 0 & 1 \\ 1 & 1 & 0 \\ 0 & 1 & 1 \end{pmatrix},\qquad \vert C\vert=2\ne0\]

故无关。若改成 $\beta_1=\alpha_1-\alpha_2$、$\beta_2=\alpha_2-\alpha_3$、$\beta_3=\alpha_3-\alpha_1$，则 $C=\begin{pmatrix} 1 & 0 & -1 \\ -1 & 1 & 0 \\ 0 & -1 & 1 \end{pmatrix}$ 的三行之和为零，$\vert C\vert=0$，故\textbf{相关}。区别只在于"符号方向是否首尾一致"，务必实际算行列式。

\textbf{题型 3：含矩阵关系型}。设 $A$ 为 $n$ 阶矩阵，$\alpha_1,\alpha_2,\alpha_3$ 是 $n$ 维列向量，且

\[A\alpha_1=\alpha_1,\qquad A\alpha_2=\alpha_1+\alpha_2,\qquad A\alpha_3=\alpha_2+\alpha_3\]

证明 $\alpha_1,\alpha_2,\alpha_3$ 线性无关。

先把矩阵关系整理成 $(A-E)\alpha_1=\mathbf{0}$、$(A-E)\alpha_2=\alpha_1$、$(A-E)\alpha_3=\alpha_2$。记 $B=A-E$，则

\[B\alpha_1=\mathbf{0},\qquad B\alpha_2=\alpha_1,\qquad B\alpha_3=\alpha_2\]

设 $k_1\alpha_1+k_2\alpha_2+k_3\alpha_3=\mathbf{0}$，两端左乘 $B^2$：由 $B^2\alpha_1=\mathbf{0}$、$B^2\alpha_2=B\alpha_1=\mathbf{0}$、$B^2\alpha_3=B\alpha_2=\alpha_1$ 得 $k_3\alpha_1=\mathbf{0}$，因 $\alpha_1\ne\mathbf{0}$（它是 $A$ 关于特征值 $1$ 的特征向量，非零），得 $k_3=0$。再左乘 $B$：得 $k_2\alpha_1=\mathbf{0}$，故 $k_2=0$。代回原式得 $k_1\alpha_1=\mathbf{0}$，故 $k_1=0$。于是 $\alpha_1,\alpha_2,\alpha_3$ 线性无关。

\textbf{这道题的解题脉络值得记住}：面对"三角状"的矩阵关系 $B\alpha_{i+1}=\alpha_i$，用 $B$ 的高次幂从最后一项开始往前"剥"，每次剥掉一项，是最直接的打法。

\begin{callwarn}{抽象证明的四个致命失误}
\textbf{失误一：忘记 $\alpha_1\ne\mathbf{0}$ 这类隐含条件}。由 $k\alpha_1=\mathbf{0}$ 推出 $k=0$ 的前提是 $\alpha_1\ne\mathbf{0}$；题目通常通过"$\alpha_1,\cdots,\alpha_s$ 线性无关"或"$\alpha$ 是非零列向量"给出，用之前必须点名。

\textbf{失误二：$C$ 的形状搞反}。$s$ 个 $\alpha$、$t$ 个 $\beta$ 时，$C$ 是 $s\times t$；若写成 $t\times s$，则行列式与秩全部算错。

\textbf{失误三：由"$\alpha$ 组无关"直接断言 "$\beta$ 组无关"而不算 $C$}。例如 $\beta_1=\alpha_1+\alpha_2$、$\beta_2=\alpha_2+\alpha_3$、$\beta_3=\alpha_1+\alpha_3$ 与 $\beta_1=\alpha_1-\alpha_2$、$\beta_2=\alpha_2-\alpha_3$、$\beta_3=\alpha_3-\alpha_1$ 前一无关后一相关，光看形式无法分辨。

\textbf{失误四：把"存在不全为零的系数使 $\sum k_j\beta_j=\mathbf{0}$"中的 $k_j$ 当成"系数为零"}。相关性的定义要求系数不全为零，证明无关要导出"全为零"，中间不能跳步。

\end{callwarn}

\begin{callexam}{抽象证明真题的采分点分布}
以"证明 $\beta_1,\cdots,\beta_t$ 线性无关"为例，标准答案的采分点通常是：

\begin{enumerate}
\item \textbf{写出设 $k_1\beta_1+\cdots+k_t\beta_t=\mathbf{0}$}（或写出 $B=AC$ 的矩阵形式）——格式分；
\item \textbf{代入并正确合并同类项}，得到关于 $\alpha_i$ 的表达式——计算分，也是大多数考生失分的地方；
\item \textbf{引用"$\alpha_1,\cdots,\alpha_s$ 线性无关"得出系数方程组}——概念分，必须明确写出"由线性无关可知，系数必须全为零"；
\item \textbf{解出 $k_1=\cdots=k_t=0$ 并下结论}——结论分。
\end{enumerate}
若采用矩阵法，则第 2、3 步合并为"计算 $\vert C\vert\ne0$（或 $r(C)=t$）"，并明确写出"因 $(\alpha_1,\cdots,\alpha_s)$ 列满秩，故 $r(\beta_1,\cdots,\beta_t)=r(C)=t$"——这句话是不能省的，它是整个推理的桥梁。

\end{callexam}

\subsection*{本章易错清单}
\addcontentsline{toc}{subsection}{本章易错清单}

\begin{enumerate}
\item \textbf{相关/无关是向量组的整体属性}，不能说"某个向量线性相关"。单个向量线性相关 $\iff$ 它是零向量。
\item $k_1\alpha_1+\cdots+k_s\alpha_s=\mathbf{0}$ 中"\textbf{不全为零}"对应相关，"\textbf{全为零才成立}"对应无关；把两个条件写反是最低级也最常见的失分。
\item 判断"$A$ 的列向量组相关性"只能用\textbf{初等行变换}；用列变换会改变列向量之间的线性关系，导致结论错误。
\item 求极大无关组时，务必回到\textbf{原矩阵}按行阶梯形中主元列的列号抽列，不能直接抄行阶梯形里的列。
\item 秩为 $r$ 的向量组中"任意 $r$ 个向量都线性无关"是\textbf{假命题}；正确说法是"任意 $r+1$ 个向量都线性相关"。
\item 秩相等 $\ne$ 向量组等价；必须加上"一个方向可线性表示"或"三秩相同"。
\item "以少表多，多必相关"：若 $(I)$ 可由 $(II)$ 线性表示，且 $(I)$ 的向量个数大于 $(II)$ 的向量个数，则 $(I)$ 必线性相关。反之"以多表少"不能保证无关。
\item $n$ 维空间中，向量个数 $s>n$ 时\textbf{必线性相关}；但 $s\le n$ 时相关与否必须实际判断。
\item 初等行变换不改变\textbf{列}向量组的线性相关性，初等列变换不改变\textbf{行}向量组的线性相关性；两者不可混用。
\item 施密特正交化的分母是 $(\beta_k,\beta_k)$（正交化后向量的内积），不是 $(\alpha_k,\alpha_k)$；单位化要除以 $\Vert\beta_k\Vert$。
\item 线性相关的向量组作正交化时会出现某个 $\beta_k=\mathbf{0}$，此时不能单位化——这说明原组秩不足，而不是公式写错。
\item $AB=O$ 的翻译方向：$B$ 的\textbf{列}向量是 $Ax=\mathbf{0}$ 的解（不是行）；$A$ 的\textbf{行}向量与解向量正交。方向搞反会连错一串结论。
\end{enumerate}
\chapter{线性方程组}
\markboth{线性方程组}{线性方程组}
\subsection*{本章考纲要求}
\addcontentsline{toc}{subsection}{本章考纲要求}

\begin{quote}
\textbf{考试内容}

\begin{enumerate}
\item 线性方程组的克拉默(Cramer)法则
\item 齐次线性方程组有非零解的充分必要条件
\item 非齐次线性方程组有解的充分必要条件
\item 线性方程组解的性质和解的结构
\item 齐次线性方程组的基础解系和通解
\item 非齐次线性方程组的通解
\end{enumerate}
\textbf{考试要求}

\begin{enumerate}
\item 会用克拉默法则。
\item 理解齐次线性方程组有非零解的充分必要条件及非齐次线性方程组有解的充分必要条件。
\item 理解齐次线性方程组的基础解系及通解的概念，掌握齐次线性方程组的基础解系和通解的求法。
\item 理解非齐次线性方程组解的结构及通解的概念。
\item 掌握用初等行变换求解线性方程组的方法。
\end{enumerate}
\end{quote}
\subsection*{本章星级总表}
\addcontentsline{toc}{subsection}{本章星级总表}

\begingroup\small\setlength{\tabcolsep}{4.0pt}
\begin{longtable}{@{}L{0.4549\linewidth}L{0.1098\linewidth}L{0.1098\linewidth}L{0.1098\linewidth}L{0.1098\linewidth}@{}}
\toprule
\textbf{考点} & \textbf{星级} & \textbf{考纲要求} & \textbf{难度} & \textbf{常见题型} \\
\midrule
\endhead
克拉默法则 & $\bigstar$$\bigstar$ & 会 & 基础 & 选择、填空 \\
齐次线性方程组有非零解的充分必要条件、基础解系与通解 & $\bigstar$$\bigstar$$\bigstar$$\bigstar$$\bigstar$ & 理解/掌握 & 中等 & 解答 \\
非齐次线性方程组有解的充分必要条件、解的结构与通解 & $\bigstar$$\bigstar$$\bigstar$$\bigstar$$\bigstar$ & 理解/掌握 & 中等 & 解答 \\
用初等行变换求解方程组（含参数讨论） & $\bigstar$$\bigstar$$\bigstar$$\bigstar$$\bigstar$ & 会 & 难 & 解答 \\
同解方程组、公共解与解的反求 & $\bigstar$$\bigstar$$\bigstar$$\bigstar$ & 会用 & 难 & 解答、证明 \\
方程组解的性质与线性组合（系数和为 1/0 的判据） & $\bigstar$$\bigstar$$\bigstar$$\bigstar$ & 理解/会 & 中等 & 选择、证明 \\
\bottomrule
\end{longtable}
\endgroup

\paragraph{学这一章要建立的三个直觉}\mbox{}{\par}

第一个直觉是"\textbf{秩 = 有效方程个数}"。一个方程组写了几行不代表提供了几条独立约束：$x_1+x_2=1$ 与 $2x_1+2x_2=2$ 其实是同一个约束。$r(A)$ 就是真正起作用的约束条数。未知数个数 $n$ 是"自由度总数"，$n-r(A)$ 就是约束之外还剩多少自由。整章所有结论都可以从"约束数"与"未知数个数"的比较中推出来。

第二个直觉是"\textbf{解集的结构只有两种形状}"。齐次方程组的解集是一个过原点的线性空间，它由 $n-r(A)$ 个"基本方向"张成，于是有基础解系与通解；非齐次方程组的解集不是一个空间，而是一个空间"搬了个家"——把齐次解空间整体平移到一个特解上，于是通解必然长成"一个特解加上任意常数乘基础解系"的样子。这个"平移"的图景一旦建立，非齐次解的结构定理就不再需要记忆。

第三个直觉是"\textbf{比较 $r(A)$ 与 $r(\bar A)$ 是唯一的万能手段}"。不管是具体的数字方程组、含参数的方程组，还是只给出解的信息要反求参数的抽象题，判断解的情况永远归结为比较这两个秩。$\vert A\vert=0$ 只是在方阵情形下对 $r(A)<n$ 的一种便捷写法，它本身\textbf{不能}区分"无解"与"无穷多解"，这一点是本章最大的坑。

\section{方程组有解的判定与克拉默法则}

线性方程组是线性代数的\textbf{出发点和归宿}：第三章的向量语言、第二章的矩阵秩，最终都服务于"解方程组"这一件事。数二试卷中本章几乎必有一道解答题（含参数讨论或抽象解的结构），分值常在 10 分以上，而且第二问往往与特征值、二次型串联出题。

\subsection{基本概念与解的判定}

\subsubsection{克拉默法则 $\bigstar$$\bigstar$}

\textbf{定理（克拉默法则）}：设线性方程组

\[\begin{cases}
a_{11}x_1+a_{12}x_2+\cdots+a_{1n}x_n=b_1 \\
a_{21}x_1+a_{22}x_2+\cdots+a_{2n}x_n=b_2 \\
\qquad\cdots\cdots\cdots\cdots \\
a_{n1}x_1+a_{n2}x_2+\cdots+a_{nn}x_n=b_n
\end{cases}\]

的系数矩阵为 $A=(a_{ij})_{n\times n}$。若系数行列式 $D=\vert A\vert\ne 0$，则方程组有\textbf{唯一解}

\[x_j=\frac{D_j}{D},\qquad j=1,2,\cdots,n\]

其中 $D_j$ 是把 $D$ 的\textbf{第 $j$ 列}换成常数列 $(b_1,b_2,\cdots,b_n)^T$、其余各列保持不动所得的行列式。

\textbf{齐次情形的直接推论}：对 $n$ 元齐次方程组 $Ax=\mathbf{0}$，

\[\vert A\vert\ne 0\implies \text{只有零解};\qquad \vert A\vert=0\implies \text{有非零解}\]

\begin{callkey}{使用前提与两条"变形"理解}
克拉默法则\textbf{只对方程个数 $=$ 未知数个数且 $\vert A\vert\ne 0$ 的方程组}成立。三个条件缺一不可。

若把 $x_j$ 看成"把常数列 $b$ 用 $A$ 的列向量组表示时的组合系数"，则克拉默法则就是线性表示唯一性的显式公式：$b=x_1\alpha_1+\cdots+x_n\alpha_n$，其中 $\alpha_j$ 是 $A$ 的第 $j$ 列。

\textbf{另一条常用等价形式}（伴随矩阵视角）：$A^{-1}=\dfrac{1}{\vert A\vert}A^*$，于是 $x=A^{-1}b$，即 $x_j=\dfrac{1}{\vert A\vert}\sum_{i=1}^{n}A_{ij}b_i$。这也给出了 $A^{-1}$ 的第 $(i,j)$ 元素为 $\dfrac{A_{ji}}{\vert A\vert}$（注意下标转置）。

\end{callkey}

\textbf{例}：用克拉默法则求解

\[\begin{cases} 2x_1-x_2=2 \\ -3x_1+2x_2=-1 \end{cases}\]

\[D=\begin{vmatrix} 2 & -1 \\ -3 & 2 \end{vmatrix}=4-3=1\ne 0,\qquad
D_1=\begin{vmatrix} 2 & -1 \\ -1 & 2 \end{vmatrix}=4-1=3,\qquad
D_2=\begin{vmatrix} 2 & 2 \\ -3 & -1 \end{vmatrix}=-2+6=4\]

\[x_1=\frac{D_1}{D}=3,\qquad x_2=\frac{D_2}{D}=4\]

验算：$2\cdot3-4=2$，$-3\cdot3+2\cdot4=-1$，均成立。

\begin{callwarn}{易错点}
忘记先验证 $\vert A\vert\ne 0$ 就直接套公式——这是克拉默法则的\textbf{唯一硬性前提}，一旦 $\vert A\vert=0$ 公式立刻失效（此时可能无解，也可能有无穷多解，必须改用秩判定）。

求 $D_j$ 时只能替换\textbf{第 $j$ 列}，千万不能顺手把常数列加到别处；而且替换后行列式的\textbf{列序不能交换}，交换一次就多一个负号。

$n\ge 4$ 时克拉默法则计算量会爆炸（要算 $n+1$ 个 $n$ 阶行列式），统考中它主要作为\textbf{理论工具}使用（例如证明 $\vert A\vert\ne 0\iff$ 唯一解），实际求解一律用初等行变换。

\end{callwarn}

\begin{callexam}{真题考法}
克拉默法则在选择题里的高频出法是\textbf{"理论判据"型}："设 $A$ 为 $n$ 阶矩阵，方程组 $Ax=b$ 对任意 $n$ 维向量 $b$ 都有唯一解"等价于 $\vert A\vert\ne 0$ 等价于 $A$ 可逆。这句话的"对任意 $b$"是关键词——去掉"任意"两个字结论立刻变弱。

另一类是\textbf{"行列式与解的联动"}：给出 $D$ 与 $D_j$ 的某个关系式反求参数，或问"$x_1$ 与 $x_2$ 的比例"。法则是桥梁，落脚点仍是行列式计算。

\end{callexam}

\subsection{齐次线性方程组的解的结构}

\subsubsection{齐次线性方程组有非零解的充分必要条件、基础解系与通解 $\bigstar$$\bigstar$$\bigstar$$\bigstar$$\bigstar$}

\textbf{符号约定}：设 $A$ 为 $m\times n$ 矩阵，则 $Ax=\mathbf{0}$ 是 $n$ 元齐次线性方程组；$r(A)$ 就是"\textbf{有效方程个数}"，$n$ 是未知数个数，$n-r(A)$ 是\textbf{自由变量个数}。

\begin{callkey}{有非零解的充分必要条件（四种等价说法）}
\[Ax=\mathbf{0}\text{ 有非零解}
\iff r(A)<n
\iff A\text{ 的列向量组线性相关}
\iff A\text{ 的 }n\text{ 个列向量"表出能力不足"}\]

当 $m=n$（方阵）时，还可写成 $\vert A\vert=0$。反之：

\[Ax=\mathbf{0}\text{ 只有零解}\iff r(A)=n \iff A\text{ 列满秩}\iff A\text{ 的列向量组线性无关}\]

\textbf{两条"无条件成立"的保底结论}：

\begin{itemize}
\item 齐次方程组\textbf{永远有解}（零解一定存在），所以讨论齐次方程组只问"是否只有零解"。
\item 若 $m<n$（方程个数少于未知数个数），则必有 $r(A)\le m<n$，故\textbf{一定有非零解}。这是"维度检验法"的秒杀依据。
\end{itemize}
\end{callkey}

\begin{callkey}{基础解系（三条件）与通解}
若 $Ax=\mathbf{0}$ 的一组解 $\xi_1,\xi_2,\cdots,\xi_{n-r}$ 满足

\begin{enumerate}
\item \textbf{是解}：每个 $\xi_i$ 都是 $Ax=\mathbf{0}$ 的解向量；
\item \textbf{无关}：$\xi_1,\cdots,\xi_{n-r}$ 线性无关；
\item \textbf{个数达标}：个数恰为 $n-r(A)$，
\end{enumerate}
则称 $\xi_1,\cdots,\xi_{n-r}$ 是 $Ax=\mathbf{0}$ 的一个\textbf{基础解系}，而

\[x=k_1\xi_1+k_2\xi_2+\cdots+k_{n-r}\xi_{n-r}\qquad(k_1,\cdots,k_{n-r}\text{ 为任意常数})\]

称为 $Ax=\mathbf{0}$ 的\textbf{通解}。

\textbf{理解要点}：基础解系就是\textbf{解向量组的极大线性无关组}；解集本身构成一个 $n-r(A)$ 维的线性空间（解空间），基础解系就是它的一组基。解空间维数公式：

\[\dim\{x\mid Ax=\mathbf{0}\}=n-r(A)\]

\end{callkey}

\textbf{为什么基础解系恰好有 $n-r(A)$ 个向量}。把行最简形写出来看：$r(A)$ 个非零行对应 $r(A)$ 个主元，也就有 $r(A)$ 个"主变量"；剩下 $n-r(A)$ 个未知数没有主元，就是"自由变量"。自由变量可以任意取值，每选定一组自由变量的值，主变量的值就由方程组唯一确定。因此"解"与"自由变量的取值"之间是一一对应的，而 $n-r(A)$ 个自由变量恰好有 $n-r(A)$ 个"基本取法"（每次让其中一个取 $1$、其余取 $0$），这就给出了 $n-r(A)$ 个解向量；又因为自由变量的取值互相独立，这些解向量线性无关，且它们的任意线性组合能覆盖所有解。

这段推理同时也解释了为什么"解空间维数 $=n-r(A)$"。

\textbf{求解基础解系与通解的标准步骤（四步）}：

\begin{enumerate}
\item 写出系数矩阵 $A$，用\textbf{初等行变换}化为\textbf{行最简形}（不是只化到行阶梯形，因为要读出解）；
\item 定位主元：每行第一个非零元所在列对应的未知数称为\textbf{主变量}，其余未知数称为\textbf{自由变量}；
\item \textbf{自由变量依次取标准基向量}（第 $1$ 个自由变量取 $1$、其余取 $0$，再第 $2$ 个取 $1$、其余取 $0$……），每次解出主变量的值，得到一个解向量；
\item 把得到的 $n-r$ 个解向量写成通解 $x=k_1\xi_1+\cdots+k_{n-r}\xi_{n-r}$。
\end{enumerate}
\begin{callkey}{快速解法口诀（行最简形直接"抄"）}
化出行最简形后，若自由变量为 $x_{j_1},\cdots,x_{j_{n-r}}$，则对\textbf{每一个自由变量} $x_{j}$：

\begin{itemize}
\item 在解向量的第 $j$ 个分量位置写 $1$；
\item 在解向量的\textbf{主变量分量位置}写该自由变量所在列的系数的\textbf{相反数}（行最简形中该列在主变量行上的数取负）。
\end{itemize}
例如行最简形为 $\begin{pmatrix}1 & 0 & -1 & -2 \\ 0 & 1 & 2 & 3\end{pmatrix}$，主变量 $x_1,x_2$，自由变量 $x_3,x_4$，则

\[\xi_1=\begin{pmatrix} 1 \\ -2 \\ 1 \\ 0 \end{pmatrix}\ (\text{对应 }x_3=1,x_4=0),\qquad
\xi_2=\begin{pmatrix} 2 \\ -3 \\ 0 \\ 1 \end{pmatrix}\ (\text{对应 }x_3=0,x_4=1)\]

口诀的道理很简单：行最简形已经给出了 $x_1-x_3-2x_4=0$、$x_2+2x_3+3x_4=0$，即 $x_1=x_3+2x_4$、$x_2=-2x_3-3x_4$，把自由变量代成标准基向量即得。

\end{callkey}

\textbf{例}：求齐次方程组 $\begin{cases} x_1+2x_2-x_3=0 \\ 2x_1+4x_2+x_3=0 \end{cases}$ 的基础解系与通解。

\[A=\begin{pmatrix} 1 & 2 & -1 \\ 2 & 4 & 1 \end{pmatrix}
\longrightarrow\begin{pmatrix} 1 & 2 & -1 \\ 0 & 0 & 3 \end{pmatrix}
\longrightarrow\begin{pmatrix} 1 & 2 & -1 \\ 0 & 0 & 1 \end{pmatrix}
\longrightarrow\begin{pmatrix} 1 & 2 & 0 \\ 0 & 0 & 1 \end{pmatrix}\]

$r(A)=2$，$n=3$，故 $n-r=1$，基础解系含 $1$ 个向量。主变量是 $x_1,x_3$，自由变量是 $x_2$。行最简形给出 $x_1+2x_2=0$ 与 $x_3=0$。取 $x_2=1$，得 $x_1=-2,x_3=0$，于是

\[\xi=\begin{pmatrix} -2 \\ 1 \\ 0 \end{pmatrix},\qquad
\text{通解 } x=k\begin{pmatrix} -2 \\ 1 \\ 0 \end{pmatrix}\ (k\text{ 为任意常数})\]

验算：$(-2)+2(1)-0=0$，$2(-2)+4(1)+0=0$，正确。

\begin{callwarn}{易错点}
\textbf{基础解系的个数由 $n-r(A)$ 决定，不能多也不能少}。写通解时漏掉某个自由变量对应的解向量、或者误把某个主变量当成自由变量，都是致命错误。

\textbf{必须化到行最简形}再读解。只化到行阶梯形时，主变量之间还纠缠着自由变量，需要逐步回代，容易算错；而行最简形的每一行主元为 $1$、该列其余元素为 $0$，系数可以直接读出。

\textbf{自由变量的取法不唯一}（也可取 $(1,1),(0,1)$ 之类），但基础解系所含向量个数 $n-r$ 是\textbf{唯一确定}的。答题时统一取标准基向量最为规范。

\textbf{只含零向量的解集}（即 $r(A)=n$）没有基础解系——此时解空间是零维的，通解就是 $x=\mathbf{0}$，不要硬凑一个解向量出来。

常见概念陷阱：若 $\alpha_1,\alpha_2,\alpha_3$ 是基础解系，那么 $\alpha_1-\alpha_2,\alpha_2-\alpha_3,\alpha_3-\alpha_1$ \textbf{不是}基础解系（它们线性相关，因为三者之和为零）；而 $\alpha_1,\alpha_1+\alpha_2,\alpha_1+\alpha_2+\alpha_3$ \textbf{是}基础解系（可验证其系数矩阵为上三角且对角元全为 $1$，行列式为 $1\ne0$）。判据是：新向量组必须\textbf{个数相同且线性无关}。

\end{callwarn}

\begin{callexam}{真题考法}
两大固定出法。$\circled{1}$ \textbf{具体计算}："求齐次方程组的基础解系与通解"，或"求一个齐次方程组，使它的基础解系为给定的两个向量"——后者的做法是设 $A$ 的每一行 $(a,b,c,d)$ 与两个已知解向量都正交，解出行的结构，再取一组线性无关的行拼成 $A$。$\circled{2}$ \textbf{抽象判定}："已知 $\xi_1,\xi_2$ 是 $Ax=\mathbf{0}$ 的基础解系，判断下列哪组向量也是基础解系"——只看两条：个数是否仍为 $n-r$，以及是否线性无关。

还有一类与伴随矩阵相结合的高频题：已知 $A$ 不可逆且某个代数余子式 $A_{12}\ne0$，则 $r(A)=n-1$，从而 $r(A^*)=1$，于是 $A^*x=\mathbf{0}$ 的基础解系含 $n-1$ 个向量。看到"$A$ 不可逆但某代数余子式非零"就要立刻反应出 "$r(A)=n-1$" 这条链。

\end{callexam}

\subsection{非齐次线性方程组}

\subsubsection{非齐次线性方程组有解的充分必要条件、解的结构与通解 $\bigstar$$\bigstar$$\bigstar$$\bigstar$$\bigstar$}

\textbf{符号约定}：设 $A=(\alpha_1,\cdots,\alpha_n)$ 为 $m\times n$ 矩阵，$\bar A=(A\mid b)$ 为\textbf{增广矩阵}，$b\ne\mathbf{0}$。称 $Ax=\mathbf{0}$ 为 $Ax=b$ 的\textbf{导出组}（导出齐次方程组）。

\begin{callkey}{有解的充分必要条件}
\[Ax=b\text{ 有解}\iff r(A)=r(\bar A)\iff b\text{ 可由 }A\text{ 的列向量组线性表示}\]

\textbf{无解的判据}：$r(A)<r(\bar A)$，且此时必有 $r(\bar A)=r(A)+1$。

\textbf{有解时的三种情形}（这是本章最核心的一张表）：

\begingroup\small\setlength{\tabcolsep}{4.0pt}
\begin{longtable}{@{}L{0.3960\linewidth}L{0.5496\linewidth}@{}}
\toprule
\textbf{条件} & \textbf{解的情况} \\
\midrule
\endhead
$r(A)\ne r(\bar A)$ & 无解 \\
$r(A)=r(\bar A)=n$ & 有唯一解 \\
$r(A)=r(\bar A)=r<n$ & 有无穷多解，通解含 $n-r$ 个自由参数 \\
\bottomrule
\end{longtable}
\endgroup

特别地，当 $m=n$ 时："有唯一解 $\iff\vert A\vert\ne0$"；"$\vert A\vert=0$ 时，或用秩判定（无解 / 无穷多解），\textbf{不能}一律说成无解"。

\end{callkey}

\begin{callkey}{解的结构定理}
若 $\eta^*$ 是 $Ax=b$ 的\textbf{某一个}特解，$\xi_1,\cdots,\xi_{n-r}$ 是导出组 $Ax=\mathbf{0}$ 的基础解系，则 $Ax=b$ 的\textbf{通解}为

\[x=\eta^*+k_1\xi_1+k_2\xi_2+\cdots+k_{n-r}\xi_{n-r}\qquad(k_i\text{ 为任意常数})\]

\textbf{等式右端的第一项可以换成任何一个特解}，通解集合不变。

\end{callkey}

\textbf{关于"特解 $+$ 齐次通解"这个结构的理解}。设 $\eta$ 是 $Ax=b$ 的任意一个解。任取另一个特解 $\eta^*$，则 $A(\eta-\eta^*)=b-b=\mathbf{0}$，说明 $\eta-\eta^*$ 是导出组的解，即 $\eta-\eta^*=k_1\xi_1+\cdots+k_{n-r}\xi_{n-r}$，移项就得到 $\eta=\eta^*+k_1\xi_1+\cdots+k_{n-r}\xi_{n-r}$。这个三行推理就完成了结构定理的证明，也是抽象证明题的模板：\textbf{要说明"某个东西是通解"，只需说明它是一个特解加上导出组的任意解}。

由此还得到两条常用的操作：已知两个特解，作差立刻得到一个齐次解；已知一个特解与导出组的通解，直接拼出非齐次的通解，不需要重新解方程。

\textbf{求解非齐次方程组的步骤（对增广矩阵操作）}：

\begin{enumerate}
\item 写出增广矩阵 $\bar A=(A\mid b)$，用初等\textbf{行}变换化为行阶梯形；
\item 比较 $r(A)$ 与 $r(\bar A)$：若不等则\textbf{无解}（行阶梯形中会出现 $\begin{pmatrix}0&\cdots&0&d\end{pmatrix}$ 而 $d\ne0$ 的"矛盾行"），结束；若相等，进入第 3 步；
\item 继续化为\textbf{行最简形}，确定主变量与自由变量；
\item \textbf{求特解 $\eta^*$}：令\textbf{所有自由变量都取 $0$}，解出主变量的值，得到的解向量就是 $\eta^*$（这是最省事的取法）；
\item \textbf{求导出组的基础解系}：把行最简形中常数列看作 $0$，按上一节的"抄系数取相反数"办法写出 $n-r$ 个解向量；
\item 拼成通解 $x=\eta^*+k_1\xi_1+\cdots+k_{n-r}\xi_{n-r}$。
\end{enumerate}
\textbf{例}：求解 $\begin{cases} x_1+x_2-x_3=1 \\ 2x_1+2x_2-x_3=3 \end{cases}$

\[\bar A=\begin{pmatrix} 1 & 1 & -1 & 1 \\ 2 & 2 & -1 & 3 \end{pmatrix}
\longrightarrow\begin{pmatrix} 1 & 1 & -1 & 1 \\ 0 & 0 & 1 & 1 \end{pmatrix}\]

$r(A)=r(\bar A)=2<3=n$，有无穷多解。继续化行最简形：

\[\longrightarrow\begin{pmatrix} 1 & 1 & 0 & 2 \\ 0 & 0 & 1 & 1 \end{pmatrix}\]

主变量 $x_1,x_3$，自由变量 $x_2$。

\begin{itemize}
\item \textbf{特解}：令 $x_2=0$，由 $x_1=2$、$x_3=1$ 得 $\eta^*=(2,0,1)^T$。
\item \textbf{导出组基础解系}：$x_1+x_2=0$，$x_3=0$，取 $x_2=1$ 得 $\xi=(-1,1,0)^T$。
\end{itemize}
故通解为

\[x=\begin{pmatrix} 2 \\ 0 \\ 1 \end{pmatrix}+k\begin{pmatrix} -1 \\ 1 \\ 0 \end{pmatrix}\qquad(k\text{ 为任意常数})\]

验算：$k=0$ 时 $2+0-1=1$、$4+0-1=3$；一般情形 $x_1=2-k,x_2=k,x_3=1$，代入第一式得 $(2-k)+k-1=1$，第二式 $2(2-k)+2k-1=3$，均成立。

\begin{callwarn}{易错点}
\textbf{求特解时"令自由变量全为 $0$"只在非齐次情形下有效}；求导出组基础解系时自由变量要取标准基向量（其中一个为 $1$）。两件事用的是\textbf{同一张行最简形}，但取值方式不同，考生常混。

\textbf{通解中特解那一项千万不能漏写}，也不能把特解与基础解系混在同一个 $k$ 里。通解的结构永远长成"一个特解 $+$ 任意常数乘基础解系"。

\textbf{"$r(A)=r(\bar A)=n$ 才有唯一解"}——这里的 $n$ 是\textbf{未知数个数}，不是方程个数 $m$。$m>n$ 时也可能有唯一解（多余方程是冗余的），$m<n$ 时\textbf{不可能}有唯一解。

用行列式判唯一解时只能用于方阵；$m\ne n$ 时 $\vert A\vert$ 根本没有定义。

\textbf{含参数时不能默认 $r(A)=r(\bar A)$}：$\vert A\vert=0$ 只是"可能无解、也可能无穷多解"，必须回代检验增广矩阵的秩。

\end{callwarn}

\begin{callexam}{真题考法}
最高频的形式是\textbf{含参数的三元方程组讨论}："问 $a$ 取何值时方程组无解 / 有唯一解 / 有无穷多解，并在有无穷多解时求通解。"标准流程：系数矩阵行列式 $\to$ 找出使 $\vert A\vert=0$ 的参数值 $\to$ 对每个值把参数代入增广矩阵化阶梯形，看 $r(A)$ 与 $r(\bar A)$ 是否相等 $\to$ 分别作答。\textbf{采分点在于"每个临界值都要单独回代"}，只算行列式就下结论会被扣掉大半分。

第二类高频形式是\textbf{抽象解的结构}：给出 $Ax=b$ 的几个具体解（形如 $A\alpha_1=b$、$A\alpha_2=b$），要求通解。做法是用解的性质作差得到导出组的解：$\alpha_1-\alpha_2$、$\alpha_1-\alpha_3$ 都是 $Ax=\mathbf{0}$ 的解，再从这些差向量中挑出 $n-r$ 个线性无关的，添上任意一个特解即得通解。

第三类是与矩阵方程结合："已知 $A^TAx=A^Tb$ 的通解"（这类方程一定相容，因为 $r(A^TA)=r(A^TA\mid A^Tb)$），解法仍是化行最简形。

\end{callexam}

\section{含参数方程组、同解与公共解}

\subsection{方程组的求解与参数讨论}

\subsubsection{用初等行变换求解方程组（含参数讨论） $\bigstar$$\bigstar$$\bigstar$$\bigstar$$\bigstar$}

\textbf{核心工具}：初等行变换求解，本质是\textbf{高斯消元法}。对方程组作初等行变换（交换两方程、某方程乘非零常数、某方程的 $k$ 倍加到另一方程）得到\textbf{同解}方程组，对应到矩阵上就是对增广矩阵作初等行变换。

\textbf{含参数问题的三条纪律}：

\begin{enumerate}
\item \textbf{不轻易除以含参数的式子}。行变换中若要用某一行除以含参数的表达式，必须先确认它不为零，否则会把某个参数值对应的情形漏掉。规避办法：把含参数项集中到同一列，让参数自然"浮"到阶梯形的下方。
\item \textbf{先算行列式找临界值（方阵情形）}，但\textbf{临界值必须逐个回代}，用秩判定到底是无解还是无穷多解。
\item \textbf{分开书写每种情形的结论}，并明确"此时 $r(A)=\cdots$，$r(\bar A)=\cdots$"，把判据写在纸上，评卷时是明确的采分点。
\end{enumerate}
\textbf{例}：讨论 $a$ 取何值时，方程组

\[\begin{cases} x_1+2x_2+3x_3=-3 \\ 4x_1+5x_2+6x_3=-6 \\ 7x_1+8x_2+9x_3=a \end{cases}\]

无解、有唯一解、有无穷多解。

\[\bar A=\begin{pmatrix} 1 & 2 & 3 & -3 \\ 4 & 5 & 6 & -6 \\ 7 & 8 & 9 & a \end{pmatrix}
\xrightarrow{\ r_2-4r_1\ }\xrightarrow{\ r_3-7r_1\ }
\begin{pmatrix} 1 & 2 & 3 & -3 \\ 0 & -3 & -6 & 6 \\ 0 & -6 & -12 & a+21 \end{pmatrix}
\xrightarrow{r_3-2r_2}
\begin{pmatrix} 1 & 2 & 3 & -3 \\ 0 & -3 & -6 & 6 \\ 0 & 0 & 0 & a-9 \end{pmatrix}\]

系数矩阵只含常数列，其秩与 $a$ 无关，恒有 $r(A)=2$（第三列可以由前两列线性表示：$(3,6,9)^T=2\cdot(2,5,8)^T-1\cdot(1,4,7)^T$）。于是：

\begin{itemize}
\item 当 $a-9\ne0$，即 $a\ne9$ 时，$r(\bar A)=3>r(A)=2$，方程组\textbf{无解}；
\item 当 $a=9$ 时，$r(A)=r(\bar A)=2<3$，方程组有\textbf{无穷多解}。
\end{itemize}
至于唯一解：因为 $r(A)$ 至多为 $2<3=n$，方程组\textbf{永远没有唯一解}。

取 $a=9$ 求通解：把 $\bar A$ 继续化为行最简形（第二行除以 $-3$，再把第一行减去第二行的 $2$ 倍）：

\[\begin{pmatrix} 1 & 2 & 3 & -3 \\ 0 & 1 & 2 & -2 \\ 0 & 0 & 0 & 0 \end{pmatrix}
\longrightarrow
\begin{pmatrix} 1 & 0 & -1 & 1 \\ 0 & 1 & 2 & -2 \\ 0 & 0 & 0 & 0 \end{pmatrix}\]

主变量 $x_1,x_2$，自由变量 $x_3$。

\begin{itemize}
\item 特解：令 $x_3=0$，得 $x_1=1,x_2=-2$，即 $\eta^*=(1,-2,0)^T$；
\item 导出组：$x_1=x_3,\ x_2=-2x_3$，取 $x_3=1$ 得 $\xi=(1,-2,1)^T$。
\end{itemize}
\[x=\begin{pmatrix} 1 \\ -2 \\ 0 \end{pmatrix}+k\begin{pmatrix} 1 \\ -2 \\ 1 \end{pmatrix}\qquad(k\text{ 为任意常数})\]

\textbf{验算说明}：取 $k=1$ 得 $x=(2,-4,1)^T$，代入前两式：$2-8+3=-3$、$8-20+6=-6$，均与原方程组吻合。第三式的左端为 $14-32+9=-9$，与右端 $a$ 比对可知：只有当 $a=-9$ 时第三式才成立，可 $a=-9$ 时 $r(\bar A)=3\ne r(A)=2$（无解），与"有解"矛盾——这正说明在 $a=9$ 时增广矩阵的第三行\textbf{已被完全消成零行}（$0=0$），第三式不再提供任何新约束，通解只由前两式决定。因此验算只需带回\textbf{原方程组的前两式}，而 $a$ 的取值由"阶梯形最后一行的常数项是否为零"单独确定。这是含参数方程组验算时的通用做法：\textbf{用不含参数（或参数已确定）的方程验算解，用阶梯形的零行条件确定参数}。

\textbf{例（经典 $2\times2$ 参数题）}：讨论 $a$ 取何值时方程组

\[\begin{cases} a x_1+x_2=1 \\ x_1+a x_2=1 \end{cases}\]

无解、有唯一解、有无穷多解。

系数行列式 $D=\begin{vmatrix} a & 1 \\ 1 & a \end{vmatrix}=a^2-1=(a-1)(a+1)$。

\begin{itemize}
\item $a\ne\pm1$ 时 $D\ne0$，有\textbf{唯一解}；
\item $a=1$ 时 $D=0$，增广矩阵 $\begin{pmatrix} 1 & 1 & 1 \\ 1 & 1 & 1 \end{pmatrix}\to\begin{pmatrix} 1 & 1 & 1 \\ 0 & 0 & 0 \end{pmatrix}$，$r(A)=r(\bar A)=1<2$，有\textbf{无穷多解} $x=k(1,-1)^T+(1,0)^T$；
\item $a=-1$ 时 $D=0$，增广矩阵 $\begin{pmatrix} -1 & 1 & 1 \\ 1 & -1 & 1 \end{pmatrix}\to\begin{pmatrix} -1 & 1 & 1 \\ 0 & 0 & 2 \end{pmatrix}$，$r(A)=1\ne r(\bar A)=2$，\textbf{无解}。
\end{itemize}
\textbf{含参数问题的三类情形的完整讨论框架}（方阵情形，$A$ 为 $n$ 阶）：

\begin{enumerate}
\item 先算 $\vert A\vert$，解出使 $\vert A\vert=0$ 的\textbf{临界参数值}（可能不止一个）；
\item 对每个临界值，把参数代回增广矩阵，用初等行变换化阶梯形，读出 $r(A)$ 与 $r(\bar A)$：
\begin{itemize}
\item 若 $r(A)\ne r(\bar A)$（阶梯形末行出现非零常数项），则\textbf{无解}；
\item 若 $r(A)=r(\bar A)=r<n$，则\textbf{有无穷多解}，继续化行最简形写出通解；
\end{itemize}
\item 对所有非临界值（$\vert A\vert\ne0$），由克拉默法则立刻得到\textbf{有唯一解}（不必求出来）。
\end{enumerate}
关键纪律是"\textbf{每个临界值都要单独回代}"。常见的失分方式是：算出 $\vert A\vert=0$ 得到 $a=1$ 或 $a=-1$，就一律写成"$a=1$ 或 $a=-1$ 时有无穷多解"——实际上这两个值里往往一个是无穷多解、另一个是无解，必须分别检验。

当系数矩阵不是方阵、或参数出现在常数列中时，$\vert A\vert$ 这条路走不通，只能直接化阶梯形，按阶梯形中"是否出现矛盾行"以及"非零行数"来分类讨论。

\begin{callwarn}{易错点}
参数位于增广矩阵中常数列时（即 $b$ 含参数），\textbf{不能}把"无解"简单归因于 $\vert A\vert=0$，必须老老实实化阶梯形看最后一列的秩贡献。

对含参数矩阵作行变换时，最常见的翻车是\textbf{用含参数的项去除分母}。例如矩阵中含 $a-1$ 这一元素，若除以 $(a-1)$ 就等于默认 $a\ne1$，把 $a=1$ 的情形整段丢了。正确做法是先把参数集中到一列，再统一消元。

"有唯一解"的条件里，$n$ 只能取\textbf{未知数个数}。许多考生写成 $m$（方程个数），在 $m\ne n$ 的题上必错。

行阶梯形化为行最简形后\textbf{务必回代验算一个特解}，这是防止整体算错的最省事手段。

\end{callwarn}

\begin{callexam}{真题考法}
含参数讨论是数二线代解答题的\textbf{头号题型}。典型设问："问 $a,b$ 取何值时方程组无解 / 有唯一解 / 有无穷多解，无穷多解时求通解。"评分细则通常把三类情形的判据、结论、通解各算一次分，\textbf{漏掉任何一种情形都会失分}。

抽象形式的考法也很常见：给出 $A$ 的具体矩阵与含参数向量 $b$，问"$Ax=b$ 无解时求参数"——由 $r(\bar A)>r(A)$ 得出矛盾行，直接锁定参数。

与后续章节的串联考法：先解 $Ax=b$ 得通解，再要求"通解中的向量都是某矩阵的特征向量"，属于拉分题，但第一步永远是"化行最简形写出通解"。

\end{callexam}

\subsection{同解与公共解}

\textbf{判别方向的一句口诀}："\textbf{谁的解集小，谁的约束就强，谁的秩就大}"。原因是解空间维数等于 $n-r$，解集越小说明维数越小，秩就越大。于是由"$Ax=\mathbf{0}$ 的解都是 $Bx=\mathbf{0}$ 的解"可以立刻判断出 $r(A)\ge r(B)$；若同时两个方程组不同解，则不等式严格，$r(A)>r(B)$。这一句口诀能在选择题里省掉大量推演。

判据的三种记忆方式也值得串起来。第一种是几何的：同解就是解集相同，两个解空间维数相同（$r(A)=r(B)$）且互相包含（一个方向的可表示性）。第二种是代数的：解集只由方程组的行向量组决定，两个方程组的解集相同当且仅当它们的行向量组等价，而"行向量组等价"又等价于"秩相同且一个能被另一个表示"，写成式子就是 $r(A)=r(B)=r\begin{pmatrix}A\\B\end{pmatrix}$。第三种是实用的：遇到具体题目就用三秩判据，遇到抽象证明就用"行向量组等价"这个等价说法。

\subsubsection{同解方程组、公共解与解的反求 $\bigstar$$\bigstar$$\bigstar$$\bigstar$}

\textbf{定义（公共解）}：方程组 $(I)$ 与 $(II)$ \textbf{公共解}就是两组解集的交集，即同时满足两组的向量。

\textbf{求公共解的三种方法（按题目给的条件选）}：

\begin{enumerate}
\item \textbf{两组方程都已知 $\to$ 联立}。把 $(I)$ 与 $(II)$ 的所有方程拼成一个大方程组，统一作行变换求解（此时"公共解"就是大方程组的解集）。
\item \textbf{$(I)$ 已知、$(II)$ 用基础解系给出 $\to$ 代入法}。先写出 $(I)$ 的通解 $x=k_1\xi_1+\cdots+k_s\xi_s$，再把它代入 $(II)$ 的方程，解出 $k_i$ 之间满足的约束关系，代回通解即得公共解。
\item \textbf{两组都用基础解系给出 $\to$ 待定系数法}。设公共解为 $\gamma$，则存在两组系数使
\end{enumerate}
\[\gamma=k_1\alpha_1+\cdots+k_s\alpha_s=l_1\beta_1+\cdots+l_t\beta_t\]

即 $k_1\alpha_1+\cdots+k_s\alpha_s-l_1\beta_1-\cdots-l_t\beta_t=\mathbf{0}$，解出 $k_i$ 与 $l_j$ 即得 $\gamma$。

\textbf{定义（同解）}：若方程组 $(I)$ 与 $(II)$ 的解集\textbf{完全相同}，则称二者\textbf{同解}。

\begin{callkey}{同解与单向归属的判据（务必分清方向）}
设 $A,B$ 均为 $m\times n$ 矩阵。

\textbf{（1）齐次的同解}

\[Ax=\mathbf{0}\text{ 与 }Bx=\mathbf{0}\text{ 同解}
\iff A\text{ 与 }B\text{ 的行向量组等价}
\iff r(A)=r(B)=r\begin{pmatrix} A \\ B \end{pmatrix}\]

\textbf{（2）单向归属}：$Bx=\mathbf{0}$ 的解\textbf{都是} $Ax=\mathbf{0}$ 的解

\[\iff A\text{ 的行向量组可由 }B\text{ 的行向量组线性表示}
\iff r(B)=r\begin{pmatrix} A \\ B \end{pmatrix}\]

由此立刻得到一条重要推论：\textbf{若 $Ax=\mathbf{0}$ 的解都是 $Bx=\mathbf{0}$ 的解，则 $r(A)\ge r(B)$}。

\textbf{（3）非齐次的同解}：$Ax=\alpha$ 与 $Bx=\beta$ 同解

\[\iff (A\mid\alpha)\text{ 与 }(B\mid\beta)\text{ 的行向量组等价}
\iff r(A\mid\alpha)=r(B\mid\beta)=r\begin{pmatrix} A & \alpha \\ B & \beta \end{pmatrix}\]

\textbf{（4）三条必背的同解结论}：

\begin{itemize}
\item 设 $P$ 为\textbf{列满秩}矩阵，则 $PAx=\mathbf{0}$ 与 $Ax=\mathbf{0}$ \textbf{同解}；
\item $A^TAx=\mathbf{0}$ 与 $Ax=\mathbf{0}$ \textbf{同解}（因为 $A^TAx=0\Rightarrow x^TA^TAx=0\Rightarrow\Vert Ax\Vert^2=0\Rightarrow Ax=0$；反方向显然）；
\item 设 $A$ 为 $n$ 阶矩阵，则 $A^nx=\mathbf{0}$ 与 $A^{n+1}x=\mathbf{0}$ \textbf{同解}。
\end{itemize}
\end{callkey}

\textbf{例（单向归属 + 不同解的标准联立）}：设

\[A=\begin{pmatrix} 1 & 0 & 1 \\ 2 & 1 & 0 \\ 0 & 1 & -1 \end{pmatrix},\qquad
B=\begin{pmatrix} 1 & 0 & 1 \\ 2 & 0 & b \end{pmatrix}\]

已知 $Ax=\mathbf{0}$ 的解都是 $Bx=\mathbf{0}$ 的解，但两者\textbf{不同解}，求 $b$。

\textbf{第一步，翻译条件}。"$Ax=\mathbf{0}$ 的解都是 $Bx=\mathbf{0}$ 的解" $\iff$ $B$ 的行向量组可由 $A$ 的行向量组线性表示 $\iff r(A)=r\begin{pmatrix} A \\ B \end{pmatrix}$。又因为两者\textbf{不同解}，包含关系必须是严格的，故还须 $r(B)<r(A)$——这一步是本题的关键，只用一个条件会漏解或错解。

\textbf{第二步，先看 $A$ 是否满秩}。设 $A=\begin{pmatrix} 1 & 0 & 1 \\ 2 & 1 & 0 \\ 1 & 1 & -1 \end{pmatrix}$。因第三行 $=(2,1,0)-(1,0,1)$（即 $r_3=r_2-r_1$），三行线性相关，故 $r(A)<3$；又第 $1,2$ 行构成的二阶子式 $\begin{vmatrix}1&0\\2&1\end{vmatrix}=1\ne0$，故 $r(A)=2$。

（若题目给的是 $A$ 满秩的三阶矩阵，如 $\begin{pmatrix} 1 & 0 & 1 \\ 2 & 1 & 0 \\ 0 & 1 & -1 \end{pmatrix}$，则 $A$ 的行向量组张成整个 $\mathbb{R}^3$，任何 $B$ 的行都能被表示，$b$ \textbf{不受任何限制}——这种"参数是摆设"的情形在真题中也会出现，要敢于下"$b$ 任意"的结论。）

\textbf{第三步，用"$B$ 的行落在 $A$ 的行空间内"定 $b$}。$A$ 的行空间是二维的 $\mathrm{span}\{(1,0,1),(2,1,0)\}$。$B$ 的第 $1$ 行 $(1,0,1)$ 恰是 $A$ 的第 $1$ 行，自动满足；只需检验第 $2$ 行 $(2,0,b)$：设

\[(2,0,b)=x(1,0,1)+y(2,1,0)\]

由第二个分量得 $0=y$；由第一个分量得 $2=x+2y=x$，即 $x=2$；再由第三个分量得 $b=x+0\cdot y=2$。故 \textbf{$b=2$}。

\textbf{第四步，验算两个条件}。$b=2$ 时 $B=\begin{pmatrix} 1 & 0 & 1 \\ 2 & 0 & 2 \end{pmatrix}$，第二行是第一行的 $2$ 倍，故 $r(B)=1$；此时 $r(A)=r\begin{pmatrix} A \\ B \end{pmatrix}=2>1=r(B)$，两个条件同时满足。

代回原始矩阵可以看清几何含义：$Ax=\mathbf{0}$ 的基础解系含 $3-2=1$ 个向量 $\eta=(-1,2,1)^T$，$Bx=\mathbf{0}$ 的解空间是二维的（基础解系含 $3-1=2$ 个向量），自然有 $\mathrm{Nul}(A)\subseteq\mathrm{Nul}(B)$ 且二者不相等，完全符合题意。

\begin{callwarn}{易错点}
\textbf{把"解都是"与"同解"混为一谈}。前者是单向包含（$Ax=\mathbf{0}$ 的解集 $\subseteq$ $Bx=\mathbf{0}$ 的解集，或反之），后者是双向相等。由包含关系只能得到\textbf{秩的不等式}，得不到秩的等式。

\textbf{"秩相等"不是同解的充分条件}。$r(A)=r(B)$ 只能说明两组解空间维数相同，还需要一个方向的包含关系（或三秩相同）才能断定同解。所以"由 $r(A)=r(B)$ 推出同解"是典型错误选项。

\textbf{方向不能反}：$Bx=\mathbf{0}$ 的解都是 $Ax=\mathbf{0}$ 的解 $\iff r(B)=r\begin{pmatrix} A \\ B \end{pmatrix}$，注意左端写的是 $B$ 的秩而不是 $A$ 的秩。记口诀："\textbf{谁的解少（约束强），谁的秩大}"——$Ax=\mathbf{0}$ 的解少说明 $A$ 的约束强，$r(A)$ 大。

\textbf{矩阵的行等价与向量组等价要区分}：$Ax=\mathbf{0}$ 与 $Bx=\mathbf{0}$ 同解 $\iff$ $A,B$ 的\textbf{行向量组}等价（不是列向量组等价，也不是矩阵等价）。

\textbf{解的反求}（"已知解求方程组/求参数"）最易丢分的地方是\textbf{漏解}：由 $r(A)=r(\bar A)$ 得到的参数值往往不止一个，必须逐一回代检验是否真的满足全部的秩关系。

\end{callwarn}

\begin{callexam}{真题考法}
"同解 / 公共解"是近几年数二线代压轴题的常客，出法有三种。

$\circled{1}$ \textbf{公共解的具体计算}：给出两组方程（一组含参数），问"当 $a$ 为何值时两组有非零公共解，并求全部非零公共解"。标准流程是先求一组的通解，代入另一组得到关于自由参数的条件方程，让该方程有非零解即得参数值。

$\circled{2}$ \textbf{同解的抽象证明}："若 $Ax=\mathbf{0}$ 与 $Bx=\mathbf{0}$ 同解，证明 $r(A)=r(B)$。"用行向量组等价 + 秩的三秩判据即可，注意两条方向都要写。

$\circled{3}$ \textbf{单向归属 + 不同解的联立条件}："已知 $Ax=\mathbf{0}$ 的解都是 $Bx=\mathbf{0}$ 的解，但两者不同解，求参数。"这一步用 $r(A)=r\begin{pmatrix} A \\ B \end{pmatrix}$ 与 $r(A)>r(B)$ \textbf{两个条件联立}，缺一不可——只用一个条件会得到多余的解或漏解。

"解的反求"的另一常见形式："求一个齐次方程组使其基础解系为 $\alpha_1,\alpha_2$"。做法是设所求矩阵的每一行 $(x_1,\cdots,x_n)$ 都满足 $x^T\alpha_1=0$ 与 $x^T\alpha_2=0$，解出这个"正交条件"的通用形式，再取一组线性无关的行拼成矩阵。这本质上是求"与已知向量组正交的向量组"的极大无关组。

\end{callexam}

\paragraph{解的反求与抽象的秩信息（补强）}\mbox{}{\par}

有一类题不给方程，只给"解的结构"的信息，要求反求矩阵的秩或参数。它们的解法都依赖一句话：\textbf{解空间的维数等于 $n-r(A)$，而"基础解系含几个向量"就是解空间维数}。

\textbf{题型一：已知基础解系的向量个数求秩}。若 $A$ 为 $m\times n$ 矩阵，$Ax=\mathbf{0}$ 的基础解系含 $t$ 个向量，则 $n-r(A)=t$，即 $r(A)=n-t$。这是最直接的一类，考场上属于送分题，但要注意 $n$ 是\textbf{未知数个数}（也就是 $A$ 的列数），不是方程个数。

\textbf{题型二：由解的具体信息反求参数}。若已知某些具体向量是 $Ax=\mathbf{0}$ 的解，把它们代入即为若干个线性约束，联立可解出参数。若已知 $Ax=\mathbf{0}$ 的解空间是某条直线，则说明 $n-r(A)=1$；若是某个平面，则 $n-r(A)=2$。\textbf{把"解集是几个自由度的东西"翻译成"$n-r$ 等于几"}，是这类题的通用第一步。

\textbf{题型三：由 $Ax=b$ 的解反求 $b$}。若已知 $Ax=b$ 的通解，把通解中的特解代回 $Ax=b$ 即得 $b=A\eta^*$。若已知通解的形式但 $A$ 未知，则利用"任意两个特解之差是齐次解"把 $A$ 的信息提取出来。

\textbf{题型四：由解反求方程组本身}。已知 $Ax=\mathbf{0}$ 的基础解系为 $\xi_1,\cdots,\xi_t$，要求写出一个满足条件的矩阵 $A$。做法是：设 $A$ 的任意一行为 $(x_1,\cdots,x_n)$，该行必须与每个 $\xi_i$ 都正交（因为 $Ax=\mathbf{0}$ 意味着每个行向量与解向量作内积为零），即 $x^T\xi_i=0$（$i=1,\cdots,t$）。解这个关于 $x$ 的齐次方程组，得到行向量的一般形式，再取 $m=n-t$ 个线性无关的行拼成 $A$。所求的 $A$ 通常不唯一，题目往往只要求写出一个。

\textbf{题型五：已知解的个数反求 $r(A)$ 的界}。若 $Ax=b$ 有三个线性无关的解，则任意两个之差都是齐次解，能得到两个线性无关的齐次解，故 $n-r(A)\ge2$，即 $r(A)\le n-2$。若已知 $Ax=b$ 有无穷多解且 $A$ 是 $n$ 阶矩阵，则 $r(A)<n$。这类"界"的估计常作为证明题的第一步，与后面的特征值、二次型内容串联。

\subsubsection{方程组解的性质与线性组合（系数和为 1/0 的判据） $\bigstar$$\bigstar$$\bigstar$$\bigstar$}

\textbf{齐次方程组解的性质}：

\[\xi_1,\cdots,\xi_s\text{ 是 }Ax=\mathbf{0}\text{ 的解}\implies k_1\xi_1+\cdots+k_s\xi_s\text{ 仍是 }Ax=\mathbf{0}\text{ 的解}\quad(k_i\text{ 任意})\]

即齐次的解集对\textbf{任意线性组合封闭}，构成解空间。

\textbf{非齐次方程组解的三条性质}：

\begin{enumerate}
\item 若 $\eta_1,\eta_2$ 都是 $Ax=b$ 的解，则 $\eta_1-\eta_2$ 是\textbf{导出组} $Ax=\mathbf{0}$ 的解；
\item 若 $\eta$ 是 $Ax=b$ 的解、$\xi$ 是 $Ax=\mathbf{0}$ 的解，则 $\eta+\xi$ 是 $Ax=b$ 的解；
\item 若 $\eta_1,\eta_2,\cdots,\eta_s$ 都是 $Ax=b$ 的解，令 $k_1\eta_1+\cdots+k_s\eta_s$，则
\end{enumerate}
\[k_1+k_2+\cdots+k_s=1\implies \text{它是 }Ax=b\text{ 的解};\qquad
k_1+k_2+\cdots+k_s=0\implies \text{它是 }Ax=\mathbf{0}\text{ 的解}\]

\begin{callkey}{系数和判据（判"是不是解"的万能钥匙）}
设 $\eta_1,\cdots,\eta_s$ 都是 $Ax=b$ 的解，$\xi_1,\cdots,\xi_t$ 都是 $Ax=\mathbf{0}$ 的解。对线性组合

\[y=\sum_{i=1}^{s}k_i\eta_i+\sum_{j=1}^{t}l_j\xi_j\]

只需计算\textbf{非齐次解的系数和} $\sigma=k_1+k_2+\cdots+k_s$，就有：

\begingroup\small\setlength{\tabcolsep}{4.0pt}
\begin{longtable}{@{}L{0.3752\linewidth}L{0.5704\linewidth}@{}}
\toprule
\textbf{系数和 $\sigma$} & \textbf{结论} \\
\midrule
\endhead
$\sigma=1$ & $y$ 是 $Ax=b$ 的解 \\
$\sigma=0$ & $y$ 是 $Ax=\mathbf{0}$ 的解 \\
$\sigma\ne0,1$ & $y$ 既不是 $Ax=b$ 的解，也不是 $Ax=\mathbf{0}$ 的解 \\
\bottomrule
\end{longtable}
\endgroup

（推导只需一行：$Ay=\sum k_iA\eta_i+\sum l_jA\xi_j=\left(\sum k_i\right)b$。）

\textbf{直接推论}：若 $\eta_1,\eta_2$ 是 $Ax=b$ 的解，则 $\dfrac{\eta_1+\eta_2}{2}$ 是 $Ax=b$ 的解（系数和 $\tfrac12+\tfrac12=1$）；$\eta_1-\eta_2$ 是 $Ax=\mathbf{0}$ 的解（系数和 $1-1=0$）；而 $2\eta_1-\eta_2$ 的系数和为 $2-1=1$，故它\textbf{也是} $Ax=b$ 的解（这是"系数和为 $1$"判据的一个漂亮应用：任何系数和为 $1$ 的组合都还是非齐次解）。

\end{callkey}

\begin{callwarn}{易错点}
上表中 $\sigma$ 只统计\textbf{非齐次解}的系数，齐次解的系数\textbf{不参与}求和（它们对 $Ay$ 的贡献恒为零）。把齐次解的系数也算进去是常见的低级错误。

$\sigma=1$ 与 $\sigma=0$ 是两个\textbf{不同的}结论档位，不能合并成"$\sigma\ne0$ 就是解"。特别地 $\sigma=2$ 时 $y$ 满足 $Ay=2b$，只有当 $b=\mathbf{0}$ 时才是解——这正说明"系数和为 $1$"的判据\textbf{只对非齐次方程组有意义}。

非齐次方程组的解集\textbf{不是线性空间}（它对加法不封闭：两个解相加一般不再满足 $Ax=b$），只能说是"仿射空间"，即"一个特解平移齐次解空间"。用"解空间"这个词描述非齐次解集是概念性错误。

\end{callwarn}

\textbf{例}：设 $\eta_1=(1,0,2)^T$、$\eta_2=(0,1,1)^T$ 都是 $Ax=b$ 的解（$b\ne\mathbf{0}$，$A$ 为 $3\times3$），判断下列向量是不是 $Ax=b$ 或 $Ax=\mathbf{0}$ 的解。

\begin{itemize}
\item $2\eta_1-2\eta_2$：非齐次系数和 $2-2=0$，是 $Ax=\mathbf{0}$ 的解；
\item $3\eta_1-2\eta_2$：系数和 $3-2=1$，是 $Ax=b$ 的解；
\item $\eta_1+\eta_2$：系数和 $1+1=2\ne0,1$，两者都不是。
\end{itemize}
\textbf{解的反求（抽象解的结构）}：若已知 $Ax=b$ 的三个解 $\eta_1,\eta_2,\eta_3$，则由 $\eta_1-\eta_2,\eta_1-\eta_3$ 得到导出组的两个解；若这两个差向量线性无关且 $n-r=2$，则通解为

\[x=\eta_1+k_1(\eta_1-\eta_2)+k_2(\eta_1-\eta_3)\]

\begin{callexam}{真题考法}
这一考点是\textbf{选择题与证明题的必考区}，考法固定：

$\circled{1}$ "设 $\eta_1,\eta_2,\eta_3$ 是 $Ax=b$ 的解，判断 $k_1\eta_1+k_2\eta_2+k_3\eta_3$ 在什么条件下仍为解"——直接套系数和判据，$k_1+k_2+k_3=1$ 为非齐次的解，$=0$ 为齐次的解。

$\circled{2}$ "设 $A$ 为 $n$ 阶矩阵，$\eta_1,\eta_2,\eta_3$ 是 $Ax=b$ 的三个线性无关解，则 $r(A)=\cdots$"——由三个线性无关解可构造出两个线性无关的齐次解（$\eta_1-\eta_2$、$\eta_1-\eta_3$），故 $n-r(A)\ge2$，即 $r(A)\le n-2$。

$\circled{3}$ 与特征值串联："已知 $\eta$ 是 $Ax=b$ 的解，且 $\eta$ 是 $A$ 的特征向量"，这类题要先把 $A\eta=\lambda\eta$ 代入 $A\eta=b$ 得到 $b=\lambda\eta$，从而把 $b$ 与特征向量联系起来。

$\circled{4}$ 抽象判定的常见坑：把"$\eta_1+\eta_2$ 是解"当成正确命题（实为错，须除以 $2$ 才是）；把"$k\eta_1$ 是解"当成对任意 $k$ 成立（实为只有 $k=1$ 才行）。

\end{callexam}

\subsection*{本章易错清单}
\addcontentsline{toc}{subsection}{本章易错清单}

\begin{enumerate}
\item 齐次方程组\textbf{永远有解}（零解），所以讨论齐次方程组只问"是否只有零解"；非齐次方程组才可能无解。
\item 非齐次有解的充要条件是 $r(A)=r(\bar A)$，\textbf{不是} $r(A)=r(\bar A)=n$；后者是"有唯一解"的条件。
\item 有唯一解的条件 $r(A)=r(\bar A)=n$ 中的 $n$ 是\textbf{未知数个数}，不是方程个数 $m$。
\item $\vert A\vert=0$ 对非齐次方程组只意味着"不能断定唯一解"，必须进一步比较 $r(A)$ 与 $r(\bar A)$ 才能区分无解与无穷多解。
\item 基础解系含 $n-r(A)$ 个向量，"个数"是硬指标，多一个少一个都不行。
\item 求基础解系必须先化到\textbf{行最简形}；只化到行阶梯形容易在回代时出错。
\item 求非齐次通解时，特解令\textbf{自由变量全取 $0$}，基础解系令\textbf{自由变量取标准基向量}，两种取值方式不能混用。
\item 通解中"特解"这一项不能漏；且特解可以换成任何一个特解，答案不唯一但必须写成"特解 $+$ 齐次通解"的形式。
\item 漏写自由参数、把 $k_1\xi_1+k_2\xi_2$ 写成 $k\xi_1+k\xi_2$（两个参数必须独立），都是常见格式错误。
\item 克拉默法则的前提是"方程个数 $=$ 未知数个数且系数行列式非零"，三个条件缺一不可；$\vert A\vert=0$ 时法则失效。
\item 含参数问题中，遇到含参数的表达式\textbf{不要随便除}，否则会漏掉使该表达式为零的参数情形；每个临界值必须单独回代检验。
\item "解都是"（单向包含）与"同解"（双向相等）是两件事；由包含关系只能得秩的\textbf{不等式}，由同解才得秩的\textbf{等式}。
\item 同解的两个方程组，其\textbf{行向量组}等价（不是列向量组）；判断时要按行看。
\item 非齐次方程组的解集对加法、数乘\textbf{均不封闭}，不能称为"解空间"；只有齐次方程组的解集才是线性空间。
\item 系数和判据只统计\textbf{非齐次解}的系数：和为 $1$ 是非齐次解，和为 $0$ 是齐次解，其余既不是。
\end{enumerate}
\chapter{特征值和特征向量}
\markboth{特征值和特征向量}{特征值和特征向量}
\subsection*{本章考纲要求}
\addcontentsline{toc}{subsection}{本章考纲要求}

\begin{quote}
\textbf{考试内容}

\begin{enumerate}
\item 矩阵的特征值和特征向量的概念、性质
\item 相似矩阵的概念及性质
\item 矩阵可相似对角化的充分必要条件及相似对角矩阵
\item 实对称矩阵的特征值、特征向量及其相似对角矩阵
\end{enumerate}
\textbf{考试要求}

\begin{enumerate}
\item 理解矩阵的特征值和特征向量的概念及性质，会求矩阵的特征值和特征向量。
\item 理解相似矩阵的概念、性质及矩阵可相似对角化的充分必要条件，掌握将矩阵化为相似对角矩阵的方法。
\item 掌握实对称矩阵的特征值和特征向量的性质。
\end{enumerate}
\end{quote}
\subsection*{本章星级总表}
\addcontentsline{toc}{subsection}{本章星级总表}

\begingroup\small\setlength{\tabcolsep}{4.0pt}
\begin{longtable}{@{}L{0.4308\linewidth}L{0.1219\linewidth}L{0.1138\linewidth}L{0.1138\linewidth}L{0.1138\linewidth}@{}}
\toprule
\textbf{考点} & \textbf{星级} & \textbf{考纲要求} & \textbf{难度} & \textbf{常见题型} \\
\midrule
\endhead
特征值与特征向量的概念、性质与求法 & $\bigstar$$\bigstar$$\bigstar$$\bigstar$$\bigstar$ & 理解/会 & 中等 & 填空、解答 \\
特征多项式与特征值的性质（$\sum\lambda_i=\mathrm{tr}A$、$\prod\lambda_i=\vert A\vert$、与秩的关系） & $\bigstar$$\bigstar$$\bigstar$$\bigstar$$\bigstar$ & 理解/会 & 中等 & 选择、填空 \\
相似矩阵的概念与性质 & $\bigstar$$\bigstar$$\bigstar$$\bigstar$ & 理解 & 中等 & 选择 \\
矩阵可相似对角化的充分必要条件与判定 & $\bigstar$$\bigstar$$\bigstar$$\bigstar$$\bigstar$ & 理解/掌握 & 难 & 解答、证明 \\
实对称矩阵的特征值、特征向量及正交相似对角化 & $\bigstar$$\bigstar$$\bigstar$$\bigstar$$\bigstar$ & 掌握 & 难 & 解答 \\
由特征值、特征向量反求矩阵或参数 & $\bigstar$$\bigstar$$\bigstar$$\bigstar$ & 会 & 难 & 解答 \\
\bottomrule
\end{longtable}
\endgroup

\section{特征值与特征向量的概念与计算}

\subsection{特征值与特征向量的定义与求法}

\subsubsection{特征值与特征向量的概念、性质与求法 $\bigstar$$\bigstar$$\bigstar$$\bigstar$$\bigstar$}

\textbf{定义}：设 $A$ 是 $n$ 阶方阵，如果存在一个数 $\lambda$ 和一个 $n$ 维\textbf{非零}列向量 $\alpha$，使得

\[A\alpha=\lambda\alpha\qquad(\alpha\neq 0),\]

则称 $\lambda$ 是 $A$ 的\textbf{特征值}，$\alpha$ 是 $A$ 的\textbf{对应于（属于）特征值 $\lambda$ 的特征向量}。

把定义式移项得 $(\lambda E-A)\alpha=0$。因为 $\alpha\neq 0$，这说明齐次线性方程组 $(\lambda E-A)x=0$ 有非零解，于是得到特征值最重要的等价刻画：\textbf{$\lambda$ 是 $A$ 的特征值，当且仅当}

\[\vert\lambda E-A\vert=0 .\]

当 $\lambda$ 是未知数时，$\vert\lambda E-A\vert$ 是关于 $\lambda$ 的 $n$ 次多项式，称为 $A$ 的\textbf{特征多项式}，记作 $f(\lambda)$；方程 $\vert\lambda E-A\vert=0$ 称为 $A$ 的\textbf{特征方程}。注意 $\vert\lambda E-A\vert$ 与 $\vert A-\lambda E\vert$ 只相差因子 $(-1)^n$，它们有完全相同的根，考试中用哪一个都行；习惯上写 $\vert\lambda E-A\vert$，因为它的 $\lambda^n$ 项系数为 $1$，展开后不容易出错。

\textbf{求特征值与特征向量的两大步}（顺序不能颠倒）：

\begin{enumerate}
\item 解特征方程 $\vert\lambda E-A\vert=0$，求出 $A$ 的全部特征值。$n$ 次方程在复数域上恰有 $n$ 个根；某个根重复出现 $k$ 次，就称 $\lambda$ 的\textbf{代数重数}为 $k$（所有特征值的代数重数之和必等于 $n$）。
\item 把每一个特征值 $\lambda_i$ 代回 $(\lambda_i E-A)x=0$，求出它的一个\textbf{基础解系} $\xi_1,\xi_2,\dots,\xi_p$，则 $A$ 的属于 $\lambda_i$ 的\textbf{全部}特征向量为
\end{enumerate}
\[k_1\xi_1+k_2\xi_2+\cdots+k_p\xi_p,\]

其中 $k_1,k_2,\dots,k_p$ 是任意常数，只要求它们\textbf{不全为零}。这里的 $p=n-r(\lambda_iE-A)$ 称为 $\lambda_i$ 的\textbf{几何重数}，它就是特征子空间 $V_{\lambda_i}=\{x\mid Ax=\lambda_ix\}$ 的维数。

\textbf{三阶数字矩阵求特征值的化简技巧（"转圈化简"）}：三阶行列式直接展开会得到一个三次方程，计算量大且容易错。可以先用行列式的性质把某一行（列）中\textbf{不含 $\lambda$ 的元素}尽量消成 $0$：沿顺时针（或逆时针）方向，用不含 $\lambda$ 的行去消其他的行（列），使某一行只剩一个含 $\lambda$ 的元素，并且这个元素里含有 $\lambda$ 的公因子。这样行列式就被拆成"$\lambda$ 的公因子 $\times$ 一个二阶行列式"，三次方程降为二次方程，特征值一眼可见。

\textbf{例 1（三阶：转圈化简）}：求 $A=\begin{pmatrix}3&1&0\\1&3&0\\0&0&2\end{pmatrix}$ 的特征值与特征向量。

\[\vert\lambda E-A\vert=\begin{vmatrix}\lambda-3&-1&0\\-1&\lambda-3&0\\0&0&\lambda-2\end{vmatrix}.\]

这个矩阵本身就是"分块对角"的形状，第一、二行三列与含 $\lambda$ 的部分容易看出公因子。把第 $2$ 列加到第 $1$ 列，第 $1$ 列变成 $(\lambda-4,\lambda-4,0)^T$，提出公因子 $(\lambda-4)$ 后第 $1$ 列成为 $(1,1,0)^T$；再用第 $1$ 行消去第 $2$ 行，得

\[\vert\lambda E-A\vert=(\lambda-4)\begin{vmatrix}1&-1&0\\0&\lambda-2&0\\0&0&\lambda-2\end{vmatrix}=(\lambda-4)(\lambda-2)^2 .\]

故特征值为 $\lambda_1=4,\ \lambda_2=\lambda_3=2$（$2$ 是二重根）。

对 $\lambda=4$：$4E-A=\begin{pmatrix}1&-1&0\\-1&1&0\\0&0&2\end{pmatrix}$，解得 $x_1=x_2$ 且 $x_3=0$，取 $\xi_1=(1,1,0)^T$；

对 $\lambda=2$：$2E-A=\begin{pmatrix}-1&-1&0\\-1&-1&0\\0&0&0\end{pmatrix}$，即 $x_1+x_2=0$，$x_3$ 自由，取基础解系 $\xi_2=(1,-1,0)^T,\ \xi_3=(0,0,1)^T$。

二重根 $2$ 恰好给出两个线性无关的特征向量，这是后面判断"可相似对角化"的关键信息。校验：$4+2+2=8=\mathrm{tr}A$，$4\times2\times2=16$；而 $\vert A\vert=3\times3\times2-1\times1\times2=16$，一致。

\textbf{例 2}：求 $A=\begin{pmatrix}3&-1\\2&0\end{pmatrix}$ 的特征值与特征向量。

\[f(\lambda)=\begin{vmatrix}\lambda-3&1\\-2&\lambda\end{vmatrix}=\lambda(\lambda-3)+2=\lambda^2-3\lambda+2=(\lambda-1)(\lambda-2),\]

故 $\lambda_1=1,\ \lambda_2=2$。

对 $\lambda_1=1$：$E-A=\begin{pmatrix}-2&1\\-2&1\end{pmatrix}$，即 $2x_1-x_2=0$，取基础解系 $\xi_1=\begin{pmatrix}1\\2\end{pmatrix}$；

对 $\lambda_2=2$：$2E-A=\begin{pmatrix}-1&1\\-2&2\end{pmatrix}$，即 $x_1=x_2$，取基础解系 $\xi_2=\begin{pmatrix}1\\1\end{pmatrix}$。

所以属于 $1$ 的全部特征向量为 $k_1\begin{pmatrix}1\\2\end{pmatrix}\ (k_1\neq0)$，属于 $2$ 的全部特征向量为 $k_2\begin{pmatrix}1\\1\end{pmatrix}\ (k_2\neq0)$。

\textbf{校验}：$1+2=3=\mathrm{tr}A$，$1\times2=2=\vert A\vert$，与矩阵本身数据吻合，说明特征值没算错。这个校验习惯每次都要养成。

\textbf{特征向量的四条性质}：

\begin{itemize}
\item 特征向量必须是\textbf{非零}向量，但特征值可以是 $0$。特别地，$\lambda=0$ 是 $A$ 的特征值 $\iff\vert A\vert=0\iff A$ 不可逆 $\iff Ax=0$ 有非零解。
\item 若 $\alpha$ 是 $A$ 的属于 $\lambda$ 的特征向量，则对任意常数 $k\neq0$，$k\alpha$ 仍是属于 $\lambda$ 的特征向量；属于同一特征值的两个特征向量的非零线性组合仍是属于该特征值的特征向量。因此属于一个特征值的特征向量总是有无穷多个，答题时不能只写一个，要写"全部特征向量为 $k\xi\ (k\neq0)$"或"$k_1\xi_1+k_2\xi_2$（$k_1,k_2$ 不全为零）"。
\item 属于\textbf{不同}特征值的特征向量一定线性无关。这是相似对角化理论的基础，必须牢记。
\item 把两个不同特征值对应的特征向量作非零线性组合（即 $k_1\xi_1+k_2\xi_2$，$k_1k_2\neq0$），得到的向量一定\textbf{不是} $A$ 的特征向量。
\end{itemize}
\textbf{与 $A$ 有关的常用矩阵的特征值、特征向量}（设 $A\alpha=\lambda\alpha$，$\alpha\neq0$）：

\begingroup\small\setlength{\tabcolsep}{4.0pt}
\begin{longtable}{@{}L{0.3095\linewidth}L{0.3095\linewidth}L{0.3095\linewidth}@{}}
\toprule
\textbf{矩阵} & \textbf{特征值} & \textbf{特征向量} \\
\midrule
\endhead
$A$ & $\lambda$ & $\alpha$ \\
$kA+bE$ & $k\lambda+b$ & $\alpha$ \\
$A^m$（$m$ 为正整数） & $\lambda^m$ & $\alpha$ \\
$f(A)=a_mA^m+\cdots+a_1A+a_0E$ & $f(\lambda)$ & $\alpha$ \\
$A^{-1}$（$\vert A\vert\neq0$） & $\frac{1}{\lambda}$ & $\alpha$ \\
$A^*$（$\vert A\vert\neq0$） & $\frac{\vert A\vert}{\lambda}$ & $\alpha$ \\
$A^T$ & $\lambda$ & 特征值相同，特征向量一般不同 \\
$P^{-1}AP$（$P$ 可逆） & $\lambda$ & $P^{-1}\alpha$ \\
\bottomrule
\end{longtable}
\endgroup

这张表是本章使用频率最高的工具：题目中一旦出现 $A^{-1}$、$A^*$、$A^k$、$f(A)$，一律先查表换成 $\lambda$ 的相应表达式。

\begin{callkey}{必背}
若矩阵 $A$ 满足多项式方程 $f(A)=O$，则 $A$ 的每个特征值 $\lambda$ 都满足 $f(\lambda)=0$；但反过来不成立，$f(\lambda)=0$ 的根不一定都是 $A$ 的特征值。例如 $A=E$ 满足 $A^2-E=O$，而 $f(\lambda)=\lambda^2-1$ 的根是 $\pm1$，$A$ 的特征值却只有 $1$。 常用特例：$A^2=A\Rightarrow\lambda=0$ 或 $1$；$A^2=E\Rightarrow\lambda=1$ 或 $-1$；$A^2=O\Rightarrow\lambda=0$。

\end{callkey}

\textbf{两个高频的特殊结构}：

（1）\textbf{秩为 $1$ 的矩阵}。$r(A)=1\iff$ 存在非零列向量 $\alpha,\beta$ 使 $A=\alpha\beta^T$。此时 $A\alpha=\alpha(\beta^T\alpha)=(\beta^T\alpha)\alpha$，且 $\mathrm{tr}A=\beta^T\alpha=\alpha^T\beta$，$A^2=(\beta^T\alpha)A=(\mathrm{tr}A)A$，从而 $A^m=(\mathrm{tr}A)^{m-1}A$（$m\geq1$）。于是 $A$ 的特征值必为

\[\mathrm{tr}A,\quad 0,\quad 0,\quad \dots,\quad 0,\]

其中 $0$ 的个数为 $n-1$（重数计入）。属于 $\mathrm{tr}A$ 的特征向量是 $\alpha$，属于 $0$ 的特征向量是 $Ax=0$ 的全部非零解。当 $\mathrm{tr}A\neq0$ 时 $A$ 可相似对角化；当 $\mathrm{tr}A=0$ 时 $A$ 不可相似对角化。

（2）\textbf{每行（列）元素之和为常数 $k$}。若 $A$ 的每一行元素之和都等于 $k$，则按行求和得

\[A\begin{pmatrix}1\\1\\\vdots\\1\end{pmatrix}=k\begin{pmatrix}1\\1\\\vdots\\1\end{pmatrix},\]

即 $k$ 是 $A$ 的一个特征值，$(1,1,\dots,1)^T$ 是属于 $k$ 的特征向量。不必算行列式就白得一个特征值，这是考场上省时间的关键。若 $A$ 的每一\textbf{列}元素之和都等于 $k$，则 $A^T$ 有特征值 $k$ 与特征向量 $(1,1,\dots,1)^T$；由于 $\vert\lambda E-A^T\vert=\vert\lambda E-A\vert$，$k$ 也是 $A$ 的特征值，但此时 $A$ 属于 $k$ 的特征向量一般不等于全 $1$ 向量，需要另行求解。

\textbf{例}：$A=\begin{pmatrix}1&1&1\\1&1&1\\1&1&1\end{pmatrix}$ 每行元素之和均为 $3$，且 $r(A)=1$，由上述两条结论立得特征值为 $3,0,0$：属于 $3$ 的特征向量是 $(1,1,1)^T$；属于 $0$ 的特征向量是 $x_1+x_2+x_3=0$ 的全部非零解。

\textbf{抽象矩阵求特征值的六个入口}：题目不给具体数字时，求特征值只能靠"翻译条件"，常用的入口有六个，按出现频率排列如下。

\begin{enumerate}
\item \textbf{用定义} $A\alpha=\lambda\alpha$：题目给出某个向量与 $A$ 的乘法结果，直接把结果写成 $\lambda\alpha$ 的形式即可读出 $\lambda$。
\item \textbf{用 $Ax=0$}：只要出现"某个非零向量 $\alpha$ 使 $A\alpha=0$"，立刻得到 $\lambda=0$ 且 $\alpha$ 是属于 $0$ 的特征向量。$AB=O$ 属于这一类（$B$ 的非零列向量都是 $A$ 属于 $0$ 的特征向量）。
\item \textbf{用行和、列和}：各行为常数 $k$ 就白得特征值 $k$ 与全 $1$ 特征向量；各列元素之和为常数则要借助 $A^T$ 转换。
\item \textbf{用与 $A$ 有关的矩阵}：由 $A^*$、$A^{-1}$、$f(A)$ 的关系式列方程，再用 $\dfrac{\vert A\vert}{\lambda}$、$\dfrac{1}{\lambda}$、$f(\lambda)$ 代换。
\item \textbf{用特征多项式}：由 $\vert\lambda E-A\vert=0$ 展开求解（抽象题常结合行列式的拆分技巧）。
\item \textbf{用多项式方程}：由 $f(A)=O$ 得 $f(\lambda)=0$，解出 $\lambda$ 的\textbf{可能}取值，再用迹、行列式、秩等条件筛掉不可能的值。
\end{enumerate}
\begin{callwarn}{易错点}
\begin{enumerate}
\item 特征向量必须非零；求"全部特征向量"时只写一个基础解系向量、不写任意常数，是最常见的丢分点。
\item $A^T$ 与 $A$ 有相同的特征值，但特征向量一般\textbf{不同}，不能直接搬用。
\item 由 $A\alpha=\lambda\alpha$ 不能"两边约去 $\alpha$"（向量没有除法），只能比较分量或移项。
\item 属于同一特征值的特征向量要写成像 $k_1\xi_1+k_2\xi_2$ 的完整形式，其中系数不能同时为 $0$。
\item 用 $\vert A\vert=0$ 求参数时，别忘了 $\lambda=0$ 是特征值这件事本身也是一种条件。
\end{enumerate}
\end{callwarn}

\begin{callexam}{真题考法}
\begin{enumerate}
\item 小题（选择、填空）求二阶、三阶具体矩阵的特征值与特征向量，重点是特征向量的\textbf{完整写法}。
\item 抽象矩阵求特征值：见到 $A\alpha=0$、$AB=O$、每行元素之和为 $k$、$r(A)=1$、$f(A)=O$ 等条件，要立刻翻译成特征值信息。
\item 求 $f(A)$、$A^{-1}$、$A^*$、$A^k$ 的特征值（查上表），常与行列式、迹结合出题。
\item 陷阱集中在"$A^T$ 的特征向量"与"特征向量必须非零"两处。
\end{enumerate}
\end{callexam}

\subsection{特征值的性质}

\subsubsection{特征多项式与特征值的性质（$\sum\lambda_i=\mathrm{tr}A$、$\prod\lambda_i=\vert A\vert$、与秩的关系）$\bigstar$$\bigstar$$\bigstar$$\bigstar$$\bigstar$}

把 $A$ 的 $n$ 个特征值（按代数重数计）记为 $\lambda_1,\lambda_2,\dots,\lambda_n$，则

\[\lambda_1+\lambda_2+\cdots+\lambda_n=\mathrm{tr}A=\sum_{i=1}^{n}a_{ii},\qquad \lambda_1\lambda_2\cdots\lambda_n=\vert A\vert .\]

这两个结论来自特征多项式的展开式

\[\vert\lambda E-A\vert=\lambda^n-(\mathrm{tr}A)\lambda^{n-1}+\cdots+(-1)^n\vert A\vert,\]

只要比较 $\lambda^{n-1}$ 的系数与常数项即可得到。它们对\textbf{任何}方阵都成立，是选择、填空题里反求参数的主力工具。

\textbf{三条常用的衍生结论}：

\begin{itemize}
\item \textbf{知二求一}：三阶矩阵已知两个特征值，第三个特征值可由 $\lambda_3=\mathrm{tr}A-\lambda_1-\lambda_2$ 或 $\lambda_3=\dfrac{\vert A\vert}{\lambda_1\lambda_2}$（当 $\lambda_1\lambda_2\neq0$）求出，再用另一个等式检验。
\item \textbf{重数要计入}：$k$ 重特征值在和式、积式中要算 $k$ 次。例如三阶矩阵特征值为 $2,2,1$，则 $\mathrm{tr}A=5$，$\vert A\vert=4$。
\item \textbf{与秩的关系}：$A$ 可逆 $\iff0$ 不是特征值。若 $A$ 可相似对角化（特别地，$A$ 是实对称矩阵），则 $r(A)$ 恰等于 $A$ 的\textbf{非零特征值的个数}（按重数计）；对一般矩阵只能保证 $r(A)\geq$ 非零特征值的个数。反例：$A=\begin{pmatrix}0&1\\0&0\end{pmatrix}$ 的特征值为 $0,0$，但 $r(A)=1$。
\end{itemize}
\begin{callkey}{必背}
若 $A$ 可逆，则 $A^*$ 的特征值为 $\dfrac{\vert A\vert}{\lambda}$；再由 $A^*=\vert A\vert A^{-1}$ 与 $\vert A^*\vert=\vert A\vert^{n-1}$ 可得 $\mathrm{tr}A^*=\vert A\vert\sum\limits_{i=1}^{n}\dfrac{1}{\lambda_i}$。此外 $r(A^*)$ 只有三种取值：$r(A)=n$ 时 $r(A^*)=n$；$r(A)=n-1$ 时 $r(A^*)=1$；$r(A)<n-1$ 时 $r(A^*)=0$。

\end{callkey}

\textbf{特征多项式的三种"眼熟"形状}：考试中出现的三阶矩阵，特征多项式大多属于下面三类，见到就能少算很多。

\begin{itemize}
\item \textbf{行和型}（每一行元素之和为常数 $k$）：$k$ 必是特征值，余下两个特征值由 $\mathrm{tr}A-k$ 与 $\dfrac{\vert A\vert}{k}$（当 $k\neq0$）确定，再用两者之和、之积交叉检验。
\item \textbf{分块对角型}：直接拆成若干小块分别求特征值，各块特征值的并集就是整体特征值。
\item \textbf{秩一型}：$A=\alpha\beta^T$ 时特征值为 $\mathrm{tr}A,0,\dots,0$，直接读出，无需展开行列式。
\end{itemize}
如果矩阵不属于以上任何一类，就用前面讲的"转圈化简"，先把含 $\lambda$ 的公因子提出来，把三次方程降成二次方程。

\textbf{例 1}：设 $A$ 为 $3$ 阶矩阵，$\mathrm{tr}A=4$，$\vert A\vert=2$，且 $\lambda=1$ 是 $A$ 的二重特征值，求第三个特征值。

设第三个特征值为 $\lambda_3$。由迹得 $1+1+\lambda_3=4$，故 $\lambda_3=2$。检验行列式：$1\times1\times2=2=\vert A\vert$，两个条件同时满足，结果自洽。

\textbf{例 2}：设 $3$ 阶矩阵 $A$ 的特征值为 $1,2,3$，求 $\mathrm{tr}A$、$\vert A\vert$、$\vert A^*\vert$、$\mathrm{tr}A^*$。

$\mathrm{tr}A=1+2+3=6$，$\vert A\vert=1\times2\times3=6$；$\vert A^*\vert=\vert A\vert^{3-1}=36$；$A^*$ 的特征值为 $\dfrac{6}{1},\dfrac{6}{2},\dfrac{6}{3}$，即 $6,3,2$，故 $\mathrm{tr}A^*=6+3+2=11$。

\textbf{例 3}（由矩阵方程确定特征值）：设 $3$ 阶实对称矩阵 $A$ 满足 $\vert A\vert=2$ 且 $A^*=A-E$，求二次型 $x^TAx$ 的规范形。

设 $\lambda$ 是 $A$ 的任一特征值。$A^*$ 的特征值为 $\dfrac{\vert A\vert}{\lambda}=\dfrac{2}{\lambda}$，由 $A^*=A-E$ 得 $\dfrac{2}{\lambda}=\lambda-1$，即 $\lambda^2-\lambda-2=0$，故 $\lambda=2$ 或 $\lambda=-1$。又 $\vert A\vert$ 等于三个特征值之积且等于 $2$，在 $\{2,-1\}$ 中取三个数使积为 $2$，只能是 $2,-1,-1$。于是 $A$ 有 $1$ 个正特征值、$2$ 个负特征值，规范形为

\[y_1^2-y_2^2-y_3^2 .\]

\textbf{例 4（用特征值算抽象行列式与迹）}：设 $3$ 阶矩阵 $A$ 的特征值为 $1,-1,2$，求 $\vert A^2+E\vert$ 与 $\mathrm{tr}(A^2+E)$。

令 $g(\lambda)=\lambda^2+1$。由 $g(A)$ 的特征值为 $g(1)=2,\ g(-1)=2,\ g(2)=5$ 得

\[\vert A^2+E\vert=2\times2\times5=20,\qquad \mathrm{tr}(A^2+E)=2+2+5=9 .\]

这类题的关键认识是：只要一个矩阵能写成 $A$ 的多项式，它的全部特征值就由 $\lambda$ 的多项式给出，于是行列式、迹、秩都可以不展开矩阵而直接算出。反过来，若题目给的是 $\vert A-\lambda_0E\vert=0$ 之类的条件，也要立刻翻译成"$\lambda_0$ 是特征值"。

\textbf{由秩反推 $0$ 的重数}：若 $A$ 可对角化，则 $0$ 作为 $A$ 的特征值的重数（代数重数）等于 $n-r(A)$。这条结论在"已知 $r(A)$ 求特征值"的题目中极其常用。例如三阶矩阵 $A$ 满足 $r(A)=1$ 且 $\mathrm{tr}A=3$，则 $A$ 有三个特征值，其中 $0$ 的重数为 $3-1=2$，剩下一个就是 $\mathrm{tr}A=3$，即特征值为 $3,0,0$。又例如已知 $\vert A\vert=0$、$r(A)=n-1$ 且 $A$ 可对角化，则可断定 $0$ 是单重特征值。

\begin{callwarn}{易错点}
\begin{enumerate}
\item $\sum\lambda_i=\mathrm{tr}A$ 与 $\prod\lambda_i=\vert A\vert$ 只是\textbf{必要条件}，无论有多少个等式成立，都不能据此断定两个矩阵相似或合同。
\item 用 $\mathrm{tr}A$ 与 $\vert A\vert$ 求出的特征值若互相矛盾，说明题目数据不可能同时成立（或自己算错），此时要停下来检查，不要硬凑。
\item 只有在 $A$ 可相似对角化（含实对称矩阵）时才有 $r(A)=$ 非零特征值的个数；一般矩阵是"$\geq$"，反例见上文的 $\begin{pmatrix}0&1\\0&0\end{pmatrix}$。
\item 用 $A^*$ 的特征值 $\dfrac{\vert A\vert}{\lambda}$ 时必须保证 $\lambda\neq0$，即 $A$ 可逆；$A$ 不可逆时 $A^*$ 的特征值要另作讨论。
\item $\vert A^*\vert=\vert A\vert^{n-1}$ 中是 $n-1$ 次幂，不要写成 $n$ 次幂。
\end{enumerate}
\end{callwarn}

\begin{callexam}{真题考法}
\begin{enumerate}
\item 选择、填空高频：给出 $\mathrm{tr}A$、$\vert A\vert$、某个特征值的重数或 $r(A)$，反求参数或求某个特征值、行列式。
\item 给出 $A^*$ 与 $A$ 的关系式（如 $A^*=A-E$、$A^*+A=O$），先用 $\dfrac{\vert A\vert}{\lambda}$ 建立关于 $\lambda$ 的方程，再用 $\vert A\vert$ 定出重数。
\item 与二次型联合：先求特征值，再数正、负个数写规范形。
\item 陷阱：把重数漏计；把 $A^*$ 的特征值写成 $\dfrac{\lambda}{\vert A\vert}$。
\end{enumerate}
\end{callexam}

\section{相似矩阵与相似对角化}

\subsection{相似矩阵}

\subsubsection{相似矩阵的概念与性质 $\bigstar$$\bigstar$$\bigstar$$\bigstar$}

\textbf{定义}：设 $A,B$ 都是 $n$ 阶方阵，如果存在\textbf{可逆}矩阵 $P$，使得

\[P^{-1}AP=B,\]

则称 $A$ 与 $B$ \textbf{相似}，记作 $A\sim B$；$P$ 称为相似变换矩阵。相似是矩阵之间最强的一种关系，后面的相似对角化问题就是围绕它展开的。

\textbf{基本性质}：相似关系满足自反性（$A\sim A$，取 $P=E$）、对称性（$A\sim B\Rightarrow B\sim A$，取 $P^{-1}$）、传递性（$A\sim B,B\sim C\Rightarrow A\sim C$，取两个 $P$ 的乘积）。

\textbf{相似不变量}（相似矩阵一定相同的量）：

\begin{itemize}
\item 特征多项式 $\vert\lambda E-A\vert=\vert\lambda E-B\vert$，从而特征值（连同重数）完全相同；
\item 迹 $\mathrm{tr}A=\mathrm{tr}B$，行列式 $\vert A\vert=\vert B\vert$，秩 $r(A)=r(B)$；
\item 对任意数 $\lambda$，$r(\lambda E-A)=r(\lambda E-B)$（因为 $P^{-1}(\lambda E-A)P=\lambda E-B$），这等价于说：同一特征值的几何重数也相同。
\end{itemize}
\textbf{相似保持的运算}：若 $A\sim B$，则

\[A^m\sim B^m,\qquad f(A)\sim f(B),\qquad A^T\sim B^T,\qquad A^*\sim B^*,\qquad A^{-1}\sim B^{-1}\ (\vert A\vert\neq0),\]

以及 $\begin{pmatrix}A&O\\O&C\end{pmatrix}\sim\begin{pmatrix}B&O\\O&D\end{pmatrix}$（当 $A\sim B$ 且 $C\sim D$）。

\textbf{判定相似的思路}：相似没有简单的充分必要条件，实战按下面的顺序走。

\begin{enumerate}
\item 先查必要条件：只要 $\mathrm{tr}A\neq\mathrm{tr}B$、$\vert A\vert\neq\vert B\vert$、$r(A)\neq r(B)$、特征值（含重数）不同、或某个 $\lambda$ 处 $r(\lambda E-A)\neq r(\lambda E-B)$，都可以直接断定\textbf{不相似}。
\item 若 $A,B$ 都可相似对角化，则 $A\sim B\iff A,B$ 的特征值完全相同（二者都相似于由特征值排成的同一个对角阵，走传递性即可）。
\item 若不可对角化，只能用定义（找出可逆 $P$）或传递性（找一个"中介矩阵"）来证明；具体题目中的"中介"常常就是对角阵。
\item 抽象小题里"见到相似就当作相同"：把与 $A$ 有关的运算全部换成 $B$ 的对应量，用特征值思想处理。
\end{enumerate}
\textbf{证明两个矩阵相似的三种具体做法}：一是\textbf{用定义}，直接构造可逆矩阵 $P$ 并验证 $P^{-1}AP=B$，题目中常给出形如 $AP=PB$ 或 $P^{-1}AP=B$ 的提示；二是\textbf{用传递性}，分别把 $A,B$ 对角化到同一个对角阵 $\Lambda$（或同一个规范形），由 $A\sim\Lambda$、$B\sim\Lambda$ 得 $A\sim B$，这是最常用的写法；三是\textbf{用中介矩阵}，例如由 $A\sim B$ 推出 $f(A)\sim f(B)$，把要证的相似关系转化成两个多项式的相似。注意第三种做法要求题目中能自然地构造出中介矩阵，不能凭空捏造。

\textbf{例 1}：设 $A\sim\mathrm{diag}(1,2,3)$，$B=(A-E)(A-2E)(A-3E)$，求 $r(B)$。

由 $f(A)\sim f(\Lambda)$ 得 $B\sim(\Lambda-E)(\Lambda-2E)(\Lambda-3E)=O$，故 $r(B)=0$，即 $B=O$。

\textbf{例 2（必要条件不能反推相似）}：$A=\begin{pmatrix}1&0\\0&1\end{pmatrix}$ 与 $B=\begin{pmatrix}1&1\\0&1\end{pmatrix}$ 的迹都是 $2$、行列式都是 $1$、秩都是 $2$，但它们\textbf{不相似}：$A$ 可对角化，而 $B$ 的特征值是二重的 $1$，$r(E-B)=r\begin{pmatrix}0&-1\\0&0\end{pmatrix}=1$，几何重数 $2-1=1<2$，$B$ 不可对角化；可对角化的矩阵不可能与非可对角化矩阵相似。

\textbf{三条常用补充结论}：

\begin{itemize}
\item $AB$ 与 $BA$ 有\textbf{相同}的特征值；若 $A$ 可逆，则还有 $AB\sim BA$（因为 $AB=A(BA)A^{-1}$）。
\item 若 $AB=aA+bB$ 且 $ab\neq0$，则 $(A-bE)(B-aE)=abE$，于是 $A-bE$ 与 $B-aE$ 都可逆且乘积为数量矩阵，可推出 $AB=BA$。
\item 若 $AB=BA$ 且 $A$ 有 $n$ 个互不相同的特征值，则 $A$ 的每个特征向量都是 $B$ 的特征向量，从而存在同一个可逆矩阵 $P$ 使 $P^{-1}AP$ 与 $P^{-1}BP$ \textbf{同时}为对角阵。
\end{itemize}
\begin{callwarn}{易错点}
\begin{enumerate}
\item 相似用的是 $P^{-1}AP$，合同用的是 $C^TAC$，绝不可混用。
\item "迹、行列式、秩相同"只是必要条件，不能作为相似的判据（上文例 2 即是反例）。
\item 若 $A\sim B$，则 $A^{-1}\sim B^{-1}$ 要求 $A$ 可逆；$\vert A\vert=0$ 时不能谈 $A^{-1}$ 的相似。
\item $A^T$ 与 $A$ 的特征值相同，但 $A^T$ 与 $A$ 一般\textbf{不相似}。
\item 相似变换矩阵 $P$ 不唯一，但必须可逆；不要写成 $P^{-1}AP=A$。
\end{enumerate}
\end{callwarn}

\begin{callexam}{真题考法}
\begin{enumerate}
\item 给"$A$ 与 $B$ 相似"，用 $\mathrm{tr}A=\mathrm{tr}B$、$\vert A\vert=\vert B\vert$ 列方程反求参数（常是解答题的第一问）。
\item 判断关于相似性质的命题正误：$A^T\sim B^T$、$A^*\sim B^*$、$A+A^{-1}\sim B+B^{-1}$ 都成立；而 $A+A^T$ 与 $B+B^T$ 一般不相似（$A+A^T$ 与 $A$ 之间没有直接的相似关系）。
\item 证明两个矩阵相似：找 $P$ 用定义，或都对角化到同一个 $\Lambda$ 走传递性。
\item 抽象题用"相似当作相同"处理，把待求量全部换成对角阵上的量。
\end{enumerate}
\end{callexam}

\subsection{相似对角化}

\subsubsection{矩阵可相似对角化的充分必要条件与判定 $\bigstar$$\bigstar$$\bigstar$$\bigstar$$\bigstar$}

\textbf{定义}：若存在可逆矩阵 $P$ 使 $P^{-1}AP=\Lambda$（$\Lambda$ 为对角阵），则称 $A$ \textbf{可相似对角化}，$\Lambda$ 称为 $A$ 的相似对角矩阵。

把 $P$ 按列分块写成 $P=(\alpha_1,\alpha_2,\dots,\alpha_n)$，则 $AP=P\Lambda$ 等价于

\[A\alpha_j=\lambda_j\alpha_j\qquad(j=1,2,\dots,n).\]

这说明 $P$ 的每一列必须是 $A$ 的特征向量，$\Lambda$ 的第 $j$ 个对角元就是 $\alpha_j$ 对应的特征值。要 $P$ 可逆，就必须要求这 $n$ 个特征向量线性无关。于是得到

\begin{callkey}{必背}
$n$ 阶矩阵 $A$ 可相似对角化 $\iff A$ 有 $n$ 个线性无关的特征向量 $\iff$ 对 $A$ 的每一个特征值，其\textbf{几何重数等于代数重数}。具体地说，若 $\lambda_0$ 是 $k$ 重特征值，则 $A$ 可对角化要求

\[n-r(\lambda_0E-A)=k .\]

\end{callkey}

\textbf{充分条件}（满足其一即可断定可对角化）：

\begin{itemize}
\item $A$ 有 $n$ 个互不相同的特征值（此时 $n$ 个特征向量自动线性无关，因为属于不同特征值的特征向量线性无关）；
\item $A$ 是实对称矩阵（下一考点详述）；
\item $A^2=A$（幂等）、$A^2=E$（对合）；
\item $r(A)=1$ 且 $\mathrm{tr}A\neq0$；
\item $A$ 的每个特征值的代数重数都是 $1$。
\end{itemize}
\textbf{必要条件}：若 $A\sim\Lambda$，则 $r(A)=r(\Lambda)=$ 非零特征值的个数，且 $\mathrm{tr}A=\mathrm{tr}\Lambda$，$\vert A\vert=\vert\Lambda\vert$。

\textbf{解题步骤}（求可逆 $P$ 使 $P^{-1}AP=\Lambda$）：

\begin{enumerate}
\item 求 $A$ 的全部特征值 $\lambda_1,\dots,\lambda_n$，并分清哪些是重根、各是几重；
\item 对每个特征值解 $(\lambda E-A)x=0$，求基础解系；若某个 $k$ 重特征值只给出少于 $k$ 个线性无关的解，则 $A$ \textbf{不可}对角化，题目一般不会这样设问；
\item 把所有线性无关的特征向量按列排成 $P$，把对应的特征值按\textbf{同样的顺序}放到 $\Lambda$ 的对角线上，则 $P^{-1}AP=\Lambda$。
\end{enumerate}
\textbf{典型不可对角化的例子}：$A=\begin{pmatrix}3&1\\0&3\end{pmatrix}$。$\lambda=3$ 是二重特征值，而 $r(3E-A)=r\begin{pmatrix}0&-1\\0&0\end{pmatrix}=1$，几何重数 $=2-1=1<2$，故 $A$ 不可对角化。一般地，主对角线元素全部相同、而非对角线上还有非零元素的上（下）三角矩阵，一定不可相似对角化。

\textbf{例}：承上节 $A=\begin{pmatrix}3&-1\\2&0\end{pmatrix}$，取 $P=\begin{pmatrix}1&1\\2&1\end{pmatrix}$，$\Lambda=\begin{pmatrix}1&0\\0&2\end{pmatrix}$，则

\[AP=\begin{pmatrix}3&-1\\2&0\end{pmatrix}\begin{pmatrix}1&1\\2&1\end{pmatrix}=\begin{pmatrix}1&2\\2&2\end{pmatrix},\qquad P\Lambda=\begin{pmatrix}1&1\\2&1\end{pmatrix}\begin{pmatrix}1&0\\0&2\end{pmatrix}=\begin{pmatrix}1&2\\2&2\end{pmatrix},\]

两边相等，故 $P^{-1}AP=\Lambda$。

\textbf{对角化的最大用途：求方阵的高次幂}。由 $A=P\Lambda P^{-1}$ 得 $A^m=P\Lambda^mP^{-1}$，而对角阵的幂只需把对角元取 $m$ 次幂。承上例，$P^{-1}=\begin{pmatrix}-1&1\\2&-1\end{pmatrix}$，于是

\[A^m=\begin{pmatrix}1&1\\2&1\end{pmatrix}\begin{pmatrix}1&0\\0&2^m\end{pmatrix}\begin{pmatrix}-1&1\\2&-1\end{pmatrix}=\begin{pmatrix}2^{m+1}-1&1-2^m\\ 2^{m+1}-2&2-2^m\end{pmatrix}.\]

代入 $m=1$ 得 $\begin{pmatrix}3&-1\\2&0\end{pmatrix}$，与 $A$ 一致，说明结果正确。

\textbf{证明可对角化的三条常用路线}：

\begin{enumerate}
\item \textbf{证有 $n$ 个互异特征值}：最简单的一条路，常用于含参数的证明题。
\item \textbf{证几何重数等于代数重数}：对重根 $\lambda_0$ 直接计算 $n-r(\lambda_0E-A)$ 并说明它等于重数。
\item \textbf{用矩阵方程}：若 $(A-aE)(A-bE)=O$ 且 $a\neq b$，则 $A$ 可对角化。理由如下：由 $(b-a)E=(A-aE)-(A-bE)$ 及秩的不等式得 $r(A-aE)+r(A-bE)\geq r((b-a)E)=n$；又由 $(A-aE)(A-bE)=O$ 知 $(A-bE)$ 的列向量都落在 $N(A-aE)$ 中，故 $r(A-bE)\leq n-r(A-aE)$，即 $r(A-aE)+r(A-bE)\leq n$。两者合起来得 $r(A-aE)+r(A-bE)=n$，即两个特征子空间的维数之和为 $n$，从而 $A$ 有 $n$ 个线性无关的特征向量。特例：$A^2=A$ 时取 $a=0,b=1$，得 $r(A)+r(A-E)=n$。
\end{enumerate}
\textbf{分块对角阵的可对角化}：若 $A=\begin{pmatrix}A_1&O\\O&A_2\end{pmatrix}$，则 $A$ 可对角化 $\iff A_1$ 与 $A_2$ 都可对角化。

\textbf{两个重要应用}：

（1）\textbf{求方阵的高次幂与数列通项}。凡是遇到 $A^m$、$f(A)$ 或由线性递推关系确定数列的问题，只要矩阵可对角化，就一律走 $A=P\Lambda P^{-1}$、$A^m=P\Lambda^mP^{-1}$ 这条路。对于递推式 $x_{n+1}=ax_n+by_n,\ y_{n+1}=cx_n+dy_n$，把它写成 $\begin{pmatrix}x_{n+1}\\y_{n+1}\end{pmatrix}=A\begin{pmatrix}x_n\\y_n\end{pmatrix}$，则 $\begin{pmatrix}x_n\\y_n\end{pmatrix}=A^n\begin{pmatrix}x_0\\y_0\end{pmatrix}$，问题化为求 $A^n$。

（2）\textbf{含参数判断可对角化}。题目若说"$A$ 有三个线性无关的特征向量"，就等价于说 $A$ 可对角化：先求出含参数的特征值，找出重根，再令重根处的 $n-r(\lambda E-A)$ 等于重数，解出参数。特别地，若重根 $\lambda_0$ 是二重根，条件是 $r(\lambda_0E-A)=n-2$（对三阶矩阵即 $r(\lambda_0E-A)=1$）。这种题几乎每年都以小题或解答第一问的形式出现。

\textbf{对角化与"同时对角化"}：若 $A,B$ 都可对角化，且满足 $AB=BA$ 与"$A$ 有 $n$ 个互异特征值"这两个条件之一，则存在同一个可逆矩阵 $P$ 使 $P^{-1}AP$ 与 $P^{-1}BP$ 同时为对角阵。做题时如果要把两个矩阵同时对角化，先验证它们可交换，再验证其中一个的特征值互异（或直接用公共特征向量构造 $P$）。

\textbf{例（含参数判断可对角化）}：设 $A=\begin{pmatrix}2&0&0\\1&a&1\\0&0&2\end{pmatrix}$，$a$ 取何值时 $A$ 可相似对角化？

按第 $1$ 行展开特征多项式得 $\vert\lambda E-A\vert=(\lambda-2)^2(\lambda-a)$，故特征值为 $2,2,a$。

当 $a\neq2$ 时，$2$ 是二重特征值、$a$ 是单根。此时

\[2E-A=\begin{pmatrix}0&0&0\\-1&2-a&-1\\0&0&0\end{pmatrix},\qquad r(2E-A)=1,\]

故 $2$ 的几何重数为 $3-1=2$，等于它的代数重数，单根 $a$ 自然提供 $1$ 个线性无关的特征向量，所以 $A$ 可对角化。

当 $a=2$ 时，$2$ 是三重特征值，而 $2E-A=\begin{pmatrix}0&0&0\\-1&0&-1\\0&0&0\end{pmatrix}$ 的秩仍为 $1$，几何重数 $3-1=2<3$，故 $A$ 不可对角化。

综上，\textbf{$a\neq2$ 时 $A$ 可相似对角化}。

需要特别提醒的是：\textbf{几何重数一定不超过代数重数}，即 $k$ 重特征值至多有 $k$ 个线性无关的特征向量，且至少有 $1$ 个。所以"能不能对角化"这件事只可能卡在重根上：单根永远不缺特征向量，重根才需要逐个检验。

\begin{callwarn}{易错点}
\begin{enumerate}
\item "有重特征值"不等于"不可对角化"。例如 $E$ 的特征值全是 $1$（$n$ 重），但 $E$ 本身就是对角阵；实对称矩阵、幂等矩阵、对合矩阵也都有重根却一定可对角化。
\item 需要的是 $n$ 个\textbf{线性无关}的特征向量，而不是 $n$ 个特征向量（同一特征值下的向量有无穷多个，但线性无关的最多只有 $k$ 个）。
\item $P$ 的列顺序必须与 $\Lambda$ 对角元的顺序\textbf{一一对应}，顺序错了整题就错。
\item 判断时不要只看重数，要看 $n-r(\lambda E-A)$ 是否等于重数。
\item $P$ 与 $\Lambda$ 都不唯一：$\Lambda$ 在不计对角元次序时唯一（由特征值决定），$P$ 可以换列、乘非零倍数，或在同一特征子空间内换基。
\end{enumerate}
\end{callwarn}

\begin{callexam}{真题考法}
\begin{enumerate}
\item 选择：判断给定的具体矩阵或抽象条件（如"有 $n$ 个两两线性无关的特征向量"、"属于不同特征值的特征向量相互正交"）能否推出可对角化，注意分清充分、必要、充要。
\item 解答第一问：判断 $A$ 能否对角化并说明理由；第二问求可逆矩阵 $P$ 与 $\Lambda$。
\item 用对角化求 $A^m$，或把递推关系写成矩阵形式后求通项。
\item 证明题：含参数的矩阵"有三个线性无关的特征向量"，据此列方程求参数。
\item 若两个矩阵相似且都可对角化，则只需比较特征值即可判定相似，不必真的求 $P$。
\end{enumerate}
\end{callexam}

\section{实对称矩阵与反求问题}

\subsection{实对称矩阵的正交相似对角化}

\subsubsection{实对称矩阵的特征值、特征向量及正交相似对角化 $\bigstar$$\bigstar$$\bigstar$$\bigstar$$\bigstar$}

设 $A$ 是 $n$ 阶\textbf{实对称矩阵}（即 $A^T=A$）。它比一般矩阵"好"得多，拥有下面三条性质：

\begin{enumerate}
\item \textbf{特征值全是实数}，且 $A$ 一定有 $n$ 个实特征值（按重数计）。
\item \textbf{属于不同特征值的特征向量必正交}（不只是线性无关）：若 $\lambda_1\neq\lambda_2$，$\xi_1,\xi_2$ 分别是属于它们的特征向量，则 $\xi_1^T\xi_2=0$。
\item \textbf{一定可以正交相似对角化}：存在\textbf{正交矩阵} $Q$，使得
\end{enumerate}
\[Q^{-1}AQ=Q^TAQ=\Lambda=\mathrm{diag}(\lambda_1,\dots,\lambda_n).\]

并且 $k$ 重特征值恰有 $k$ 个线性无关的特征向量，即 $n-r(\lambda_0E-A)=k$（几何重数等于代数重数）。

反过来也成立：\textbf{如果实矩阵 $A$ 可正交相似对角化，那么 $A$ 必是实对称矩阵}。也就是"实对称"与"可正交相似对角化"互为充要条件。

\textbf{求正交矩阵 $Q$ 的四步法}：

\begin{enumerate}
\item 解特征方程 $\vert\lambda E-A\vert=0$，求全部特征值；
\item 对每个特征值解 $(\lambda E-A)x=0$，求基础解系（合起来得到 $n$ 个线性无关的特征向量）；
\item \textbf{正交化}：只需对\textbf{同一特征值}的多个特征向量作施密特正交化（不同特征值的特征向量已经自动正交，不要再动它）；
\item \textbf{单位化}：把所有正交化后的向量化为单位向量 $\gamma_1,\dots,\gamma_n$，按列拼成 $Q=(\gamma_1,\dots,\gamma_n)$（列的顺序与 $\Lambda$ 对角元一致），则 $Q^TAQ=Q^{-1}AQ=\Lambda$。
\end{enumerate}
施密特正交化与单位化的公式（已有线性无关向量 $\eta_1,\eta_2,\eta_3$）：

\[\beta_1=\eta_1,\qquad \beta_2=\eta_2-\frac{(\eta_2,\beta_1)}{(\beta_1,\beta_1)}\beta_1,\qquad \beta_3=\eta_3-\frac{(\eta_3,\beta_1)}{(\beta_1,\beta_1)}\beta_1-\frac{(\eta_3,\beta_2)}{(\beta_2,\beta_2)}\beta_2,\]

\[\gamma_i=\frac{1}{\Vert\beta_i\Vert}\beta_i\qquad(i=1,2,3).\]

其中 $(\ ,\ )$ 表示内积，$\Vert\beta_i\Vert=\sqrt{(\beta_i,\beta_i)}$。

\textbf{例}：设 $A=\begin{pmatrix}0&1&1\\1&0&1\\1&1&0\end{pmatrix}$，求正交矩阵 $Q$ 使 $Q^TAQ$ 为对角阵。

$A$ 的各行元素之和均为 $2$，由行和结论立刻得 $\lambda=2$ 是特征值，$\eta_3=(1,1,1)^T$ 是对应特征向量。再算特征多项式：把第 $2,3$ 行加到第 $1$ 行，第 $1$ 行变成 $(\lambda-2,\lambda-2,\lambda-2)$，提出公因子 $\lambda-2$ 后第 $1$ 行全为 $1$，再用它消去第 $2,3$ 行，得

\[\vert\lambda E-A\vert=(\lambda-2)(\lambda+1)^2,\]

故特征值为 $\lambda_1=\lambda_2=-1,\ \lambda_3=2$。

对 $\lambda=-1$：$-E-A=\begin{pmatrix}-1&-1&-1\\-1&-1&-1\\-1&-1&-1\end{pmatrix}$，即 $x_1+x_2+x_3=0$，取基础解系 $\eta_1=(1,-1,0)^T,\ \eta_2=(1,0,-1)^T$。

对 $\lambda=2$：$2E-A=\begin{pmatrix}2&-1&-1\\-1&2&-1\\-1&-1&2\end{pmatrix}$，解得 $x_1=x_2=x_3$，取 $\eta_3=(1,1,1)^T$。

$\eta_1,\eta_2$ 都与 $\eta_3$ 正交，只需把 $\eta_1,\eta_2$ 正交化：

\[\beta_1=\eta_1=(1,-1,0)^T,\qquad \beta_2=\eta_2-\frac{(\eta_2,\beta_1)}{(\beta_1,\beta_1)}\beta_1=(1,0,-1)^T-\frac12(1,-1,0)^T=\left(\frac12,\frac12,-1\right)^T .\]

单位化得

\[\gamma_1=\frac{1}{\sqrt2}\begin{pmatrix}1\\-1\\0\end{pmatrix},\qquad \gamma_2=\frac{1}{\sqrt6}\begin{pmatrix}1\\1\\-2\end{pmatrix},\qquad \gamma_3=\frac{1}{\sqrt3}\begin{pmatrix}1\\1\\1\end{pmatrix}.\]

取 $Q=(\gamma_1,\gamma_2,\gamma_3)$，则 $Q^TAQ=\mathrm{diag}(-1,-1,2)$。可以核对 $Q^TQ=E$，并用 $A\gamma_2=-\gamma_2$ 检验。

\textbf{三个实战技巧}：

\begin{itemize}
\item \textbf{知二求一}：若已知两个线性无关的特征向量 $\xi_1,\xi_2$ 属于同一特征值，要求另一个特征值对应的特征向量，只需解方程组 $\begin{pmatrix}\xi_1^T\\ \xi_2^T\end{pmatrix}x=0$，得到的非零解的方向就是所需特征向量。
\item \textbf{知一求二}：若已知 $\lambda=2$ 的一个特征向量是 $(1,1,1)^T$，而 $2$ 是单根、$1$ 是二重根，则属于 $1$ 的两个特征向量只需与 $(1,1,1)^T$ 正交，即解 $x_1+x_2+x_3=0$，取两个线性无关的解即可。
\item \textbf{谱分解}：若 $\gamma_1,\dots,\gamma_n$ 是 $A$ 的标准正交特征向量，则
\end{itemize}
\[A=\lambda_1\gamma_1\gamma_1^T+\lambda_2\gamma_2\gamma_2^T+\cdots+\lambda_n\gamma_n\gamma_n^T .\]

这条公式在"反求实对称矩阵"时非常好用，可以完全避免求逆矩阵。

\textbf{与二次型的衔接（下一章要用）}：由 $Q^TAQ=\Lambda$ 可知，实对称矩阵 $A$ 的\textbf{正特征值的个数}与\textbf{负特征值的个数}（都按重数计）是由 $A$ 唯一决定的，它们分别称为 $A$ 对应的二次型的正惯性指数与负惯性指数；$A$ 的秩等于非零特征值的个数。于是只要把 $A$ 的特征值全部求出，就能立刻写出二次型的规范形与正、负惯性指数，从而判断正定性、判断两个实对称矩阵是否合同。这条桥梁把本章与下一章连成一体，是数二线代解答题最常见的组合考法。

\textbf{例（已知 $Q$ 的一列反求参数）}：设 $A=\begin{pmatrix}0&-1&4\\-1&3&a\\4&a&0\end{pmatrix}$ 是实对称矩阵，正交矩阵 $Q$ 使 $Q^TAQ$ 为对角阵，且 $Q$ 的第 $1$ 列为 $\dfrac{1}{\sqrt6}(1,2,1)^T$，求 $a$ 与相应的特征值。

因为 $Q$ 的第 $1$ 列必须是 $A$ 的特征向量，所以 $A\begin{pmatrix}1\\2\\1\end{pmatrix}=\lambda\begin{pmatrix}1\\2\\1\end{pmatrix}$。先看第 $1$ 个分量：$0\times1+(-1)\times2+4\times1=2$，故 $\lambda=2$（否则第一个分量应当是 $\lambda\cdot1$）。再看第 $2$ 个分量：$-1+6+a=5+a$，应等于 $2\times2=4$，得 $a=-1$。最后用第 $3$ 个分量检验：$4+2a+0=4-2=2$，恰等于 $2\times1=2$，自洽。

这类题的要害是：\textbf{正交矩阵的每一列都是 $A$ 的特征向量}，所以不必先算特征多项式，直接用 $A\gamma=\lambda\gamma$ 列方程最快。

\begin{callwarn}{易错点}
\begin{enumerate}
\item 题目要求"正交矩阵 $Q$"时，既要正交化又要\textbf{单位化}；只正交不单位化是最常见的丢分点。
\item 只需对\textbf{同一特征值}的特征向量正交化；若把不同特征值的特征向量也拿去施密特正交化，会破坏"它们是特征向量"这一性质，答案就错了。
\item $Q$ 不唯一：特征向量的正负号、重根子空间中基的选取、列的顺序都可以变，但都必须满足 $Q^TAQ=\Lambda$，且列与 $\Lambda$ 的对角元一一对应。
\item 只有 $A^T=A$ 才能真正做到"不同特征值的特征向量必正交"，非对称矩阵没有这个性质。
\item 用"行和 $=$ 特征值"技巧时，只对\textbf{每行}元素之和为常数有效；若是每列之和为常数，得到的特征向量不是全 $1$ 向量。
\end{enumerate}
\end{callwarn}

\begin{callexam}{真题考法}
\begin{enumerate}
\item 数二解答题的常客："求正交矩阵 $Q$，使 $Q^TAQ$ 为对角阵"，通常与下一章二次型标准化合并考查。
\item 变式：题目直接给出 $Q$ 的某一列是某个向量，要求先由"它必须是特征向量"反求矩阵中的参数，再求完整的 $Q$。此时不要急着算特征多项式，先用 $A\gamma=\lambda\gamma$ 列方程。
\item 由"实对称矩阵二重特征值 $+$ 一个已知特征向量"出发，用正交性"知一求二"。
\item 陷阱：把非对称矩阵当实对称处理；忘记单位化；重根的特征向量忘记正交化。
\end{enumerate}
\end{callexam}

\subsection{反求矩阵}

\subsubsection{由特征值、特征向量反求矩阵或参数 $\bigstar$$\bigstar$$\bigstar$$\bigstar$}

这类题的核心是\textbf{条件翻译}：题目给出的每一句话，都要翻译成"某个向量是某个特征值的特征向量"。常见的翻译如下。

\begingroup\small\setlength{\tabcolsep}{4.0pt}
\begin{longtable}{@{}L{0.3359\linewidth}L{0.6098\linewidth}@{}}
\toprule
\textbf{题目条件} & \textbf{翻译结果} \\
\midrule
\endhead
$A\alpha=0$，$\alpha\neq0$ & $0$ 是 $A$ 的特征值，$\alpha$ 是属于 $0$ 的特征向量 \\
$AB=O$ & $B$ 的每一个非零列向量都是 $A$ 的属于 $0$ 的特征向量 \\
$AB=C$，把 $B,C$ 按列分块后 $C$ 的第 $i$ 列 $=\lambda_i\cdot B$ 的第 $i$ 列 & $B$ 的第 $i$ 列是 $A$ 的属于 $\lambda_i$ 的特征向量 \\
$A$ 的每行元素之和均为 $k$ & $k$ 是特征值，$(1,1,\dots,1)^T$ 是属于 $k$ 的特征向量 \\
$A$ 与 $B$ 相似 & 特征值（含重数）、迹、行列式、秩全部相同 \\
$\vert A\vert=0$ & $0$ 是 $A$ 的特征值 \\
$r(A)=k$ 且 $A$ 可对角化 & $A$ 恰有 $k$ 个非零特征值（按重数计） \\
\bottomrule
\end{longtable}
\endgroup

\textbf{反求参数的套路}：若已知 $\alpha$ 是 $A$ 的属于 $\lambda$ 的特征向量，就直接把 $A\alpha=\lambda\alpha$ 写成分量等式，得到一个方程组，从中解出参数与 $\lambda$。

\textbf{例 1}：设 $A=\begin{pmatrix}1&b\\0&2\end{pmatrix}$ 的一个特征向量为 $(1,1)^T$，求 $b$ 与对应的特征值。

由 $A\begin{pmatrix}1\\1\end{pmatrix}=\begin{pmatrix}1+b\\2\end{pmatrix}=\lambda\begin{pmatrix}1\\1\end{pmatrix}$，比较第二个分量得 $\lambda=2$，再比较第一个分量得 $1+b=2$，故 $b=1$，对应特征值为 $2$。

\textbf{反求矩阵的套路}：

\begin{itemize}
\item 若已知 $A$ 可对角化，特征值 $\lambda_1,\dots,\lambda_n$ 与对应的线性无关特征向量 $\alpha_1,\dots,\alpha_n$，则取 $P=(\alpha_1,\dots,\alpha_n)$，$\Lambda=\mathrm{diag}(\lambda_1,\dots,\lambda_n)$，由
\end{itemize}
\[A=P\Lambda P^{-1}\]

回代即可（这一步往往还要算一个逆矩阵）。

\begin{itemize}
\item 若 $A$ 是\textbf{实对称矩阵}，则已知标准正交特征向量 $\gamma_i$ 后可直接用谱分解
\end{itemize}
\[A=\lambda_1\gamma_1\gamma_1^T+\lambda_2\gamma_2\gamma_2^T+\cdots+\lambda_n\gamma_n\gamma_n^T\]

把它写成矩阵形式就是 $A=Q\Lambda Q^T$（其中 $Q=(\gamma_1,\dots,\gamma_n)$）。这样写可以避免求逆矩阵，计算量大幅下降。

\begin{itemize}
\item 若只给出部分信息（例如"$A$ 是 $3$ 阶实对称矩阵，二重特征值 $\lambda=1$，$r(A)=2$"），先用迹、行列式、秩把\textbf{全部特征值}确定下来，再定特征向量。
\end{itemize}
\textbf{例 2}：设 $A$ 为 $3$ 阶实对称矩阵，$A$ 的二重特征值 $\lambda=1$ 对应的特征向量为 $\alpha_1=(1,1,0)^T,\ \alpha_2=(1,-1,0)^T$，另一个特征值为 $\lambda=-1$，求 $A$。

先求属于 $-1$ 的特征向量 $\alpha_3$：它必须与 $\alpha_1,\alpha_2$ 都正交，即满足 $\begin{pmatrix}1&1&0\\1&-1&0\end{pmatrix}x=0$，解得 $x_1=x_2=0$，$x_3$ 任意，取 $\alpha_3=(0,0,1)^T$。

把 $\alpha_1,\alpha_2,\alpha_3$ 单位化：$\gamma_1=\dfrac{1}{\sqrt2}(1,1,0)^T$，$\gamma_2=\dfrac{1}{\sqrt2}(1,-1,0)^T$，$\gamma_3=(0,0,1)^T$。用谱分解

\[A=1\cdot\gamma_1\gamma_1^T+1\cdot\gamma_2\gamma_2^T+(-1)\cdot\gamma_3\gamma_3^T =\begin{pmatrix}1&0&0\\0&1&0\\0&0&-1\end{pmatrix}.\]

代回检验：$\gamma_1,\gamma_2$ 的第三个分量为 $0$，故 $A\gamma_1=\gamma_1$，$A\gamma_2=\gamma_2$；$A\gamma_3=-\gamma_3$，与题意完全一致。

\textbf{例 3}（由条件反求特征值）：设 $3$ 阶矩阵 $A$ 满足 $A^2=A$ 且 $r(A)=2$，求 $A$ 的特征值。

由 $A^2=A$ 知 $A$ 的特征值只可能是 $0$ 或 $1$，且 $A$ 必可对角化（幂等矩阵）。对角化后 $r(A)$ 等于非零特征值的个数，故 $A$ 恰有两个特征值 $1$、一个特征值 $0$，即特征值为 $1,1,0$（进而 $\mathrm{tr}A=2$）。

\textbf{例 4（由相似反求参数）}：设 $A=\begin{pmatrix}2&0\\0&a\end{pmatrix}$ 与 $B=\begin{pmatrix}1&1\\0&2\end{pmatrix}$ 相似，求 $a$。

$B$ 是上三角矩阵，其特征值为对角线元素 $1$ 与 $2$；$A$ 是对角阵，特征值为 $2$ 与 $a$。相似矩阵的特征值（连同重数）完全相同，故 $\{2,a\}=\{1,2\}$，从而 $a=1$。检验：此时 $\mathrm{tr}A=3=\mathrm{tr}B$，$\vert A\vert=2=\vert B\vert$，条件自洽。

\textbf{反求问题的完整流程}可以概括成四句话：先由"哪些向量是特征向量"确定特征向量的方向（个数必须够 $n$ 个且线性无关）；再由"迹、行列式、秩、相似关系"确定特征值；实对称情形用正交性把缺的特征向量补齐；最后用 $A=P\Lambda P^{-1}$ 或谱分解写出矩阵并回代检验。

\begin{callwarn}{易错点}
\begin{enumerate}
\item 由特征值、特征向量反求矩阵的前提是这 $n$ 个特征向量\textbf{线性无关}（一般矩阵还要求 $A$ 可对角化）；实对称矩阵总有 $n$ 个线性无关的特征向量，可以直接用。
\item $A=P\Lambda P^{-1}$ 与 $A=Q\Lambda Q^T$ 中的顺序必须严格对应；用谱分解时每个外积 $\gamma_i\gamma_i^T$ 前面的系数是它自己的特征值。
\item 由 $AB=O$ 得到的是"$B$ 的\textbf{非零}列向量是 $A$ 属于 $0$ 的特征向量"；$B$ 的零列向量不是特征向量。
\item 只给特征值不给特征向量，或只给特征向量不给特征值，都无法确定矩阵。
\item 用 $r(A)$ 推非零特征值个数时，必须先确认 $A$ 可对角化（如实对称、幂等、对合）。
\end{enumerate}
\end{callwarn}

\begin{callexam}{真题考法}
\begin{enumerate}
\item 解答题常见两问：第一问由 $A\alpha=\lambda\alpha$ 或其变形反求参数；第二问求正交矩阵 $Q$ 并写出 $A=Q\Lambda Q^T$。
\item 给出 $A$ 满足的矩阵方程（如 $A^*=A-E$、$A^2-3A+2E=O$）与少量数值信息，要求先定出全部特征值，再反求 $A$ 或求其行列式、迹。
\item 已知 $A$ 为实对称矩阵、某个重根的特征向量、另一个特征值，要求补出全部特征向量并写出 $A$；此时"知二求一"用正交条件最快。
\item 陷阱：漏掉"实对称"这一前提而直接说特征向量正交；把 $A=P\Lambda P^{-1}$ 写成 $A=P^{-1}\Lambda P$。
\end{enumerate}
\end{callexam}

\subsection*{本章易错清单}
\addcontentsline{toc}{subsection}{本章易错清单}

\begin{enumerate}
\item 特征向量必须是\textbf{非零}向量，特征值可以是 $0$。求"全部特征向量"时只写一个基础解系向量、不写任意常数，或漏写"系数不全为零"，都是丢分点。
\item $A$ 与 $A^T$ 的特征值相同，但特征向量一般不同，$A$ 与 $A^T$ 一般也不相似。凡出现 $A^T$ 就要重新算特征向量。
\item $\sum\lambda_i=\mathrm{tr}A$、$\prod\lambda_i=\vert A\vert$ 只是必要条件，不能用来证明相似；重根在求和、求积时要按重数计入。
\item $r(A)=$ 非零特征值个数只在 $A$ 可相似对角化（含实对称矩阵）时成立；一般只有 $r(A)\geq$ 非零特征值个数，反例是 $\begin{pmatrix}0&1\\0&0\end{pmatrix}$。
\item "有重特征值"不等于"不可对角化"：$E$、实对称矩阵、幂等矩阵都有重根却一定可对角化。判断的唯一标准是每个 $k$ 重特征值都有 $k$ 个线性无关的特征向量，即 $n-r(\lambda E-A)=k$。
\item 需要的是 $n$ 个\textbf{线性无关}的特征向量，而不是 $n$ 个特征向量；同一特征值最多只能选出 $k$ 个线性无关的特征向量。
\item 用对角化求 $P$ 时，$P$ 的列顺序必须与 $\Lambda$ 对角元的顺序一致；$P$ 不唯一，但每列都必须是特征向量且整体线性无关。
\item 实对称矩阵求正交矩阵 $Q$ 时\textbf{必须单位化}；不同特征值的特征向量已经自动正交，只需对重根的特征向量作施密特正交化。
\item 相似、合同、等价是三个不同的概念：相似是 $P^{-1}AP$，合同是 $C^TAC$，等价是 $PAQ$，绝不能混用记号。
\item 由 $A\alpha=\lambda\alpha$ 反求参数时不要"约去 $\alpha$"，要比较分量；由 $AB=O$ 得到的是 $B$ 的非零列向量为 $0$ 特征向量。
\item 若 $A\sim B$，则 $A^{-1}\sim B^{-1}$ 额外要求 $A$ 可逆；$\vert A\vert=0$ 时不能对 $A^*$ 使用 $\dfrac{\vert A\vert}{\lambda}$。
\item 主对角线元素全相同而仍有非零非对角元的上（下）三角矩阵一定不可相似对角化，这是构造"不可对角化"反例最快的来源。
\end{enumerate}
\chapter{二次型}
\markboth{二次型}{二次型}
\subsection*{本章考纲要求}
\addcontentsline{toc}{subsection}{本章考纲要求}

\begin{quote}
\textbf{考试内容}

\begin{enumerate}
\item 二次型及其矩阵表示
\item 合同变换与合同矩阵
\item 二次型的秩
\item 惯性定理
\item 二次型的标准形和规范形
\item 用正交变换和配方法化二次型为标准形
\item 二次型及其矩阵的正定性
\end{enumerate}
\textbf{考试要求}

\begin{enumerate}
\item 掌握二次型及其矩阵表示，了解二次型秩的概念，了解合同变换与合同矩阵的概念，了解二次型的标准形、规范形的概念以及惯性定理。
\item 掌握用正交变换化二次型为标准形的方法，会用配方法化二次型为标准形。
\item 理解正定二次型、正定矩阵的概念，并掌握其判别法。
\end{enumerate}
\end{quote}
\subsection*{本章星级总表}
\addcontentsline{toc}{subsection}{本章星级总表}

\begingroup\small\setlength{\tabcolsep}{4.0pt}
\begin{longtable}{@{}L{0.4080\linewidth}L{0.1215\linewidth}L{0.1215\linewidth}L{0.1215\linewidth}L{0.1215\linewidth}@{}}
\toprule
\textbf{考点} & \textbf{星级} & \textbf{考纲要求} & \textbf{难度} & \textbf{常见题型} \\
\midrule
\endhead
二次型及其矩阵表示、二次型的秩 & $\bigstar$$\bigstar$$\bigstar$$\bigstar$ & 掌握/了解 & 基础 & 选择、填空 \\
合同变换与合同矩阵、惯性定理 & $\bigstar$$\bigstar$$\bigstar$$\bigstar$ & 了解 & 中等 & 选择、填空 \\
用正交变换化二次型为标准形 & $\bigstar$$\bigstar$$\bigstar$$\bigstar$$\bigstar$ & 掌握 & 难 & 解答 \\
用配方法化二次型为标准形与规范形 & $\bigstar$$\bigstar$$\bigstar$$\bigstar$ & 会用 & 中等 & 解答 \\
正定二次型、正定矩阵的概念与判别法 & $\bigstar$$\bigstar$$\bigstar$$\bigstar$$\bigstar$ & 理解/掌握 & 难 & 选择、解答、证明 \\
二次型的最值、瑞利商与特征值的应用 & $\bigstar$$\bigstar$$\bigstar$ & 会用 & 难 & 解答 \\
\bottomrule
\end{longtable}
\endgroup

\section{二次型的表示与合同关系}

\subsection{二次型及其矩阵}

\subsubsection{二次型及其矩阵表示、二次型的秩 $\bigstar$$\bigstar$$\bigstar$$\bigstar$}

\textbf{定义}：含有 $n$ 个变量 $x_1,x_2,\dots,x_n$ 的二次齐次多项式

\[f(x_1,x_2,\dots,x_n)=\sum_{i=1}^{n}\sum_{j=1}^{n}a_{ij}x_ix_j\qquad(a_{ij}=a_{ji})\]

称为 $n$ 元\textbf{二次型}。"二次齐次"意味着每一项都是二次的：可以是平方项 $x_i^2$，也可以是乘积项 $x_ix_j$，但绝不能出现一次项、常数项或三次以上的项。数学二只讨论实系数二次型。

取 $A=(a_{ij})_{n\times n}$，则 $A$ 是\textbf{实对称矩阵}，且

\[f(x_1,\dots,x_n)=x^TAx,\qquad x=(x_1,\dots,x_n)^T .\]

此时称 $A$ 为该二次型的\textbf{矩阵}。二次型的矩阵一定是对称的，而且这样的矩阵是\textbf{唯一}的——这正是"二次型"与"对称矩阵"可以互相翻译的原因。

\textbf{由二次型写矩阵的规则}（这一步错了整题都错，必须练到条件反射）：

\begin{itemize}
\item $x_i^2$ 的系数直接放在对角元 $a_{ii}$ 上；
\item 乘积项 $x_ix_j$（$i\neq j$）的系数要\textbf{取一半}放在 $a_{ij}$ 与 $a_{ji}$ 上，即若 $x_ix_j$ 的系数是 $2c$，则 $a_{ij}=a_{ji}=c$。
\end{itemize}
\textbf{例 1}：写出 $f(x_1,x_2,x_3)=x_1^2+2x_2^2+3x_3^2+4x_1x_2+2x_1x_3-6x_2x_3$ 的矩阵。

平方项系数直接写对角元：$a_{11}=1,a_{22}=2,a_{33}=3$；乘积项系数取半：$a_{12}=a_{21}=2$，$a_{13}=a_{31}=1$，$a_{23}=a_{32}=-3$。故

\[A=\begin{pmatrix}1&2&1\\2&2&-3\\1&-3&3\end{pmatrix}.\]

回代检验：$x^TAx=x_1^2+2x_2^2+3x_3^2+2\cdot2x_1x_2+2\cdot1x_1x_3+2\cdot(-3)x_2x_3$，与题给二次型完全一致。

\textbf{二次型的秩}：规定 $r(f)=r(A)$，即二次型矩阵的秩。由惯性定理（下一考点）还可知道，二次型化为标准形后\textbf{非零平方项的个数}就等于 $r(A)$，规范形中 $\pm1$ 的个数之和也等于 $r(A)$。

若题目给出的是\textbf{平方和形式} $f=\sum\limits_{i=1}^{m}(b_{i1}x_1+b_{i2}x_2+\cdots+b_{in}x_n)^2$，则把各线性形式的系数按行排成 $m\times n$ 矩阵 $B$，有

\[f=x^T(B^TB)x,\qquad r(f)=r(B^TB)=r(B).\]

于是"平方和形式的二次型的秩"就等于"线性形式组系数矩阵的秩"，比直接乘开求矩阵快得多。

\textbf{例 2}：设 $f(x_1,x_2,x_3)=(x_1-x_2+x_3)^2+(x_1+x_2+x_3)^2+(x_1+2x_2+ax_3)^2$ 的秩为 $2$，求 $a$。

取 $B=\begin{pmatrix}1&-1&1\\1&1&1\\1&2&a\end{pmatrix}$，则 $r(f)=r(B)=2$，故 $\vert B\vert=0$。计算

\[\vert B\vert=\begin{vmatrix}1&-1&1\\1&1&1\\1&2&a\end{vmatrix}=(a-2)+(a-1)+1=2a-2,\]

由 $2a-2=0$ 得 $a=1$。检验 $a=1$：此时第 $3$ 行为 $(1,2,1)$，第 $1$ 行与第 $2$ 行线性无关，第 $3$ 行 $-$ 第 $1$ 行 $=(0,3,0)$ 与第 $1$ 行也线性无关（第一、二个分量不成比例），故 $r(B)=2$，符合要求。

\textbf{例 3（重要结构）}：设 $f(x_1,\dots,x_n)=\sum\limits_{i=1}^{n}(x_i-\bar x)^2$，其中 $\bar x=\dfrac{1}{n}(x_1+\cdots+x_n)$，求 $f$ 的矩阵、秩与规范形。

记 $a=(1,1,\dots,1)^T$。把 $f$ 展开整理可得

\[f=x^T\left(E-\frac1naa^T\right)x,\qquad A=E-\frac1naa^T .\]

$Aa=a-\dfrac1na(a^Ta)=a-a=0$（因为 $a^Ta=n$）；而对任何与 $a$ 正交的 $x$，$Ax=x$。所以 $A$ 的特征值为 $0$（对应 $a$ 方向）与 $1$（$n-1$ 重），$r(A)=n-1$，$f$ 半正定，规范形为

\[y_1^2+y_2^2+\cdots+y_{n-1}^2 .\]

这个结论告诉我们：$n$ 个数的离差平方和只含 $n-1$ 个"独立方向"，这也是"自由度是 $n-1$"的代数来源。

\textbf{二次型与实对称矩阵的双向翻译}：本章所有问题都可以在"二次型语言"与"矩阵语言"之间来回切换。见到二次型要能立刻写出它的对称矩阵，用于求秩、求特征值、算顺序主子式；见到实对称矩阵与向量 $x$，要能立刻写出二次型 $f=x^TAx$，用于配方、求最值、判断正定。特别地，"二次型的秩""正负惯性指数""正定"这三个概念，用矩阵语言说分别就是"矩阵的秩""正负特征值的个数""矩阵正定"，两边说的是同一件事。养成"看到二次型先写矩阵"的习惯，本章的题目就都转化成第五章已经练熟的题目。

\textbf{为什么一定要化标准形}：标准形只含平方项，二次型取值的正负一眼可辨——$d_i>0$，在该方向上 $f$ 取正值；$d_i<0$，在该方向上 $f$ 取负值；系数为零的方向上 $f$ 恒为零。于是"判断正定性""求惯性指数""求最值""解方程 $f=0$"这四类问题，都要先把二次型化成标准形或规范形。本章后两节的两套方法（正交变换法与配方法），本质上就是为这个目标服务的两种工具。

\begin{callwarn}{易错点}
\begin{enumerate}
\item 乘积项系数必须\textbf{取一半}再放到非对角元上，这是本章最高频的低级错误。
\item 二次型的矩阵必须是对称矩阵；若题目给的是一个不对称的 $A$ 而要求 $x^TAx$ 的二次型矩阵，答案是 $\dfrac12(A+A^T)$，因为 $x^TAx=x^T\dfrac{A+A^T}{2}x$。
\item 二次型的秩是矩阵的秩，不是矩阵的阶数；秩可能小于阶数（退化情形）。
\item 平方和形式的二次型矩阵是 $B^TB$，不是 $B$；但求秩时用 $r(B)$ 即可。
\item 二次型不允许出现一次项和常数项；题目中若出现，要么是印刷错误，要么根本不是二次型。
\end{enumerate}
\end{callwarn}

\begin{callexam}{真题考法}
\begin{enumerate}
\item 选择、填空：写出给定二次型的矩阵；或给出二次型矩阵，求二次型的秩、正负惯性指数。
\item 平方和形式求参数：由"秩为 $k$"列出行列式为零的方程，注意最后要检验秩确实等于 $k$（不能只保证不满秩）。
\item 抽象二次型 $f=x^TAx$ 的秩，常与 $A$ 的特征值、$A$ 的结构（如 $A=E-\frac1naa^T$）结合。
\item 陷阱：忘记乘积项取半；把非对称矩阵直接当成二次型矩阵。
\end{enumerate}
\end{callexam}

\subsection{合同变换与惯性定理}

\subsubsection{合同变换与合同矩阵、惯性定理 $\bigstar$$\bigstar$$\bigstar$$\bigstar$}

\textbf{坐标变换（可逆线性变换）}：设 $x=(x_1,\dots,x_n)^T,\ y=(y_1,\dots,y_n)^T$，若 $x=Cy$ 且 $C$ 是\textbf{可逆}矩阵，就称它为由 $x$ 到 $y$ 的坐标变换。若进一步要求 $C$ 是正交矩阵，则称为\textbf{正交变换}。做代换

\[x=Cy\ \Longrightarrow\ f=x^TAx=(Cy)^TA(Cy)=y^T(C^TAC)y,\]

即二次型的矩阵从 $A$ 变成了 $C^TAC$。这引出合同的概念。

\textbf{定义}：设 $A,B$ 是同阶方阵，若存在\textbf{可逆}矩阵 $C$ 使

\[C^TAC=B,\]

则称 $A$ 与 $B$ \textbf{合同}，称这个变换为合同变换。注意必须是对称矩阵才能谈二次型，因此实对称矩阵只可能与实对称矩阵合同，不可能与非对称矩阵合同。

\textbf{合同的性质}：自反性（取 $C=E$）、对称性（$A$ 与 $B$ 合同 $\Rightarrow$ $B$ 与 $A$ 合同，取 $C^{-1}$）、传递性（两次合同变换复合仍是一次合同变换，因为 $C_2^T(C_1^TAC_1)C_2=(C_1C_2)^TA(C_1C_2)$）。并且合同矩阵的\textbf{秩相同}。特别地，若 $Q$ 是正交矩阵，则 $Q^TAQ=B$ 既表示 $A$ 与 $B$ 合同，又表示 $A$ 与 $B$ 相似。

\textbf{惯性定理}：无论用什么样的可逆线性变换把二次型 $f=x^TAx$ 化为标准形，标准形中\textbf{正平方项的个数 $p$} 与\textbf{负平方项的个数 $q$} 都是固定不变的，它们分别称为二次型的\textbf{正惯性指数}与\textbf{负惯性指数}；$r=p+q$ 就是二次型的秩。也就是说：

\begin{itemize}
\item 二次型的\textbf{规范形是唯一的}（不计平方项的次序）；
\item $p$ 等于 $A$ 的正特征值的个数，$q$ 等于 $A$ 的负特征值的个数（都按重数计）；
\item 把标准形 $d_1y_1^2+\cdots+d_ny_n^2$ 中非零的系数按正负分别数一下，就得到 $p$ 与 $q$；符号差 $s=p-q$ 也是不变量。
\end{itemize}
\textbf{判别两个同阶实对称矩阵是否合同}：设 $A,B$ 都是 $n$ 阶实对称矩阵，下面几条彼此等价——$A$ 与 $B$ 合同 $\iff p_A=p_B$ 且 $q_A=q_B\iff$ 秩相同且正惯性指数相同 $\iff$ 二者的规范形相同。一句话概括就是："正特征值个数相同、负特征值个数相同。"注意：\textbf{只要求个数相同，不要求特征值本身相同}。例如 $\mathrm{diag}(1,1,0)$ 与 $\mathrm{diag}(2,2,0)$ 的正特征值个数都是 $2$、负特征值个数都是 $0$，二者合同，但特征值不同，故不相似。

\textbf{符号差与秩的关系}：符号差 $s=p-q$ 与秩 $r=p+q$ 一起可以唯一确定 $p$ 与 $q$：

\[p=\frac{r+s}{2},\qquad q=\frac{r-s}{2}.\]

由此立刻得到两个常用的必要条件：$s$ 与 $r$ 必须\textbf{同奇偶}（否则 $p,q$ 不是整数），且 $\vert s\vert\leq r\leq n$。填空题里若给出"二次型的秩为 $r$、符号差为 $s$"，用这两个式子即可一步反求出正、负惯性指数。

\textbf{求正、负惯性指数的三条路径}：一是\textbf{配方}，把二次型化成标准形后数正负项的个数，计算量最小，适合具体数字题；二是\textbf{求特征值}，数正特征值与负特征值的个数，适合矩阵结构清楚、特征值容易求的题（分块对角、秩一、行和型等）；三是\textbf{看主元符号}，用合同变换把矩阵化成对角阵后数主元的正负，适合含参数的讨论题。三条路径的结论必须一致，因此可以互相检验：如果配方得到的 $p$ 与由特征值数出的 $p$ 不相等，一定有哪一步算错了。

\textbf{二次型的三种形态}：

\begingroup\small\setlength{\tabcolsep}{4.0pt}
\begin{longtable}{@{}L{0.1416\linewidth}L{0.1416\linewidth}L{0.1416\linewidth}L{0.2302\linewidth}L{0.2390\linewidth}@{}}
\toprule
\textbf{名称} & \textbf{形式} & \textbf{使用的变换} & \textbf{系数特点} & \textbf{是否唯一} \\
\midrule
\endhead
二次型 & $x^TAx$ & 无 & 含乘积项，乘积项系数取半 & 矩阵唯一 \\
标准形 & $d_1y_1^2+\cdots+d_ny_n^2$ & $x=Cy$，$C$ 可逆 & 只含平方项 & 与 $C$ 有关，一般不唯一 \\
规范形 & $z_1^2+\cdots+z_p^2-z_{p+1}^2-\cdots-z_{p+q}^2$ & $x=Cz$，$C$ 可逆 & 系数只取 $1,-1,0$ & 唯一（不计次序） \\
\bottomrule
\end{longtable}
\endgroup

\textbf{三种矩阵关系的对比}：

\begingroup\small\setlength{\tabcolsep}{4.0pt}
\begin{longtable}{@{}L{0.1122\linewidth}L{0.1823\linewidth}L{0.3014\linewidth}L{0.3155\linewidth}@{}}
\toprule
\textbf{关系} & \textbf{定义} & \textbf{保持不变的性质} & \textbf{判别方法} \\
\midrule
\endhead
等价 & 存在可逆 $P,Q$ 使 $PAQ=B$ & 秩 & $r(A)=r(B)$ \\
相似 & 存在可逆 $P$ 使 $P^{-1}AP=B$ & 特征值（含重数）、迹、行列式、秩、特征多项式 & 没有简单充要条件；二者都可对角化时看特征值是否相同 \\
合同 & 存在可逆 $C$ 使 $C^TAC=B$ & 秩、正惯性指数、负惯性指数 & 实对称时，正、负惯性指数分别相同 \\
\bottomrule
\end{longtable}
\endgroup

对\textbf{实对称矩阵}还有两条要记住的结论：相似一定合同（特征值相同自然正负个数也相同），但合同不一定相似（例如 $\mathrm{diag}(1,1,0)$ 与 $\mathrm{diag}(2,2,0)$ 合同——正惯性指数都是 $2$、负惯性指数都是 $0$；但特征值不同，故不相似）。此外，对实对称矩阵，"相似"与"正交相似"是同一件事。

\textbf{例}：判断 $A=\begin{pmatrix}2&-1&-1\\-1&2&-1\\-1&-1&2\end{pmatrix}$ 与 $B=\begin{pmatrix}1&0&0\\0&1&0\\0&0&0\end{pmatrix}$ 是否合同、是否相似。

把 $A$ 写成 $A=3E-J$（$J$ 为全 $1$ 矩阵），而 $J$ 的特征值为 $3,0,0$，故 $A$ 的特征值为 $0,3,3$，于是 $p_A=2,q_A=0$，$r(A)=2$。$B$ 的特征值为 $1,1,0$，$p_B=2,q_B=0$，$r(B)=2$。正、负惯性指数分别相同，所以 $A$ 与 $B$ \textbf{合同}。但 $A$ 的特征值是 $0,3,3$，$B$ 的特征值是 $0,1,1$，特征值不完全相同，所以二者\textbf{不相似}。这个例子正是"实对称矩阵合同但不相似"的标准反例。

\textbf{用初等变换求合同变换矩阵 $C$}：若要求出把 $A$ 化为对角阵的具体可逆矩阵 $C$，可对 $A$ 作一次初等行变换的同时，作\textbf{同一种}初等列变换（即用同一个初等矩阵 $E_i$ 左乘又右乘），这样每一步都相当于 $A\to E_i^TAE_i$；同时对单位矩阵 $E$ 只作\textbf{相同的列变换}（即 $E\to EE_i$）。当 $A$ 被化成对角阵 $D$ 时，$E$ 就变成了 $C$，此时 $C^TAC=D$。这个方法的优点是能同时得到对角形与变换矩阵，缺点是操作时要"行、列成对"，稍不留神就会漏掉一边。

\textbf{例（用合同变换把二次型化为标准形）}：把 $f=2x_1x_2$ 化为标准形。它的矩阵是 $A=\begin{pmatrix}0&1\\1&0\end{pmatrix}$。作可逆线性变换 $x_1=y_1+y_2,\ x_2=y_1-y_2$，即 $C=\begin{pmatrix}1&1\\1&-1\end{pmatrix}$，代入得

\[f=2(y_1+y_2)(y_1-y_2)=2y_1^2-2y_2^2,\]

即标准形为 $2y_1^2-2y_2^2$，于是 $p=1,q=1,r=2$。用矩阵验证：$C^TAC=\begin{pmatrix}2&0\\0&-2\end{pmatrix}$，与直接代入的结果完全一致。可以看到，作一次合同变换与作一次代换是同一件事的两种写法。

\begin{callwarn}{易错点}
\begin{enumerate}
\item 合同用的是 $C^TAC$，相似用的是 $P^{-1}AP$，两张皮不能混；只有 $C$ 为正交矩阵时两者才同时成立。
\item 合同\textbf{不要求}特征值相同，只要求正、负惯性指数相同；把"特征值相同"当作合同的条件是常见错误。
\item 标准形不唯一（配方顺序不同、所用可逆变换不同，系数就不同），但正负项的\textbf{个数}一定不变；规范形唯一（不计次序）。
\item 惯性指数是"个数"，不是"具体的系数值"；把 $2y_1^2$ 说成"正惯性指数是 $2$"是典型错误（应为"正惯性指数是 $1$"）。
\item 实对称矩阵不可能与非对称矩阵合同；讨论合同时先确认是实对称矩阵。
\item 用初等变换法求 $C$ 时，对 $A$ 必须"行、列同时"作同一种变换，对 $E$ 只作列变换，三者不能错位。
\end{enumerate}
\end{callwarn}

\begin{callexam}{真题考法}
\begin{enumerate}
\item 选择：给出两个实对称矩阵，判断是否合同、是否相似，或给出"合同且相似""合同但不相似"等选项。
\item 填空：求二次型的正、负惯性指数与符号差；给参数使两个矩阵合同。
\item 计算 $p,q$ 的三条路：化标准形后数正负项（配方或正交变换）、求特征值数正负个数、对简单矩阵直接看主元符号。
\item 陷阱：把"秩相同"当成合同的充要条件（还应加上正惯性指数相同）。
\end{enumerate}
\end{callexam}

\section{化二次型为标准形}

\subsection{正交变换法}

\subsubsection{用正交变换化二次型为标准形 $\bigstar$$\bigstar$$\bigstar$$\bigstar$$\bigstar$}

\textbf{定理}：任一 $n$ 元实二次型 $f=x^TAx$（$A$ 为实对称矩阵），必存在\textbf{正交变换} $x=Qy$，把它化为标准形

\[f=\lambda_1y_1^2+\lambda_2y_2^2+\cdots+\lambda_ny_n^2,\]

其中 $\lambda_1,\dots,\lambda_n$ 恰好是 $A$ 的 $n$ 个特征值。用矩阵语言说，就是存在正交矩阵 $Q$ 使

\[Q^TAQ=Q^{-1}AQ=\Lambda=\mathrm{diag}(\lambda_1,\dots,\lambda_n).\]

\textbf{这说明：用正交变换得到的标准形，其系数一定是特征值（因而在"系数就是特征值"这个意义下是唯一的），并且标准形中平方项的顺序与 $Q$ 的列、$\Lambda$ 的对角元顺序一一对应。}

\textbf{解题步骤}（与第五章求正交矩阵 $Q$ 完全一致，务必背熟）：

\begin{enumerate}
\item 写出二次型矩阵 $A$（乘积项系数记着取半）；
\item 解 $\vert\lambda E-A\vert=0$，求 $A$ 的全部特征值；
\item 对每个特征值解 $(\lambda E-A)x=0$，求基础解系；
\item 把\textbf{同一特征值}的多个特征向量作施密特正交化，再全部单位化，按列排成 $Q$；
\item 写出结果：$x=Qy$ 使 $f=\lambda_1y_1^2+\cdots+\lambda_ny_n^2$。
\end{enumerate}
\textbf{正交变换的几何意义}：因为 $Q^TQ=E$，所以 $\Vert x\Vert^2=x^Tx=(Qy)^T(Qy)=y^Ty=\Vert y\Vert^2$，即正交变换\textbf{保持向量的长度不变}，几何上只作旋转与反射，因此二次曲面的形状不变。这既是正交变换法与配方法的本质区别，也是它在几何问题中更受青睐的原因。而配方法所用的可逆变换会改变长度，只保持二次型的"正负结构"。

\textbf{例 1}：用正交变换把 $f(x_1,x_2,x_3)=x_1^2+x_2^2+x_3^2+4x_1x_2+4x_1x_3+4x_2x_3$ 化为标准形。

乘积项系数都是 $4$，取半得 $2$，故

\[A=\begin{pmatrix}1&2&2\\2&1&2\\2&2&1\end{pmatrix}=2J-E,\]

其中 $J$ 是元素全为 $1$ 的矩阵。

由 $2J$ 的特征值为 $6,0,0$ 得 $A$ 的特征值为 $\lambda_1=\lambda_2=-1,\ \lambda_3=5$。也可直接看出：$A$ 的每行元素之和均为 $5$，故 $5$ 是特征值，$\dfrac1{\sqrt3}(1,1,1)^T$ 是相应特征向量。再由 $\mathrm{tr}A=3$ 与 $\vert A\vert=5$，剩下两个特征值满足"和为 $-2$、积为 $-5$"，即 $-1,-1$。

对 $\lambda=5$：解 $5E-A=6E-2J$，得特征向量 $(1,1,1)^T$。

对 $\lambda=-1$：解 $-E-A=-2J$，即 $x_1+x_2+x_3=0$，取 $\eta_1=(1,-1,0)^T,\eta_2=(1,0,-1)^T$，正交化后得 $\beta_2=\left(\dfrac12,\dfrac12,-1\right)^T$。

单位化后取

\[Q=\left(\frac{1}{\sqrt2}\begin{pmatrix}1\\-1\\0\end{pmatrix},\ \frac{1}{\sqrt6}\begin{pmatrix}1\\1\\-2\end{pmatrix},\ \frac{1}{\sqrt3}\begin{pmatrix}1\\1\\1\end{pmatrix}\right),\]

则 $Q^TAQ=\mathrm{diag}(-1,-1,5)$，于是在正交变换 $x=Qy$ 下

\[f=-y_1^2-y_2^2+5y_3^2 .\]

由这个标准形立刻读出正惯性指数 $p=1$、负惯性指数 $q=2$、秩 $r=3$。

\textbf{例 2（已知标准形反求参数）}：设 $f=x_1^2+x_2^2+x_3^2+2ax_1x_2+2ax_1x_3+2ax_2x_3$ 的秩为 $1$，求 $a$，并写出正交变换下的标准形。

\[A=(1-a)E+aJ,\]

$J$ 的特征值为 $3,0,0$，故 $A$ 的特征值为 $1+2a$（对应 $(1,1,1)^T$）与 $1-a$（二重）。秩为 $1$ 意味着非零特征值只有 $1$ 个，故 $1-a=0$，即 $a=1$。此时 $A=J$，特征值为 $3,0,0$，正交变换下的标准形为 $3y_1^2$，规范形为 $y_1^2$。回代检验：$a=1$ 时 $f=(x_1+x_2+x_3)^2$，确实秩为 $1$。

\textbf{常见变形与应对}：一是\textbf{已知 $A$ 求标准形}，直接走上面四步；二是\textbf{已知标准形或已知 $Q$ 的某一列反求参数}，用"正交矩阵的每一列都是特征向量"以及 $\mathrm{tr}A=\sum\lambda_i$、$\vert A\vert=\prod\lambda_i$ 列方程；三是\textbf{先定参数再求 $Q$}，题目常给一个含参数的矩阵并说"它能用正交变换化为标准形"，此时隐含的条件是 $A^T=A$，先由对称性解出参数；四是\textbf{与第五章的特征值问题合并考查}，本质上只有一道题——求出特征值与标准正交特征向量，同时得到 $Q$ 与标准形。无论哪种变形，落笔的第一件事都是写出对称矩阵 $A$。

\begin{callwarn}{易错点}
\begin{enumerate}
\item 必须是\textbf{正交矩阵} $Q$（$Q^{-1}=Q^T$），不能只求一个可逆矩阵；题目若说"用正交变换"，答案里必须出现单位化后的向量。
\item 同一特征值的多个特征向量要作施密特正交化，不同特征值的不要动；最后所有向量都要单位化。
\item $Q$ 的列顺序与 $\Lambda$ 对角元顺序、与标准形中 $\lambda_i y_i^2$ 的次序必须三者一致。
\item 正交变换得到的标准形系数就是 $A$ 的特征值，因此结果（在对应顺序下）唯一；如果算出系数不是特征值，一定错了。
\item $f$ 与 $A$ 的对应关系要写对：$f=x^TAx$，变换是 $x=Qy$，代入得 $y^TQ^TAQy$，不要漏掉 $Q^T$。
\end{enumerate}
\end{callwarn}

\begin{callexam}{真题考法}
\begin{enumerate}
\item 数二解答题的核心考点之一：求正交变换 $x=Qy$ 把二次型化为标准形，并写出标准形；有时还要求写出二次型的矩阵或判断正定性。
\item 变式一：先由"正交变换后标准形为 $\lambda_1y_1^2+\lambda_2y_2^2+\lambda_3y_3^2$"或"$Q$ 的某一列"反求参数 $a,b$，再求 $Q$；利用 $\mathrm{tr}A=\sum\lambda_i$、$\vert A\vert=\prod\lambda_i$ 可快速定参。
\item 变式二：已知二次型经正交变换化为 $2y_1^2+y_2^2-y_3^2$ 之类的具体形式，反求原二次型中的参数并求 $Q$。
\item 陷阱：忘记单位化；重根的特征向量忘记正交化；标准形系数与 $Q$ 列顺序不匹配。
\end{enumerate}
\end{callexam}

\subsection{配方法}

\subsubsection{用配方法化二次型为标准形与规范形 $\bigstar$$\bigstar$$\bigstar$$\bigstar$}

\textbf{配方法（拉格朗日配方法）}的目标是找一个\textbf{可逆}线性变换 $x=Cy$，把二次型化成只含平方项的标准形

\[f=d_1y_1^2+d_2y_2^2+\cdots+d_ny_n^2 .\]

\textbf{步骤}：

\begin{enumerate}
\item \textbf{有平方项时}：先挑一个系数不为零的平方项，例如 $a_{kk}\neq0$，把二次型中所有含 $x_k$ 的项集中在一起，配成完全平方
\end{enumerate}
\[f=a_{kk}\left(x_k+\sum_{j\neq k}b_jx_j\right)^2+g(x_1,\dots,x_{k-1},x_{k+1},\dots,x_n),\]

其中 $g$ 是一个不再含 $x_k$ 的 $n-1$ 元二次型。配制时要注意"加了什么就要减去什么"，保证二次型恒等变形。

\begin{enumerate}
\item \textbf{对 $g$ 重复第一步}，直到剩下的部分也是完全平方；没有平方项可配时配方结束。
\item \textbf{没有平方项时}（所有 $a_{ii}=0$ 但某个 $a_{ij}\neq0,\ i\neq j$）：先作可逆线性变换
\end{enumerate}
\[x_i=y_i+y_j,\qquad x_j=y_i-y_j,\qquad x_k=y_k\ (k\neq i,j),\]

这个变换是可逆的（行列式为 $-2$），代入后会出现平方项，再按第 1、2 步配方。

\begin{enumerate}
\item \textbf{由标准形化规范形}：对标准形 $f=d_1y_1^2+\cdots+d_ny_n^2$（$d_i\neq0$），再作伸缩变换
\end{enumerate}
\[y_i=\frac{1}{\sqrt{\vert d_i\vert}}z_i\qquad(i=1,\dots,r),\qquad y_i=z_i\ (i>r),\]

就把系数全部变成 $\pm1$，得到规范形 $z_1^2+\cdots+z_p^2-z_{p+1}^2-\cdots-z_{p+q}^2$。这两步都是可逆变换，复合起来仍是一个可逆线性变换。

\textbf{例}：用配方法把 $f(x_1,x_2,x_3)=x_1x_2+x_1x_3+x_2x_3$ 化为标准形，并写出所用的可逆线性变换。

二次型没有平方项，先作变换

\[x_1=y_1+y_2,\qquad x_2=y_1-y_2,\qquad x_3=y_3,\]

代入得

\[f=(y_1+y_2)(y_1-y_2)+(y_1+y_2)y_3+(y_1-y_2)y_3=y_1^2-y_2^2+2y_1y_3 .\]

现在有平方项了，对 $y_1$ 配方：

\[y_1^2+2y_1y_3=(y_1+y_3)^2-y_3^2,\]

故 $f=(y_1+y_3)^2-y_2^2-y_3^2$。令

\[z_1=y_1+y_3,\qquad z_2=y_2,\qquad z_3=y_3,\]

则 $f=z_1^2-z_2^2-z_3^2$，这已经是规范形。把两步变换复合（$y_1=z_1-z_3,\ y_2=z_2,\ y_3=z_3$ 代回 $x$）：

\[x_1=z_1+z_2-z_3,\qquad x_2=z_1-z_2-z_3,\qquad x_3=z_3,\]

即 $x=Cz$，其中 $C=\begin{pmatrix}1&1&-1\\1&-1&-1\\0&0&1\end{pmatrix}$，$\vert C\vert=-2\neq0$，$C$ 可逆。检验：取 $z=(1,0,0)^T$，则 $x=(1,1,0)^T$，$f=1\cdot1+0+0=1=z_1^2$；取 $z=(0,1,0)^T$，则 $x=(1,-1,0)^T$，$f=-1=-z_2^2$，都对。

由此也读出 $p=1,q=2,r=3$，与用正交变换求特征值（$A$ 的特征值为 $1,-\dfrac12,-\dfrac12$）得到的结论一致。

\textbf{两种方法的对比}：

\begingroup\small\setlength{\tabcolsep}{4.0pt}
\begin{longtable}{@{}L{0.1663\linewidth}L{0.4019\linewidth}L{0.3603\linewidth}@{}}
\toprule
\textbf{项目} & \textbf{正交变换法} & \textbf{配方法} \\
\midrule
\endhead
变换类型 & $x=Qy$，$Q$ 正交 & $x=Cy$，$C$ 可逆 \\
标准形系数 & 恰好是 $A$ 的特征值（唯一） & 由配方过程决定（不唯一） \\
是否保持长度 & 保持 & 一般不保持 \\
计算量 & 大（要求特征值、特征向量、正交化、单位化） & 小（只作代数配方） \\
主要用途 & 几何问题、二次曲面、求最值 & 快速求标准形、规范形、惯性指数、判正定 \\
\bottomrule
\end{longtable}
\endgroup

\textbf{配方法的一个常见应用：求 $f(x)=0$ 的全部解}。若 $f$ 是平方和形式 $f=\sum(b_{i1}x_1+\cdots+b_{in}x_n)^2$，则 $f=0\iff$ 每个线性形式都等于 $0$，于是问题化为解齐次线性方程组。若 $f$ 已经化为标准形 $d_1y_1^2+\cdots+d_ry_r^2$（$d_i\neq0$），则 $f=0$ 等价于 $y_1=\cdots=y_r=0$，而 $y_{r+1},\dots,y_n$ 是自由的，再回代 $x=Cy$ 就是全部解。这两种做法都要注意：\textbf{系数为零的那些平方项不给出任何限制}，相应变量是自由变量，不能漏。

\textbf{配方时的三个细节}：第一，优先选\textbf{只与少数变量耦合}或系数较大的平方项所在变量来配，可以让后续式子更简单；第二，只要某个变量的平方项系数不为零，就直接配方，不要再绕道使用 $x_i=y_i\pm y_j$ 的技巧；第三，配方完毕务必\textbf{回代检验}——把 $x=Cy$ 代回原二次型，看是否恰好得到所写的标准形。配方题的失分大多不是配不出来，而是所用的变换写错、或者"加了什么就要减什么"没有对齐。

\textbf{配方法与合同变换法的关系}：两种方法都能找到可逆矩阵 $C$ 使 $C^TAC$ 为对角阵。区别在于：配方法直接对二次型的表达式变形，"人算"更灵活，也容易把 $C$ 一并写出来；合同变换法直接对矩阵作成对的初等行、列变换，过程更机械、更适合程序化操作。考试中大题一般会指定"用正交变换"或"用配方法"，按要求做即可；如果题目只说"求一个可逆线性变换 $x=Cy$"，那么两种方法都算对，用配方法通常更快。

\begin{callwarn}{易错点}
\begin{enumerate}
\item 配方法必须用\textbf{可逆}线性变换；判断 $x_i=y_i+y_j,\ x_j=y_i-y_j$ 是否可逆（行列式 $-2\neq0$）是必备检查。
\item 配方是恒等变形，"加了什么就要减回去"，漏减会让二次型被改变。
\item 标准形要用\textbf{原变量}表示出完整变换 $x=Cy$，只写 $f=$ 标准形而不给 $C$，在解答题里通常拿不到分。
\item 配方法与正交变换法得到的标准形一般不同（系数不同），但正、负惯性指数一定相同，这可以用来互相检验。
\item 化规范形的伸缩变换中，$d_i<0$ 时 $y_i=z_i/\sqrt{\vert d_i\vert}$ 会带来一个负号，答案要写成 $-z_i^2$。
\end{enumerate}
\end{callwarn}

\begin{callexam}{真题考法}
\begin{enumerate}
\item 解答题：用配方法（可逆线性变换）化二次型为标准形与规范形，并写出所作变换；常与"求正负惯性指数"连考。
\item 含参数的二次型：先用配方或惯性指数确定参数，再写标准形。
\item 求 $f(x_1,x_2,x_3)=0$ 的全部解，注意自由变量的个数与 $n-r(A)$ 的关系。
\item 陷阱：变换不可逆；漏写变换矩阵；用正交变换的答案去回答"配方法"的题目（要求指出所用方法是正交还是可逆）。
\end{enumerate}
\end{callexam}

\section{正定性与二次型的最值}

\subsection{正定二次型}

\subsubsection{正定二次型、正定矩阵的概念与判别法 $\bigstar$$\bigstar$$\bigstar$$\bigstar$$\bigstar$}

\textbf{定义}：设 $f=x^TAx$ 是 $n$ 元实二次型（$A$ 为实对称矩阵）。若对任意\textbf{非零}实向量 $x$，都有

\[f(x)=x^TAx>0,\]

则称 $f$ 为\textbf{正定二次型}，称 $A$ 为\textbf{正定矩阵}（记作 $A>0$）。

必须强调：\textbf{正定矩阵一定是实对称矩阵}，只有当 $A$ 实对称时，$x^TAx>0$ 这样的性质才被称为正定。若题目给出的矩阵不对称，先把它对称化为 $\dfrac12(A+A^T)$ 再讨论。

\textbf{五个等价的判别法}（$n$ 元实二次型 $f=x^TAx$ 正定 $\iff$ 下面任意一条成立）：

\begin{enumerate}
\item 正惯性指数 $p=n$（等价地 $q=0$、$r(A)=n$），即规范形为 $z_1^2+\cdots+z_n^2$，没有负项也没有零项；
\item $A$ 的\textbf{全部特征值都大于零}；
\item $A$ 的\textbf{各阶顺序主子式都大于零}（Sylvester 准则）：记 $D_k$ 为 $A$ 左上角 $k$ 阶子式，则 $D_1>0,D_2>0,\dots,D_n>0$；
\item $A$ 与单位矩阵 $E$ \textbf{合同}，即存在可逆矩阵 $C$ 使 $A=C^TEC=C^TC$；
\item 存在可逆矩阵 $D$ 使 $A=D^TD$（与第 4 条是同一件事的另一种写法）。
\end{enumerate}
此外，配方法也是一条实用判据：配方后若各平方项的系数全大于零，则二次型正定。

\textbf{必要条件}（可以用来快速排除）：$A$ 的主对角元全大于零（$a_{ii}>0$），且 $\vert A\vert>0$。但这两条\textbf{不充分}。反例就是前面用过的

\[A=\begin{pmatrix}1&2&2\\2&1&2\\2&2&1\end{pmatrix},\qquad a_{ii}=1>0,\ \ \vert A\vert=5>0,\]

但 $A$ 的特征值为 $5,-1,-1$，有负特征值，$f$ \textbf{不正定}（例如取 $x=(1,-1,0)^T$，$f=1+1-4=-2<0$）。所以想证明正定不能只验证对角元和行列式，必须用五条等价判据之一。

\textbf{例 1（用顺序主子式求参数范围）}：求 $t$ 的取值范围，使 $f=x_1^2+2x_2^2+3x_3^2+2tx_1x_2$ 正定。

\[A=\begin{pmatrix}1&t&0\\t&2&0\\0&0&3\end{pmatrix},\qquad D_1=1>0,\qquad D_2=\begin{vmatrix}1&t\\t&2\end{vmatrix}=2-t^2,\qquad D_3=3(2-t^2) .\]

要求全部顺序主子式大于零，即 $2-t^2>0$，故 $-\sqrt2<t<\sqrt2$。

\textbf{例 2（用"矩阵秩"判正定）}：$f=(ax_1+2x_2-3x_3)^2+(x_2-2x_3)^2+(x_1+ax_2-x_3)^2$ 是正定二次型，求 $a$ 满足的条件。

平方和形式 $f=x^T(B^TB)x$，其中 $B=\begin{pmatrix}a&2&-3\\0&1&-2\\1&a&-1\end{pmatrix}$。因为 $B^TB$ 正定 $\iff B$ 可逆，故只需 $\vert B\vert\neq0$。计算

\[\vert B\vert=a(2a-1)-2(2)+(-3)(-1)=2a^2-a-1=(2a+1)(a-1),\]

所以条件是 $a\neq1$ 且 $a\neq-\dfrac12$。这一类题的通用结论是：\textbf{平方和形式的二次型正定 $\iff$ 各线性形式线性无关 $\iff$ 系数矩阵可逆（满秩）}。

\textbf{半正定、负定与不定}：

\begin{itemize}
\item \textbf{半正定}：对一切 $x$ 有 $x^TAx\geq0$，且存在非零 $x$ 使 $x^TAx=0$。判别：特征值全 $\geq0$；$p=r(A)$（即 $q=0$）；存在矩阵 $D$（不要求可逆）使 $A=D^TD$；\textbf{所有主子式}（不只是顺序主子式）都 $\geq0$。
\item \textbf{负定}：对一切非零 $x$ 有 $x^TAx<0$。判别：特征值全 $<0$；$-A$ 正定；顺序主子式满足符号交替 $(-1)^kD_k>0$，即 $D_1<0,D_2>0,D_3<0,\dots$。
\item \textbf{不定}：$f$ 既能取到正值也能取到负值，等价于特征值有正有负。
\end{itemize}
\textbf{正定矩阵的运算性质}（设 $A,B$ 都是 $n$ 阶正定矩阵）：

\begin{itemize}
\item $A^T$、$A^{-1}$、$A^*$、$kA\ (k>0)$、$A^m$（$m$ 为正整数）都正定；
\item $A+B$ 正定；分块对角阵 $\begin{pmatrix}A&O\\O&B\end{pmatrix}$ 正定；
\item 若 $C$ 可逆，则 $C^TAC$ 正定（合同的传递性）；
\item 存在正定矩阵 $S$ 使 $A=S^2$；存在可逆矩阵 $D$ 使 $A=D^TD$；
\item $AB$ \textbf{未必}正定（甚至未必对称），这是最常见的错误选项。
\end{itemize}
\textbf{证明抽象矩阵正定的四条路线}：一是\textbf{定义法}，直接验证对任意 $x\neq0$ 有 $x^TAx>0$（常把 $x^TAx$ 配成平方和）；二是\textbf{构造 $A=D^TD$}，只要找到可逆的 $D$ 即可；三是\textbf{特征值法}，证明 $A$ 实对称且特征值全正；四是\textbf{合同法}，证明 $A$ 与 $E$ 合同。具体题目用哪一条，取决于题目给出的条件形式：给平方和用第一条，给 $A=B^TB$ 或 $A=P^TBP$ 用第二条或第四条，给特征值信息用第三条。

\textbf{判别方法的选择策略}：如果是\textbf{具体数字矩阵}，首选顺序主子式——最多算 $n$ 个左上角行列式，含参数时也最方便讨论；如果矩阵具有\textbf{分块对角、秩一、行和型}等特殊结构，直接算特征值更快；如果是\textbf{抽象矩阵}（只给出 $A=B^TB$、$A=P^TBP$、$A$ 与某个正定矩阵合同等信息），就用"与 $E$ 合同""$A=D^TD$""特征值全正"这三条之一来论证，必要时回到定义验证 $x^TAx>0$。做选择题时还常用排除法：先看主对角元是否全正、行列式是否大于零，不满足就直接排除。

\textbf{例 3（抽象矩阵正定的证明）}：设 $A$ 是 $m\times n$ 实矩阵，$B=\lambda E+A^TA$，证明当 $\lambda>0$ 时 $B$ 正定。

先验证对称性：$B^T=\lambda E^T+(A^TA)^T=\lambda E+A^TA=B$，所以 $B$ 是实对称矩阵。再对任意 $x\neq0$ 计算

\[x^TBx=\lambda x^Tx+x^TA^TAx=\lambda\Vert x\Vert^2+\Vert Ax\Vert^2 .\]

因为 $\lambda>0$ 且 $x\neq0$，第一项 $\lambda\Vert x\Vert^2>0$；第二项是向量长度的平方，恒有 $\Vert Ax\Vert^2\geq0$。两项之和大于零，故 $B$ 正定。这个例子展示了"定义法"的标准写法：先说明实对称，再对一般 $x$ 把 $x^TBx$ 化成"正项 $+$ 非负项"。

\begin{callwarn}{易错点}
\begin{enumerate}
\item 正定的前提是\textbf{实对称矩阵}。题目若给出非对称矩阵，必须先说明"$x^TAx=x^T\dfrac{A+A^T}{2}x$"，再对对称部分讨论。
\item 判正定看的是\textbf{顺序}主子式；判半正定才要看\textbf{全部}主子式。二者不能互换。
\item "$a_{ii}>0$ 且 $\vert A\vert>0$"只是必要条件，不能作为判据（反例见上文 $2J-E$）。
\item $AB$ 在 $A,B$ 都正定时\textbf{未必}正定，也未必对称；只有 $A+B$、$A^{-1}$、$A^*$、$A^m$ 这些才保证正定。
\item 用特征值判别时，必须保证 $A$ 实对称，否则可能有复特征值，"全为正"就无从谈起。
\item 平方和形式正定要求各线性形式\textbf{线性无关}（系数矩阵满秩），不是"每个平方项都非零"就够了。
\end{enumerate}
\end{callwarn}

\begin{callexam}{真题考法}
\begin{enumerate}
\item 选择：判断给定的二次型或矩阵是否正定、求参数范围（常给含参数的三阶矩阵用顺序主子式）。
\item 解答、证明：证明抽象矩阵正定，常用 $A=D^TD$、$A$ 与 $E$ 合同、或 $x^TAx$ 配方。常见素材：$E+\alpha\alpha^T$（$\alpha$ 为实向量）、$A^TA+\lambda E\ (\lambda>0)$、$A^*$ 与 $A^{-1}$ 的组合。
\item 判断正定与二次型最值、特征值联合出题：先证正定，再求最值。
\item 陷阱：把非对称矩阵直接判正定；把半正定的判据（全部主子式）用到正定上；误以为 $AB$ 正定。
\end{enumerate}
\end{callexam}

\subsection{二次型的最值}

\subsubsection{二次型的最值、瑞利商与特征值的应用 $\bigstar$$\bigstar$$\bigstar$}

\textbf{瑞利商结论}：设 $A$ 是 $n$ 阶实对称矩阵，特征值按大小排列为

\[\lambda_1\leq\lambda_2\leq\cdots\leq\lambda_n,\]

则对任意\textbf{非零}实向量 $x$ 都有

\[\lambda_1\leq\frac{x^TAx}{x^Tx}\leq\lambda_n .\]

比值 $\dfrac{x^TAx}{x^Tx}$ 称为关于 $A$ 的\textbf{瑞利商}。等号成立的条件是：$x$ 取为属于 $\lambda_1$ 的特征向量时左边取等号，$x$ 取为属于 $\lambda_n$ 的特征向量时右边取等号。

\textbf{推论（单位球面上的最值）}：当 $\Vert x\Vert=1$（即 $x^Tx=1$）时，

\[\min_{\Vert x\Vert=1}x^TAx=\lambda_1,\qquad \max_{\Vert x\Vert=1}x^TAx=\lambda_n .\]

更一般地，在约束 $x^Tx=k\ (k>0)$ 下，$\min x^TAx=k\lambda_1$，$\max x^TAx=k\lambda_n$。

这两个结论的证明思路很直接：取 $A$ 的一组标准正交特征向量 $\gamma_1,\dots,\gamma_n$，把 $x$ 表成 $x=\sum c_i\gamma_i$，则 $x^Tx=\sum c_i^2$，$x^TAx=\sum\lambda_ic_i^2$，于是

\[\lambda_1=\frac{\lambda_1\sum c_i^2}{\sum c_i^2}\leq\frac{\sum\lambda_ic_i^2}{\sum c_i^2}\leq\frac{\lambda_n\sum c_i^2}{\sum c_i^2}=\lambda_n .\]

\textbf{一个常用加强}：若进一步限制 $x$ 与属于 $\lambda_n$ 的特征向量 $\gamma_n$ 正交（即在上式的展开中 $c_n=0$），则上界可缩小为 $\lambda_{n-1}$：

\[\frac{x^TAx}{x^Tx}\leq\lambda_{n-1}.\]

同理，若 $x$ 与 $\gamma_1$ 正交，则下界可提高到 $\lambda_2$。这个加强在求"在给定子空间上的最值"时非常好用。

\textbf{例 1}：设 $f(x_1,x_2,x_3)=3x_1^2+4x_2^2+3x_3^2+2x_1x_3$，证明 $\min\limits_{x\neq0}\dfrac{f(x)}{x^Tx}=2$。

\[A=\begin{pmatrix}3&0&1\\0&4&0\\1&0&3\end{pmatrix}.\]

从第 $1,3$ 行（列）组成的二阶块得特征方程 $(\lambda-3)^2-1=0$，即 $\lambda=4$ 或 $2$；再由第 $2$ 个分量单独给出 $\lambda=4$。所以 $A$ 的特征值为 $2,4,4$，最小特征值 $\lambda_1=2$。由瑞利商结论

\[\frac{f(x)}{x^Tx}=\frac{x^TAx}{x^Tx}\geq\lambda_1=2,\]

故比值的最小值为 $2$，且在属于 $2$ 的特征向量处取到。求特征向量：解线性方程组 $(2E-A)x=0$，即 $\begin{pmatrix}-1&0&-1\\0&-2&0\\-1&0&-1\end{pmatrix}x=0$，得 $x_1=-x_3,\ x_2=0$，取 $(1,0,-1)^T$，代入 $f$ 得 $3+3-2=4$，$x^Tx=2$，比值恰为 $2$，验证成立。

\textbf{例 2（用正定半正定证明不等式）}：证明柯西-施瓦茨不等式 $\left(\sum\limits_{i=1}^{n}a_ib_i\right)^2\leq\left(\sum\limits_{i=1}^{n}a_i^2\right)\left(\sum\limits_{i=1}^{n}b_i^2\right)$。

构造二次型 $f(x,y)=\sum\limits_{i=1}^{n}(a_ix+b_iy)^2$。它是平方和，故对一切 $(x,y)$ 有 $f(x,y)\geq0$，即 $f$ 半正定。另一方面把 $f$ 写成矩阵形式得 $f=x^TAx$，其中

\[A=\begin{pmatrix}\sum a_i^2&\sum a_ib_i\\ \sum a_ib_i&\sum b_i^2\end{pmatrix}.\]

半正定矩阵的行列式大于等于零，即 $\vert A\vert=\left(\sum a_i^2\right)\left(\sum b_i^2\right)-\left(\sum a_ib_i\right)^2\geq0$，整理即得所证不等式。

\textbf{例 3（由不等式求参数）}：设 $3$ 阶实对称矩阵 $A$ 的特征值为 $2,3,4$，若对任意 $3$ 维实列向量 $x$ 都有 $x^TAx\leq a\,x^Tx$，求 $a$ 的最小值。

由瑞利商结论，$\dfrac{x^TAx}{x^Tx}\leq\lambda_{\max}=4$，且在属于 $4$ 的特征向量处取到等号。因此要使不等式对一切 $x$ 成立，必须且只需 $a\geq4$，$a$ 的最小值为 $4$。推广开来："对一切 $x$ 有 $x^TAx\leq a\,x^Tx$"等价于 $a\geq\lambda_{\max}$；"$x^TAx\geq a\,x^Tx$"等价于 $a\leq\lambda_{\min}$；"$b\,x^Tx\leq x^TAx\leq a\,x^Tx$"则要求 $b\leq\lambda_{\min}$ 且 $a\geq\lambda_{\max}$。把矩阵不等式与"最大、最小特征值"互相翻译，是处理这类参数题的通法。

\textbf{最值结论的常见考法}：

\begin{itemize}
\item 求 $\min\limits_{x\neq0}\dfrac{f(x)}{x^Tx}$ 或 $\max\limits_{x\neq0}\dfrac{f(x)}{x^Tx}$：先求矩阵的特征值，最小值是最小特征值，最大值是最大特征值。
\item 求"在 $\Vert x\Vert=1$ 条件下 $f$ 的最大值"或"满足条件的最大值点"：最大值点为最大特征值对应的单位特征向量。
\item 证明关于 $x^TAx$ 的不等式：把 $x^TAx$ 夹在 $\lambda_1x^Tx$ 与 $\lambda_nx^Tx$ 之间，或把某个矩阵写成两个二次型之差再用正定性。
\item 求参数范围：由"对一切 $x$ 都有 $x^TAx\geq0$（或大于零）"转化为 $A$ 半正定（或正定），用特征值或顺序主子式解参数。
\end{itemize}
\textbf{二次型最值与正定性的联系}：这两件事天然相通。若 $f$ 正定，则比值 $\dfrac{f(x)}{x^Tx}$ 的最小值就是最小特征值 $\lambda_{\min}$，且 $\lambda_{\min}>0$；若 $f$ 半正定，则 $\lambda_{\min}=0$；若 $f$ 不定，则比值既能取到正值也能取到负值，取值范围恰好是 $[\lambda_{\min},\lambda_{\max}]$。所以"求二次型的最值"与"判断二次型的正定性"往往只是同一个问题的两种问法，命题人也经常把它们安排在同一道题的两问里。

\begin{callwarn}{易错点}
\begin{enumerate}
\item 瑞利商结论要求 $A$ 是\textbf{实对称矩阵}；非对称矩阵没有这个结论。
\item 最值是在 $x\neq0$（或 $\Vert x\Vert=1$）的约束下取的；无约束时二次型一般无界（不定或半正定时可取到正无穷）。
\item 求"最大值点"时，若最大特征值是重根，最大值点不唯一（对应的整个特征子空间中的非零向量都取到最大值），要写全。
\item 用"行列式大于等于零"证不等式时，前提是该矩阵\textbf{半正定}，不是随便什么矩阵都行；构造出平方和后要说明半正定。
\item 瑞利商的比值 $\dfrac{x^TAx}{x^Tx}$ 与 $\dfrac{x^TAx}{\Vert x\Vert^2}$ 是同一个量，不要多除一次或漏除一次。
\end{enumerate}
\end{callwarn}

\begin{callexam}{真题考法}
\begin{enumerate}
\item 解答题：求正交变换化二次型为标准形后，进一步证明 $\min\limits_{x\neq0}\dfrac{f(x)}{x^Tx}$ 等于最小特征值（或求最大值）。这是二次型与特征值结合的标准考法。
\item 求 $\Vert x\Vert=1$ 条件下二次型的最大值及其取到点。
\item 由"对任意 $x$ 都有 $x^TAx\leq a\,x^Tx$"求参数 $a$ 的最小值，转化为 $a\geq\lambda_{\max}$。
\item 用半正定构造证明柯西-施瓦茨不等式或其它二次型不等式。
\item 陷阱：忘记说明 $A$ 实对称；把最大（小）值写成最大（小）特征值乘以自由常数；重根时漏写最值点。
\end{enumerate}
\end{callexam}

\subsection*{本章易错清单}
\addcontentsline{toc}{subsection}{本章易错清单}

\begin{enumerate}
\item 由二次型写矩阵时，乘积项 $x_ix_j$ 的系数必须\textbf{取一半}放在 $a_{ij}$ 与 $a_{ji}$ 上；这是本章出现频率最高的计算错误。
\item 二次型的矩阵一定是实对称矩阵；若题目给的是非对称矩阵 $A$，则 $x^TAx$ 的二次型矩阵是 $\dfrac12(A+A^T)$，讨论正定性、惯性指数前必须先对称化。
\item 二次型的秩 $=$ 矩阵的秩 $=p+q=$ 标准形中非零项的个数；不要把阶数当秩。
\item 合同用 $C^TAC$，相似用 $P^{-1}AP$；只有 $C$ 为正交矩阵时两者同时成立。实对称矩阵相似一定合同，但合同不一定相似。
\item 惯性指数是"正负项的\textbf{个数}"，不是"平方项前面的系数值"。$2y_1^2-3y_2^2$ 的正惯性指数是 $1$、负惯性指数是 $1$、符号差是 $0$。
\item 标准形不唯一（配方方式不同系数就不同），规范形唯一（不计平方项次序）；判断两个二次型是否等价（合同）只需比较正、负惯性指数。
\item "秩相同"只是合同与相似的必要条件，不是充分条件；判别合同必须比较正惯性指数（或规范形），判别相似必须比较特征值。
\item 正交变换法必须\textbf{正交}（要单位化，重根的特征向量要正交化）；配方法必须用\textbf{可逆}变换，且要写出完整的 $x=Cy$。
\item 正交变换保持向量长度，配方法的可逆变换一般不保持；两种方法得到的标准形系数不同，但正、负惯性指数一定相同，可用来互相检验。
\item 判正定看\textbf{顺序}主子式全大于零；判半正定才看\textbf{全部}主子式非负。$a_{ii}>0$ 与 $\vert A\vert>0$ 只是必要条件，不充分（反例 $2J-E$：对角元全为 $1$、行列式为 $5$，却有负特征值）。
\item 正定矩阵必须实对称；$A,B$ 正定时 $A+B$、$A^{-1}$、$A^*$、$A^m$、$kA\ (k>0)$ 都正定，但 $AB$ 未必正定（甚至未必对称）。
\item 平方和形式 $\sum(b_{i1}x_1+\cdots+b_{in}x_n)^2$ 正定 $\iff$ 系数矩阵 $B$ 满秩（可逆），不是"每一项都非零"；求秩时用 $r(B)$。
\item 瑞利商 $\lambda_1\leq\dfrac{x^TAx}{x^Tx}\leq\lambda_n$ 只对实对称矩阵成立，且最值在对应的特征向量处取到；无约束的二次型一般没有最值。
\end{enumerate}
\chapter*{附录一 五星考点冲刺路线}
\addcontentsline{toc}{chapter}{附录一 五星考点冲刺路线}
\markboth{附录一 五星考点冲刺路线}{附录一 五星考点冲刺路线}
\subsection{这份路线怎么用}

线性代数只占约 20\textbackslash{}\% 的分值，却是\textbf{整张卷子里最容易拿满}的部分：考点少、题型固定、套路单一。 本附录把六章考点按星级重排，给出投入顺序与达标标准。

线代复习有一条铁律：\textbf{先背结论，再练计算}。线代的失分几乎都是"结论记错"或"前提漏掉"， 而不是"想不出来"。

\subsection{冲刺优先级总表}

\begingroup\small\setlength{\tabcolsep}{4.0pt}
\begin{longtable}{@{}L{0.1291\linewidth}L{0.2758\linewidth}L{0.0878\linewidth}L{0.0878\linewidth}L{0.3136\linewidth}@{}}
\toprule
\textbf{优先级} & \textbf{考点} & \textbf{星级} & \textbf{建议投入} & \textbf{达标标准} \\
\midrule
\endhead
第 1 优先 & 矩阵的秩（求法、不等式、常用结论） & $\bigstar$$\bigstar$$\bigstar$$\bigstar$$\bigstar$ & 6 小时 & 抽象矩阵秩的讨论题能写出不等式链并说明取等条件 \\
第 1 优先 & 可逆矩阵、伴随矩阵及其公式 & $\bigstar$$\bigstar$$\bigstar$$\bigstar$$\bigstar$ & 6 小时 & $AA^*=\lvert A\rvert E$ 系列公式与 $r(A^*)$ 的三种情形脱口而出 \\
第 1 优先 & 向量组线性相关与线性无关（含抽象证明） & $\bigstar$$\bigstar$$\bigstar$$\bigstar$$\bigstar$ & 8 小时 & 抽象向量组线性无关的证明能独立写出两种套路 \\
第 1 优先 & 线性方程组解的结构、基础解系与通解 & $\bigstar$$\bigstar$$\bigstar$$\bigstar$$\bigstar$ & 8 小时 & 含参数方程组能完整讨论并写出通解 \\
第 1 优先 & 特征值与特征向量、相似对角化 & $\bigstar$$\bigstar$$\bigstar$$\bigstar$$\bigstar$ & 8 小时 & 判定能否对角化并求出可逆矩阵 $P$ \\
第 1 优先 & 实对称矩阵的正交相似对角化 & $\bigstar$$\bigstar$$\bigstar$$\bigstar$$\bigstar$ & 6 小时 & 正交化、单位化、拼矩阵三步不丢分 \\
第 1 优先 & 二次型的标准形（正交变换法与配方法） & $\bigstar$$\bigstar$$\bigstar$$\bigstar$$\bigstar$ & 6 小时 & 两种方法都能求出标准形与所用变换 \\
第 1 优先 & 正定二次型与正定矩阵的判别 & $\bigstar$$\bigstar$$\bigstar$$\bigstar$$\bigstar$ & 4 小时 & 四种判别法（顺序主子式、特征值、定义、合同于 $E$）能用 \\
第 1 优先 & 抽象行列式（$\lvert kA\rvert,\lvert AB\rvert,\lvert A^*\rvert$） & $\bigstar$$\bigstar$$\bigstar$$\bigstar$$\bigstar$ & 4 小时 & 数字与抽象两类行列式都能算 \\
第 2 优先 & 数字行列式的计算技巧 & $\bigstar$$\bigstar$$\bigstar$$\bigstar$ & 5 小时 & 化三角形、降阶、范德蒙德、爪型都会 \\
第 2 优先 & 极大线性无关组与向量组的秩 & $\bigstar$$\bigstar$$\bigstar$$\bigstar$ & 5 小时 & 把向量拼成矩阵求秩并写出答案 \\
第 2 优先 & 初等变换与初等矩阵（左行右列） & $\bigstar$$\bigstar$$\bigstar$$\bigstar$ & 4 小时 & 已知变换求矩阵类题不丢分 \\
第 2 优先 & 矩阵方程 $AX=B$ 的求解 & $\bigstar$$\bigstar$$\bigstar$$\bigstar$ & 3 小时 & 用初等行变换解 $[A\mid B]$ 一步到位 \\
第 2 优先 & 用特征值、特征向量反求矩阵 & $\bigstar$$\bigstar$$\bigstar$$\bigstar$ & 4 小时 & 由 $P^{-1}AP=\Lambda$ 反解 $A$ 熟练 \\
第 2 优先 & 合同变换、惯性定理与规范形 & $\bigstar$$\bigstar$$\bigstar$$\bigstar$ & 4 小时 & 会写规范形并判断是否合同 \\
第 2 优先 & 同解方程组、公共解问题 & $\bigstar$$\bigstar$$\bigstar$$\bigstar$ & 4 小时 & 能联立两方程组求公共解 \\
第 2 优先 & 解的线性组合判据（系数和为 0 或 1） & $\bigstar$$\bigstar$$\bigstar$$\bigstar$ & 3 小时 & 一眼判断某向量是特解还是齐次解 \\
第 3 优先 & 伴随矩阵的秩、代数余子式线性组合 & $\bigstar$$\bigstar$$\bigstar$$\bigstar$ & 3 小时 & 会构造行列式求 $\sum A_{ij}$ 型问题 \\
第 3 优先 & 矩阵的线性运算与转置规律 & $\bigstar$$\bigstar$$\bigstar$ & 2 小时 & 知道 $(AB)^T=B^TA^T$、不可交换等基本点 \\
第 3 优先 & 方阵的幂与秩一矩阵技巧 & $\bigstar$$\bigstar$$\bigstar$ & 3 小时 & 秩一矩阵的幂、每行和为 $k$ 的矩阵的幂会算 \\
第 3 优先 & 克拉默法则 & $\bigstar$$\bigstar$ & 1 小时 & 记住形式与适用条件 \\
第 3 优先 & 内积与施密特正交规范化 & $\bigstar$$\bigstar$ & 2 小时 & 会算正交化与单位化 \\
第 3 优先 & 分块矩阵及其运算 & $\bigstar$$\bigstar$ & 2 小时 & 会算分块乘法与分块对角阵的逆 \\
第 4 优先 & 对称、反对称、正交矩阵的性质 & $\bigstar$$\bigstar$ & 2 小时 & 记住定义与两条常用性质 \\
\bottomrule
\end{longtable}
\endgroup

\subsection{三轮时间表（以 90 天为例）}

\textbf{第一轮：通读六章（第 1--15 天，每天 2 小时）}

\begin{itemize}
\item 顺序：行列式 $\to$ 矩阵 $\to$ 向量 $\to$ 线性方程组 $\to$ 特征值 $\to$ 二次型。
\item 边读边做一张 A4 速查表（对照本册附录二），把每条结论的\textbf{前提}写在结论旁边。
\item 每章读完立刻做 10--15 道基础题，特别是"求秩、求逆、求通解、求特征值"四类计算。
\end{itemize}
\textbf{第二轮：攻四类大题（第 16--60 天，穿插在高数复习中）}

每年数二线代的重头戏就四类：

\begin{enumerate}
\item \textbf{抽象矩阵的秩与可逆性}（证明题）：常用 $AB=O\Rightarrow r(A)+r(B)\le n$、$r(A)=n-1$ 时 $A^*\ne O$。
\item \textbf{向量组线性相关性的证明}：设 $\sum k_i\alpha_i=0$，用已知条件推出所有 $k_i=0$。
\item \textbf{含参数方程组的解的讨论}：对增广矩阵行变换，按参数分情况讨论。
\item \textbf{相似对角化与二次型标准化}：一条龙"求特征值 $\to$ 求特征向量 $\to$ 正交化 $\to$ 写出标准形 $\to$ 判断正定"。
\end{enumerate}
\textbf{第三轮：考前（最后 10 天）}

\begin{itemize}
\item 每天过一遍本册附录二"结论速查与易混对照"。
\item 重点核对易混项：相似不相似、合同不合同、等价、正交矩阵、伴随的秩。
\item 把真题里线代的两三个小题与一个大题重做一遍，限时 30 分钟。
\end{itemize}
\subsection{考前 7 天清单}

\begin{enumerate}
\item 默写：$\lvert kA\rvert=k^n\lvert A\rvert$，$\lvert AB\rvert=\lvert A\rvert\lvert B\rvert$，$\lvert A^*\rvert=\lvert A\rvert^{n-1}$，$\lvert A^{-1}\rvert=\lvert A\rvert^{-1}$。
\item 默写：$AA^*=A^*A=\lvert A\rvert E$，$A^{-1}=\frac{1}{\lvert A\rvert}A^*$，$(A^*)^{-1}=(A^{-1})^*$，$\lvert A^*\rvert$，$(A^*)^*=\lvert A\rvert^{n-2}A$。
\item 默写：$r(A^*)$ 的三种情形（$n$、$1$、$0$）。
\item 默写：$A$ 可逆的等价条件 6 条（$\lvert A\rvert\ne0$、$r(A)=n$、$A$ 的列向量组线性无关、$Ax=0$ 只有零解、$A$ 与 $E$ 等价、$A$ 可表示为有限个初等矩阵之积）。
\item 默写：相似判定的必要条件（同迹、同行列式、同秩、同特征多项式）与可对角化条件（$n$ 个线性无关特征向量）。
\item 默写：$A$ 为实对称矩阵时的三件事——特征值全为实数、不同特征值的特征向量正交、一定可正交相似对角化。
\item 默写：正定的五个等价说法与四种判别法（顺序主子式、特征值、定义、合同于单位阵）。
\end{enumerate}
\chapter*{附录二 结论速查与易混对照}
\addcontentsline{toc}{chapter}{附录二 结论速查与易混对照}
\markboth{附录二 结论速查与易混对照}{附录二 结论速查与易混对照}
\subsection{行列式公式表}

\begingroup\small\setlength{\tabcolsep}{4.0pt}
\begin{longtable}{@{}L{0.3911\linewidth}L{0.5546\linewidth}@{}}
\toprule
\textbf{公式} & \textbf{前提 / 说明} \\
\midrule
\endhead
$\lvert A^T\rvert=\lvert A\rvert$ & 对任意方阵 \\
$\lvert kA\rvert=k^n\lvert A\rvert$ & $A$ 为 $n$ 阶方阵（\textbf{不是} $k\lvert A\rvert$） \\
$\lvert AB\rvert=\lvert A\rvert\lvert B\rvert$ & $A,B$ 同阶方阵 \\
$\lvert A^{-1}\rvert=\lvert A\rvert^{-1}$ & $A$ 可逆 \\
$\lvert A^*\rvert=\lvert A\rvert^{n-1}$ & $A$ 为 $n$ 阶方阵 \\
$\lvert A\rvert=\prod_{i=1}^{n}\lambda_i$ & $\lambda_i$ 为 $A$ 的全部特征值 \\
$\lvert A+B\rvert\ne\lvert A\rvert+\lvert B\rvert$ & 一般\textbf{不等于}，只有特殊情形可拆 \\
上（下）三角行列式 $=$ 主对角线元素之积 & — \\
范德蒙德行列式 $\prod_{1\le i<j\le n}(x_j-x_i)$ & — \\
两行（列）成比例 $\Rightarrow$ 行列式为 $0$ & 含某行全为 $0$ 同理 \\
\bottomrule
\end{longtable}
\endgroup

\textbf{代数余子式的两个"自交 / 互交"结论}：

\[\sum_{j=1}^{n}a_{ij}A_{ij}=\lvert A\rvert,\qquad \sum_{j=1}^{n}a_{ij}A_{kj}=0\ (i\ne k).\]

\subsection{伴随矩阵公式表}

\[AA^*=A^*A=\lvert A\rvert E,\qquad A^{-1}=\frac{1}{\lvert A\rvert}A^*,\qquad A^*=\lvert A\rvert A^{-1}\ (A\ \text{可逆}),\]

\[(A^*)^{-1}=(A^{-1})^*,\qquad (A^*)^T=(A^T)^*,\qquad (kA)^*=k^{n-1}A^*,\]

\[(AB)^*=B^*A^*,\qquad (A^*)^*=\lvert A\rvert^{n-2}A,\qquad \lvert A^*\rvert=\lvert A\rvert^{n-1}.\]

\textbf{伴随矩阵的秩（必背）}：

\[r(A^*)=\begin{cases}n,& r(A)=n,\\ 1,& r(A)=n-1,\\ 0,& r(A)\le n-2.\end{cases}\]

\begin{callwarn}{易错}
$(AB)^*=B^*A^*$ 的顺序与 $(AB)^{-1}=B^{-1}A^{-1}$ 一致，别写反。 另外 $A^*$ 只有在 $A$ 可逆时才等于 $\lvert A\rvert A^{-1}$。

\end{callwarn}

\subsection{秩的不等式与常用结论}

\begingroup\small\setlength{\tabcolsep}{4.0pt}
\begin{longtable}{@{}L{0.4272\linewidth}L{0.5185\linewidth}@{}}
\toprule
\textbf{结论} & \textbf{说明} \\
\midrule
\endhead
$0\le r(A)\le\min\{m,n\}$ & $A$ 为 $m\times n$ \\
$r(A)=r(A^T)=r(kA)\allowbreak\ (k\ne0)$ & 转置、数乘不变 \\
$r(A+B)\le r(A)+r(B)$ & 越"拼"秩越大 \\
$r(AB)\le\min\{r(A),r(B)\}$ & 越"乘"秩越小 \\
$AB=O\allowbreak\ \Rightarrow\allowbreak\ r(A)+r(B)\le n$ & $A$ 为 $m\times n$，$B$ 为 $n\times s$ \\
$P,Q$ 可逆 $\Rightarrow r(PAQ)=r(A)$ & 乘可逆阵秩不变 \\
$r(A)=n\iff A$ 可逆 & $A$ 为 $n$ 阶 \\
$n$ 阶矩阵 $A$：$r(A)<n\iff\lvert A\rvert=0\iff Ax=0$ 有非零解 & 三条等价 \\
$r\begin{pmatrix}A&O\\O&B\end{pmatrix}=r(A)+r(B)$ & 分块对角 \\
若 $A$ 的列向量组可由 $B$ 的列向量组线性表示，则 $r(A)\le r(B)$ & 表示关系 \\
\bottomrule
\end{longtable}
\endgroup

\subsection{可逆矩阵的等价条件（$A$ 为 $n$ 阶方阵）}

以下每一条都\textbf{等价于"$A$ 可逆"}：

\begin{enumerate}
\item $\lvert A\rvert\ne0$；
\item $r(A)=n$；
\item $A$ 的列（行）向量组线性无关；
\item $Ax=0$ 只有零解；
\item 对任意 $b$，$Ax=b$ 有唯一解；
\item $A$ 与 $E$ 等价，即 $A$ 可经有限次初等变换化为 $E$；
\item $A$ 可表示为有限个初等矩阵的乘积；
\item $A$ 的特征值全不为 $0$。
\end{enumerate}
\subsection{等价、相似、合同对照表}

\begingroup\small\setlength{\tabcolsep}{4.0pt}
\begin{longtable}{@{}L{0.1236\linewidth}L{0.2008\linewidth}L{0.2858\linewidth}L{0.3012\linewidth}@{}}
\toprule
\textbf{关系} & \textbf{定义} & \textbf{必要条件} & \textbf{充分条件} \\
\midrule
\endhead
等价 & 存在可逆 $P,Q$ 使 $PAQ=B$ & 同型、同秩 & 同型且同秩 \\
相似 & 存在可逆 $P$ 使 $P^{-1}AP=B$ & 同秩、同行列式、同迹、同特征多项式 & $n$ 阶方阵有 $n$ 个线性无关特征向量 \\
合同 & 存在可逆 $C$ 使 $C^TAC=B$ & 同秩、正负惯性指数相同 & 实对称同阶且正负惯性指数相同 \\
\bottomrule
\end{longtable}
\endgroup

\begin{callkey}{必背}
三条铁律：\textbf{相似 $\Rightarrow$ 等价}；\textbf{实对称矩阵相似 $\iff$ 合同}； 合同与相似都要求同阶，而等价只要求同型。

\end{callkey}

\subsection{特征值的性质}

\begingroup\small\setlength{\tabcolsep}{4.0pt}
\begin{longtable}{@{}L{0.6482\linewidth}L{0.2975\linewidth}@{}}
\toprule
\textbf{结论} & \textbf{说明} \\
\midrule
\endhead
$\sum_{i=1}^{n}\lambda_i=\operatorname{tr}(A)$，$\prod_{i=1}^{n}\lambda_i=\lvert A\rvert$ & 检验答案的两把尺子 \\
不同特征值对应的特征向量线性无关 & 常考证明 \\
$k$ 重特征值至多有 $k$ 个线性无关的特征向量 & 判定能否对角化的核心 \\
$A$ 可对角化 $\iff$ 每个 $k$ 重特征值恰有 $k$ 个线性无关特征向量 $\iff A$ 有 $n$ 个线性无关特征向量 & 充要条件 \\
$A$ 的特征值为 $\lambda$ $\Rightarrow$ $f(A)$ 的特征值为 $f(\lambda)$ & $f$ 为多项式 \\
$A$ 可逆时：$A^{-1}$ 的特征值为 $\frac1\lambda$，$A^*$ 的特征值为 $\frac{\lvert A\rvert}{\lambda}$ & 特征向量不变 \\
$A$ 与 $A^T$ 有相同特征值 & 但特征向量一般不同 \\
实对称矩阵的特征值全为实数，不同特征值的特征向量必正交 & 数二高频 \\
每行元素之和均为 $k$ $\Rightarrow$ $k$ 是一个特征值，$(1,1,\dots,1)^T$ 是对应特征向量 & 秒杀型结论 \\
\bottomrule
\end{longtable}
\endgroup

\subsection{二次型与正定性}

\textbf{标准形、规范形与惯性定理}：二次型 $f=x^TAx$（$A$ 实对称）经正交变换 $x=Py$ 化为 $\lambda_1y_1^2+\dots+\lambda_ny_n^2$（$\lambda_i$ 为 $A$ 的特征值），再经可逆线性变换化为规范形 $y_1^2+\dots+y_p^2-y_{p+1}^2-\dots-y_{r}^2$，其中 $p$ 为正惯性指数，$r$ 为二次型的秩，$p$ 与所作变换无关。

\textbf{正定的判别（$A$ 实对称，$n$ 阶）}：

\begin{enumerate}
\item 定义法：对任意 $x\ne0$ 有 $x^TAx>0$；
\item 特征值法：$A$ 的全部特征值 $>0$；
\item 顺序主子式法：$A$ 的各阶顺序主子式全 $>0$；
\item 合同法：$A$ 与单位阵 $E$ 合同，即存在可逆 $C$ 使 $A=C^TC$。
\end{enumerate}
\textbf{常用结论}：

\begin{itemize}
\item $A$ 正定 $\Rightarrow$ $\lvert A\rvert>0$，且主对角元全 $>0$（反之不然）；
\item $A$ 正定 $\Rightarrow$ $A^{-1}$、$A^*$、$A^k$ 均正定；
\item $A$ 正定 $\Rightarrow$ $A$ 的顺序主子阵全部正定；
\item 若 $A$ 正定、$B$ 实对称，则 $B^TAB$ 正定 $\iff B$ 可逆。
\end{itemize}
\begin{callwarn}{易错}
"正定"必须先说明 $A$ 是\textbf{实对称矩阵}——严格的定义只对实对称矩阵给出。 考试中若给定一般矩阵要求判断正定，先说明 $A^TA$ 之类的对称化处理。

\end{callwarn}

\subsection{高频秒杀结论与易混对照}

\begingroup\small\setlength{\tabcolsep}{4.0pt}
\begin{longtable}{@{}L{0.3483\linewidth}L{0.5974\linewidth}@{}}
\toprule
\textbf{题面信号} & \textbf{直接用} \\
\midrule
\endhead
$A$ 每行元素之和都是 $k$ & $k$ 是特征值，$(1,\dots,1)^T$ 是特征向量 \\
$r(A)=1$ & $A=\alpha\beta^T$，$\operatorname{tr}(A)=\beta^T\alpha$ 是非零特征值，其余为 $0$ \\
$A^2=A$（幂等） & 特征值只能是 $0$ 或 $1$，且 $A$ 可对角化 \\
$A^2=E$（对合） & 特征值只能是 $\pm1$，且 $A$ 可对角化 \\
$A$ 为 $n$ 阶正交矩阵 & $A^TA=E$，$\lvert A\rvert=\pm1$，列向量组标准正交 \\
已知 $A$ 的特征值求 $A^*$ & 特征值 $\frac{\lvert A\rvert}{\lambda}$，特征向量不变 \\
$\sum A_{ij}$ 型求和 & 构造新行列式（把求和列替换为全 $1$ 列） \\
\bottomrule
\end{longtable}
\endgroup

\textbf{最容易记错的 10 组对照}：

\begin{enumerate}
\item $\lvert kA\rvert=k^n\lvert A\rvert$，\textbf{不是} $k\lvert A\rvert$。
\item $(AB)^T=B^TA^T$，$(AB)^{-1}=B^{-1}A^{-1}$，$(AB)^*=B^*A^*$——顺序全部相反。
\item $A$ 与 $A^T$ 特征值相同，但特征向量不同。
\item 相似要求同阶可逆 $P$，合同要求同阶可逆 $C$ 且为 $C^TAC$。
\item 等价只要求同秩，合同还要求正负惯性指数相同。
\item 正交矩阵的行列式是 $\pm1$；行列式为 $\pm1$ 的矩阵不一定是正交矩阵。
\item 实对称矩阵必可对角化；一般矩阵不一定。
\item $A$ 可对角化 $\iff$ 有 $n$ 个线性无关特征向量，\textbf{不是}"有 $n$ 个特征值"。
\item 正定矩阵必须是实对称矩阵。
\item $r(A)=n-1$ 时 $A^*\ne O$ 且 $r(A^*)=1$；$r(A)\le n-2$ 时 $A^*=O$。
\end{enumerate}
\end{document}
